\documentclass[11pt]{article}
\usepackage[margin=1in]{geometry}
\usepackage{amsmath,amssymb,amsthm,amscd,mathtools,array,booktabs,enumitem,mathrsfs}
\usepackage[T1]{fontenc}
\usepackage{lmodern}
\usepackage[hidelinks]{hyperref}
\usepackage[nameinlink,noabbrev]{cleveref}

\newtheorem{theorem}{Theorem}[section]
\newtheorem{proposition}[theorem]{Proposition}
\newtheorem{lemma}[theorem]{Lemma}
\newtheorem{corollary}[theorem]{Corollary}
\theoremstyle{definition}
\newtheorem{definition}[theorem]{Definition}
\newtheorem{remark}[theorem]{Remark}

\newcommand{\C}{\mathbb C}
\newcommand{\R}{\mathbb R}
\newcommand{\Z}{\mathbb Z}
\newcommand{\Hh}{\mathbb H}
\newcommand{\Oo}{\mathbb O}
\newcommand{\PP}{\mathbb P}
\newcommand{\T}{\mathbb T}
\newcommand{\Herm}{\operatorname{Herm}}
\newcommand{\Mat}{\operatorname{Mat}}
\newcommand{\tr}{\operatorname{tr}}
\newcommand{\im}{\operatorname{im}}
\newcommand{\rank}{\operatorname{rank}}
\newcommand{\diag}{\operatorname{diag}}
\newcommand{\Span}{\operatorname{span}}
\newcommand{\adj}{\operatorname{adj}}
\newcommand{\id}{\operatorname{id}}
\newcommand{\one}{\mathbf 1}

\title{Retained cusp topology, triality morphisms, and the
marked Niemeier--Leech construction}
\author{Standalone theorem note for the higher-rung search}
\date{6 September 2026}

\begin{document}
\maketitle

\begin{abstract}
For the retained modulo-$840$ lattice $N$ with roots $A_5^4D_4$ and
its specified index-six Leech neighbour $\Lambda$, we compute the
arithmetic of the original Albert cubic, retaining its coefficient $2$.
Fix the Cayley--Dickson frame $(1,i,j,k,\ell,i\ell,j\ell,k\ell)$ and
the explicit isometry $\mathcal L:N_{\mathbb R}\to\mathbb O^3$.
Writing $\tau(q)=\operatorname{Re}((q_1q_2)q_3)$ and
$F=\tau\circ\mathcal L$ on the original Niemeier span,
let $\mathcal V_F(L)$ be the additive group generated by its values.
Then
\[
\begin{aligned}
 \mathcal V_F(N)&=\tfrac1{27}\mathbb Z\oplus
       \tfrac{\sqrt2}{54}\mathbb Z\oplus\tfrac{\sqrt3}{18}\mathbb Z,\\
 \mathcal V_F(\Lambda)&=\tfrac1{27}\mathbb Z\oplus
       \tfrac{\sqrt2}{216}\mathbb Z\oplus\tfrac{\sqrt3}{54}\mathbb Z,
 \qquad \mathcal V_F(\Lambda)/\mathcal V_F(N)\cong\mathbb Z/12.
\end{aligned}
\]
The determinant is $\lambda^3-\lambda(y,y)+2F(y)$; this value-group
inclusion is not a lattice inclusion $N\subset\Lambda$.
The residue operator $r=\rho(361)$ cycles the three neighbours
$\Lambda_j=r^j\Lambda$; all their pairwise intersections equal $J$,
with $\det J=81$ and $[\Lambda_j:J]=9$. A fixed fourth isotropic
plane gives the even unimodular $A_2^{12}$ lattice $M$, with $72$
roots, $M\cap\Lambda_j=J$ and $[M:J]=9$.
Its cubic-value group already equals that of the common lattice:
\[
 \mathcal V_F(J)=\mathcal V_F(M)
 =\tfrac1{27}\mathbb Z\oplus\tfrac{\sqrt2}{216}\mathbb Z
                         \oplus\tfrac{\sqrt3}{18}\mathbb Z.
\]
The two successive value-group inclusions have matrices
$\operatorname{diag}(1,4,1)$ and $\operatorname{diag}(1,1,3)$,
with quotients $\mathbb Z/4$ and $\mathbb Z/3$.

Wilson's Leech construction \cite{esn:Wilson2009} is embedded in the
same original metric by the specified isometry $x\mapsto x/\sqrt2$.
Its cyclic fixed Gram matrix is $3G_{E_8}$, its orthogonal complement
is explicitly isometric to $A_2\otimes E_8$, and their sum has index
$3^8$ with quotient $\mathbb F_3^8$. Retaining
$c(q_1,q_2,q_3)=(q_3,q_1,q_2)$, Wilson's opposite cycle is
$c_{\rm W}=c^{-1}$. The literal intersections with both
$\mathcal L(N)$ and $\mathcal L(\Lambda)$ equal
$\{\sqrt2(a,a,a):a\in D_{\mathbb H}\}$, where $D_{\mathbb H}$
is $D_4$ in the first four fixed Cayley coordinates. This rank-four
lattice has Gram matrix $6G_{D_4}$, determinant $5184$, minimum squared
norm $12$, and value-generated subgroup $4\sqrt2\mathbb Z$ for $\tau$
($8\sqrt2\mathbb Z$ for $2\tau$). Its complete image is
$2\sqrt2\bigl(18\mathbb Z\cup(6\mathbb Z+2)\cup(6\mathbb Z-2)\bigr)$,
with an explicit integral section. Writing $a=a_0+v$, $t=|v|$,
the corresponding Albert matrix retains the scalar $\mu$ and has
eigenvalues $\mu+2\sqrt2a_0$ and
$\mu-\sqrt2a_0\pm\sqrt6t$. Its three labelled projectors recover
$v$ exactly. The complete orbit-family proof keeps mean, radius and
Cartan angle, including the two focal projective-plane strata.
The marked umbral shadow is $(5I_{12}+J)S_6$, with
$J_{rs}=\mathbf1_{s=5r\ (12)}$; its coefficient map gives the
eta quotient $T_6$ and $T_6(-1/(6\tau))=72/T_6(\tau)$,
with the exact multiplier and original glue-action comparisons.

The higher target is three-directional, not required to carry a
complex structure. The ordered triality products retain the Jordan
factor $1/2$, mixed determinant derivative $2\tau$, and symmetric
polarization $\tau/3$, where $\tau(a,b,c)=\operatorname{Re}((ab)c)$.
The spectral graph identifies $\mathbb O^3$ with the equal-diagonal
positive determinant boundary. An explicit flag-cell chart and its
inverse give the degree-$+1$ collapse
$F_4/\operatorname{Spin}(8)\to S^{24}$; every reverse sphere map has
degree zero. The angle-forgetting cusp map bounds an explicit framed
$D^2\times T^2$. The retained class $[H]=\eta_4\eta_5\ne0$, under
the specified based $X\cong S^6$ identification, instead gives the
nontrivial real three-plane reduction inside the triality carrier and
the Hopf--cusp class $12\nu_6\ne0$ in $\pi_9(S^6)$.
The complex-toroidal comparison retains its original periods, kernel,
monodromy and fillings; its Euler contribution $6$ prevents the unchanged
wrap from closing on $S^8$. These comparisons do not by themselves
certify the claimed complex $S^6$ or an interacting quantum spectrum.
\end{abstract}

\section{The source construction and the toroidal comparison}

The higher-rung investigation asks how a structure with three interacting
directions can extend the source geometry. The present algebraic datum is
the vector and two half-spin sectors of triality, together with their
invariant trilinear tensor. The complex-torus extension studied first below
is one specified comparison with the source construction. Its complex
structure is a hypothesis of those toroidal theorems; it is not a
requirement on the higher-rung target. The retained cusp-angle map and the
marked $A_5^4D_4$ lattice are retained as inputs to the subsequent comparison
maps.

We use the following exact input from the reconstructed candidate
construction \cite{AlpogeS6_2026}.  Let
\[
 B=\PP^1,\qquad \Sigma=\{p_0,p_1,p_2\},\qquad
 B^\circ=B\setminus\Sigma,
\]
where $p_0$ is a cusp and $p_1,p_2$ have orbifold orders $3,4$.
There is a claimed compact connected complex threefold and proper
holomorphic map
\[
 f:X\longrightarrow B,
 \qquad X\cong_{\mathrm{diff}}S^6,
\]
whose regular fibres are complex two-tori.  In the basis
$(\gamma,u,w,\delta)$ of $V=\Z^4$ the integral monodromies are generated by
\begin{equation}
 T_1=
 \begin{pmatrix}1&0&-6&2\\0&-1&1&1\\0&-1&0&1\\0&0&0&1\end{pmatrix},
 \qquad
 T_2=
 \begin{pmatrix}1&6&0&-3\\0&0&-1&1\\0&1&0&0\\0&0&0&1\end{pmatrix},
 \label{eq:original-T}
\end{equation}
with $T_1^3=T_2^4=I$ and
$T_0=(T_1T_2)^{-1}=I+N$, $N^2=0$.  On the dual lattice
$\Lambda=V^*$ the monodromies are
$A_j=(T_j^{-1})^{\mathsf T}$.  The period matrix is
\begin{equation}
 \Pi(z)=
 \begin{pmatrix}6\mu(z)&\tau(z)&1&0\\
                 \beta(z)&\mu(z)&0&1\end{pmatrix},
 \qquad F_z=\C^2/\Pi(z)\Lambda .
 \label{eq:original-period}
\end{equation}
At $p_0$ the filling is a Mumford toric quotient with reduced irreducible
normal-crossing central fibre $W$: its normalization is the degree-six
del Pezzo surface, opposite sides of its anticanonical hexagon are
identified, and $\chi(W)=2$.  At $p_j$, $j=1,2$, the filling is a
logarithmic transform with
\[
 f^*p_1=3S_1,\qquad f^*p_2=4S_2,
\]
where the $S_j$ are smooth bielliptic surfaces and hence have Euler
characteristic zero.  The primitive affine translations have
\begin{equation}
 v_1=\widehat\gamma+2\widehat u-4\widehat w,\qquad
 v_2=-\widehat\gamma-3\widehat u+3\widehat w,\qquad
 (\ell_0,\ell_1,\ell_2)=(0,1,-1).
 \label{eq:original-twists}
\end{equation}
They give
$\pi_1(X)=\Z/|12\ell_0-4\ell_1-3\ell_2|=0$ and the integral homology of
$S^6$.

The later marked observation uses the integral difference matrix and its
fixed right inverse
\begin{equation}
 P=\begin{pmatrix}1&-1&0\\0&1&-1\end{pmatrix},\quad
 u=\begin{pmatrix}1\\1\\1\end{pmatrix},\quad
 C=PP^*=\begin{pmatrix}2&-1\\-1&2\end{pmatrix},\quad
 R=P^*C^{-1}=\frac13
 \begin{pmatrix}2&1\\-1&1\\-1&-2\end{pmatrix}.
 \label{eq:marked-P}
\end{equation}
Thus $PR=I_2$, $Pu=0$, $u^*R=0$, and for $X_P=(R\;u)$,
\begin{equation}
 X_P^{-1}=\begin{pmatrix}P\\u^*/3\end{pmatrix},
 \qquad \det X_P=1,
 \qquad X_P^*X_P=\diag(C^{-1},3).
 \label{eq:marked-inverse}
\end{equation}

\begin{definition}[Data for the complex-toroidal comparison]
\label[definition]{def:feature-ledger}
A \emph{full higher toroidal wrap} consists of all the following data and
proofs.  Neither its dimension nor its homeomorphism type is prescribed.
\begin{enumerate}[label=\textup{(F\arabic*)},leftmargin=3.2em]
\item A compact connected smooth carrier obtained from the stated pieces,
      with its smooth and homotopy type computed from the attaching maps.
\item A geometric tangent structure on the whole carrier.  For a literal
      toroidal continuation this is an integrable complex structure; a
      different category must state its replacement rather than inherit the
      word ``complex''.
\item A proper holomorphic map to the same one-dimensional base, with a
      complex torus of the dimension forced by total dimension as regular
      fibre.
\item Specified local spaces over the cusp and the two finite points,
      including their multiplicities, reduced fibres, attaching maps, and
      Euler characteristics.
\item An integral period lattice, monodromy representation, affine
      translations, and actual specialization maps at every filling.
\item The van Kampen and integral Mayer--Vietoris maps proving the claimed
      smooth sphere topology.
\item A marked observation whose domain, codomain, lattice map, and
      compatibility with periods and fillings are specified.
\item Rigidity: after all markings above are fixed, no auxiliary metric,
      phase section, fan identification, or gluing choice is left
      unrecorded, and existence and uniqueness up to the stated notion of
      isomorphism are proved.
\end{enumerate}
This definition specifies the complex-toroidal comparison studied in the
following sections. The three-sector triality construction is developed
through its own explicitly defined products, bundles, boundary, and lattice
maps below; the list above does not constrain it to be another complex
structure.
\end{definition}

For a rank-three successor over the regular base, require a complex-linear
surjection $P_z:\C^3\to\C^2$ with $K_{\C,z}=\ker P_z$, period embeddings
$\widetilde\Pi_z:\widetilde\Lambda\hookrightarrow\C^3$ and
$\Pi_z:\Lambda\hookrightarrow\C^2$, and
\[
 P_z\widetilde\Pi_z=\Pi_z\Phi,\qquad
 K_{\C,z}\cap\widetilde\Pi_z(\widetilde\Lambda)
 =\widetilde\Pi_z\kappa(K_\Z).
\]
Put
\[
 E_z=K_{\C,z}/\widetilde\Pi_z\kappa(K_\Z),\qquad
 \widetilde F_z=\C^3/\widetilde\Pi_z(\widetilde\Lambda),\qquad
 F_z=\C^2/\Pi_z(\Lambda).
\]
The desired portion of the map diagram is then
\begin{equation}
\begin{CD}
0@>>>K_\Z@>{\kappa}>>\widetilde\Lambda@>{\Phi}>>\Lambda@>>>0\\
@.@V{\widetilde\Pi_z\kappa}VV@V{\widetilde\Pi_z}VV@V{\Pi_z}VV@.\\
0@>>>K_{\C,z}@>>>\C^3@>{P_z}>>\C^2@>>>0\\
@.@V{q_K}VV@V{q_3}VV@V{q_2}VV@.\\
0@>>>E_z@>>>\widetilde F_z@>{\overline P_z}>>F_z@>>>0.
\end{CD}
\label{eq:required-diagram}
\end{equation}
All three rows are exact.  The maps $\overline P_z$ must assemble to a map
$\widetilde X^\circ\to X^\circ$ over $B^\circ$, together with commuting
specialization squares at all three deleted points.
Sections \ref{sec:torus-successor} and \ref{sec:sphere-ladder} prove exactly
which parts of this diagram exist.

\section{Sphere and Euler obstructions}

\begin{theorem}[The $S^8$ tangent obstruction with all hypotheses]
\label[theorem]{thm:s8-obstruction}
Let $M$ be a closed smooth manifold homotopy equivalent to $S^8$.  There is
no bundle endomorphism $J:TM\to TM$ satisfying $J^2=-\id$.  Consequently
no smooth manifold in that homotopy type is a complex fourfold, and in
particular no proper holomorphic complex-three-torus fibration over
$\PP^1$ can have such a total space.
\end{theorem}

\begin{proof}
Suppose $J$ exists and let $E=(TM,J)$, a complex rank-four bundle.  Orient
$M$ by $J$ and choose $a\in H^8(M;\Z)$ with
$\langle a,[M]\rangle=1$.  Since the cohomology below the top degree
vanishes,
\[
 c_1(E)=c_2(E)=c_3(E)=0.
\]
The top Chern class of a complex bundle is the Euler class of its underlying
oriented real bundle.  Poincar\'e--Hopf and $\chi(M)=2$ therefore give
\[
 c_4(E)=e(TM)=2a.
\]
For the underlying real bundle,
\[
 p_2(E_{\R})=c_2(E)^2-2c_1(E)c_3(E)+2c_4(E)=4a,
\]
whereas $p_1(TM)=0$ because $H^4(M;\Z)=0$.  The middle cohomology of $M$
vanishes, so $\operatorname{sign}(M)=0$.  Hirzebruch's degree-eight
signature formula now gives
\[
 0=\operatorname{sign}(M)
 =\left\langle\frac{7p_2(TM)-p_1(TM)^2}{45},[M]\right\rangle
 =\frac{28}{45},
\]
a contradiction.  The last displayed rational number is not being asserted
to be a possible signature: its nonintegrality is precisely the
contradiction forced by the incompatible characteristic classes.
\end{proof}


\begin{theorem}[Parity and the complete higher-sphere exclusion]
\label[theorem]{thm:all-higher-spheres}
Let $M$ be a smooth homotopy $m$-sphere with $m>6$.
Then $M$ admits no almost-complex structure and hence no integrable complex
structure.  In particular $S^9$ cannot be a complex carrier, and there is
no different higher sphere dimension to which the literal complex part of
the candidate can be moved.
\end{theorem}

\begin{proof}
If $m$ is odd, an endomorphism $J$ with $J^2=-I$ would make each real tangent
space a complex vector space, whose real dimension is even.  Hence $m$ must
be even.  Write $m=2n$; then $m>6$ gives $n\ge4$.

If $E=(TM,J)$ existed, $c_i(E)=0$ for $i<n$ and, after using the complex
orientation, $c_n(E)=2a$ for the integral top generator $a$.  Pull $E$ back
along a homotopy equivalence $S^{2n}\to M$.  Bott integrality says that the
top Chern class of every complex bundle on $S^{2n}$ is divisible by
 $(n-1)!$ \cite[Chapter~17, Section~7, Theorem~7.2,
 pp.~274--275]{Husemoller1994}.  Since $(n-1)!\ge6$ does not divide $2$,
 this contradicts $c_n(E)=2a$.  Borel and Serre prove the same conclusion for every
differentiable structure on $S^{2n}$ in \cite[Section~15,
Proposition~15.1 and the following remark, p.~434]{BorelSerre1953}; their
universal reduced-power input is \cite[Section~12, Propositions~12.4 and
12.6 and Lemma~12.5, p.~430]{BorelSerre1953}.
\end{proof}

\begin{proposition}[The direct rank-three Euler mismatch]
\label[proposition]{prop:euler-mismatch}
Consider the direct rank-three toroidal local model obtained by coning the
periodic Kuhn triangulation of $\R^3$ and dividing by $\Z^3$.  Its central
fibre $W_3$ has
\[
 \chi(W_3)=6.
\]
If a proper triangulable map $\widetilde f:\widetilde X^8\to S^2$ has this
as one special fibre, has precisely two other special fibres that are free
finite quotients of $T^6$, and is a $T^6$-bundle elsewhere, then
$\chi(\widetilde X)=6$.  Therefore the direct three-fibre pattern cannot
have total space $S^8$ even after forgetting complex structures.
\end{proposition}

\begin{proof}
A fundamental cube of the Kuhn triangulation is divided into the six
tetrahedra indexed by permutations of its three coordinate directions.
The central toric fibre is stratified by algebraic tori; every
positive-dimensional torus stratum has Euler characteristic zero, while
the zero-dimensional strata are indexed by the maximal tetrahedra modulo
$\Z^3$.  There are exactly six, so additivity over the finite torus
stratification gives $\chi(W_3)=6$.

Let $\Sigma_f$ be the three critical values.  Euler integration, or a cut
into the regular bundle and three filling pieces, gives
\[
 \chi(\widetilde X)
 =\chi(S^2\setminus\Sigma_f)\chi(T^6)
  +\sum_{p\in\Sigma_f}\chi(\widetilde f^{-1}(p)).
\]
The regular term vanishes since $\chi(T^6)=0$.  A free finite quotient of a
torus has Euler characteristic zero: the covering formula multiplies its
Euler characteristic by the positive degree and the covering torus has
Euler characteristic zero.  Only $W_3$ remains, giving six.  Since
$\chi(S^8)=2$, the total spaces cannot be homeomorphic.
\end{proof}

\begin{corollary}[The exact singular-fibre repair equation]
\label[corollary]{cor:euler-repair}
Let $g:M\to S^2$ be a proper triangulable map which is locally trivial
with fibre $T^6$ off a finite set $\Delta$, and suppose the direct Kuhn
cusp fibre $W_3$ occurs over $p_0\in\Delta$.  Then
\begin{equation}
 \chi(M)=6+\sum_{p\in\Delta\setminus\{p_0\}}
                    \chi\bigl(g^{-1}(p)\bigr).
 \label{eq:euler-repair-general}
\end{equation}
Consequently an $S^8$ carrier requires, exactly at the Euler-characteristic
level,
\begin{equation}
 \sum_{p\in\Delta\setminus\{p_0\}}
                    \chi\bigl(g^{-1}(p)\bigr)=-4.
 \label{eq:euler-repair-minus-four}
\end{equation}
If there are only the original two finite values, their replacement fibres
must have Euler characteristics summing to $-4$.  If those two fibres remain
free finite quotients of $T^6$, they still contribute zero, so the cusp
fibre itself must instead be replaced by a filling with Euler
characteristic $2$.  Without either change, topology is an output and every
closed carrier of the unchanged wrap necessarily has Euler characteristic
$6$.
\end{corollary}

\begin{proof}
Euler integration gives
$\chi(M)=\chi(S^2\setminus\Delta)\chi(T^6)
+\sum_{p\in\Delta}\chi(g^{-1}(p))$.  The first term is zero and
$\chi(W_3)=6$ by \cref{prop:euler-mismatch}, giving
\eqref{eq:euler-repair-general}.  Substitution of $\chi(S^8)=2$ gives
\eqref{eq:euler-repair-minus-four}.  The remaining assertions are its two
specializations, using the covering calculation in the proof of
\cref{prop:euler-mismatch}.
\end{proof}

\begin{remark}[What remains possible topologically]
A merely real singular map $S^8\to S^2$ with regular fibre $T^6$ is not
excluded by \cref{thm:s8-obstruction}; it must have special-fibre Euler sum
$2$.  It cannot be a fibre bundle: the homotopy sequence would require the
connecting map $\pi_2(S^2)=\Z\to\pi_1(T^6)=\Z^6$ to be surjective.  A real
$S^9\to S^2$ model with regular fibre $T^7$ would likewise require genuine
singularities and special-fibre Euler sum $0$, but neither its carrier nor
its odd-dimensional regular fibre can be complex in the required sense.
\end{remark}

\section{The exact toroidal successor over the cusp}
\label{sec:torus-successor}

We now retain, rather than discard, every exact map that survives the
obstructions.  Put
\[
 A_F=\Z\langle\widehat w,\widehat\delta\rangle,\qquad
 B_F=\Z\langle\widehat\gamma,\widehat u\rangle,\qquad
 B_0=\begin{pmatrix}0&1\\-1&0\end{pmatrix}:B_F\to A_F.
\]
In angular--deck order the candidate cusp monodromy and filling map are
\begin{equation}
 M_F(a,b)=(a+B_0b,b),\qquad
 c_F:A_F\oplus B_F\longrightarrow B_F,\quad c_F(a,b)=b.
 \label{eq:candidate-cusp-module}
\end{equation}
Let $L_3=\Z^3$, let $k=(1,1,1)^{\mathsf T}$, and put
\begin{equation}
 P=\begin{pmatrix}1&-1&0\\0&1&-1\end{pmatrix},\qquad
 Q=B_0^{-1}P=\begin{pmatrix}0&-1&1\\1&-1&0\end{pmatrix}.
 \label{eq:PQ}
\end{equation}
Both maps are onto, $\ker P=\ker Q=\Z k$, and $B_0Q=P$.
The rank-three cusp module and its specialization are
\begin{equation}
 M_3(a,b)=(a+b,b),\qquad
 c_3:L_3\oplus L_3\longrightarrow L_3,\quad c_3(a,b)=b.
 \label{eq:rank3-cusp-module}
\end{equation}

\begin{theorem}[Primitive lattice, monodromy, and filling quotient]
\label[theorem]{thm:lattice-quotient}
The homomorphism
\[
 \Phi=P\oplus Q:L_3\oplus L_3\longrightarrow A_F\oplus B_F,
 \qquad \Phi(a,b)=(Pa,Qb),
\]
fits into the exact sequence
\begin{equation}
 0\longrightarrow\Z^2
 \xrightarrow{\ \kappa\ }\Z^6
 \xrightarrow{\ \Phi\ }\Z^4\longrightarrow0,
 \qquad \kappa(m,n)=(mk,nk),
 \label{eq:lattice-exact}
\end{equation}
and satisfies
\begin{equation}
 \Phi M_3=M_F\Phi,\qquad c_F\Phi=Qc_3.
 \label{eq:mono-fill-compatibility}
\end{equation}
The kernel in \eqref{eq:lattice-exact} is primitive.
\end{theorem}

\begin{proof}
The integer matrix
\[
 S=\begin{pmatrix}1&1\\0&1\\0&0\end{pmatrix}
\]
satisfies $PS=I_2$ and $QSB_0=I_2$.  Thus $P$ and $Q$, and hence $\Phi$,
are integrally surjective.  Their displayed kernels give
$\ker\Phi=\Z(k,0)\oplus\Z(0,k)$, proving exactness.  A kernel of a
surjection between free abelian groups is primitive.

Directly,
\[
 \Phi M_3(a,b)=(P(a+b),Qb),
\]
while
\[
 M_F\Phi(a,b)=(Pa+B_0Qb,Qb)=(Pa+Pb,Qb).
\]
The filling identity is equally literal:
$c_F\Phi(a,b)=Qb=Qc_3(a,b)$.
\end{proof}

Let $t=e^{2\pi i s}$ on a logarithmic cusp cover.  Retain the complete
bounded holomorphic correction $C_2(t)$ of the candidate and define
\begin{equation}
 Z_2=sB_0+C_2(t),\qquad
 C_3(t)=S C_2(t)Q,\qquad
 Z_3=sI_3+C_3(t).
 \label{eq:period-extension}
\end{equation}
We use the angular--deck period convention
$\Omega_r=(I_r\;Z_r)$, and write $\Omega_r(t)$ when the dependence through
$s$ and $C_r(t)$ must be displayed.  The manuscript convention
$(Z_r\;I_r)$ is related by the displayed integral block permutation
$\left(\begin{smallmatrix}0&I_r\\I_r&0\end{smallmatrix}\right)$; no
period, sign, or lattice coordinate is silently changed.

\begin{theorem}[Holomorphic elliptic extension and its exact fibre]
\label[theorem]{thm:elliptic-extension}
After shrinking the cusp disc, $\operatorname{Im}Z_3$ is invertible and
\begin{equation}
 PZ_3=Z_2Q,\qquad P\Omega_3=\Omega_2\Phi.
 \label{eq:period-intertwining}
\end{equation}
Consequently $P:\C^3\to\C^2$ descends to a short exact sequence of complex
tori
\begin{equation}
 0\longrightarrow E_s\longrightarrow
 \frac{\C^3}{\Z^3+Z_3\Z^3}
 \xrightarrow{\ \overline P_s\ }
 \frac{\C^2}{\Z^2+Z_2\Z^2}
 \longrightarrow0,
 \qquad E_s\cong\frac{\C}{\Z+s\Z}.
 \label{eq:torus-short-exact}
\end{equation}
The correction $C_2$ is recoverable from its prescribed lift by the exact
left inverse
\begin{equation}
 C_2=P C_3 S B_0.
 \label{eq:correction-inverse}
\end{equation}
Thus the extension neither deletes the original periods nor replaces them
by an identity period matrix.
\end{theorem}

\begin{proof}
Since $C_3$ is bounded while $\operatorname{Im}s\to+\infty$,
$\operatorname{Im}Z_3=(\operatorname{Im}s)I_3+\operatorname{Im}C_3$ is
invertible on a sufficiently small cusp.  Using $PS=I_2$ and $B_0Q=P$,
\[
 PZ_3=sP+PSC_2Q=sB_0Q+C_2Q=Z_2Q,
\]
which is exactly \eqref{eq:period-intertwining}.  It makes the complex
linear map $P$ a homomorphism of the displayed quotient tori.  It is
surjective on vector spaces and its lattice image is the whole target
lattice because $\Phi$ is onto; therefore its kernel is connected.

The vector-space kernel is $\C k$.  Since $Qk=0$, one has $C_3k=0$ and
$Z_3k=sk$.  If a source period lies in $\C k$, its image under $P$ is the
zero target period.  The real period map $\Omega_2:\R^4\to\C^2$ is an
isomorphism, so \eqref{eq:period-intertwining} and
$\ker\Phi=\Z(k,0)\oplus\Z(0,k)$ show that all such periods are precisely
$\Z k+s\Z k$.  This proves the fibre formula.

Finally $QS=B_0^{-1}$, whence
\[
 PC_3SB_0=PSC_2QSB_0=C_2.
\]
\end{proof}

\subsection{The cusp filling map, not merely its boundary lattice}

For $m\in\Z^3$ and a permutation
$\pi=(\pi_1,\pi_2,\pi_3)\in S_3$, let
\begin{equation}
 \sigma_{m,\pi}=\operatorname{conv}
 \{m,m+e_{\pi_1},m+e_{\pi_1}+e_{\pi_2},m+e_1+e_2+e_3\}.
 \label{eq:kuhn-tetrahedron}
\end{equation}
These are the tetrahedra of the periodic Kuhn triangulation; denote that
triangulation by $\mathcal T_3$.  Let $\mathcal T_2$ be the periodic
triangulation of $\R^2$ whose triangles are
\[
 \operatorname{conv}\{u,u+e_1,u+e_2\},\qquad
 \operatorname{conv}\{u+e_1,u+e_2,u+e_1+e_2\},
 \qquad u\in\Z^2.
\]
For $r=2,3$, let $\mathfrak S_r$ be the fan in $\Z^r\oplus\Z$ consisting
of all faces of the cones over $\tau\times\{1\}$ for
$\tau\in\mathcal T_r$, and put $Y_r=X_{\mathfrak S_r}$.  In the target
lattice write
\[
 n_1=\widehat w,\qquad n_2=\widehat\delta-\widehat w,\qquad
 n_3=-\widehat\delta;\qquad n_1+n_2+n_3=0.
\]
The map $P$ in \eqref{eq:PQ} is exactly $P(e_i)=n_i$.

\begin{proposition}[Exact fan and quotient-filling morphism]
\label[proposition]{prop:toric-filling-map}
Coning the Kuhn triangulation at height one gives a smooth toric fourfold
$Y_3$ with height map locally equal to
\[
 t=z_0z_1z_2z_3.
\]
The integral map
\[
 \widehat P:\Z^3\oplus\Z\longrightarrow A_F\oplus\Z,\qquad
 \widehat P(x,r)=(Px,r),
\]
is a map from this fan to the candidate's periodic $A_2$ fan and hence
induces a toric morphism $F:Y_3\to Y_2$ over the $t$-line.

On the dense torus define
\begin{align}
 \Psi^{(3)}_\lambda(x,t)
  &=\bigl(\exp(2\pi iC_3(t)\lambda)t^\lambda x,t\bigr),
       &&\lambda\in\Z^3,\label{eq:source-toric-action}\\
 \Psi^{(2)}_\mu(y,t)
  &=\bigl(\exp(2\pi iC_2(t)\mu)t^{B_0\mu}y,t\bigr),
       &&\mu\in\Z^2.\label{eq:target-toric-action}
\end{align}
Then
\begin{equation}
 F\Psi^{(3)}_\lambda=\Psi^{(2)}_{Q\lambda}F.
 \label{eq:toric-action-intertwining}
\end{equation}
For a sufficiently small open cusp disc
$\Delta_\epsilon=\{t\in\C:|t|<\epsilon\}$, set
\begin{equation}
 Y_{r,\epsilon}=t^{-1}(\Delta_\epsilon),\qquad
 N_3^C=Y_{3,\epsilon}/\{\Psi^{(3)}_\lambda:\lambda\in\Z^3\},\qquad
 N_0=Y_{2,\epsilon}/\{\Psi^{(2)}_\mu:\mu\in\Z^2\}.
 \label{eq:analytic-filling-quotients}
\end{equation}
After shrinking $\Delta_\epsilon$, both actions are free and properly
discontinuous, and $F$ descends to a holomorphic filling morphism
\begin{equation}
 \overline F:N_3^C\longrightarrow N_0
 \label{eq:filling-morphism}
\end{equation}
whose map on smooth-fibre $H_1$ is $\Phi$ and whose monodromy compatibility
is \eqref{eq:mono-fill-compatibility}.  Under
$H_1(T^6)=L_3\oplus L_3$ and $H_1(T^4)=A_F\oplus B_F$, the angular
generators bound the toric coordinate discs while the deck generators
survive.  Hence the inclusion-induced specialization maps are exactly
$c_3(a,b)=b$ and $c_F(a,b)=b$, and they form the commutative square
\begin{equation}
\begin{CD}
H_1(T^6;\Z)@>{\Phi}>>H_1(T^4;\Z)\\
@V{c_3}VV@VV{c_F}V\\
\Z^3@>{Q}>>\Z^2.
\end{CD}
\label{eq:filling-square}
\end{equation}
\end{proposition}

\begin{proof}
The four vertices of \eqref{eq:kuhn-tetrahedron} map to
\[
 Pm,\quad Pm+n_{\pi_1},\quad Pm-n_{\pi_3},\quad Pm.
\]
In target coordinates
$(n_1,n_2,n_3)=(e_1,e_2-e_1,-e_2)$.  For the six permutations the three
distinct images are the six $A_2$ triangles incident to $Pm$:
\[
\begin{array}{c|c@{\qquad}c|c}
(1,2,3)&\{v,v+e_1,v+e_2\}&(2,1,3)&\{v,v-e_1+e_2,v+e_2\}\\
(1,3,2)&\{v,v+e_1,v+e_1-e_2\}&(3,1,2)&\{v,v-e_2,v+e_1-e_2\}\\
(2,3,1)&\{v,v-e_1+e_2,v-e_1\}&(3,2,1)&\{v,v-e_2,v-e_1\}.
\end{array}
\]
Thus every source cone maps into a target cone.  The toric fan criterion
gives $F$, and preservation of the height coordinate gives $t\circ F=t$.
The maximal Kuhn cones are unimodular, giving the displayed local monomial
equation and smoothness of $Y_3$.

On the dense torus, $PC_3=C_2Q$ and $P=B_0Q$ give
\[
 F\Psi^{(3)}_\lambda(x,t)
 =\bigl(\exp(2\pi iC_2Q\lambda)t^{B_0Q\lambda}F(x),t\bigr),
\]
which proves \eqref{eq:toric-action-intertwining}; the fan map extends it
over every toric orbit.

For the analytic quotient, put
$y=-\log|t|/(2\pi)$ and
\[
 \operatorname{trop}_t(x)=\frac{1}{\log|t|}
       (\log|x_1|,\log|x_2|,\log|x_3|).
\]
The exact transformation law is
\[
 \operatorname{trop}_t(\Psi^{(3)}_\lambda x)
 =\operatorname{trop}_t(x)+A(t)\lambda,\qquad
 A(t)=I_3+\frac{\operatorname{Im}C_3(t)}{y}.
\]
Boundedness of $C_3$ permits the cusp to be shrunk until
\[
\tfrac12\|\lambda\|\le\|A(t)\lambda\|
 \le\tfrac32\|\lambda\|.
\]
For the target use
\[
 \operatorname{trop}^{(2)}_t(y)
 =\frac{1}{\log|t|}(\log|y_1|,\log|y_2|).
\]
Its exact transformation law is
\[
 \operatorname{trop}^{(2)}_t(\Psi^{(2)}_\mu y)
 =\operatorname{trop}^{(2)}_t(y)+A_2(t)\mu,
 \qquad A_2(t)=B_0+\frac{\operatorname{Im}C_2(t)}{y}.
\]
Since $B_0$ is orthogonal, after one further shrinking the same bounds hold
with $A_2(t),\mu$ in place of $A(t),\lambda$.  Therefore a compact set can
meet only finitely many of its translates for either action.
On a boundary orbit the bounded cell label is translated by $\lambda$, so
the same conclusion holds there; the target argument uses the invertible
lattice map $B_0$ in exactly the same way.  A fixed point in a smooth fibre would
give $\operatorname{Im}Z_3\lambda=0$, hence $\lambda=0$; equality of a
boundary cell with its translate likewise forces $\lambda=0$.  The target
equalities give $\operatorname{Im}Z_2\mu=0$ in the smooth fibre and
$B_0\mu=0$ on the boundary, hence $\mu=0$.  Rounding
the coordinates of $A(t)^{-1}\operatorname{trop}_t(x)$ gives a uniformly
 bounded set of representatives.  The closures of these representatives lie
 in the finite union of closed toric polydiscs belonging to the six Kuhn
 tetrahedra; that finite union is compact.  This proves proper
discontinuity.  The same construction with $A_2(t)$ supplies target
representatives.  Over every smaller closed disc
$\overline\Delta_{\epsilon'}$, $0<\epsilon'<\epsilon$, the displayed
finite unions are compact, proving properness there.  The equivariant $F$
therefore descends, and
\eqref{eq:filling-square} is exactly the second identity in
\eqref{eq:mono-fill-compatibility}.
\end{proof}

\subsection{Finite peripheral data and the precise globalization gap}

In the manuscript lattice order the candidate matrices are
\begin{equation}
 A_1=\begin{pmatrix}1&0&0&0\\6&0&1&0\\-6&-1&-1&0\\-2&1&0&1\end{pmatrix},
 \qquad
 A_2=\begin{pmatrix}1&0&0&0\\0&0&-1&0\\-6&1&0&0\\3&0&1&1\end{pmatrix}.
 \label{eq:finite-candidate-matrices}
\end{equation}
The matrices $A_1,A_2$ are written in the manuscript order
$(\widehat\gamma,\widehat u,\widehat w,\widehat\delta)=B_F\oplus A_F$,
whereas $M_F$ in \eqref{eq:candidate-cusp-module} is written in the
angular--deck order $A_F\oplus B_F$.  With the explicit block swap
\begin{equation}
 J=\begin{pmatrix}0&I_2\\ I_2&0\end{pmatrix},\qquad
 M_F^{\rm man}=JM_FJ=
 \begin{pmatrix}I_2&0\\ B_0&I_2\end{pmatrix},
 \label{eq:manuscript-cusp-order}
\end{equation}
one has
$A_1^3=A_2^4=I$ and $A_1A_2M_F^{\rm man}=I$.  Define
\begin{equation}
 C_1=\begin{pmatrix}-1&-1\\1&0\end{pmatrix},\quad
 C_2'=\begin{pmatrix}0&1\\-1&0\end{pmatrix},\quad
 C_0=\begin{pmatrix}1&1\\0&1\end{pmatrix}.
\label{eq:elliptic-peripheral-matrices}
\end{equation}
For uniform notation below set $D_1=C_1$ and $D_2=C_2'$.
Then $(C_1)^3=(C_2')^4=I$, $C_1C_2'=C_0^{-1}$, and
$(C_0-I)^2=0$ with rank one.

The kernel block in this presentation is not an unrelated direct summand of
the rank-three cusp lattice.  Put
\begin{equation}
 \begin{gathered}
 r=\begin{pmatrix}0&0&1\end{pmatrix},\qquad
 A=\begin{pmatrix}P\\r\end{pmatrix}
 =\begin{pmatrix}1&-1&0\\0&1&-1\\0&0&1\end{pmatrix},\qquad
 B=\begin{pmatrix}Q\\r\end{pmatrix}
 =\begin{pmatrix}0&-1&1\\1&-1&0\\0&0&1\end{pmatrix},\\
 D=B_0\oplus[1]
 =\begin{pmatrix}0&1&0\\-1&0&0\\0&0&1\end{pmatrix},\qquad
 U_6=\diag(A,B)\in SL(6,\Z).
 \end{gathered}
 \label{eq:rank-six-basis-change}
\end{equation}

\begin{lemma}[The identical cusp lattice in quotient--kernel coordinates]
\label[lemma]{lem:cusp-basis-identification}
Let
\[
 \rho(Pa,ra;Qb,rb)=(Pa,Qb;ra,rb),\qquad
 \rho_{\rm man}=(J\oplus I_2)\rho.
\]
Then
\begin{align}
 U_6M_3U_6^{-1}
 &=\begin{pmatrix}I_3&D\\0&I_3\end{pmatrix},
 \label{eq:rank-six-conjugacy}\\
 \rho U_6M_3U_6^{-1}\rho^{-1}
 &=M_F\oplus C_0,
 \label{eq:angular-kernel-splitting}\\
 \rho_{\rm man}U_6M_3U_6^{-1}\rho_{\rm man}^{-1}
 &=M_F^{\rm man}\oplus C_0.
 \label{eq:manuscript-kernel-splitting}
\end{align}
Projection onto the candidate block is exactly
$\operatorname{pr}_{1,2}\rho U_6(a,b)=(Pa,Qb)=\Phi(a,b)$ in
angular--deck coordinates and $J\Phi(a,b)$ in manuscript coordinates.
\end{lemma}

\begin{proof}
Both $A$ and $B$ have determinant one, and direct multiplication gives
$A=DB$, so $AB^{-1}=D$.  Conjugating
$M_3=\left(\begin{smallmatrix}I_3&I_3\\0&I_3\end{smallmatrix}\right)$
by $U_6$ proves \eqref{eq:rank-six-conjugacy}.  The first two rows of
$A,B$ are exactly $P,Q$, while $r(1,1,1)^{\mathsf T}=1$; regrouping the
first two and last coordinates therefore gives
\eqref{eq:angular-kernel-splitting}.  Applying the already displayed block
swap $J$ proves \eqref{eq:manuscript-kernel-splitting} and both projection
formulas.
\end{proof}

\begin{proposition}[Exact finite-order affine lift]
\label[proposition]{prop:finite-lift}
On $\widetilde\Lambda=\Lambda\oplus\Z^2$, identified with the original
$\Z^3\oplus\Z^3$ by \cref{lem:cusp-basis-identification}, set
\[
 \widetilde A_1=A_1\oplus C_1,\qquad
 \widetilde A_2=A_2\oplus C_2',\qquad
 \widetilde M_0=M_F^{\rm man}\oplus C_0.
\]
Then
\[
 \widetilde A_1^3=I,\qquad \widetilde A_2^4=I,\qquad
 \widetilde A_1\widetilde A_2\widetilde M_0=I.
\]
The affine transformations
\begin{equation}
 \widetilde g_j(x,y)=(A_jx+v_j/m_j,D_jy),\qquad
 (m_1,m_2)=(3,4),
 \label{eq:affine-lifts}
\end{equation}
of $T^4\times T^2$ have exact orders $m_j$, act freely, and project to
the candidate finite-fibre actions.
\end{proposition}

\begin{proof}
All three matrix identities follow blockwise from the identities just
displayed and the candidate identities.  The vector $(v_j,0)$ is fixed by
$\widetilde A_j$.  The $m_j$-th power of \eqref{eq:affine-lifts} is
translation by the integral vector $(v_j,0)$ and hence is the identity on
the torus.  If a smaller positive power fixed a point, its projection to
$T^4$ would fix a point under the corresponding smaller power of the
candidate affine transformation.  The logarithmic-transform action in the
candidate is free, so this is impossible.  Projection to the first block is
equivariant by construction.
\end{proof}

\begin{remark}[Exact endpoint]
Theorems \ref{thm:lattice-quotient}--\ref{thm:elliptic-extension} and
Propositions \ref{prop:toric-filling-map}--\ref{prop:finite-lift} prove the
regular cusp, the entire toric cusp filling, and compatible finite
\emph{integral affine} data.  They do not construct one holomorphic
rank-three period family on the full $(3,4,\infty)$ cover whose two finite
analytic fillings glue to the cusp family.  Even if that analytic gluing is
constructed, \cref{prop:euler-mismatch} proves that the unchanged
three-fibre pattern has Euler characteristic six and hence is not $S^8$.
This is a proved morphism with a proved endpoint, not an assertion that the
rank-three and rank-two objects are unrelated.
\end{remark}

\section[The forced S9, S10, and rank-loss S8]
{The forced $S^9$, $S^{10}$, and rank-loss $S^8$}
\label{sec:sphere-ladder}

Put
\[
 \mathscr V=\Mat_{3\times3}(\C),\qquad
 \mathcal P_3=\{H\in\Herm_3(\C):H\succeq0\},\qquad
 \mathcal K_3=\{H\in\mathcal P_3:\tr H=1\}.
\]
The superscript $+$ denotes one-point compactification.

\begin{theorem}[Complete $U(3)$ and $SU(3)$ orbit invariants]
\label[theorem]{thm:orbit-invariants}
The maps
\begin{align}
 \Gamma_U(A)&=A^*A,\label{eq:U-gram}\\
 \Gamma_{SU}(A)&=(A^*A,\det A)\label{eq:SU-gram}
\end{align}
induce homeomorphisms
\begin{align}
 \mathscr V/U(3)&\xrightarrow{\sim}\mathcal P_3,\label{eq:U-quotient}\\
 \mathscr V/SU(3)&\xrightarrow{\sim}
 \mathcal Y_3:=\{(H,z)\in\mathcal P_3\times\C:|z|^2=\det H\}.
 \label{eq:SU-quotient}
\end{align}
For a rank-$r$ matrix $A$, the two orbit fibres and stabilizers are
\[
 U(3)\cdot A\cong U(3)/U(3-r),\qquad
 SU(3)\cdot A\cong SU(3)/SU(3-r),
\]
with $SU(1)$ interpreted as the trivial group.
\end{theorem}

\begin{proof}
Both maps are invariant under the indicated left actions, and
$|\det A|^2=\det(A^*A)$.  Every $H\succeq0$ is the Gram matrix of
$H^{1/2}$.  If $A^*A=B^*B$, the rule $Ax\mapsto Bx$ is well-defined because
the two matrices have the same kernel, and is isometric because
\[
 \langle Bx,By\rangle=x^*B^*By=x^*A^*Ay=\langle Ax,Ay\rangle.
\]
Extending it orthogonally gives $U\in U(3)$ with $UA=B$.  Thus
\eqref{eq:U-gram} separates $U(3)$-orbits.

Given $(H,z)\in\mathcal Y_3$, if $H$ is singular then $z=0$ and
$H^{1/2}$ realizes it.  If $H$ is positive definite, put
\[
 \omega=z/\sqrt{\det H},\qquad
 A=\diag(\omega,1,1)H^{1/2}.
\]
Then $(A^*A,\det A)=(H,z)$.  If two matrices have both invariants equal,
the unitary $U$ above has determinant one when they are invertible.  In
rank at most two, multiply $U$ on the left by a unitary that fixes
$\im B$ pointwise and corrects the determinant on its nonzero orthogonal
complement.  Hence \eqref{eq:SU-gram} separates $SU(3)$-orbits.

Both invariant maps are proper because $\tr(A^*A)=\|A\|_F^2$.  The induced
continuous bijections to the displayed closed Hausdorff sets are therefore
homeomorphisms.  A left stabilizer fixes $\im A$ pointwise and is arbitrary
unitary, respectively special unitary, on $(\im A)^\perp$, proving the
fibre statements.
\end{proof}

Let
\[
 E=\{X\in\Herm_3(\C):\tr X=0\},\qquad \dim_{\R}E=8,\qquad c=I_3/3.
\]
Equip $E$ with the Frobenius inner product
$\langle X,Y\rangle_F=\tr(XY)$ and its norm; $B(E)$ and $S(E)$ below
denote its closed unit ball and unit sphere.
For $v\in S(E)$ define
\begin{equation}
 R_*(v)=-\frac{1}{3\lambda_{\min}(v)}.
 \label{eq:radial-length}
\end{equation}
Since a nonzero traceless Hermitian matrix has a negative eigenvalue,
$R_*(v)$ is finite and positive.

\begin{lemma}[Explicit eight-ball chart]
\label[lemma]{lem:K3-ball}
The maps
\begin{align}
 \chi(c)&=0,&
 \chi(c+s v)&=\frac{s}{R_*(v)}v,\quad0<s\le R_*(v),
 \label{eq:chi}\\
 \chi^{-1}(rv)&=c+rR_*(v)v,\quad0<r\le1
 \label{eq:chi-inverse}
\end{align}
are inverse homeomorphisms $\mathcal K_3\cong B(E)\cong B^8$.  They send
the rank-deficient locus onto $S(E)$.
\end{lemma}

\begin{proof}
The eigenvalues of $c+sv$ are $1/3+s\lambda_j(v)$.  Hence this matrix is
positive semidefinite exactly for $0\le s\le R_*(v)$ and is singular
exactly at the upper endpoint.  Every $H\ne c$ in $\mathcal K_3$ lies on
one unique such ray.  Continuity of the least eigenvalue makes
$v\mapsto-\lambda_{\min}(v)$ a positive continuous function on the compact
space $S(E)$, hence bounded above and bounded away from zero.  Thus
$R_*(v)$ is continuous and uniformly bounded above and away from zero.
The formulas \eqref{eq:chi}--\eqref{eq:chi-inverse} are consequently
continuous, including at $c$ and at the origin, and prove all assertions.
\end{proof}

\begin{theorem}[The forced sphere ladder and explicit inverses]
\label[theorem]{thm:sphere-ladder}
Let $S_F^{17}=\{A\in\mathscr V:\|A\|_F=1\}$.  Then
\begin{equation}
 S_F^{17}/SU(3)\cong
 \mathcal Y_{3,1}:=\{(H,z)\in\mathcal Y_3:\tr H=1\}\cong S^9.
 \label{eq:S9}
\end{equation}
An explicit second homeomorphism is
\begin{equation}
 \Omega(H,z)=
 \begin{cases}
 \left(\chi(H),\sqrt{1-\|\chi(H)\|^2}\,z/|z|\right),&H\succ0,\\
 (\chi(H),0),&H\text{ singular},
 \end{cases}
 \label{eq:Omega-S9}
\end{equation}
from $\mathcal Y_{3,1}$ to $S(E\oplus\C)$.

Furthermore the weighted radial map
\begin{equation}
 \Phi_{SU}(H,z)=
 \begin{cases}
 0,&(H,z)=(0,0),\\
 s\,\Omega(H/s^2,z/s^3),&s=\sqrt{\tr H}>0
 \end{cases}
 \label{eq:weighted-radial-map}
\end{equation}
is a homeomorphism $\mathcal Y_3\cong E\oplus\C\cong\R^{10}$.  It extends
to one-point compactifications and gives a homeomorphism of pairs
\begin{equation}
 \left(\mathscr V^+/SU(3),\{\det A=0\}^+/SU(3)\right)
 \cong\left(S^{10},S^8_{\mathrm{lin}}\right).
 \label{eq:S10-S8-pair}
\end{equation}
The complement is
\begin{equation}
 S^{10}\setminus S^8\cong E\times(\C\setminus\{0\})
 \simeq S^1,
 \label{eq:rank-complement}
\end{equation}
and $\arg\det A$ is the angular deformation-retraction coordinate.
\end{theorem}

\begin{proof}
The first identification in \eqref{eq:S9} follows from
\cref{thm:orbit-invariants}.  If $H\succ0$, then $z\ne0$; if $H$ is
singular, then $z=0$ and \cref{lem:K3-ball} gives
$\|\chi(H)\|=1$.  Thus \eqref{eq:Omega-S9} lands on the unit sphere and is
continuous even at rank loss because the coefficient of $z/|z|$ tends to
zero.  Given $(x,w)\in S(E\oplus\C)$, put $H=\chi^{-1}(x)$.  If
$\|x\|<1$, then $w\ne0$ and the unique inverse determinant is
\begin{equation}
 z=\sqrt{\det H}\,w/|w|.
 \label{eq:Omega-inverse}
\end{equation}
If $\|x\|=1$, then $w=0$, $H$ is singular, and $z=0$.  This proves the
$S^9$ homeomorphism with its inverse.

For $(H,z)\ne0$, the normalized pair
$(H/s^2,z/s^3)$ belongs to $\mathcal Y_{3,1}$ because
$|z/s^3|^2=\det(H/s^2)$.  Conversely, for $v\ne0$ in $E\oplus\C$, set
\[
 s=\|v\|,\qquad (h,\zeta)=\Omega^{-1}(v/\|v\|),\qquad
 H=s^2h,\qquad z=s^3\zeta.
\]
This is the inverse of \eqref{eq:weighted-radial-map}.  The norm identity
$\|\Phi_{SU}(H,z)\|=s$ proves continuity of both maps at zero and properness,
so the map extends to compactifications.

Rank loss is exactly $z=0$.  Under $\Omega$, the normalized rank-loss set
maps to $S(E\oplus0)$, and radialization maps the entire affine rank-loss
set to $E\oplus0$.  Its one-point compactification is the standard linear
$S^8\subset S^{10}$.  Removing it gives
$E\times(\C\setminus0)$; projection of the second factor to its phase is a
deformation retraction.  By construction this phase is
$z/|z|=\det A/|\det A|$.
\end{proof}

Put
\[
 \Sigma_3=\{H\in\Herm_3(\C):H\succeq0,\ \det H=0\}.
\]

\begin{proposition}[The $U(3)$ ball and why its boundary is the $S^8$]
\label[proposition]{prop:U3-ball}
The map
\begin{equation}
 E\times\R_{\ge0}\longrightarrow\mathcal P_3,\qquad
 (K,t)\longmapsto K+(t-\lambda_{\min}(K))I_3
 \label{eq:halfspace-map}
\end{equation}
is a homeomorphism with inverse
\[
 H\longmapsto
 \left(H-\frac{\tr H}{3}I_3,\lambda_{\min}(H)\right).
\]
Consequently
\begin{equation}
 \mathscr V^+/U(3)\cong\mathcal P_3^+\cong B^9,\qquad
 \partial B^9=\Sigma_3^+\cong S^8.
 \label{eq:U3-ball-boundary}
\end{equation}
By contrast, the quotient of the source unit sphere is
\begin{equation}
 S_F^{17}/U(3)\cong\mathcal K_3\cong B^8,
 \label{eq:U3-unit-ball}
\end{equation}
whose boundary is the trace-one rank-loss $S^7$.  Thus ``take the source
sphere quotient'' and ``take the boundary of the compactified full quotient''
are different operations.
\end{proposition}

\begin{proof}
The least eigenvalue of the image of \eqref{eq:halfspace-map} is $t$.
Conversely, subtracting one third of the trace from $H$ produces a traceless
matrix $K$ whose least eigenvalue is
$\lambda_{\min}(H)-\tr(H)/3$; substitution recovers $H$.  Eigenvalue
continuity proves that the formulas are inverse homeomorphisms.  The
one-point compactification of a closed half-space is a closed ball, and its
boundary is the one-point compactification of the boundary hyperplane.
Now use \cref{thm:orbit-invariants}.  Equation \eqref{eq:U3-unit-ball}
follows from $\|A\|_F^2=\tr(A^*A)$ and \cref{lem:K3-ball}.
\end{proof}

\begin{corollary}[Residual determinant circle and its fibres]
\label[corollary]{cor:residual-circle}
The residual group $U(3)/SU(3)\cong S^1$ acts on the pair in
\eqref{eq:S10-S8-pair} by
\[
 e^{i\theta}\cdot(x,w)=(x,e^{i\theta}w).
\]
Its fixed set is precisely the rank-loss $S^8$, and the orbit map has
\begin{equation}
 S^{10}/S^1\cong B^9,\qquad
 q^{-1}(q(p))\cong
 \begin{cases}S^1,&p\notin S^8,\\\{p\},&p\in S^8.\end{cases}
 \label{eq:circle-fibres}
\end{equation}
On the normalized section, $S^9/S^1\cong B^8$ and the fixed set is the
rank-deficient $S^7=\partial B^8$.
\end{corollary}

\begin{proof}
Left multiplication by $U\in U(3)$ fixes $H=A^*A$ and sends
$z=\det A$ to $(\det U)z$.  Every phase is a determinant.  Equations
\eqref{eq:Omega-S9} and \eqref{eq:weighted-radial-map} therefore turn the
residual action into scalar rotation of the complex coordinate.  Its orbit
space on $E\oplus\C$ is $E\times\R_{\ge0}$, whose one-point compactification
is a closed $9$-ball; the zero complex coordinate is its boundary.  The
stabilizer is trivial for $w\ne0$ and all of $S^1$ for $w=0$, giving the
fibres.  On the unit sphere, replacing $w$ by $|w|$ gives the closed unit
ball in $E$ and the last statement.
\end{proof}

The $S^8$ in \eqref{eq:S10-S8-pair} can also be defined without first
forming an orbit quotient.

\begin{proposition}[The exact residual-circle degeneration and join]
\label[proposition]{prop:S9-join}
Under \eqref{eq:S9}, forgetting determinant phase is exactly
\begin{equation}
 q_1:S(E\oplus\C)=S^9\longrightarrow B(E)=B^8,\qquad
 q_1(x,w)=x.
 \label{eq:S9-to-B8}
\end{equation}
Its fibre is $S^1$ for $\|x\|<1$ and a singleton for
$\|x\|=1$.  More precisely,
\begin{equation}
 S^9=S(E\oplus\C)\cong S(E)*S(\C)\cong S^7*S^1,
 \label{eq:S9-join}
\end{equation}
and $q_1$ collapses the $S^1$ join factor at the $S^7$ endpoint.  On the
full compactification the analogous map is
\begin{equation}
 q_+:S^{10}\longrightarrow B^9,
 \label{eq:S10-to-B9}
\end{equation}
with circle fibre off the rank-loss $S^8$ and point fibre on that $S^8$.
The integral class of the circle on either complement is generated by
$\arg\det A$.
\end{proposition}

\begin{proof}
For fixed $H\succ0$ of trace one, the equation
$|z|^2=\det H>0$ leaves exactly its phase, an $S^1$; for singular $H$ it
forces $z=0$.  Formula \eqref{eq:Omega-S9} converts $H$ into $x=\chi(H)$
and the determinant phase into the phase of $w$, so the orbit map is
literally \eqref{eq:S9-to-B8}.  Every unit pair has the unique form
\[
 (x,w)=(\cos\alpha\,u,\sin\alpha\,v),\qquad
 u\in S(E),\ v\in S(\C),\ 0\le\alpha\le\pi/2,
\]
away from the endpoints, with the usual endpoint identifications.  This is
the defining quotient of the join and proves \eqref{eq:S9-join}.  The full
statement is \cref{cor:residual-circle,prop:U3-ball}.  The generator claim
follows from the loop
\[
 A_\theta=3^{-1/2}\diag(e^{i\theta},1,1),\qquad
 0\le\theta\le2\pi,
\]
whose Gram matrix is $I_3/3$ and whose determinant phase winds once.
\end{proof}

\begin{proposition}[Spectral-floor chart and the marked fibres]
\label[proposition]{prop:spectral-floor-observation}
The maps
\begin{equation}
 \mathcal F:E\longrightarrow\Sigma_3,\qquad
 \mathcal F(K)=K-\lambda_{\min}(K)I_3,\qquad
 \mathcal F^{-1}(H)=H-\frac{\tr H}{3}I_3
 \label{eq:spectral-floor}
\end{equation}
are inverse homeomorphisms.  Hence
$\Sigma_3^+\cong S^8$, and it is exactly the $S^8$ of
\eqref{eq:S10-S8-pair}.

The fixed marked compression
\begin{equation}
 \mathscr O:\Sigma_3\longrightarrow\Herm_2(\C)_{\succeq0},\qquad
 \mathscr O(H)=PHP^*
 \label{eq:marked-compression}
\end{equation}
is onto.  For a fixed $G\succeq0$, put
\[
 X_P^{-1}HX_P^{-*}=\begin{pmatrix}G&a\\a^*&\lambda\end{pmatrix}.
\]
Its complete fibre in $\Sigma_3$ is
\begin{equation}
 \mathscr O^{-1}(G)=
 \left\{
 X_P\begin{pmatrix}G&a\\a^*&a^*G^\dagger a+\tau\end{pmatrix}X_P^*:
 a\in\operatorname{Ran}G,\ \tau\ge0,\
 \bigl(\rank G<2\ \text{or}\ \tau=0\bigr)
 \right\}.
 \label{eq:marked-complete-fibre}
\end{equation}
Consequently the fibres are homeomorphic to
\begin{equation}
 \begin{cases}
 \C^2,&G\succ0,\\
 \C\times\R_{\ge0},&\rank G=1,\\
 \R_{\ge0},&G=0.
 \end{cases}
 \label{eq:marked-fibre-types}
\end{equation}
The positive section is $G\mapsto RGR^*$.  The affine map
\eqref{eq:marked-compression} is not proper, so these formulas do not assert
an unproved extension through the one-point compactification point.
\end{proposition}

\begin{proof}
For $K\in E$, the matrix $\mathcal F(K)$ is positive semidefinite and has
least eigenvalue zero.  If $H\in\Sigma_3$, then
$K=H-(\tr H/3)I_3$ is traceless and has least eigenvalue $-\tr H/3$;
therefore the two formulas in \eqref{eq:spectral-floor} are inverse.
Continuity of the least eigenvalue proves the homeomorphism.  Its one-point
compactification is $E^+\cong S^8$.  The proof of
\cref{thm:sphere-ladder} identifies the same rank-loss affine space with
$E\oplus0$, so the two $S^8$ descriptions agree.

Equations \eqref{eq:marked-P}--\eqref{eq:marked-inverse} show that the
upper-left block of $X_P^{-1}HX_P^{-*}$ is $PHP^*=G$.  Congruence by the
invertible $X_P$ reduces positivity to positivity of the displayed block
matrix.  The exact generalized Schur-complement criterion is
\[
 G\succeq0,\qquad a\in\operatorname{Ran}G,\qquad
 \tau:=\lambda-a^*G^\dagger a\ge0.
\]
Indeed, multiplying on both sides by the invertible triangular matrices
with off-diagonal entry $-G^\dagger a$ reduces it to
$\diag(G,\tau)$; if $a$ had a component on $\ker G$, testing the quadratic
form on a large signed multiple of a kernel vector would violate
positivity.  The same congruence gives
\[
 \rank H=\rank G+\one_{\tau>0}.
\]
Thus $H$ is singular precisely when $\rank G<2$ or $\tau=0$, proving
\eqref{eq:marked-complete-fibre}.  The three fibre types follow by using
$\dim_\C\operatorname{Ran}G=2,1,0$.  Setting $a=0$ and $\tau=0$ gives
$H=RGR^*$ and proves surjectivity and the section.  Finally, for fixed
$G\succ0$ the parameter $a\in\C^2$ is unbounded, so the inverse image of
the compact singleton $\{G\}$ is not compact; $\mathscr O$ is not proper.
\end{proof}

\subsection[A global marked observation on the S9 quotient]
{A global marked observation on the $S^9$ quotient}

Although compression on the affine rank-loss $S^8$ is nonproper, the
trace-one $S^9$ has a global compact observation with exactly computable
fibres.  Define
\begin{equation}
 \mathscr O_9:\mathcal Y_{3,1}\longrightarrow
 \mathcal B_P:=\{G\succeq0:\tr(C^{-1}G)\le1\},\qquad
 \mathscr O_9(H,z)=PHP^*.
 \label{eq:O9}
\end{equation}
The set $\mathcal B_P$ is a compact convex body with nonempty interior in
$\Herm_2(\C)$, hence is homeomorphic to $B^4$ by radial projection from an
interior point.

\begin{theorem}[Every fibre of the $S^9$ marked observation]
\label[theorem]{thm:O9-fibres}
The map \eqref{eq:O9} is a continuous surjection.  For
$G\in\mathcal B_P$ set
\begin{equation}
 L(G)=\frac{1-\tr(C^{-1}G)}{3}.
 \label{eq:L-G}
\end{equation}
Its fibre is exactly the set of pairs
\begin{equation}
 a\in\operatorname{Ran}G,\qquad
 a^*G^\dagger a\le L(G),\qquad
 z\in\C,\qquad
 |z|^2=\det G\bigl(L(G)-a^*G^\dagger a\bigr),
 \label{eq:O9-complete-fibre}
\end{equation}
where $(H,z)$ is recovered uniquely from
\begin{equation}
 H=X_P\begin{pmatrix}G&a\\a^*&L(G)\end{pmatrix}X_P^*.
 \label{eq:O9-H-recovery}
\end{equation}
The fibre is $S^5$ when $G\succ0$ and $L(G)>0$; it is a point when
$G\succ0$ and $L(G)=0$; it is a closed two-disc when $\rank G=1$ and
$L(G)>0$, degenerating to a point when $L(G)=0$; and it is a point at
$G=0$.  These are orbit-space fibres, not torus fibres.
\end{theorem}

\begin{proof}
Since $\tr H=\tr(X_P^*X_PK)$ for
$K=X_P^{-1}HX_P^{-*}$, \eqref{eq:marked-inverse} gives
\[
 1=\tr H=\tr(C^{-1}G)+3\lambda.
\]
Thus $\lambda=L(G)$.  The Schur-complement argument in
\cref{prop:spectral-floor-observation} gives exactly the first two
conditions in \eqref{eq:O9-complete-fibre}.  Since $\det X_P=1$,
\[
 \det H=\det G\bigl(L(G)-a^*G^\dagger a\bigr)
\]
when $G$ is invertible, and the same identity follows polynomially, or from
rank, when $G$ is singular.  The defining equation
$|z|^2=\det H$ of $\mathcal Y_{3,1}$ gives the last condition.  Conversely
these conditions construct a unique point by \eqref{eq:O9-H-recovery};
taking $a=0$ proves surjectivity for every $G\in\mathcal B_P$.

If $G\succ0$, put $b=G^{-1/2}a$ and
$w=z/\sqrt{\det G}$.  Equation \eqref{eq:O9-complete-fibre} becomes
$\|b\|^2+|w|^2=L(G)$ in $\C^2\oplus\C$.  It is $S^5$ for $L(G)>0$ and a
point for $L(G)=0$.  If $\rank G=1$, then $z=0$ and the allowed $a$ form
the closed Hermitian norm disc in the one-dimensional complex space
$\operatorname{Ran}G$.  If $G=0$, then $a=z=0$ and
$H=uu^*/3$ is the unique point.
\end{proof}

\subsection[The exact period-to-S9 and period-to-S8 maps]
{The exact period-to-$S^9$ and period-to-$S^8$ maps}

For $t$ in the punctured cusp and
$G_6\in\operatorname{Sym}^+_6(\R)$, retain both arguments and put
\begin{equation}
 G_4=\Phi G_6\Phi^{\mathsf T},\qquad
 H_3(t;G_6)=\Omega_3(t)G_6\Omega_3(t)^*,\qquad
 H_2(t;G_6)=\Omega_2(t)G_4\Omega_2(t)^*.
 \label{eq:period-Grams}
\end{equation}

\begin{proposition}[Commuting marked Gram square]
\label[proposition]{prop:Gram-square}
The matrices $G_4,H_3(t;G_6)$, and $H_2(t;G_6)$ are positive definite, and
\begin{equation}
 PH_3(t;G_6)P^*=H_2(t;G_6).
 \label{eq:Gram-observation}
\end{equation}
Equivalently, for every fixed $t$ the square
\begin{equation}
\begin{CD}
\operatorname{Sym}^+_6(\R)@>{G\mapsto\Omega_3(t)G\Omega_3(t)^*}>>
\Herm_3(\C)_{\succ0}\\
@V{G\mapsto\Phi G\Phi^{\mathsf T}}VV
@VV{H\mapsto PHP^*}V\\
\operatorname{Sym}^+_4(\R)@>{G\mapsto\Omega_2(t)G\Omega_2(t)^*}>>
\Herm_2(\C)_{\succ0}
\end{CD}
\label{eq:Gram-square}
\end{equation}
commutes.  Normalize without suppressing either argument:
\[
 h_3(t;G_6)=\frac{H_3(t;G_6)}{\tr H_3(t;G_6)},\qquad
 z_3(t;G_6)=+\sqrt{\det h_3(t;G_6)}.
\]
Then the explicit point
\begin{equation}
 \jmath(t;G_6)=\bigl(h_3(t;G_6),z_3(t;G_6)\bigr)
 \in\mathcal Y_{3,1}\cong S^9,
 \qquad
 \mathscr O_9\bigl(\jmath(t;G_6)\bigr)
 =\frac{H_2(t;G_6)}{\tr H_3(t;G_6)}.
 \label{eq:period-to-S9}
\end{equation}
The positive square root is a specified phase-zero section; it does not
erase the existence of the other determinant phases in $S^9$.

The spectral-floor map gives a second exact, generally noncompact-chart map
\begin{equation}
 \begin{aligned}
 \sigma(t;G_6)
 &=H_3(t;G_6)-\lambda_{\min}\bigl(H_3(t;G_6)\bigr)I_3
   \in\Sigma_3\hookrightarrow\Sigma_3^+\cong S^8,\\
 \mathscr O\bigl(\sigma(t;G_6)\bigr)
 &=H_2(t;G_6)-\lambda_{\min}\bigl(H_3(t;G_6)\bigr)C.
 \end{aligned}
 \label{eq:period-to-S8}
\end{equation}
\end{proposition}

\begin{proof}
Surjectivity of $\Phi$ makes $\Phi^{\mathsf T}$ injective; hence
$G_4=\Phi G_6\Phi^{\mathsf T}$ is positive definite.  The identity blocks
in $\Omega_r(t)$ make $\Omega_r(t)^*$ injective, so both $H_r(t;G_6)$ are
positive definite.  Using $P\Omega_3(t)=\Omega_2(t)\Phi$ from
\eqref{eq:period-intertwining},
\[
 PH_3(t;G_6)P^*=P\Omega_3(t)G_6\Omega_3(t)^*P^*
 =\Omega_2(t)\Phi G_6\Phi^{\mathsf T}\Omega_2(t)^*=H_2(t;G_6).
\]
This proves the square and \eqref{eq:period-to-S9}.  Subtracting the least
eigenvalue makes a positive matrix singular without altering its
eigenspaces; applying $P$ and $PI_3P^*=C$ proves
\eqref{eq:period-to-S8}.
\end{proof}

There is one useful cusp endpoint for a fixed, explicitly chosen metric.
Take $G_6=I_6$ in the real period frame, choose the standard logarithmic
strip $0\le x<1$ for $s=x+iy$, and abbreviate
$H_{3,I}(t)=H_3(t;I_6)$.  The exact formula is
\begin{equation}
 \begin{split}
 H_{3,I}(t)
 &=I_3+Z_3(t)Z_3(t)^*\\
 &=(1+|s|^2)I_3+sC_3(t)^*+\overline{s}C_3(t)
   +C_3(t)C_3(t)^*.
 \end{split}
 \label{eq:identity-period-metric}
\end{equation}
Put
$h_{3,I}(t)=H_{3,I}(t)/\tr H_{3,I}(t)$,
$z_{3,I}(t)=+\sqrt{\det h_{3,I}(t)}$, and
$\jmath_I(t)=\jmath(t;I_6)$.

\begin{corollary}[A proved cusp limit in the $S^9$ carrier]
\label[corollary]{cor:S9-cusp-limit}
For \eqref{eq:identity-period-metric}, boundedness of $C_3(t)$ gives
\begin{equation}
 h_{3,I}(t)\longrightarrow I_3/3,\qquad
 z_{3,I}(t)\longrightarrow1/(3\sqrt3),\qquad
 \mathscr O_9\bigl(\jmath_I(t)\bigr)\longrightarrow C/3
 \quad(t\to0).
 \label{eq:S9-limit}
\end{equation}
Thus the period metric has an exact continuous marked limit in the interior
of the $S^9$ orbit carrier.  This is a map of period--metric data, not a
compactification of every point of the torus family and not a fourth
singular filling.
\end{corollary}

\begin{proof}
Equation \eqref{eq:identity-period-metric} and boundedness of $C_3$ give
$y^{-2}H_{3,I}\to I_3$ as $y\to\infty$: indeed $0\le x<1$ gives
$|s|^2/y^2\to1$, while the two terms linear in $s$ are $O(y)$ and the
remaining correction is $O(1)$.  Taking traces and determinants yields
$h_{3,I}\to I_3/3$ and the positive square-root limit for $z_{3,I}$.
Compressing by $P$ and using \eqref{eq:period-to-S9} gives the last limit.
\end{proof}

\begin{theorem}[Which Gram carrier is strongest]
\label[theorem]{thm:strongest-gram}
Within the specified left-action invariant theory, the information order is
strict:
\[
 \mathcal Y_3\xrightarrow{\text{forget }\arg z}\mathcal P_3,
 \qquad
 \mathcal Y_{3,1}\xrightarrow{\text{forget }\arg z}\mathcal K_3.
\]
The maps have circle fibres on full-rank matrices and point fibres at rank
loss.  Hence:
\begin{enumerate}[label=\textup{(\roman*)}]
\item $S^9=S_F^{17}/SU(3)$ is the strongest \emph{normalized single-sphere}
carrier: it retains the complete Gram matrix, determinant phase, its
integral winding, the exact circle degeneration \eqref{eq:S9-to-B8}, and
the global marked observation \eqref{eq:O9}.
\item $(S^{10},S^8)$ is the strongest \emph{unnormalized compact pair}: it
additionally retains the radial scale, and its $S^8$ is exactly the fixed
rank-loss locus.
\item The standalone $S^8=\partial(\mathscr V^+/U(3))$ is the strongest
rank-loss sphere, but it has forgotten determinant phase and does not carry
the global compact extension of the affine compression
\eqref{eq:marked-compression}.
\end{enumerate}
None of this ordering supplies a complex tangent structure, a $T^6$ regular
fibre, the base $\PP^1$, the three special values, or their integral
specialization maps.
\end{theorem}

\begin{proof}
The orbit-separation theorem shows that $(H,z)$ is the complete $SU(3)$
invariant and $H$ the complete $U(3)$ invariant.  For $H\succ0$, changing
the phase of $z$ changes the $SU(3)$ orbit but not the $U(3)$ orbit; at rank
loss $z=0$, so no phase remains.  This proves strict information loss and
the fibre claim.  Propositions \ref{prop:U3-ball} and \ref{prop:S9-join},
Theorem \ref{thm:O9-fibres}, and Equation \eqref{eq:S10-S8-pair} prove
(i)--(iii).  The absent structures cannot be inferred from the quotient
map: its base has dimensions eight or nine, its generic fibre is $S^1$, its
singular values form the full boundary $S^7$ or $S^8$, and its complement
has fundamental group $\Z$.  In contrast, since
$f|_{X^\circ}:X^\circ\to B^\circ$ is a $T^4$-bundle and
$\pi_2(B^\circ)=0$, its homotopy exact sequence gives
\[
 1\longrightarrow\Z^4\longrightarrow\pi_1(X^\circ)
 \longrightarrow\pi_1(B^\circ)=F_2\longrightarrow1,
\]
which surjects onto the nonabelian free group $F_2$.  Thus even the regular
complements cannot be homotopy equivalent.  The tangent obstructions in
\cref{thm:s8-obstruction,thm:all-higher-spheres} finish the claimed scope.
\end{proof}

\section{The 24-dimensional octonionic flag}
\label{sec:octonionic-flag}

We now construct the exceptional Jordan object instead of selecting it by
dimension.  Let
\[
 \mathfrak h_3(\Oo)=
 \left\{
 \begin{pmatrix}
 x_1&p&q\\ \overline p&x_2&r\\ \overline q&\overline r&x_3
 \end{pmatrix}:
 x_1,x_2,x_3\in\R,\ p,q,r\in\Oo
 \right\},
 \qquad
 a\circ b=\frac12(ab+ba).
\]
Define
\[
 \Oo\PP^2=\{a\in\mathfrak h_3(\Oo):a^2=a,\ \tr a=1\}
\]
and, without replacing the incidence equation by a mnemonic,
\begin{equation}
 \operatorname{Fl}(\Oo)=
 \{(a,b)\in\Oo\PP^2\times\Oo\PP^2:
                 \operatorname{Re}\tr(ab)=0\}.
 \label{eq:octonionic-flag-definition}
\end{equation}
Let
\[
 F_4=\operatorname{Aut}_{\R}\bigl(\mathfrak h_3(\Oo),\circ\bigr),\qquad
 d_1=\diag(1,0,0),\qquad d_2=\diag(0,1,0),\qquad
 x_{\Oo,0}=(d_1,d_2).
\]
Here \(F_4\) denotes the compact real form.  These definitions and all
claims in the next theorem are recorded in the primary calculation of
 Mare and Willems \cite[Definition~2.2, p.~488;
 Proposition~2.4, p.~489, and Proposition~2.7, p.~491, and the
 root-space discussion on pp.~491--492]{MareWillems2013}.

\begin{theorem}[The exact triality tangent module]
\label[theorem]{thm:octonionic-triality}
The diagonal \(F_4\)-action on \(\operatorname{Fl}(\Oo)\) is transitive,
the stabilizer of \(x_{\Oo,0}\) is \(\operatorname{Spin}(8)\), and hence
\[
 \operatorname{Fl}(\Oo)\cong F_4/\operatorname{Spin}(8).
\]
Fix pairwise distinct real numbers \(\lambda_1,\lambda_2,\lambda_3\)
with \(\lambda_1+\lambda_2+\lambda_3=0\), and set
\(d_3=I_3-d_1-d_2\) and \(d_\lambda=\diag(\lambda_1,\lambda_2,\lambda_3)\).
The exact comparison with the diagonal-orbit model is the equivariant
diffeomorphism
\begin{equation}
 \begin{aligned}
 \Xi_\lambda:\operatorname{Fl}(\Oo)&\longrightarrow
                 F_4d_\lambda\subset\mathfrak h_3^0(\Oo),\\
 (a,b)&\longmapsto
       \lambda_1a+\lambda_2b+\lambda_3(I_3-a-b).
 \end{aligned}
 \label{eq:flag-diagonal-orbit-map}
\end{equation}
For \(z\in F_4d_\lambda\), its inverse is \((p_1(z),p_2(z))\), where
\begin{equation}
 p_i(z)=\frac{(z-\lambda_j I_3)\circ(z-\lambda_k I_3)}
                  {(\lambda_i-\lambda_j)(\lambda_i-\lambda_k)},
 \qquad\{i,j,k\}=\{1,2,3\}.
 \label{eq:flag-spectral-idempotents}
\end{equation}
Define the three coordinate subspaces of \(\mathfrak h_3^0(\Oo)\) by
\begin{align}
 \mathfrak h_{\gamma_1}
 &=\left\{
 \begin{pmatrix}0&0&0\\0&0&r\\0&\overline r&0\end{pmatrix}:r\in\Oo
 \right\},\nonumber\\
 \mathfrak h_{\gamma_2}
 &=\left\{
 \begin{pmatrix}0&0&q\\0&0&0\\\overline q&0&0\end{pmatrix}:q\in\Oo
 \right\},\label{eq:octonionic-root-spaces}\\
 \mathfrak h_{\gamma_3}
 &=\left\{
 \begin{pmatrix}0&p&0\\\overline p&0&0\\0&0&0\end{pmatrix}:p\in\Oo
 \right\}.\nonumber
\end{align}
As real \(\operatorname{Spin}(8)\)-modules,
\begin{equation}
 T_{x_{\Oo,0}}\operatorname{Fl}(\Oo)
 \xrightarrow[\cong]{\ d\Xi_\lambda\ }
 T_{d_\lambda}(F_4d_\lambda)
 =\mathfrak h_{\gamma_1}\oplus
  \mathfrak h_{\gamma_2}\oplus
  \mathfrak h_{\gamma_3}
 \cong V_8\oplus S_8^+\oplus S_8^-.
 \label{eq:triality-tangent}
\end{equation}
In particular,
\[
 \dim_{\R}\operatorname{Fl}(\Oo)=8+8+8=24.
\]
\end{theorem}

\begin{proof}
Mare and Willems prove that the stabilizer of \((d_1,d_2)\) consists
exactly of the automorphisms fixing the diagonal Jordan subalgebra
pointwise, which is \(\operatorname{Spin}(8)\), and that the action is
transitive.  Their restricted-root decomposition of
\(\mathfrak h_3^0(\Oo)\) has the three spaces in
\eqref{eq:octonionic-root-spaces}; the tangent space to the orbit is their
direct sum.  To verify the model comparison, write
\((a,b)=(h d_1,h d_2)\) with \(h\in F_4\).  Since \(hI_3=I_3\),
\(I_3-a-b=hd_3\), so \(\Xi_\lambda(a,b)=hd_\lambda\).
The three orthogonal diagonal idempotents satisfy the polynomial
identities \(p_i(d_\lambda)=d_i\).  Jordan automorphisms preserve these
identities, whence \(p_i(hd_\lambda)=hd_i\).  This proves both inverse
compositions in \eqref{eq:flag-diagonal-orbit-map}--
\eqref{eq:flag-spectral-idempotents}.  All formulas are polynomial on
the ambient real vector spaces with nonzero constant denominators, so the
maps are smooth.  Their differentials are inverse, which gives the
displayed tangent identification at the correctly typed basepoints.
Under the stabilizer, the three summands are respectively
the vector representation and the two real half-spin representations
\cite[pp.~491--492, especially the discussion around
Equation~(2.10)]{MareWillems2013}.  Each is the displayed copy of the
eight-dimensional real vector space \(\Oo\), which proves the dimension
calculation without suppressing any summand.
\end{proof}

For \(k=1,2,3\), let
\[
 \mathcal E_k=F_4\times_{\operatorname{Spin}(8)}
                         \mathfrak h_{\gamma_k}.
\]
These are rank-eight real subbundles, and
\begin{equation}
 T\operatorname{Fl}(\Oo)=
 \mathcal E_1\oplus\mathcal E_2\oplus\mathcal E_3.
 \label{eq:octonionic-tangent-bundles}
\end{equation}
They are associated bundles; the notation does not assert that any
\(\mathcal E_k\) is trivial.

\begin{theorem}[Integral ring, cells, and Euler class]
\label[theorem]{thm:octonionic-cohomology}
There are orientations and a basis
\(\beta_1,\beta_2\in H^8(\operatorname{Fl}(\Oo);\Z)\) for which
\begin{equation}
 \begin{aligned}
 e_1:=e(\mathcal E_1)&=2\beta_1-\beta_2,\\
 e_2:=e(\mathcal E_2)&=-\beta_1+2\beta_2,\\
 e_3:=e(\mathcal E_3)&=e_1+e_2=\beta_1+\beta_2,
 \end{aligned}
 \label{eq:octonionic-Euler-generators}
\end{equation}
and
\begin{equation}
 H^*(\operatorname{Fl}(\Oo);\Z)
 \cong
 \frac{\Z[\beta_1,\beta_2]}
 {\bigl(\beta_1^2-\beta_1\beta_2+\beta_2^2,\,
         \beta_1\beta_2(\beta_1-\beta_2)\bigr)},
 \qquad |\beta_1|=|\beta_2|=8.
 \label{eq:octonionic-cohomology-ring}
\end{equation}
Writing
\[
 u_{\Oo}=\beta_1^2\beta_2=\beta_1\beta_2^2
 \in H^{24}(\operatorname{Fl}(\Oo);\Z)
\]
with the induced orientation, \(u_{\Oo}\) is a primitive top generator and
\begin{equation}
 e\bigl(T\operatorname{Fl}(\Oo)\bigr)
 =e_1e_2e_3=6u_{\Oo},\qquad
 \left\langle e(T\operatorname{Fl}(\Oo)),
          [\operatorname{Fl}(\Oo)]\right\rangle=6.
 \label{eq:octonionic-Euler-six}
\end{equation}
The integral cohomology is free and has Poincar\'e polynomial
\begin{equation}
 1+2t^8+2t^{16}+t^{24}.
 \label{eq:octonionic-Poincare}
\end{equation}

More precisely, on
\(\mathfrak d^0=\{\diag(x_1,x_2,x_3):x_1+x_2+x_3=0\}\), put
\[
 \gamma_1(x)=x_3-x_2,\qquad
 \gamma_2(x)=x_1-x_3,\qquad
 \gamma_3(x)=x_1-x_2=\gamma_1(x)+\gamma_2(x).
\]
In the symmetric-pair normalization the actual restricted roots are
\(\{\pm\gamma_1/2,\pm\gamma_2/2,\pm\gamma_3/2\}\), of type \(A_2\).
The doubled labels \(\gamma_i\) name the same root spaces and reflection
hyperplanes; this factor of two is the one stated in the source's
root-normalization footnote \cite[footnote~3, p.~491]{MareWillems2013}.
There is a cell decomposition
\begin{equation}
 \operatorname{Fl}(\Oo)=\bigsqcup_{\pi\in S_3}C_\pi,\qquad
 C_\pi\cong\R^{8\ell(\pi)}.
 \label{eq:octonionic-six-cells}
\end{equation}
For the stabilizer action,
\begin{equation}
 \operatorname{Fl}(\Oo)^{\operatorname{Spin}(8)}
 =\{\pi x_{\Oo,0}:\pi\in S_3\}.
 \label{eq:octonionic-six-fixed-points}
\end{equation}
Thus both the cell calculation and the Euler-class calculation give
\(\chi(\operatorname{Fl}(\Oo))=6\).
\end{theorem}

\begin{proof}
Mare and Willems' integral theorem gives a free ring generated by
\(\beta_1,\beta_2\), with relations
\[
 S_i(\beta_1,\beta_2-\beta_1,-\beta_2)=0,\qquad i=2,3,
\]
and gives \eqref{eq:octonionic-Euler-generators}
\cite[Theorem~1.1, p.~484, and Equation~(3.2),
p.~496]{MareWillems2013}.  Expanding without changing the generators,
\[
 S_2=-\beta_1^2+\beta_1\beta_2-\beta_2^2,\qquad
 S_3=\beta_1\beta_2(\beta_1-\beta_2),
\]
which is exactly \eqref{eq:octonionic-cohomology-ring}.  The second
relation gives
\(\beta_1^2\beta_2=\beta_1\beta_2^2\).  Multiplying the first relation by
\(\beta_1\) and by \(\beta_2\) then gives
\(\beta_1^3=\beta_2^3=0\).  Consequently
\[
 (2\beta_1-\beta_2)(-\beta_1+2\beta_2)(\beta_1+\beta_2)
 =3\beta_1^2\beta_2+3\beta_1\beta_2^2=6u_{\Oo}.
\]
The primary calculation states that one sixth of this Euler product is a
basis of the top integral cohomology
\cite[Equation~(3.3) and the following paragraph,
p.~497]{MareWillems2013}; hence \(u_{\Oo}\) is primitive and
\eqref{eq:octonionic-Euler-six} follows.

The same source constructs the \(S_3\)-indexed cells, with one
eight-dimensional real root space for each inversion of \(\pi\), and
identifies the fixed-point set
\cite[Equations~(2.11)--(2.12) and Proposition~2.10,
pp.~493--494]{MareWillems2013}.  Their real dimensions are
\[
 0,\quad8,8,\quad16,16,\quad24.
\]
All cells are even-dimensional.  The cellular chain groups therefore have
zero differentials, yielding \eqref{eq:octonionic-Poincare} and a second
calculation of the Euler characteristic as the number of cells, namely
six.
Equivalently, the cellwise complex coordinates have dimensions
\(0,4,4,8,8,12\), since \(C_\pi\cong\C^{4\ell(\pi)}\) as real
cells.  This cell notation supplies no global complex tangent structure.
\end{proof}

\begin{proposition}[The exact relation between the two Weyl data]
\label[proposition]{prop:two-Coxeter-systems}
The ordinary Coxeter number of the ambient root system \(F_4\) is
\[
 h(F_4)=12.
\]
The restricted root system that governs
\eqref{eq:octonionic-six-cells}, however, is \(A_2\); its Weyl group is
\(S_3\) and its Coxeter number is \(3\).  For a common maximal torus of
\(\operatorname{Spin}(8)\) and \(F_4\), there is a split exact sequence
\[
 1\longrightarrow W(D_4)\longrightarrow W(F_4)
   \longrightarrow S_3\longrightarrow1,
 \qquad [W(F_4):W(D_4)]=1152/192=6.
\]
The quotient permutes the three triality representations.  It gives an
exact Weyl-group relation behind the number six; the scalar
\(h(F_4)=12\) alone does not determine the cells or Euler characteristic.
\end{proposition}

\begin{proof}
For a finite root system the invariant degrees satisfy \(d_i=e_i+1\),
and the largest degree is the Coxeter number.  The authoritative \(F_4\)
table gives
\[
 (d_1,d_2,d_3,d_4)=(2,6,8,12),
 \qquad(e_1,e_2,e_3,e_4)=(1,5,7,11),
\]
so \(h(F_4)=12\)
\cite[Section~10.1, p.~213, and the table on p.~216]
{KomoriMatsumotoTsumura2023}.  By contrast, the roots
\(\gamma_1,\gamma_2,\gamma_1+\gamma_2\) in the proof of
\cref{thm:octonionic-cohomology} form \(A_2\).  Its two simple reflections
act as the adjacent transpositions of three coordinates, and their product
is a three-cycle; hence its Coxeter element has order \(3\).
Mare and Willems explicitly identify this restricted Weyl group with
\(S_3\) \cite[pp.~492--493]{MareWillems2013}.

For completeness the extension can be checked in root coordinates.
Take an orthonormal basis \(\varepsilon_1,\ldots,\varepsilon_4\).
The long roots of \(F_4\) are
\(\{\pm\varepsilon_i\pm\varepsilon_j:i<j\}\), the \(D_4\) roots.
Partition the short roots into
\[
 \mathcal V=\{\pm\varepsilon_i\},\qquad
 \mathcal S_\pm=
 \left\{\tfrac12\sum_{i=1}^4\delta_i\varepsilon_i:
           \delta_i\in\{1,-1\},\ \prod_i\delta_i=\pm1\right\}.
\]
These are exactly the three connected components of the graph joining
orthogonal short roots: distinct components have inner products
\(\pm1/2\), and within each component every non-antipodal distinct pair
is orthogonal.  Thus \(W(F_4)\) permutes the three sets.  An element in
the kernel preserves \(\mathcal V\), so is a signed permutation matrix;
preserving both half-spin sets forces an even number of sign changes.
These are exactly \(W(D_4)\), of order \(2^3\cdot4!=192\).
The reflections in \(\varepsilon_1\) and
\(h=(\varepsilon_1+\varepsilon_2+\varepsilon_3+\varepsilon_4)/2\)
induce two distinct transpositions: the first exchanges
\(\mathcal S_+,\mathcal S_-\), and the second exchanges
\(\mathcal V,\mathcal S_-\).  Their product has order three, since
\(\langle\varepsilon_1,h\rangle=1/2\).  They generate a copy of
\(S_3\) mapping isomorphically onto the permutation quotient.  This proves
surjectivity and splitting, and gives \(|W(F_4)|=6\cdot192=1152\).
The three displayed short-root sets are the weights of the three triality
modules, so this is the asserted permutation of those modules.
\end{proof}

\subsection{The exact six-chamber comparison}

Let \(\mathscr K_3^{\max}/\Z^3\) be the translation classes of the maximal
Kuhn tetrahedra \(\sigma_{m,\pi}\) in
\eqref{eq:kuhn-tetrahedron}.  Translation by \(-m\) gives the exact
identification
\[
 \mathscr K_3^{\max}/\Z^3
 =\{[\sigma_{0,\pi}]:\pi\in S_3\}.
\]
The zero-dimensional torus strata of \(W_3\) are indexed by this set; call
the stratum corresponding to \([\sigma_{0,\pi}]\) by \(z_\pi\).

\begin{theorem}[The exact marked map of finite indexing sets]
\label[theorem]{thm:Kuhn-octonionic-comparison}
After fixing the coordinate order \((e_1,e_2,e_3)\) and the base flag
\(x_{\Oo,0}\), the formula
\begin{equation}
 \Theta:W_3^{(0)}\longrightarrow
 \operatorname{Fl}(\Oo)^{\operatorname{Spin}(8)},\qquad
 \Theta(z_\pi)=\pi x_{\Oo,0}
 \label{eq:Kuhn-flag-bijection}
\end{equation}
is an \(S_3\)-equivariant bijection.  Here
\(W_3^{(0)}=\{z_\pi:\pi\in S_3\}\) is a finite set of torus zero-strata,
and \(S_3\) acts on the source by
permuting coordinate directions and on the target as the restricted Weyl
group.  It identifies the six summands contributing to the two Euler
characteristics, but it does not extend to an asserted identification of
the spaces.
\end{theorem}

\begin{proof}
Every translation class of a maximal Kuhn tetrahedron has one and only one
representative \(\sigma_{0,\pi}\), so the source of \(\Theta\) has exactly
the six displayed elements.  Equation
\eqref{eq:octonionic-six-fixed-points} gives exactly the six displayed
target elements.  Thus \(\Theta\) is bijective.  For
\(\tau\in S_3\), permutation of coordinate directions gives
\[
 \tau\cdot[\sigma_{0,\pi}]=[\sigma_{0,\tau\pi}],
\]
while the restricted Weyl action gives
\(\tau\cdot(\pi x_{\Oo,0})=(\tau\pi)x_{\Oo,0}\).  Hence \(\Theta\) is
equivariant.

For \(W_3\), every positive-dimensional algebraic-torus stratum has Euler
characteristic zero, so its six zero-strata produce \(\chi(W_3)=6\).
For the flag, the six cells in
\eqref{eq:octonionic-six-cells} are all even-dimensional, so each
contributes \(+1\).  The bijection matches the index sets; the two
stratifications and their positive-dimensional pieces are not identified
by \(\Theta\).
More formally, let \(I_{\mathrm{cell}}=\{\operatorname{label}(C_\pi):
\pi\in S_3\}\) be a set of six labels and put
\(\tau\operatorname{label}(C_\pi)=\operatorname{label}(C_{\tau\pi})\).
Then \(z_\pi\mapsto\pi x_{\Oo,0}\mapsto
\operatorname{label}(C_\pi)\) is a chain of equivariant bijections of
finite sets.  The last action is defined on labels only: cells of different
dimensions occur, so no action carrying each cell homeomorphically to
the cell with its translated label is asserted.
\end{proof}

\begin{proposition}[A genuine embedded lower flag, and the degree boundary]
\label[proposition]{prop:complex-flag-embedding}
Choose \(\iota\in\Oo\) with \(\iota^2=-1\) and put
\(\C_\iota=\Span_{\R}\{1,\iota\}\).  Entrywise inclusion in
\(\mathfrak h_3(\Oo)\) gives a topological embedding
\begin{equation}
 \jmath_\iota:\operatorname{Fl}_3(\C_\iota)
 \hookrightarrow\operatorname{Fl}(\Oo)
 \label{eq:complex-flag-embedding}
\end{equation}
which sends the six coordinate flags to the six points in
\eqref{eq:octonionic-six-fixed-points}.

There is also an ungraded-ring isomorphism
\begin{equation}
 \rho:
 H^*(\operatorname{Fl}_3(\C);\Z)
 \xrightarrow{\ \cong\ }
 H^*(\operatorname{Fl}(\Oo);\Z),\qquad
 x_i\longmapsto\beta_i,\quad i=1,2,
 \label{eq:degree-regraded-ring-map}
\end{equation}
where the source has the presentation in
\eqref{eq:octonionic-cohomology-ring} with \(|x_i|=2\), while
\(|\beta_i|=8\).  No continuous map in either direction can induce this
degree regrading.
\end{proposition}

\begin{proof}
The subalgebra \(\C_\iota\) is associative and isomorphic to \(\C\).
Rank-one Hermitian projectors in
\(\mathfrak h_3(\C_\iota)\) remain idempotent trace-one elements after
entrywise inclusion, and orthogonality is the unchanged equation
\(\operatorname{Re}\tr(ab)=0\).  The resulting map is injective.
Its domain is compact and its target Hausdorff, so it is a topological
embedding.  The diagonal projectors are unchanged, which proves the
fixed-point assertion.

The standard complete-flag presentation is
\[
 H^*(\operatorname{Fl}_3(\C);\Z)
 \cong
 \frac{\Z[x_1,x_2]}
 {(x_1^2-x_1x_2+x_2^2,\,
   x_1x_2(x_1-x_2))},\qquad |x_i|=2;
\]
Mare and Willems derive the octonionic presentation precisely by comparison
with this \(A_2\) flag ring
\cite[Theorem~1.1, p.~484, and the proof on
pp.~496--497]{MareWillems2013}.  Thus \(\rho\) is an isomorphism after
forgetting degrees.  If \(f:\operatorname{Fl}(\Oo)\to
\operatorname{Fl}_3(\C)\) were to induce \(\rho\), then
\(f^*(x_i)\in H^2(\operatorname{Fl}(\Oo);\Z)=0\).  If a map in the other
direction were to induce \(\rho^{-1}\), then the image of each
\(\beta_i\) would lie in
\(H^8(\operatorname{Fl}_3(\C);\Z)=0\).  Either conclusion contradicts the
nonzero generators in \eqref{eq:degree-regraded-ring-map}.
\end{proof}

\subsection{The octonionic determinant-boundary sphere}

For \(\lambda\in\R\) and \(q_1,q_2,q_3\in\Oo\), retain the ordered
coordinates
\begin{equation}
 Q_{\Oo}(\lambda;q_1,q_2,q_3)=
 \begin{pmatrix}
 \lambda&q_3&\overline q_2\\
 \overline q_3&\lambda&q_1\\
 q_2&\overline q_1&\lambda
 \end{pmatrix}
 \in\mathfrak h_3(\Oo).
 \label{eq:octonionic-equal-diagonal}
\end{equation}
The equal-diagonal slice and its positive cone are
\begin{equation}
 \mathcal T(\Oo)=
 \{Q_{\Oo}(\lambda;q_1,q_2,q_3)\}
 \cong\R\oplus\Oo^3,\qquad
 \mathcal C(\Oo)=\mathcal T(\Oo)\cap
                 \mathfrak h_3(\Oo)_{\succeq0}.
 \label{eq:octonionic-equal-diagonal-cone}
\end{equation}
Thus \(\dim_{\R}\mathcal T(\Oo)=1+3\cdot8=25\).  The Jordan determinant
on this slice is the exact cubic
\begin{equation}
 \det_{\Oo}Q_{\Oo}
 =\lambda^3-\lambda\bigl(
 |q_1|^2+|q_2|^2+|q_3|^2\bigr)
 +2\operatorname{Re}\bigl((q_1q_2)q_3\bigr).
 \label{eq:octonionic-equal-diagonal-determinant}
\end{equation}
The real part is unchanged if \(q_1(q_2q_3)\) is used, because the
octonionic associator has zero real part.

\begin{theorem}[The exact octonionic \(S^{24}\) determinant carrier]
\label[theorem]{thm:octonionic-determinant-S24}
The relative boundary of \(\mathcal C(\Oo)\) inside
\(\mathcal T(\Oo)\) is
\[
 \partial\mathcal C(\Oo)
 =\{Q\in\mathcal C(\Oo):\det_{\Oo}Q=0\}.
\]
There are homeomorphisms of pairs
\begin{equation}
 \bigl(\mathcal C(\Oo)^+,
       (\partial\mathcal C(\Oo))^+\bigr)
 \cong(B^{25},S^{24}).
 \label{eq:octonionic-determinant-pair}
\end{equation}
In particular
\[
 S_{\det,\Oo}^{24}:=(\partial\mathcal C(\Oo))^+\cong S^{24}
\]
is an exact determinant-boundary sphere.  It is not the
\(24\)-dimensional flag \(F_4/\operatorname{Spin}(8)\).
\end{theorem}

\begin{proof}
The nonnegative cone of the Euclidean Jordan algebra
\(\mathfrak h_3(\Oo)\) is closed and convex.  Its intersection
\[
 K_{\Oo}=\{(q_1,q_2,q_3)\in\Oo^3:
              Q_{\Oo}(1;q_1,q_2,q_3)\succeq0\}
\]
with the affine hyperplane \(\lambda=1\) is therefore closed and convex.
Every positive semidefinite \(2\times2\) principal corner gives
\(|q_i|\le1\), so \(K_{\Oo}\) is bounded and hence compact.
The origin of \(\Oo^3\) corresponds to \(I_3\), which is positive
definite, so it is an interior point of \(K_{\Oo}\).

For \(v\in S(\Oo^3)\), let
\[
 R_{\Oo}(v)=\max\{r\ge0:rv\in K_{\Oo}\}.
\]
Convexity, compactness, and the interior-point property make
\(R_{\Oo}:S(\Oo^3)\to\R_{>0}\) continuous.  Indeed, choose
\(0<\epsilon<M\) with
\(\epsilon B^{24}\subset K_{\Oo}\subset M B^{24}\).  The Minkowski
functional
\[
 p_K(x)=\inf\{t>0:x\in tK_{\Oo}\}
\]
is positive for \(x\ne0\), positively homogeneous, and subadditive:
if \(x\in aK_{\Oo}\) and \(y\in bK_{\Oo}\), convexity gives
\(x+y\in(a+b)K_{\Oo}\).  Taking infima proves subadditivity.
As \(p_K(w)\le\|w\|/\epsilon\), subadditivity in both directions gives
\(|p_K(x)-p_K(y)|\le\|x-y\|/\epsilon\).  Hence
\(R_{\Oo}(v)=1/p_K(v)\) is continuous.  Radial rescaling
\[
 0\longmapsto0,\qquad rv\longmapsto
 \frac{r}{R_{\Oo}(v)}v
\]
is consequently a homeomorphism \(K_{\Oo}\cong B^{24}\) carrying
\(\partial K_{\Oo}\) onto \(S^{23}\).  Denote it by \(F_K\); its
inverse sends \(sv\), \(0<s\le1\), to \(sR_{\Oo}(v)v\), and sends
zero to zero.  Both maps are continuous at zero because the radial
function is bounded above and below by positive constants.

A nonzero element of \(\mathcal C(\Oo)\) has \(\lambda>0\): if
\(\lambda=0\), the same principal-corner inequalities force
\(q_1=q_2=q_3=0\).  Positive scaling therefore reconstructs the whole
cone as
\[
 \mathcal C(\Oo)=\{0\}\cup
 \{\lambda Q_{\Oo}(1;q_1,q_2,q_3):
       \lambda>0,\ (q_1,q_2,q_3)\in K_{\Oo}\}.
\]
Here are explicit maps through the Lorentz cone, with
\(q=(q_1,q_2,q_3)\):
\[
 \begin{aligned}
 \mathcal C(\Oo)&\longrightarrow
       L_{25}:=\{(s,x)\in\R\oplus\Oo^3:s\ge\|x\|\},\\
 Q_{\Oo}(\lambda;q)&\longmapsto
       (\lambda,\lambda F_K(q/\lambda))\quad(\lambda>0),\\
 0&\longmapsto(0,0),\\
 L_{25}&\longrightarrow[0,\infty)\times\Oo^3,\\
 (s,x)&\longmapsto(s-\|x\|,x).
 \end{aligned}
\]
The inverses are \((s,x)\mapsto Q_{\Oo}(s;sF_K^{-1}(x/s))\) for
\(s>0\), with zero sent to zero, and
\((t,x)\mapsto(t+\|x\|,x)\), respectively.  The boundedness of
\(F_K\) and \(F_K^{-1}\) proves continuity at the cone vertices.
The boundary maps to \(s=\|x\|\), then to \(t=0\), proving an explicit
pair homeomorphism with a closed half-space and its boundary.

To identify this boundary with determinant zero in the relative topology,
use the spectral decomposition in the Euclidean Jordan algebra.
A positive definite element has three positive eigenvalues and is an
interior point even in the full algebra, hence in the slice.  Conversely,
if a positive semidefinite \(Q\) has a zero eigenvalue, then
\(Q-\delta I_3\in\mathcal T(\Oo)\), for every \(\delta>0\), has a
negative eigenvalue.  Such points lie outside \(\mathcal C(\Oo)\) and
converge to \(Q\); thus \(Q\) is not a relative interior point.
Since the cone is closed, these are exactly its relative boundary points.
The determinant is the product of the three nonnegative eigenvalues, so
its zero set in the cone is exactly that boundary.

Fixing an orthonormal identification \(\Oo^3\cong\R^{24}\), the map
\[
 (t,x)\longmapsto
 \left(\frac{2x}{D},\frac{t^2+\|x\|^2-1}{D}\right)\in B^{25},
 \qquad D=1+t^2+\|x\|^2,
\]
extends at infinity to \((0,1)\in\partial B^{25}\).
It is inverse stereographic projection to the closed hemisphere in
\(S^{25}\) with first coordinate nonnegative, followed by deletion of
that first coordinate.  The hemisphere projection is bijective, with
inverse adjoining \(\sqrt{1-\|y\|^2}\); hence the displayed map is a
homeomorphism after one-point compactification.  At \(t=0\) it is
\[
 x\longmapsto
 \left(\frac{2x}{1+\|x\|^2},
       \frac{\|x\|^2-1}{1+\|x\|^2}\right)\in S^{24},
\]
with infinity sent to \((0,1)\).  This proves
\eqref{eq:octonionic-determinant-pair}, including its specified boundary.
\end{proof}

% Staged inclusion. No canonical file was edited.
\subsection{The ordered Albert triality tensor and its global descent}
\label{att:sec:ordered-tensor}

\subsubsection{Original coordinates and octonion identities}

Let $\mathbb O$ be the real octonion division algebra, with its original
conjugation, norm $N(a)=a\bar a=\bar a a$, and real inner product
$h(a,b)=\operatorname{Re}(a\bar b)$. Put $J=\mathfrak h_3(\mathbb O)$,
with its Albert Jordan product $X\circ Y=(XY+YX)/2$. Matrix multiplication
here involves only products of two octonions in each entry; no product
of three octonionic matrices is left unparenthesized. Set
$e_i=\operatorname{diag}(\delta_{1i},\delta_{2i},\delta_{3i})$,
$I=e_1+e_2+e_3$, and
\[
 E_1(a)=\begin{pmatrix}0&0&0\\0&0&a\\0&\bar a&0\end{pmatrix},\quad
 E_2(b)=\begin{pmatrix}0&0&\bar b\\0&0&0\\b&0&0\end{pmatrix},\quad
 E_3(c)=\begin{pmatrix}0&c&0\\\bar c&0&0\\0&0&0\end{pmatrix}.
\]
The specified equal-diagonal object remains exactly
\begin{equation}
 Q_{\mathbb O}(\lambda;q_1,q_2,q_3)
 =\lambda I+E_1(q_1)+E_2(q_2)+E_3(q_3)
 =\begin{pmatrix}\lambda&q_3&\bar q_2\\
 \bar q_3&\lambda&q_1\\q_2&\bar q_1&\lambda\end{pmatrix}.
 \label{att:eq:Q}
\end{equation}
Let $V_i$ be three labelled copies of $\mathbb O$ and $J_i=E_i(V_i)$.
Labelling does not identify their group actions. The original trace
metric on $J$ is
\begin{equation}
\begin{split}
 G(X,Y)&=\operatorname{tr}(X\circ Y),\\
 G\left(\sum_i\xi_i e_i+\sum_iE_i(a_i),
        \sum_i\eta_i e_i+\sum_iE_i(b_i)\right)
 &=\sum_i\xi_i\eta_i+2\sum_i h(a_i,b_i).
\end{split}
\label{att:eq:metrics}
\end{equation}
In particular $E_i^*G=2h$, not $h$.

We use alternativity, $N(ab)=N(a)N(b)$, and conjugation reversing
products. The parenthesis changes used below follow as follows.
Because $\bar a=2\operatorname{Re}(a)-a$, alternativity gives
\begin{equation}
 \bar a(ab)=N(a)b,\quad a(\bar a b)=N(a)b,\quad
 (ba)\bar a=N(a)b,\quad (b\bar a)a=N(a)b.
 \label{att:eq:inverse-property}
\end{equation}
Polarizing norm composition gives
$h(ab,ac)=N(a)h(b,c)$ and $h(ba,ca)=N(a)h(b,c)$.
For $a\ne0$, these equations and \eqref{att:eq:inverse-property}
imply $L_a^*=L_{\bar a}$ and $R_a^*=R_{\bar a}$; for $a=0$
the identities are immediate. Consequently
\begin{align}
 h(ab,c)&=h(b,\bar a c)=h(a,c\bar b),\nonumber\\
 \operatorname{Re}((ab)c)&=\operatorname{Re}(a(bc))
                         =\operatorname{Re}((bc)a).
 \label{att:eq:real-associativity}
\end{align}
For the first real-part equality, both expressions equal
$h(c,\bar b\bar a)$ by the adjoint identities. For the second, use
$\operatorname{Re}(xy)=h(x,\bar y)=h(y,\bar x)=\operatorname{Re}(yx)$.
These are identities of real parts; they do not equate the octonions
$(ab)c$ and $a(bc)$.

\subsubsection{The complete Peirce product table}

A cyclic triple means exactly $(i,j,k)=(1,2,3),(2,3,1),(3,1,2)$.

\begin{proposition}\label{att:prop:peirce}
The decomposition $J=\bigoplus_i\mathbb R e_i\oplus\bigoplus_iJ_i$
has the complete product table
\begin{align}
 e_i\circ e_j&=\delta_{ij}e_i,\nonumber\\
 e_i\circ E_i(a)&=0,&
 e_j\circ E_i(a)&=e_k\circ E_i(a)=\tfrac12E_i(a),\nonumber\\
 E_i(a)\circ E_i(b)&=h(a,b)(e_j+e_k),\label{att:eq:peirce-table}\\
 E_i(a)\circ E_j(b)&=\tfrac12E_k(\overline{ab}).\nonumber
\end{align}
Reversing the Jordan arguments uses commutativity of $\circ$,
without reversing the octonion order in the cyclic formula.
\end{proposition}

\begin{proof}
The first three formulas follow entry by entry. For example the two
nonzero diagonal entries of $E_1(a)\circ E_1(b)$ are
$(a\bar b+b\bar a)/2$ and $(\bar a b+\bar b a)/2$, both $h(a,b)$.
The mixed products have precisely the following nonzero entries:
\[
\begin{array}{c|c|c}
\text{product}&\text{first entry}&\text{second entry}\\\hline
E_1(a)\circ E_2(b)&(1,2)=\bar b\bar a/2&(2,1)=ab/2\\
E_2(b)\circ E_3(c)&(2,3)=\bar c\bar b/2&(3,2)=bc/2\\
E_3(c)\circ E_1(a)&(3,1)=\bar a\bar c/2&(1,3)=ca/2
\end{array}
\]
These are respectively $E_3(\overline{ab})/2$,
$E_1(\overline{bc})/2$, and $E_2(\overline{ca})/2$, by the original
definitions of $E_1,E_2,E_3$.
\end{proof}

\subsubsection{Three typed products, adjoints, and inverses}

Define the trilinear form and all three bilinear maps by
\begin{align}
 \tau:V_1\times V_2\times V_3&\longrightarrow\mathbb R,&
 \tau(a,b,c)&=\operatorname{Re}((ab)c),\label{att:eq:tau}\\
 \mu_{12}:V_1\times V_2&\longrightarrow V_3,&
 \mu_{12}(a,b)&=\overline{ab}=\bar b\bar a,\nonumber\\
 \mu_{23}:V_2\times V_3&\longrightarrow V_1,&
 \mu_{23}(b,c)&=\overline{bc}=\bar c\bar b,\nonumber\\
 \mu_{31}:V_3\times V_1&\longrightarrow V_2,&
 \mu_{31}(c,a)&=\overline{ca}=\bar a\bar c.\nonumber
\end{align}
Every $V_i$ carries $h$, with its label retained. Equation
\eqref{att:eq:real-associativity} and conjugation give
\begin{equation}
 \tau(a,b,c)=h(\mu_{12}(a,b),c)
 =h(\mu_{23}(b,c),a)=h(\mu_{31}(c,a),b).
 \label{att:eq:riesz}
\end{equation}
For example $h(\overline{ab},c)=
\operatorname{Re}(\overline{ab}\bar c)=
\operatorname{Re}(\overline{c(ab)})=\operatorname{Re}((ab)c)$.
Thus these are exactly the three Riesz-curried maps for $h$.
In the original matrix metric,
\begin{equation}
 \tau(a,b,c)=G(E_1(a)\circ E_2(b),E_3(c))
 =\operatorname{tr}((E_1(a)\circ E_2(b))\circ E_3(c)).
 \label{att:eq:tau-jordan}
\end{equation}

\begin{proposition}\label{att:prop:norm-inverse}
For cyclic $(i,j,k)$, $a\in V_i$, and $b\in V_j$,
\begin{align}
 N(\mu_{ij}(a,b))&=N(a)N(b),\label{att:eq:composition}\\
 \mu_{ki}(\mu_{ij}(a,b),a)&=N(a)b,\label{att:eq:inverse-one}\\
 \mu_{jk}(b,\mu_{ij}(a,b))&=N(b)a.\label{att:eq:inverse-two}
\end{align}
For fixed $a\ne0$, the real-linear isomorphism
$M_a:V_j\to V_k$, $M_a(b)=\mu_{ij}(a,b)$, has
\[
 M_a^*(c)=\mu_{ki}(c,a),\quad
 M_a^*M_a=N(a)\operatorname{id}_{V_j},\quad
 M_aM_a^*=N(a)\operatorname{id}_{V_k},\quad
 M_a^{-1}(c)=\frac{\mu_{ki}(c,a)}{N(a)}.
\]
For fixed $b\ne0$, the inverse of $a\mapsto\mu_{ij}(a,b)$ is
$c\mapsto\mu_{jk}(b,c)/N(b)$, and its numerator is also the adjoint.
\end{proposition}

\begin{proof}
Norm composition proves the first formula. The inverse numerators are,
with their parentheses preserved,
\[
 \overline{\overline{ab}\,a}=\bar a(ab)=N(a)b,\qquad
 \overline{b\,\overline{ab}}=(ab)\bar b=N(b)a.
\]
Their final equalities are \eqref{att:eq:inverse-property}. Relabelling
by the three stated cyclic triples proves all six inverse identities.
The adjoint statements are \eqref{att:eq:riesz}. Both adjoint-composition
identities follow from the two inverse identities, including with the
other argument nonzero. Division by $N(a)$ or $N(b)$ gives the inverses.
No unit-length hypothesis has been imposed.
\end{proof}

Retaining $G$ gives the following different numerical identities.
For $X=E_i(a)\in J_i,Y=E_j(b)\in J_j$ with $i\ne j$, put
$B_{ij}(X,Y)=2X\circ Y=E_k(\mu_{ij}(a,b))$, using cyclic argument order.
Then
\begin{equation}
 G(B_{ij}(X,Y),B_{ij}(X,Y))=\tfrac12G(X,X)G(Y,Y),\qquad
 G(X\circ Y,X\circ Y)=\tfrac18G(X,X)G(Y,Y).
 \label{att:eq:trace-norm-factors}
\end{equation}
Indeed the left sides are $2N(a)N(b)$ and $N(a)N(b)/2$,
whereas $G(X,X)G(Y,Y)=4N(a)N(b)$. The coefficient-one composition
identity uses the explicitly stated $h$, not the trace metric $G$.

\subsubsection{The determinant and the factor six in polarization}

\begin{proposition}\label{att:prop:determinant}
For arbitrary real $\xi_i$ and the original octonionic coordinates
$q_i$, let $X=\sum_i\xi_i e_i+\sum_iE_i(q_i)$. Its Jordan determinant is
\begin{equation}
 d(X)=\xi_1\xi_2\xi_3-\sum_{i=1}^3\xi_iN(q_i)
                    +2\operatorname{Re}((q_1q_2)q_3).
 \label{att:eq:determinant-general}
\end{equation}
In particular,
\[
 d(Q_{\mathbb O}(\lambda;q_1,q_2,q_3))
 =\lambda^3-\lambda\bigl(N(q_1)+N(q_2)+N(q_3)\bigr)
                  +2\operatorname{Re}((q_1q_2)q_3).
\]
\end{proposition}

\begin{proof}
The Albert cubic norm is given by the Jordan trace formula
\[
 d(X)=\tfrac13\operatorname{tr}(X\circ(X\circ X))
 -\tfrac12\operatorname{tr}(X)\operatorname{tr}(X\circ X)
 +\tfrac16\operatorname{tr}(X)^3.
\]
The product table computes the diagonal entries and three coordinates
of $X^2$:
\[
 (X^2)_{ii}=\xi_i^2+N(q_j)+N(q_k),\qquad
 (X^2)_i=(\xi_j+\xi_k)q_i+\overline{q_jq_k},
 \quad(i,j,k)\text{ cyclic}.
\]
With $s=\xi_1+\xi_2+\xi_3$, it follows that
\begin{align*}
 \operatorname{tr}(X)&=s,\\
 \operatorname{tr}(X^2)&=\sum_i\xi_i^2+2\sum_iN(q_i),\\
 \operatorname{tr}(X\circ X^2)
 &=\sum_i\xi_i^3+3\sum_i(\xi_j+\xi_k)N(q_i)
                     +6\operatorname{Re}((q_1q_2)q_3).
\end{align*}
Each scalar product of $q_i$ with $\overline{q_jq_k}$ equals
$\tau(q_1,q_2,q_3)$ by \eqref{att:eq:riesz}; the trace metric supplies
the factor $2$. Substitution into the cubic gives coefficient
$(\xi_j+\xi_k)-s=-\xi_i$ for $N(q_i)$ and coefficient $6/3=2$ for
$\tau$. Direct expansion gives the scalar part
\[
 \tfrac13\sum_i\xi_i^3-\tfrac12s\sum_i\xi_i^2+\tfrac16s^3
 =\xi_1\xi_2\xi_3.
\]
This proves \eqref{att:eq:determinant-general}. Setting all three
$\xi_i=\lambda$ retains exactly the original equal-diagonal slice.
\end{proof}

Define the symmetric polarization $\mathcal D$ by
$\mathcal D(X,X,X)=d(X)$, explicitly
\begin{align}
 6\mathcal D(X,Y,Z)={}&d(X+Y+Z)-d(X+Y)-d(X+Z)-d(Y+Z)
 \nonumber\\&+d(X)+d(Y)+d(Z).
 \label{att:eq:polarization}
\end{align}
Expansion of homogeneous cubic monomials proves symmetry and
trilinearity: inclusion--exclusion removes every term missing an
argument, and degree three forces each remaining argument to occur
once. On the diagonal the multiplier is $27-3\cdot8+3=6$.
The same expansion proves
\[
 \left.\frac{\partial^3}{\partial r\partial s\partial t}
 d(rX+sY+tZ)\right|_{r=s=t=0}=6\mathcal D(X,Y,Z).
\]
For the three distinct sectors one obtains
\begin{equation}
\begin{split}
 \left.\partial_r\partial_s\partial_t
 d(rE_1(a)+sE_2(b)+tE_3(c))\right|_0&=2\tau(a,b,c),\\
 \mathcal D(E_1(a),E_2(b),E_3(c))&=\tfrac13\tau(a,b,c).
\end{split}
 \label{att:eq:factor-six}
\end{equation}
Indeed that determinant is $2rst\tau(a,b,c)$. Relative to the original
$G$, the three Riesz representatives for fixed $E_1(a),E_2(b)$ are
\[
\begin{array}{c|c}
\text{scalar form on }J_1\times J_2\times J_3&
             \text{representative in }J_3\\\hline
\tau&E_3(\overline{ab})/2=E_1(a)\circ E_2(b)\\
\partial_r\partial_s\partial_t d&
                  E_3(\overline{ab})=2E_1(a)\circ E_2(b)\\
\mathcal D&E_3(\overline{ab})/6=(E_1(a)\circ E_2(b))/3
\end{array}
\]
Thus the determinant coefficient $2$ and polarization factor $6$
are separate, fully specified operations.

\subsubsection{The exact stabilizer and its three actions}

Let $\kappa(a)=\bar a$ and define
\begin{equation}
\begin{split}
H_\tau=\{(g_1,g_2,g_3)\in O(V_1,h)\times O(V_2,h)\times O(V_3,h):\quad&\\
\tau(g_1a,g_2b,g_3c)=\tau(a,b,c)\text{ for every }a,b,c&\}.
\end{split}
\label{att:eq:metric-stabilizer}
\end{equation}
By \eqref{att:eq:riesz}, orthogonality and nondegeneracy of $h$,
this is equivalent to each of the equations
\begin{align}
 \mu_{12}(g_1a,g_2b)&=g_3\mu_{12}(a,b),\nonumber\\
 \mu_{23}(g_2b,g_3c)&=g_1\mu_{23}(b,c),\label{att:eq:equivariance}\\
 \mu_{31}(g_3c,g_1a)&=g_2\mu_{31}(c,a).\nonumber
\end{align}
With every conjugation retained, the first equation reads
\begin{equation}
 (g_1a)(g_2b)=\kappa g_3\kappa(ab).
 \label{att:eq:related-triple}
\end{equation}
Consequently $(g_1,g_2,g_3)\mapsto(g_1,g_2,\kappa g_3\kappa)$
is a group isomorphism onto the ordinary product-related triples
$\mathcal R=\{(A,B,C)\in O(8)^3:(Ax)(By)=C(xy)\text{ for all }x,y\}$.

\begin{theorem}\label{att:thm:stabilizer}
Every element of $H_\tau$ has $g_i\in SO(8)$, and $H_\tau$ is
isomorphic to the compact simply connected $\operatorname{Spin}(8)$.
The map
\[
\begin{split}
 \Phi:H_\tau&\longrightarrow
 \{g\in\operatorname{Aut}(J,\circ):g(e_i)=e_i\ (i=1,2,3)\},\\
 \Phi(g)\left(\sum_i\xi_i e_i+\sum_iE_i(a_i)\right)
 &=\sum_i\xi_i e_i+\sum_iE_i(g_i a_i)
\end{split}
\]
is an exact group isomorphism. Thus \eqref{att:eq:equivariance} is
$\operatorname{Spin}(8)$ equivariance for the actions on the original
Albert coordinates.
\end{theorem}

\begin{proof}
For unit $u,v\in\mathbb O$, reflection in $u^\perp$ is
\[
 D_u(x)=x-2h(x,u)u=-u\bar x u.
\]
To verify the second expression, multiply
$x\bar u+u\bar x=2h(x,u)$ on the right by $u$ and use
$(x\bar u)u=x$. Define
\[
 A_{u,v}(x)=v(\bar u x),\quad
 B_{u,v}(y)=(y\bar u)v,\quad
 C_{u,v}(z)=v(\bar u z\bar u)v=D_vD_u(z).
\]
Repeated outside letters are parenthesized by flexibility. The
Moufang identity $(vx)(yv)=v(xy)v$, first for $v$ and then for $\bar u$,
proves
\[
 (v(\bar u x))((y\bar u)v)=v(\bar u(xy)\bar u)v.
\]
This is a product-related triple. All its coordinates lie in $SO(8)$:
multiplication by a unit octonion is orthogonal, and its determinant
is continuously $+1$ on the connected unit sphere because it is $+1$
at $1$. Every element of $SO(8)$ is a product of an even number of
reflections. To obtain such a decomposition, reflect the image of the
first vector of an orthonormal basis onto that vector, then repeat
inside its orthogonal complement. After finitely many steps the map is
the identity; the determinant records the parity of the reflections.
Projection of $\mathcal R\cap SO(8)^3$ to $C$ is therefore surjective.

Its kernel contains exactly two elements. If $(Ax)(By)=xy$, put
$p=A1$, $q=B1$. Then $pq=1$, $By=\bar p y$, and $Ax=xp$.
The equation becomes $(xp)(\bar p y)=xy$. Replacing $y$ by $py$
gives $(xp)y=x(py)$ for every $x,y$. Hence $p$ lies in the middle
nucleus of $\mathbb O$, which is $\mathbb R$. The latter fact can be
checked without an implicit associativity assumption. In the
Cayley--Dickson presentation $\mathbb O=\mathbb H\oplus\mathbb H\ell$
with
\[
 (a+b\ell)(c+d\ell)=(ac-\bar d b)+(da+b\bar c)\ell,
\]
write $p=a+b\ell$ and use quaternion units
$\mathsf i\mathsf j=\mathsf k$. The $\ell$ coefficient of
$(\mathsf i p)\mathsf j-\mathsf i(p\mathsf j)$ is $-2b\mathsf k$,
so $b=0$. The $\ell$ coefficient of
$(\ell a)x-\ell(ax)$ is $\bar a\bar x-\bar x\bar a$.
Its vanishing for $x=\mathsf i,\mathsf j$ forces $a$ real.
Since $|p|=1$, the kernel is
$\{(I,I,I),(-I,-I,I)\}$.

There are no orientation-reversing related triples. First, the cyclic
invariance of $\tau$ gives surjectivity onto $SO(8)$ in each of its
three coordinates, and hence the same holds for the three projections
of $\mathcal R\cap SO(8)^3$. If a product-related triple had
$\det A=-1$, compose it with an orientation-preserving product-related
triple whose first coordinate is $\kappa A^{-1}$. The result has first
coordinate $\kappa$, so satisfies $\bar x(By)=C(xy)$.
Setting $x=1$ yields $B=C$. Setting $y=1$, and writing $p=B1$, yields
$Bx=\bar x p$. Hence
\[
 \bar x(\bar y p)=\overline{xy}\,p\qquad\text{for all }x,y.
\]
Set $y=p$ and multiply on the right by $\bar p$. The inverse property
gives $\bar x\bar p=\bar p\bar x$ for all $x$. The center of
$\mathbb O$ is $\mathbb R$: in the displayed Cayley--Dickson
presentation, commuting $a+b\ell$ with $\mathsf i$ forces $b=0$,
and commuting $a$ with $\mathsf i,\mathsf j$ forces $a$ real.
Thus $p=\pm1$. The preceding equation would imply
$\bar x\bar y=\overline{xy}=\bar y\bar x$ for all $x,y$,
contradicting $\mathsf i\mathsf j=-\mathsf j\mathsf i$.
Thus $\det A=1$. Cyclic permutation of $H_\tau$ excludes negative
determinants in the other two coordinates too.

The group is connected. The reflection-pair triples are continuously
parametrized by $S^7\times S^7$, contain the identity at $u=v=1$,
and contain $(-I,-I,I)$ at $u=1,v=-1$. Their third coordinates
generate $SO(8)$, so every related triple is a product of these
triples and a kernel element. Every such element therefore lies in
the identity component. Projection to $C$ is a two-sheeted covering:
it is a surjective homomorphism of compact Lie groups with finite
kernel, hence a local diffeomorphism with those two-element fibres.
Here $\pi_1(SO(8))=\mathbb Z/2$. One proof starts with the unit
quaternion conjugation action on imaginary quaternions, which
identifies $SO(3)$ with $S^3/\{\pm1\}$, and uses the fibrations
$SO(n-1)\to SO(n)\to S^{n-1}$ for $n\ge4$; their bases have zero
$\pi_1$ and $\pi_2$, so the fundamental group remains $\mathbb Z/2$.
A connected two-sheeted cover of this group is simply connected.
It is $\operatorname{Spin}(8)$, the simply connected cover of $SO(8)$.

The map $\Phi$ preserves every product in
\eqref{att:eq:peirce-table}, by orthogonality and
\eqref{att:eq:equivariance}; bilinearity proves preservation of all
Jordan products. Its inverse is $\Phi(g^{-1})$. Conversely, a Jordan
automorphism fixing all $e_i$ preserves
\[
 J_i=\{X\in J:e_i\circ X=0,\quad e_j\circ X=e_k\circ X=X/2\}.
\]
Its restrictions define linear maps $g_i$ in the original $E_i$
coordinates. Applying the automorphism to the same-sector Peirce
product proves $h(g_i a,g_i b)=h(a,b)$; applying it to the mixed
product proves \eqref{att:eq:related-triple}. Thus the restrictions
belong to $H_\tau$. The direct-sum decomposition of $J$ proves
injectivity and surjectivity of $\Phi$ with no further kernel.
\end{proof}

\begin{proposition}\label{att:prop:linear-stabilizer}
Without the metrics, the stabilizer of the same typed tensor in
$GL(V_1)\times GL(V_2)\times GL(V_3)$ is exactly
\[
 \{(s_1g_1,s_2g_2,s_3g_3):s_i>0,\quad s_1s_2s_3=1,\quad
                          (g_1,g_2,g_3)\in H_\tau\}.
\]
Thus the metric-preserving stabilizer is $\operatorname{Spin}(8)$;
the unrestricted typed linear stabilizer has two additional
independent positive scale parameters.
\end{proposition}

\begin{proof}
Tensor invariance for $(A,B,C)$ and the $h$ adjoint imply
$\mu_{12}(Aa,Bb)=C^{-T}\mu_{12}(a,b)$.
For nonzero $a$, Proposition~\ref{att:prop:norm-inverse} gives
$|\det M_a|=N(a)^4$ on dimension eight. Taking determinants yields
$N(Aa)^4=|\det B\det C|^{-1}N(a)^4$. Positivity therefore gives
$A=s_1g_1$, where $s_1>0$ and $g_1$ is orthogonal. The cyclic equations
give $B=s_2g_2$, $C=s_3g_3$ likewise. For unit $a,b$, the two sides
of the adjoint equation have norms $s_1s_2$ and $s_3^{-1}$, so
$s_1s_2s_3=1$. Dividing by those explicit scalars proves
$(g_1,g_2,g_3)\in H_\tau$. The converse follows by trilinearity.
\end{proof}

The three actions can also be read directly as the vector and two
real half-spin actions. On $S=V_2\oplus V_3$ define
\begin{equation}
 c(a)(b,c_0)=(-\mu_{31}(c_0,a),\mu_{12}(a,b)),\qquad a\in V_1.
 \label{att:eq:clifford}
\end{equation}
The adjoint and inverse identities prove
$c(a)^*=-c(a)$ and $c(a)^2=-N(a)I$. Polarization gives
\[
 c(a)c(a')+c(a')c(a)=-2h(a,a')I.
\]
This is the real Clifford action with negative norm convention.
Even products preserve each of the two summands.
Equation \eqref{att:eq:equivariance} gives the exact intertwining map
\[
 \operatorname{diag}(g_2,g_3)c(a)
 =c(g_1a)\operatorname{diag}(g_2,g_3).
\]
Conversely, an even product of unit Clifford operators covers the
corresponding product of reflections on $V_1$, is orthogonal on $S$,
and satisfies this equation. Its two diagonal blocks, together with
its vector action, therefore form a triple in $H_\tau$. This map
from the spin double cover is onto: it covers $SO(V_1)$ and includes
both elements over the identity, because the Clifford scalar $-1$
acts as $-I$ on both summands. The two groups are connected double
covers and the map is an isomorphism. This identifies $V_1$ as the
vector representation and $V_2,V_3$ as the two half-spin modules in
the displayed Clifford realization, without conjugating any $q_i$.
Their names $S^+,S^-$ can be fixed by the parity summands; when the
names are instead defined by an oriented volume element, that
orientation must be specified as additional notation.

There is no nonzero $\operatorname{Spin}(8)$ equivariant linear map
$V_i\to V_j$ for $i\ne j$. The four triples
\[
 (I,I,I),\quad(I,-I,-I),\quad(-I,I,-I),\quad(-I,-I,I)
\]
belong to $H_\tau$ by the product equation. For any two different
sectors, one of these acts as $+I$ on the source and $-I$ on the
target. An equivariant linear map then equals its own negative and
vanishes. Their exact positive relation is the three nonzero
equivariant bilinear maps \eqref{att:eq:tau}; fixing any nonzero
argument gives the intersector isomorphisms and inverses already
proved in Proposition~\ref{att:prop:norm-inverse}.

\subsubsection{An explicit tangent comparison and global bundle descent}

Put $F_4=\operatorname{Aut}(J,\circ)$ and use
Theorem~\ref{att:thm:stabilizer} to identify $H_\tau$ with the ordered
frame stabilizer. Then $M=F_4/H_\tau$ parametrizes the orbit of the
ordered frame $(e_1,e_2,e_3)$.

Every Jordan automorphism preserves $\operatorname{tr}$ and $G$.
Indeed the Peirce table gives
$\operatorname{Tr}_{\operatorname{End}(J)}L_{e_i}=1+16/2=9$ and
$\operatorname{Tr}_{\operatorname{End}(J)}L_{E_i(a)}=0$.
The latter operator has zero diagonal blocks in the Peirce
decomposition, as each product in its table changes summand.
Thus $\operatorname{Tr}L_X=9\operatorname{tr}X$.
Conjugating $L_X$ by a Jordan automorphism gives $L_{gX}$, proving
$\operatorname{tr}(gX)=\operatorname{tr}X$. Product preservation now
gives $G(gX,gY)=G(X,Y)$.

Write $L_X(Y)=X\circ Y$. In any Jordan algebra $[L_X,L_Y]$ is a
derivation, with the following algebraic proof. The coefficient of
three distinct parameters in the Jordan identity gives
\begin{align*}
x\circ((z\circ w)\circ y)+z\circ((w\circ x)\circ y)
 +w\circ((x\circ z)\circ y)
={}&(z\circ w)\circ(y\circ x)+(w\circ x)\circ(y\circ z)\\
&+(x\circ z)\circ(y\circ w).
\end{align*}
Subtracting this identity with $x,y$ interchanged cancels the right
sides and gives exactly
\[
 [L_x,L_y](z\circ w)
 =([L_x,L_y]z)\circ w+z\circ([L_x,L_y]w).
\]
For cyclic $(i,j,k)$ define
\begin{equation}
 D_i(a)=4[L_{E_i(a)},L_{e_j}].
 \label{att:eq:derivations}
\end{equation}
Its complete values follow from the Peirce product table:
\begin{align}
 D_i(a)e_i&=0,& D_i(a)e_j&=E_i(a),&
 D_i(a)e_k&=-E_i(a),\nonumber\\
 D_i(a)E_i(b)&=2h(a,b)(e_k-e_j),\label{att:eq:D-table}\\
 D_i(a)E_j(b)&=-E_k(\overline{ab}),&
 D_i(a)E_k(c)&=E_j(\overline{ca}).\nonumber
\end{align}
For example on $e_j$ the commutator before multiplication by $4$ is
$E_i(a)/2-E_i(a)/4=E_i(a)/4$, and on $E_j(b)$ it is
$0-E_k(\overline{ab})/4$. The same product table supplies the other
entries. Exponentiating a derivation gives Jordan automorphisms,
since the derivation rule and the exponential power series give
$e^{tD}(X\circ Y)=(e^{tD}X)\circ(e^{tD}Y)$.

Let $\mathfrak f_4$ be the derivation algebra and $\mathfrak h$
its subalgebra fixing each $e_i$. Put
$\mathfrak m_i=\{D_i(a):a\in V_i\}$ and
$\mathfrak m=\bigoplus_i\mathfrak m_i$. Then
\begin{equation}
 \mathfrak f_4=\mathfrak h\oplus\mathfrak m,\qquad
 \mathfrak m_i\cong V_i,\qquad
 \operatorname{Ad}(h)D_i(a)=D_i(h_i a).
 \label{att:eq:reductive}
\end{equation}
To check that no tangent directions are missing, differentiate
$e_i^2=e_i$ to get $De_i=2e_i\circ De_i$. Thus $De_i$ has only the
components $J_j\oplus J_k$. Differentiating $I^2=I$ gives $DI=0$,
so the two possible occurrences of any $J_i$ component in the triple
$(De_1,De_2,De_3)$ are opposite. Equation \eqref{att:eq:D-table}
therefore supplies unique $a_i$ such that $D-\sum_iD_i(a_i)$ fixes
all three $e_i$. That equation also proves the sum direct.
Conjugation of the commutator defining $D_i$ by an automorphism
fixing $e_j$ proves the displayed equivariance.

Retain arbitrary distinct original spectral parameters
$\lambda_1,\lambda_2,\lambda_3\in\mathbb R$, with
$\lambda_1+\lambda_2+\lambda_3=0$, and set
\[
 d_\lambda=\sum_i\lambda_i e_i,\qquad
 \Delta_1=\lambda_2-\lambda_3,\quad
 \Delta_2=\lambda_3-\lambda_1,\quad
 \Delta_3=\lambda_1-\lambda_2.
\]
Every $\Delta_i$ is nonzero and
\begin{equation}
 D_i(a)d_\lambda=\Delta_i E_i(a).
 \label{att:eq:spectral-differential}
\end{equation}
The equivariant map
$\Xi_\lambda:M\to F_4d_\lambda$, $[g]\mapsto gd_\lambda$,
is a diffeomorphism. Its inverse recovers the ordered frame from the
Jordan polynomials
\[
 p_i(z)=\frac{(z-\lambda_jI)\circ(z-\lambda_kI)}
              {(\lambda_i-\lambda_j)(\lambda_i-\lambda_k)}.
\]
For $z=gd_\lambda$ these equal $ge_i$: evaluate on $d_\lambda$ first,
then apply the Jordan automorphism $g$. In particular the stabilizer
of $d_\lambda$ is exactly $H_\tau$, and both maps are smooth.
Equation \eqref{att:eq:spectral-differential} is the differential of
this comparison with all scale factors retained.

Use the associated-bundle convention
$(g,v)\sim(gh,h^{-1}v)$ for $F_4\times_{H_\tau}V_i$, and denote this
rank-eight real vector bundle by $\mathcal V_i$. Define
\begin{align}
 K_i:\mathcal V_i&\longrightarrow TM,&
 K_i[g,a]&=\left.\frac{d}{dt}\right|_0
                   g\exp(tD_i(a))H_\tau,\label{att:eq:Kmap}\\
 P_i:\mathcal V_i&\longrightarrow M\times J,&
 P_i[g,a]&=(gH_\tau,gE_i(a)).\nonumber
\end{align}
Both maps are well defined by the proved equivariance.
Equation \eqref{att:eq:reductive} shows that
$\bigoplus_iK_i:\bigoplus_i\mathcal V_i\to TM$ is an isomorphism.
The image of $P_i$ above the frame $(p_1,p_2,p_3)=(ge_1,ge_2,ge_3)$
is precisely the intrinsic Peirce space
\[
 \{X\in J:p_i\circ X=0,\quad p_j\circ X=p_k\circ X=X/2\}.
\]
Suppressing the recorded base-point component of $P_i$ when viewing
it as a tangent vector to the orbit, the exact comparison is
\begin{equation}
 d\Xi_\lambda\circ K_i=\Delta_iP_i.
 \label{att:eq:bundle-comparison}
\end{equation}
Thus the original Albert slots and frame tangent slots are related
by an explicit map, rather than an unrecorded change of coordinates.

The tensor and products descend by the formulas
\begin{align}
 \boldsymbol\tau([g,a],[g,b],[g,c])&=\tau(a,b,c),\nonumber\\
 \boldsymbol\mu_{ij}([g,a],[g,b])&=[g,\mu_{ij}(a,b)]
                                  \in\mathcal V_k.
 \label{att:eq:descent}
\end{align}
Representatives over one point can always be chosen with common $g$.
Replacing it by $gh$ replaces the arguments by
$h_i^{-1}a,h_j^{-1}b,h_k^{-1}c$.
Tensor invariance and \eqref{att:eq:equivariance} give exactly the
same scalar and the same associated class. This proves well-defined
global descent, beyond invariance at a single fibre. The associated
metrics induced by $h$ carry all norm, adjoint, and inverse formulas
of Proposition~\ref{att:prop:norm-inverse} to identities of bundle maps.

Transport by $K_i$ gives the three subbundles
$\mathcal T_i=K_i(\mathcal V_i)$, the direct sum
$TM=\mathcal T_1\oplus\mathcal T_2\oplus\mathcal T_3$, and global
bilinear maps $\mathcal T_i\otimes\mathcal T_j\to\mathcal T_k$.
Define a second, orbit-coordinate realization explicitly by
$S_i=K_i/\Delta_i$. Then $d\Xi_\lambda S_i=P_i$.
Transport of $\boldsymbol\mu$ by $S_i$ is exactly the original
coordinate map $E_i(a),E_j(b)\mapsto E_k(\overline{ab})$.
These transports obey the following exact comparison. For
$u_i\in J_i,u_j\in J_j$ at the reference orbit point, the
$K$-transported product has orbit-coordinate value
\begin{equation}
 \frac{\Delta_k}{\Delta_i\Delta_j}\,B_{ij}(u_i,u_j),
 \label{att:eq:tangent-scale-products}
\end{equation}
whereas the $S$-transported product has value $B_{ij}(u_i,u_j)$.
Indeed the $K$-coordinate arguments are
$E_i^{-1}u_i/\Delta_i$ and $E_j^{-1}u_j/\Delta_j$; bilinearity and
the output multiplier $\Delta_k$ give the displayed coefficient.
At every other point the same identity is transported by $g\in F_4$.
The $K$-transported scalar in orbit coordinates is
\[
 \frac{\tau(E_1^{-1}u_1,E_2^{-1}u_2,E_3^{-1}u_3)}
             {\Delta_1\Delta_2\Delta_3},
\]
and the $S$-transported scalar is its numerator.
For the pullback $G_\lambda=\Xi_\lambda^*G$ of the original trace
metric on the orbit, trace invariance and \eqref{att:eq:metrics} give
\[
 G_\lambda(K_i[g,a],K_i[g,b])=2\Delta_i^2h(a,b),\qquad
 G_\lambda(S_i[g,a],S_i[g,b])=2h(a,b).
\]
No $\Delta_i$ is set to $1$ and no coordinate is conjugated to absorb
a sign or a metric factor.

There is a specified connection for which these descended tensors
are parallel. On $F_4\to F_4/H_\tau$ take
\[
 \omega=\operatorname{pr}_{\mathfrak h}(g^{-1}dg)
 \quad\text{for the decomposition }\eqref{att:eq:reductive}.
\]
Because the projection commutes with $\operatorname{Ad}(H_\tau)$,
this form is right equivariant and sends each fundamental vertical
vector to its generating element of $\mathfrak h$. These are exactly
the principal-connection conditions. Upstairs, the descended tensor
and products are represented by the constant $H_\tau$-invariant
tensors $\tau,\mu_{ij}$. Their covariant derivatives are the
ordinary derivatives of those constant coefficients plus the
infinitesimal $\mathfrak h$ action. Both terms vanish.
They are therefore parallel for this explicitly defined homogeneous
connection and for the connections transported to $TM$ by $K$ or $S$.
No assertion about Levi--Civita parallelism follows from this argument.

\subsubsection{Source conventions and scope}

The matrix convention, related-triple equation, and Jordan frame
action agree entry for entry with Ichiro Yokota,
\emph{Exceptional Lie groups}, arXiv:0902.0431,
Theorem~1.14.2, Lemmas~1.14.3--1.14.4, Theorem~1.16.2, and
Theorem~2.7.1 \cite{Yokota2009}. The present proof supplies the coordinate calculations,
norm and inverse maps, determinant coefficients, and global bundle
maps in that convention. Its algebraic starting objects are the
real alternative composition algebra $\mathbb O$ and its Albert
Jordan algebra $J$; those defining structures are retained.

The proved result is an ordered three-sector tangent structure on
$F_4/\operatorname{Spin}(8)$ with a normed triality tensor and three
exact bilinear morphisms. It does not assert toroidal filling,
an integrable complex structure, or physical dynamics.
The obstruction to equivariant unary identifications and the
positive bilinear relations between sectors were both proved.
The dependence on the diagonal-orbit parameters is the exact
morphism \eqref{att:eq:bundle-comparison}, with its product and metric
consequences; it is not a missing comparison.


\subsection{A coordinate-preserving triality map onto the determinant boundary}
\label{tri:sec:spectral-graph}

Retain the ordered matrix in \eqref{eq:octonionic-equal-diagonal} and write
\[
 E(q)=Q_{\mathbb O}(0;q_1,q_2,q_3),\qquad
 s(q)=|q_1|^2+|q_2|^2+|q_3|^2,\qquad
 \tau(q)=\operatorname{Re}((q_1q_2)q_3).
\]
Here the eigenvalues are the three real Jordan eigenvalues, with their
multiplicities. They are the roots of
\(z^3-s(q)z-2\tau(q)\). Define
\[
 a(q)=-\lambda_{\min}(E(q)).
\]
Thus \(a(q)\) is exactly the largest real root of
\(\lambda^3-s(q)\lambda+2\tau(q)=0\). This prescription keeps the
complete cubic, including its mixed term and its factor two.

\begin{theorem}[The spectral graph and its exact inverse]
\label{tri:thm:spectral-graph}
The map
\begin{equation}
 \mathcal B:\mathbb O^3\longrightarrow\partial\mathcal C(\mathbb O),
 \qquad q\longmapsto Q_{\mathbb O}(a(q);q)
 \label{tri:eq:spectral-graph}
\end{equation}
is a proper homeomorphism. Its inverse is the coordinate projection
\[
 \mathcal P\left(
 \begin{pmatrix}
 \lambda&q_3&\overline q_2\\
 \overline q_3&\lambda&q_1\\
 q_2&\overline q_1&\lambda
 \end{pmatrix}\right)=(q_1,q_2,q_3).
\]
It extends to a based homeomorphism
\[
 \mathcal B^+:(\mathbb O^3)^+
       \xrightarrow{\ \cong\ }S_{\det,\mathbb O}^{24}.
\]
Let \(K=\operatorname{Spin}(8)\subset F_4\) be the pointwise stabilizer
of the ordered diagonal Jordan frame. On \(\mathbb O^3\), use the
three representations defined by
\(kE_i(q_i)=E_i(\rho_i(k)q_i)\), with \(E_i\) the three ordered
off-diagonal inclusions. Then \(\mathcal B\) and \(\mathcal P\) are
\(K\)-equivariant. The norm and the cubic interaction are recovered
on the boundary without discarding any coordinate:
\begin{equation}
 2\tau(q)=a(q)s(q)-a(q)^3.
 \label{tri:eq:boundary-triality}
\end{equation}
\end{theorem}

\begin{proof}
Use the Euclidean Jordan trace inner product
\(\langle u,v\rangle_J=\operatorname{tr}(u\circ v)\).
For \(q,p\in\mathbb O^3\), direct multiplication gives
\[
 \|E(q)-E(p)\|_J^2=2\sum_{i=1}^3|q_i-p_i|^2.
\]
Let \(\mathcal I\) be the compact set of primitive idempotents of
trace one. The spectral theorem for the Euclidean Albert algebra gives
\[
 \lambda_{\min}(u)=\min_{e\in\mathcal I}\langle u,e\rangle_J.
\]
For completeness, in a spectral Jordan frame write
\(u=\sum_i\mu_i e_i\). For every \(e\in\mathcal I\), the numbers
\(\langle e_i,e\rangle_J\) are nonnegative and have sum one: positivity
uses self-duality of the cone of squares, and the sum is
\(\langle I_3,e\rangle_J=1\). Thus
\(\langle u,e\rangle_J\geq\min_i\mu_i\), with equality at an
idempotent for a least eigenvalue. This proves the displayed formula.
Every such \(e\) has \(\|e\|_J^2=\operatorname{tr}(e^2)=1\).
Cauchy--Schwarz and the minimum formula now give
\begin{equation}
 |a(q)-a(p)|\leq\|E(q)-E(p)\|_J
       =\sqrt{2\sum_{i=1}^3|q_i-p_i|^2}.
 \label{tri:eq:spectral-Lipschitz}
\end{equation}
In particular \(a\) is continuous.

Since \(E(q)\) has trace zero, its least eigenvalue is nonpositive.
If \(q\ne0\), all eigenvalues cannot be zero, since the spectral theorem
would then give \(E(q)=0\). Their zero sum consequently forces a
strictly negative least eigenvalue. Hence \(a(0)=0\) and
\(a(q)>0\) for \(q\ne0\).
Adding \(\lambda I_3\) adds \(\lambda\) to each of the three Jordan
eigenvalues. Therefore
\[
 Q_{\mathbb O}(\lambda;q)\succeq0
       \quad\Longleftrightarrow\quad\lambda\geq a(q),
\]
and within this positive cone the determinant vanishes exactly at
\(\lambda=a(q)\). This is the relative boundary established in
\cref{thm:octonionic-determinant-S24}. It proves both inverse
compositions in \eqref{tri:eq:spectral-graph}.
Continuity of the graph and coordinate projection proves the
homeomorphism. Moreover
\[
 \|Q_{\mathbb O}(a(q);q)\|_J^2=3a(q)^2+2s(q)\geq2s(q),
\]
so unbounded \(q\) leaves every compact subset of the boundary.
This proves properness directly and the extension at infinity.

A member of \(K\) fixes \(I_3\), preserves the Jordan product and trace,
and sends each \(E_i(q_i)\) to \(E_i(\rho_i(k)q_i)\). It preserves
the characteristic polynomial and the ordered eigenvalues. Thus
\(a(\rho_1(k)q_1,\rho_2(k)q_2,\rho_3(k)q_3)=a(q)\), which proves
both equivariance identities. Substitution of \(\lambda=a(q)\) into
the unchanged determinant formula gives
\eqref{tri:eq:boundary-triality}.
\end{proof}

\paragraph{The exact tangent map from the ordered flag.}
For pairwise different real \(\lambda_1,\lambda_2,\lambda_3\) with
zero sum, retain \(\Xi_\lambda\) of
\eqref{eq:flag-diagonal-orbit-map}. At \(x_0=(d_1,d_2)\), the inverse
of its differential, written in the original off-diagonal coordinates, is
\begin{equation}
 \begin{aligned}
 A(q)&=E_2\left(\frac{q_2}{\lambda_1-\lambda_3}\right)
       +E_3\left(\frac{q_3}{\lambda_1-\lambda_2}\right),\\
 B(q)&=E_1\left(\frac{q_1}{\lambda_2-\lambda_3}\right)
       -E_3\left(\frac{q_3}{\lambda_1-\lambda_2}\right),\\
 (d\Xi_\lambda)_{x_0}(A(q),B(q))&=E(q).
 \end{aligned}
 \label{tri:eq:flag-tangent-inverse}
\end{equation}
Indeed, differentiate the three idempotents in a Jordan frame. The
idempotence equations and the identity that their sum is \(I_3\)
force \(A\) to have only sectors two and three, \(B\) to have only
sectors one and three, and their sector-three parts to be negatives.
To check that these are actual tangent vectors, the differential of the
already proved diffeomorphism \(\Xi_\lambda\) is an isomorphism onto
\(E_1(\mathbb O)\oplus E_2(\mathbb O)\oplus E_3(\mathbb O)\).
Its formula is
\((\lambda_1-\lambda_3)A+(\lambda_2-\lambda_3)B\).
Solving its three sector equations gives precisely
\eqref{tri:eq:flag-tangent-inverse}, so no extra tangent solutions or
denominators are suppressed. Composing this tangent coordinate
isomorphism with \(\mathcal B\) gives an explicit equivariant
homeomorphism from the complete tangent space at \(x_0\) onto
\(\partial\mathcal C(\mathbb O)\).


% Staged proof fragment. No canonical file has been edited.
% Parent integration prefix: fc:
\subsection{An explicit equivariant flag cell and its degree-one collapse}
\label{fc:section}

Write $J=\mathfrak h_3(\mathbb O)$, $J_0=\ker\operatorname{tr}$,
$e=I_3$, and $d_i=\operatorname{diag}(\delta_{i1},\delta_{i2},\delta_{i3})$.
The inner product on $J$ is
\[
 \langle x,y\rangle_J=\operatorname{tr}(x\circ y),\qquad
 L_x(y)=x\circ y.
\]
All operator products and brackets below belong to
$\operatorname{End}_{\mathbb R}(J)$; thus $[A,B]=AB-BA$ is a commutator
of real linear operators. An adjoint $A^*$ in these formulas is the
adjoint for $\langle\ ,\ \rangle_J$.

We preserve the ordered coordinates of the equal-diagonal slice by putting
\begin{equation}
 \begin{aligned}
 E_1(r)&=\begin{pmatrix}0&0&0\\0&0&r\\0&\bar r&0\end{pmatrix},\qquad
 E_2(s)=\begin{pmatrix}0&0&\bar s\\0&0&0\\s&0&0\end{pmatrix},\\
 E_3(t)&=\begin{pmatrix}0&t&0\\\bar t&0&0\\0&0&0\end{pmatrix},\qquad
 W=\mathbb O\oplus\mathbb O\oplus\mathbb O,\\
 E(q)&=E_1(q_1)+E_2(q_2)+E_3(q_3)
      =Q_{\mathbb O}(0;q_1,q_2,q_3).
 \end{aligned}
 \label{fc:coordinate-operators}
\end{equation}
In particular the second coordinate is the lower $(3,1)$ entry, and
conjugating it in the upper $(1,3)$ entry remains part of every formula.
Let $K=\operatorname{Aut}(J)=F_4$ and
$M=\{m\in K:m d_i=d_i\ (i=1,2,3)\}=\operatorname{Spin}(8)$.
The action of $M$ on $W$ is defined, without choosing three unrelated
orthogonal actions, by
\begin{equation}
 E_i((m\cdot q)_i)=mE_i(q_i),\qquad i=1,2,3.
 \label{fc:triality-action}
\end{equation}
The three summands are the vector and the two real half-spin modules in
this specified coordinate convention.

We use the real symmetric pair
\[
 G=E_{6(-26)},\qquad
 \mathfrak g=\operatorname{Der}(J)\oplus L_{J_0},\qquad
 \mathfrak k=\operatorname{Der}(J),\qquad
 \mathfrak s=L_{J_0}.
\]
Here $G$ is the connected reduced structure group in its real linear
representation on $J$. Its Cartan involution on the Lie algebra is
$\theta(A)=-A^*$; hence
$\theta(D+L_x)=D-L_x$ for $D\in\operatorname{Der}(J)$ and $x\in J_0$.
The symmetric-pair realization and restricted-root spaces are as in
\cite[Section~2.6, pp.~491--493]{MareWillems2013}.
The real-flag theorem we use is precisely
\cite[Appendix~B, Theorem~B.2, Proposition~B.3 and Corollary~B.4,
p.~519]{MareWillems2013}: for $P=MAN$ the space $G/P=K/M$
has cells $Nw$, and $X\mapsto\exp(X)wP$ is a diffeomorphism from
the sum of the positive root spaces made negative by $w^{-1}$ onto
that cell. Its Cartan comparison is $\Theta(X)=X-\theta X$.
We now determine that comparison, the exponential, its flag coordinates,
its inverse, and the collapse with every coefficient specified.

\begin{lemma}[The Cartan inverse with its exact coefficient]
\label{fc:cartan-inverse}
For $h=\operatorname{diag}(h_1,h_2,h_3)\in J_0$, set
\[
 \rho_1(h)=\frac{h_3-h_2}{2},\qquad
 \rho_2(h)=\frac{h_1-h_3}{2},\qquad
 \rho_3(h)=\frac{h_1-h_2}{2}=\rho_1(h)+\rho_2(h).
\]
Fix any
$h_*=\operatorname{diag}(h_{*,1},h_{*,2},h_{*,3})\in J_0$
with $h_{*,1}>h_{*,3}>h_{*,2}$. For $u\in E_i(\mathbb O)$ define
\begin{equation}
 \nu_i(u)=\frac12\left(L_u+
             \frac{[L_{h_*},L_u]}{\rho_i(h_*)}\right).
 \label{fc:nu}
\end{equation}
Then $\nu_i(u)$ is independent of the permitted choice of $h_*$,
belongs to $\mathfrak g_{\rho_i}$, and satisfies
\begin{equation}
 [L_h,\nu_i(u)]=\rho_i(h)\nu_i(u),\qquad
 \nu_i(u)-\theta\nu_i(u)=L_u,\qquad
 \nu_i(u)e=\frac12u.
 \label{fc:nu-identities}
\end{equation}
Thus $\nu_i$ is exactly the inverse of $\Theta$ on the $i$th summand;
the factor $1/2$ in \eqref{fc:nu} is necessary.
The isomorphism
\[
 Z:W\longrightarrow\mathfrak n
       =\mathfrak g_{\rho_1}\oplus\mathfrak g_{\rho_2}
                               \oplus\mathfrak g_{\rho_3},\qquad
 Z(q)=\sum_{i=1}^3\nu_i(E_i(q_i))
\]
is $M$-equivariant.
\end{lemma}

\begin{proof}
We first record the Jordan identity needed to ensure that the operators
are in the correct Lie algebra. For any $a,b\in J$,
$D_{a,b}=[L_a,L_b]$ is a derivation, and therefore
\begin{equation}
 [D_{a,b},L_c]=L_{D_{a,b}c}.
 \label{fc:inner-derivation}
\end{equation}
Here is the polarization proving this assertion. The coefficient of the
product of three distinct real parameters in
$(x^2\circ y)\circ x=x^2\circ(y\circ x)$ gives
\[
 \begin{split}
 &d\circ\bigl(b\circ(a\circ c)\bigr)
 +a\circ\bigl(b\circ(c\circ d)\bigr)
 +c\circ\bigl(b\circ(d\circ a)\bigr)\\
 &\hspace{1em}=(a\circ c)\circ(b\circ d)
 +(c\circ d)\circ(b\circ a)
 +(d\circ a)\circ(b\circ c).
 \end{split}
\]
Subtract the same identity with $a$ and $b$ exchanged. The three
terms on the right cancel, by commutativity of the Jordan product.
The result, rearranged, is
\[
 D_{a,b}(c\circ d)=(D_{a,b}c)\circ d+c\circ(D_{a,b}d),
\]
which proves the assertion and \eqref{fc:inner-derivation}.

Let $u$ have its two off-diagonal entries in positions $(j,k)$ and
$(k,j)$, and write $D_h=[L_h,L_u]$ for real diagonal $h$.
If $h'$ is another real diagonal element, direct use of the binary
Jordan product gives
\begin{equation}
 \begin{split}
 D_hh'
 &=\left(\frac{(h_j+h_k)(h'_j+h'_k)}4
                 -\frac{h_jh'_j+h_kh'_k}2\right)u\\
 &=-\frac{(h_j-h_k)(h'_j-h'_k)}4u
  =-\rho_i(h)\rho_i(h')u.
 \end{split}
 \label{fc:diagonal-calculation}
\end{equation}
For completeness, the proportionality of $D_h$ in $h$ can also be
checked on every $z\in J$ with no associativity assumption on three
octonions. Expanding the two Jordan products, the fact that every
entry of $h$ is real gives the restricted identity
\[
 4[L_h,L_u]z=(hu-uh)z-z(hu-uh).
\]
In this displayed identity only, the products are binary matrix
products; the expansion is valid because every reassociation it uses
contains a real entry from $h$. Since $u$ has only the indicated two
off-diagonal entries, $hu-uh$ depends linearly on the single real
scalar $h_j-h_k$. Consequently
\begin{equation}
 [L_h,L_u]=\frac{\rho_i(h)}{\rho_i(h_*)}[L_{h_*},L_u].
 \label{fc:proportionality}
\end{equation}
This assertion includes $\rho_i(h)=0$. We use no corresponding
identity with arbitrary octonionic $h$, and no octonionic matrix
conjugation is used to define a group action.

Put $D_*=[L_{h_*},L_u]$. Equations
\eqref{fc:inner-derivation} and
\eqref{fc:diagonal-calculation} give
\[
 [L_h,D_*]=-L_{D_*h}=\rho_i(h_*)\rho_i(h)L_u.
\]
Combining this with \eqref{fc:proportionality} proves
$[L_h,\nu_i(u)]=\rho_i(h)\nu_i(u)$ and also proves independence of
$h_*$. Since $\theta L_u=-L_u$ and $\theta D_*=D_*$, subtracting
$\theta\nu_i(u)$ from \eqref{fc:nu} gives $L_u$ exactly.
Every derivation kills the unit, so $\nu_i(u)e=u/2$.
The map $u\mapsto L_u$ is injective because $L_ue=u$.
The positive restricted-root spaces have the three stated
eight-dimensional Cartan images; thus these injective maps exhaust
them. Finally, an $m\in M$ fixes $h_*$ and satisfies
$mL_um^{-1}=L_{mu}$ and
$m[L_{h_*},L_u]m^{-1}=[L_{h_*},L_{mu}]$.
These formulas prove the asserted equivariance.
\end{proof}

\begin{lemma}[Finite exponential and the coefficient in the nilpotent law]
\label{fc:finite-exponential}
For every $q\in W$,
\begin{equation}
 Z(q)^5=0,\qquad
 \exp Z(q)=I+Z(q)+\frac{Z(q)^2}{2}
                   +\frac{Z(q)^3}{6}+\frac{Z(q)^4}{24}.
 \label{fc:finite-exp}
\end{equation}
The same formula holds with $Z(q)$ replaced by $-Z(q)^*$.
Writing $\nu_i(r)$ temporarily for $\nu_i(E_i(r))$, the brackets are
\begin{equation}
 [\nu_1(r),\nu_2(s)]=\nu_3\left(-\frac12\overline{rs}\right),
 \qquad [\nu_i(r),\nu_i(s)]=0,
 \qquad[\nu_3(t),\nu_j(r)]=0.
 \label{fc:nil-bracket}
\end{equation}
The last equality holds for $j=1,2,3$. In exponential coordinates the
group law on $N$ is
\begin{equation}
 \begin{split}
 (q_1,q_2,q_3)\star(q'_1,q'_2,q'_3)
 =\biggl(&q_1+q'_1,\ q_2+q'_2,\\
         &q_3+q'_3-\frac14\overline{q_1q'_2}
                         +\frac14\overline{q'_1q_2}\biggr).
 \end{split}
 \label{fc:nil-group-law}
\end{equation}
\end{lemma}

\begin{proof}
The auxiliary operator $H_0=L_{\operatorname{diag}(2,-2,0)}$ has
the following eigenspaces in its representation on $J$:
\[
 \begin{array}{c|ccccc}
 \text{eigenvalue}&-2&-1&0&1&2\\ \hline
 \text{eigenspace}&\mathbb R d_2&E_1(\mathbb O)&
 \mathbb R d_3\oplus E_3(\mathbb O)&E_2(\mathbb O)&\mathbb R d_1.
 \end{array}
\]
This table follows from
$L_h d_j=h_jd_j$ and
$L_hu=(h_j+h_k)u/2$ for an off-diagonal $(j,k)$ element.
The actual positive roots evaluated on $\operatorname{diag}(2,-2,0)$
are $1,1,2$. Therefore every term of $Z(q)$ strictly raises these
weights by at least one. Five successive applications give zero.
Adjoints strictly lower the weights, proving the analogous assertion
for $-Z(q)^*$. This proves both finite exponential formulas.

Brackets add restricted roots. There is no root $2\rho_i$,
$\rho_1+\rho_3$, or $\rho_2+\rho_3$, so the brackets claimed to
vanish in \eqref{fc:nil-bracket} do vanish. To compute the
remaining coefficient, put $u=E_1(r)$, $v=E_2(s)$. Binary matrix
multiplication gives, in the retained coordinate order,
\[
 u\circ v=\frac12E_3(\overline{rs})=:w.
\]
The weight of $u$, $v$, and $w$ under $L_h$ is respectively
$-h_1/2$, $-h_2/2$, and $-h_3/2$. Hence
\[
 [L_h,L_u]v=-\rho_1(h)w,\qquad
 [L_h,L_v]u=\rho_2(h)w,
\]
and \eqref{fc:nu} gives
$\nu_1(r)v=0$ and $\nu_2(s)u=w$. Applying the commutator to $e$,
\[
 [\nu_1(r),\nu_2(s)]e
       =\frac12(\nu_1(r)v-\nu_2(s)u)
       =-\frac14E_3(\overline{rs}).
\]
On the third root space the map $A\mapsto2Ae$ is inverse to $\nu_3$.
This proves the first formula of \eqref{fc:nil-bracket}.
All iterated brackets of length three vanish. The exact
Baker--Campbell--Hausdorff formula is therefore
$\log(\exp X\exp Y)=X+Y+[X,Y]/2$; substituting the just-computed
bracket gives \eqref{fc:nil-group-law}, including both signs and
the coefficient $1/4$.
\end{proof}

\begin{theorem}[A finite operator formula for the entire open cell]
\label{fc:open-cell}
For $q\in W$ define
\begin{equation}
 \begin{aligned}
 A(q)&=\left(I+Z+\frac{Z^2}{2}+\frac{Z^3}{6}
                       +\frac{Z^4}{24}\right)d_2,\\
 B(q)&=\left(I-Z^*+\frac{(Z^*)^2}{2}-\frac{(Z^*)^3}{6}
                       +\frac{(Z^*)^4}{24}\right)d_1,
       \qquad Z=Z(q),\\
 \Psi(q)&=\left(\frac{A(q)}{\operatorname{tr}A(q)},
               \frac{B(q)}{\operatorname{tr}B(q)}\right).
 \end{aligned}
 \label{fc:explicit-chart}
\end{equation}
Both traces in these formulas are strictly positive. The map $\Psi$
is an $M$-equivariant diffeomorphism
\begin{equation}
 W\xrightarrow{\ \Psi\ }
 U=\{(a,b)\in\operatorname{Fl}(\mathbb O):a_{22}>0, b_{11}>0\}
   =C_{w_0},\qquad w_0=(12).
 \label{fc:open-cell-locus}
\end{equation}
The inverse $q=\Psi^{-1}(a,b)$ is the exact formula
\begin{equation}
 \boxed{\quad
 q_1=\frac{2a_{23}}{a_{22}},\qquad
 q_2=-\frac{2b_{31}}{b_{11}},\qquad
 q_3=\frac{2a_{12}}{a_{22}}-\frac14\overline{q_1q_2}.
 \quad}
 \label{fc:explicit-inverse}
\end{equation}
Every denominator in \eqref{fc:explicit-inverse} is a positive
real scalar. In particular $\Psi(0)=(d_2,d_1)$.
\end{theorem}

\begin{proof}
Let $\widehat w_0$ be the Jordan automorphism that interchanges
the first and second real diagonal positions, obtained by conjugation
with the real permutation matrix for $(12)$. For this real matrix
the conjugation is a well-defined Jordan automorphism. Its action on
the retained positive roots is
\[
 w_0\rho_1=-\rho_2,\qquad
 w_0\rho_2=-\rho_1,\qquad
 w_0\rho_3=-\rho_3.
\]
It is consequently the longest Weyl element, and the real-flag
cell theorem gives a diffeomorphism
\begin{equation}
 W\longrightarrow N\widehat w_0P/P,\qquad
 q\longmapsto\exp Z(q)\widehat w_0P.
 \label{fc:group-chart}
\end{equation}

Here are the exact Iwasawa factors that identify the flag pair in
\eqref{fc:group-chart}. Write $g=kan$ uniquely with
$k\in K$, $a=\exp L_h\in A$, $n\in N$. Each positive root
operator kills $d_1$. This follows either by applying
\eqref{fc:nu} to $d_1$, or from the $H_0$ weight table,
whose highest weight is $\mathbb R d_1$. Each operator
$-\nu_i^*$ kills $d_2$ by the lowest-weight statement. Thus
\[
 nd_1=d_1,\qquad (n^*)^{-1}d_2=d_2,
\]
and consequently
\begin{equation}
 gd_1=e^{h_1}kd_1,\qquad
 (g^*)^{-1}d_2=e^{-h_2}kd_2.
 \label{fc:iwasawa-flag-factors}
\end{equation}
Since Jordan automorphisms preserve the trace,
the traces on the right are exactly $e^{h_1}>0$ and $e^{-h_2}>0$.
It follows that the flag corresponding to $gP$ is
\[
 \left(\frac{gd_1}{\operatorname{tr}(gd_1)},
       \frac{(g^*)^{-1}d_2}{\operatorname{tr}((g^*)^{-1}d_2)}\right)
       =(kd_1,kd_2).
\]
For $g=\exp Z(q)\widehat w_0$, its two numerators are
$\exp Z(q)d_2$ and $\exp(-Z(q)^*)d_1$, respectively.
The finite exponential identities prove
\eqref{fc:explicit-chart}, all its positivity assertions,
and the identification with \eqref{fc:group-chart}.

We next verify the image description, so that no membership criterion
is left implicit. For any $n\in N$, the $H_0$ weight table shows
that the coefficient of $d_2$ in $nd_j$ is $1$ if $j=2$ and $0$
otherwise. Indeed $n$ strictly raises weights apart from its identity
term. Similarly the coefficient of $d_1$ in $(n^*)^{-1}d_j$ is
$1$ if $j=1$ and $0$ otherwise. A flag in $N\widehat wP/P$ is
represented by the positive-trace multiples of
$nd_{w(1)}$ and $(n^*)^{-1}d_{w(2)}$. It satisfies
$a_{22}>0$ precisely when $w(1)=2$, and $b_{11}>0$ precisely
when $w(2)=1$. Both conditions hold precisely for $w=w_0=(12)$.
The six-cell decomposition covers the flag, proving
\eqref{fc:open-cell-locus}.

For the inverse we record the required coefficients of the finite
exponentials. The same weights show
\begin{equation}
 \begin{gathered}
 A(q)_{22}=1,\qquad B(q)_{11}=1,\qquad
 A(q)_{23}=\frac12q_1,\qquad
 B(q)_{31}=-\frac12q_2,\\
 A(q)_{12}=\frac12q_3+\frac18\overline{q_1q_2}.
 \label{fc:coefficient-extraction}
 \end{gathered}
\end{equation}
To see also the last coefficient directly, the relevant first-order
terms are
\[
 Zd_2=\frac12E_1(q_1)+\frac12E_3(q_3),\qquad
 -Z^*d_1=-\frac12E_2(q_2)-\frac12E_3(q_3).
\]
An $E_3$ contribution to $Z^2d_2$ can arise only from
$\nu_2(q_2)\nu_1(q_1)d_2$ or
$\nu_1(q_1)\nu_2(q_2)d_2$; the latter is zero, and the former
is $\frac14E_3(\overline{q_1q_2})$ by the calculation in
Lemma~\ref{fc:finite-exponential}. The other weight-two
contribution $\nu_1(q_1)^2d_2$ lies in $\mathbb R d_3$,
and $\nu_2(q_2)^2d_2=0$. The coefficient $1/2$ in the
exponential gives $1/8$. Terms of degree three or four raise
weight too far to have an $E_3$ component. The other four
identities in \eqref{fc:coefficient-extraction} follow from
the unique lowest, highest, and adjacent weights.

For $(a,b)=\Psi(q)$, the first two identities imply
\[
 A(q)=\frac{a}{a_{22}},\qquad
 B(q)=\frac{b}{b_{11}}.
\]
The remaining identities solve for $q_1,q_2,q_3$ in precisely
the order \eqref{fc:explicit-inverse}. Since
\eqref{fc:group-chart} is onto $U$, this proves both inverse
compositions everywhere on $U$, not only a formal local inverse.
The displayed formulas show their smoothness.

Finally, $Z(mq)=mZ(q)m^{-1}$ and $m$ is orthogonal for the
trace inner product. It fixes $d_1,d_2$ and preserves the trace.
Each expression in \eqref{fc:explicit-chart} therefore
commutes with the $M$ action. This proves equivariance of $\Psi$
and, by its proved bijectivity, of its inverse.
\end{proof}

\begin{proposition}[The orientation sign, including the conjugated coordinate]
\label{fc:orientation}
Fix an orientation of $\mathbb O$, and give $W$ the ordered product
orientation $(q_1,q_2,q_3)$. Give the flag the orientation specified
in the integral Euler-class calculation: orient its first root
summand by the chosen orientation of $E_1(\mathbb O)$, its third
by $n_2E_1(\mathbb O)$, and its second by
$n_1^{-1}n_2E_1(\mathbb O)$, where $n_1=(23)$ and $n_2=(13)$
are real-permutation lifts. The total orientation is the ordered
sum of root summands $(1,2,3)$, transported by $K$.
Then $\Psi:W\to U$ is orientation preserving.
\end{proposition}

\begin{proof}
These are the orientation transport conventions of
\cite[Section~3, pp.~495--496]{MareWillems2013}.
They are independent of the lifts within their cosets because
$M$ is connected and each of its orthogonal representations has
positive real determinant.
We compute the sign explicitly. For this computation only, call
the product orientation of the three upper entries $(r,q,p)$ in
positions $(23),(13),(12)$ the upper-entry orientation. Octonionic
conjugation fixes the real line and acts by $-1$ on the seven
imaginary directions, so its real determinant is $(-1)^7=-1$.
The lift $n_2=(13)$ takes an upper $(23)$ entry $r$ to the upper
$(12)$ entry $\bar r$. The lift $n_1^{-1}=(23)$ then takes that
upper $(12)$ entry to the upper $(13)$ entry with the same value.
The source orientations of summands $1,2,3$ thus have signs
$+1,-1,-1$ relative to the upper-entry coordinates. Their product
is $+1$, so the flag tangent orientation at $d_\lambda$ agrees
with the upper-entry orientation. Here $d_\lambda$ has the original
pairwise distinct parameters $\lambda_i$ and the trace-zero relation;
no ordering or rescaling of these parameters is needed.

The lift $\widehat w_0=(12)$ interchanges the upper-entry
$(23)$ and $(13)$ blocks and conjugates the $(12)$ block.
Interchanging two eight-dimensional blocks has determinant
$(-1)^{8\cdot8}=+1$, and conjugation has determinant $-1$.
Thus the transported flag orientation at
$\widehat w_0d_\lambda$ is the negative of its upper-entry
orientation. This transport preserves the flag orientation because
it is the action of an element of connected $K$.

In the retained $q$ coordinates, $E$ conjugates exactly the second
upper-entry block; its determinant relative to the upper-entry
orientation is therefore $-1$. Differentiate
\eqref{fc:explicit-chart} and the original orbit comparison
$\Xi_\lambda(a,b)=\lambda_1a+\lambda_2b+\lambda_3(e-a-b)$
at $q=0$. The traces have zero first derivative, and the result is
\begin{equation}
 \begin{split}
 d(\Xi_\lambda\Psi)_0(q)
 ={}&\frac{\lambda_1-\lambda_3}{2}E_1(q_1)
    +\frac{\lambda_3-\lambda_2}{2}E_2(q_2)
    +\frac{\lambda_1-\lambda_2}{2}E_3(q_3)\\
 ={}&-\sum_{i=1}^3\rho_i(\widehat w_0d_\lambda)E_i(q_i).
 \end{split}
 \label{fc:chart-differential}
\end{equation}
Its three real scalar factors are nonzero; each acts on eight
real coordinates, so their determinant contribution is positive.
The total differential has sign $-1$ relative to upper-entry
coordinates, exactly the sign of the transported flag orientation
at this point. Therefore its sign as an oriented map is $+1$.
The domain is connected and the map is a diffeomorphism, so this
sign is $+1$ everywhere.
\end{proof}

\begin{theorem}[The exact equivariant quotient onto the representation sphere]
\label{fc:collapse}
Let $F=\operatorname{Fl}(\mathbb O)$ and put
\[
 A=F\setminus U=\{(a,b)\in F:a_{22}=0\ \text{or}\ b_{11}=0\}.
\]
It is the union of the five cells other than $C_{w_0}$, and is a
closed $16$-dimensional CW subcomplex. Define
\begin{equation}
 c:F\longrightarrow W^+,
 \qquad
 c(a,b)=
 \begin{cases}
   \Psi^{-1}(a,b),&(a,b)\in U,\\
   \infty,&(a,b)\in A.
 \end{cases}
 \label{fc:collapse-map}
\end{equation}
Then $c$ is continuous, surjective, and $M$-equivariant and induces
a based $M$-equivariant homeomorphism
\begin{equation}
 F/A\xrightarrow{\ \cong\ }W^+.
 \label{fc:quotient-homeomorphism}
\end{equation}
With the orientations in Proposition~\ref{fc:orientation},
$c$ has degree $+1$. If $\omega_W\in\widetilde H^{24}(W^+;\mathbb Z)$
is its positive generator, then
\begin{equation}
 c_*[F]=[W^+],\qquad c^*\omega_W=u_{\mathbb O},
 \qquad u_{\mathbb O}=\beta_1^2\beta_2=\beta_1\beta_2^2.
 \label{fc:degree-one}
\end{equation}
On integral homology $c_*$ is the identity in degree zero, an
isomorphism in degree $24$, and the zero map on the two copies
of $\mathbb Z$ in each of degrees $8$ and $16$.
On fixed points its exact formula is
\[
 c(w_0x_{\mathbb O,0})=0,\qquad
 c(wx_{\mathbb O,0})=\infty\quad(w\ne w_0).
\]
These are the six points of $F^M$ mapping onto the two points
$(W^+)^M=\{0,\infty\}$.
\end{theorem}

\begin{proof}
All diagonal entries of primitive positive idempotents are
nonnegative. The open-cell locus of
Theorem~\ref{fc:open-cell} consequently gives the stated
closed-set formula for $A$. The restricted roots have multiplicity
eight. The six cell dimensions are $0,8,8,16,16,24$, with the
unique top cell $C_{w_0}$. The CW property of the real-flag
decomposition implies that their union below degree $24$ is the
closed $16$-skeleton, which proves the subcomplex assertion.

Continuity of $c$ on $U$ is already proved. Let $K_0\subset W$
be compact. Its image $\Psi(K_0)$ is compact in the Hausdorff
compact manifold $F$, and hence closed in $F$. Thus
\[
 c^{-1}\bigl(W^+\setminus K_0\bigr)=F\setminus\Psi(K_0)
\]
is open and contains $A$. Complements of compact subsets of $W$
form a neighborhood basis of $\infty$. This proves continuity
at every point of $A$, including all its lower-dimensional strata.
It also proves the exact escape statement: a net in $U$ approaching
$A$ eventually leaves $\Psi(K_0)$ for every compact $K_0\subset W$.
Surjectivity and the description of the fibers are immediate from
the two inverse compositions of the chart: every finite point has
one preimage, and $c^{-1}(\infty)=A$. Therefore the induced map
$F/A\to W^+$ is a continuous bijection from a compact space to a
Hausdorff space and hence a homeomorphism.
The open cell is $M$-invariant by the proved chart equivariance;
its complement is invariant too. The formula for $c$ is therefore
$M$-equivariant, with $M$ fixing $\infty$.

Since $A$ has dimension $16$, its integral homology in degrees
$23$ and $24$ is zero. The exact sequence of the pair gives an
isomorphism
\[
 H_{24}(F;\mathbb Z)\xrightarrow{\ \cong\ }H_{24}(F,A;\mathbb Z).
\]
Because $A$ is a subcomplex, quotienting identifies this relative
group with $\widetilde H_{24}(F/A;\mathbb Z)$.
At any point of $U$, the local map is the inverse of the
orientation-preserving chart. It has local degree $+1$.
There is exactly one preimage of that finite regular value.
This fixes the sign of the preceding homology isomorphism and
proves $c_*[F]=[W^+]$. Pairing with cohomology proves
$c^*\omega_W=u_{\mathbb O}$ for the already fixed positive
top generator of the flag. The remaining homology assertions
follow because $W^+$ is a $24$-sphere and $F$ is connected with
the stated six-cell homology.
Each of the three nontrivial irreducible eight-dimensional
$M$-modules has zero invariant subspace, so $W^M=\{0\}$.
The point $w_0x_{\mathbb O,0}$ is $\Psi(0)$ and belongs to the
open cell; the other five Weyl points lie in $A$. Substitution
in the defining formula for $c$ proves the fixed-point assertions.
\end{proof}

\begin{corollary}[Composition with the retained spectral boundary graph]
\label{fc:boundary-composite}
For $q\in W$, let
\[
 \alpha(q)=-\lambda_{\min}(E(q)),\qquad
 \Gamma(q)=Q_{\mathbb O}(\alpha(q);q_1,q_2,q_3).
\]
The exact boundary graph has inverse the off-diagonal coordinate
projection $Q_{\mathbb O}(\lambda;q)\mapsto q$. It extends to an
$M$-equivariant based homeomorphism
\[
 \Gamma^+:W^+\longrightarrow
 (\partial\mathcal C(\mathbb O))^+,
 \qquad\Gamma^+(\infty)=\infty.
\]
Consequently the fully specified map
\begin{equation}
 C_{\det}=\Gamma^+c:F\longrightarrow
           S_{\det,\mathbb O}^{24}
 \label{fc:determinant-collapse}
\end{equation}
is continuous, $M$-equivariant, and has degree $+1$ when the
boundary sphere carries the orientation transported from its
retained ordered coordinates by $\Gamma^+$. Its exact fibers are
\[
 C_{\det}^{-1}\bigl(\Gamma(q)\bigr)=\{\Psi(q)\}
 \quad(q\in W),\qquad C_{\det}^{-1}(\infty)=A.
\]
\end{corollary}

\begin{proof}
The eigenvalues of $\lambda e+E(q)$ are
$\lambda+\lambda_j(E(q))$, with their multiplicities retained.
It is positive semidefinite precisely when
$\lambda\ge-\lambda_{\min}(E(q))$, and lies on the relative
boundary precisely when equality holds. This proves both
compositions with the stated coordinate projection. The minimum
eigenvalue is continuous, so $\Gamma$ is continuous, and the
projection is continuous. A homeomorphism is proper, because the
inverse maps compact sets to compact sets; therefore it extends
to the indicated one-point compactifications.
An element of $M$ fixes $e$, preserves eigenvalues as a Jordan
automorphism, and sends $E(q)$ to $E(mq)$. Thus
$\alpha(mq)=\alpha(q)$ and $\Gamma(mq)=m\Gamma(q)$.
This proves equivariance also at infinity. The asserted
properties, degree, and fibers of $C_{\det}$ now follow from
Theorem~\ref{fc:collapse} by composition with the explicitly
oriented homeomorphism $\Gamma^+$.
\end{proof}

\begin{corollary}[The exact limit on a reverse sphere map]
\label{fc:no-reverse-degree}
Every continuous map $f:S^{24}\to F$ has degree zero.
In particular neither $c$ nor $C_{\det}$ admits a continuous
section or a homotopy section. These assertions coexist with
the degree-one equivariant maps just constructed.
\end{corollary}

\begin{proof}
The cohomology groups $H^8(S^{24};\mathbb Z)$ vanish, so
$f^*\beta_1=f^*\beta_2=0$. Hence
\[
 f^*u_{\mathbb O}=f^*(\beta_1^2\beta_2)
                 =(f^*\beta_1)^2f^*\beta_2=0.
\]
This is exactly the statement that the degree is zero.
A homotopy section of $c$ would satisfy
$f^*c^*\omega_W=\omega_W$, contradicting the equality just
proved and \eqref{fc:degree-one}. The same argument applies
after the homeomorphism $\Gamma^+$.
\end{proof}

\begin{proposition}[Exact maps among the three marked 24-dimensional constructions]
\label[proposition]{prop:three-24d-objects}
The following constructions are related by the explicit spectral graph,
flag-cell collapse, and smash-coordinate maps proved here.
\begin{enumerate}[label=\textup{(\roman*)}]
\item The octonionic flag is the homogeneous \(24\)-manifold
      \(F_4/\operatorname{Spin}(8)\), with
      \(H^8\cong H^{16}\cong\Z^2\) and \(\chi=6\).
\item After a based diffeomorphism \(X\cong S^6\), the fourth smash power
      \(X^{\wedge4}\) is a based sphere \(S_{\wedge}^{24}\); it has
      \(\widetilde H^k=0\) for \(k\ne24\) and \(\chi=2\), and its quotient
      maps depend on the chosen based candidate and cusp collapses.
\item The octonionic equal-diagonal cone has the determinant-boundary
      compactification
      \(S_{\det,\Oo}^{24}=(\partial\mathcal C(\Oo))^+\).  It is produced by
      the explicit cubic \eqref{eq:octonionic-equal-diagonal-determinant},
      not by a smash product.
\end{enumerate}
The first package is not homotopy equivalent to either sphere.  The two
sphere carriers in (ii) and (iii) are related by the based homeomorphism
displayed below. The flag maps to the determinant sphere by the explicit
degree-one map of Corollary~\ref{fc:boundary-composite}; its entire
infinite fibre is the stated $16$-skeleton. The obstruction to a reverse
homotopy section is Corollary~\ref{fc:no-reverse-degree}.

The lower complex determinant orbit pair \((S^{10},S^8)\) from
\cref{thm:sphere-ladder} is a lower rung: its \(S^8\) is the compactified
complex rank-loss locus and its complement has determinant-phase
fundamental group \(\Z\).  It is not a third \(24\)-dimensional object.
\end{proposition}

\begin{proof}
Item (i) is \cref{thm:octonionic-triality,thm:octonionic-cohomology}.
The based smash identity
\((S^6)^{\wedge4}\cong S^{24}\) proves the sphere and cohomology assertions
in (ii).  Item (iii) is
\cref{thm:octonionic-determinant-S24}.  The nonzero degree-eight and
degree-sixteen cohomology of the flag precludes a homotopy equivalence with
either sphere; the Euler characteristics give an independent obstruction.
For an explicit choice-dependent comparison, let
\(\eta:(X,x_{\mathrm{base}})\to(\R^6)^+\) be a based homeomorphism,
write the four six-coordinate vectors consecutively as
$(x_0,\ldots,x_{23})$, and fix the explicit linear isometry
$L:(\mathbb R^6)^4\to\mathbb O^3$ given by
$q_j=\sum_{a=0}^7x_{8(j-1)+a}u_a$ in the fixed ordered orthonormal
octonion basis. Let $b=\mathcal P$ be the off-diagonal coordinate inverse
of the spectral graph in Theorem~\ref{tri:thm:spectral-graph}.
The smash compactification map
\[
 \kappa:((\R^6)^+)^{\wedge4}\longrightarrow((\R^6)^4)^+
\]
sends a tuple of finite vectors to their ordered tuple and sends the
collapsed wedge to infinity.  It is a continuous bijection from a compact
space to a Hausdorff space, hence a homeomorphism.  Thus
\[
 (b^+)^{-1}\circ L^+\circ\kappa\circ\eta^{\wedge4}:
 X^{\wedge4}\xrightarrow{\ \cong\ }
       (\partial\mathcal C(\Oo))^+
\]
is a based homeomorphism with inverse obtained by reversing these four
maps.  It depends explicitly on \(\eta\) and \(L\); neither marked
construction specifies them.  The final
lower-rung assertions are \eqref{eq:S10-S8-pair} and
\eqref{eq:rank-complement}.
\end{proof}

\section{What the mod--12 and mod--24 arithmetic actually is}
\label{sec:mod-12-24}

\subsection{The integral peripheral representation}

The matrices \(C_1,C_2',C_0\) in
\eqref{eq:elliptic-peripheral-matrices} are not merely of orders three,
four, and infinite.  Since \(C_2'\) and \(C_0\) generate
\(SL_2(\Z)\), their abelianization detects one exact modulus.

\begin{theorem}[Peripheral modulus and the mod--24 boundary]
\label[theorem]{thm:peripheral-moduli}
Let
\[
 \operatorname{ab}:SL_2(\Z)\twoheadrightarrow SL_2(\Z)_{\mathrm{ab}}
\]
be the abelianization map.  There is a unique isomorphism
\[
 \overline\nu:SL_2(\Z)_{\mathrm{ab}}
       \xrightarrow{\ \cong\ }\Z/12
\]
for which the marked character
\begin{equation}
 \chi=\overline\nu\circ\operatorname{ab}:
 SL_2(\Z)\longrightarrow\Z/12
 \label{eq:marked-peripheral-character}
\end{equation}
sends \(C_0\) to \(-1\).  Its values on the three peripheral matrices are
\begin{equation}
 \chi(C_1)=4,\qquad \chi(C_2')=-3,\qquad \chi(C_0)=-1
 \pmod {12}.
 \label{eq:peripheral-mod12-values}
\end{equation}
The residues \(4\) and \(-3=9\) are complementary idempotents in
\(\Z/12\); additively they have exact orders \(3\) and \(4\).

Let
\[
 j:\Z/12\hookrightarrow\Z/24,\qquad j(a)=2a.
\]
Every group homomorphism \(\psi:SL_2(\Z)\to\Z/24\) has image in the even
subgroup and factors uniquely through \(j\).  In particular
\begin{equation}
 \chi_{24}=j\chi,\qquad
 \bigl(\chi_{24}(C_1),\chi_{24}(C_2'),\chi_{24}(C_0)\bigr)
 =(8,-6,-2)\pmod {24}
 \label{eq:peripheral-mod24-values}
\end{equation}
is the unique mod--\(24\) re-encoding with
\[
 \exp\!\left(\frac{2\pi i\,\chi_{24}(M)}{24}\right)
 =
 \exp\!\left(\frac{2\pi i\,\chi(M)}{12}\right)
 \qquad(M\in SL_2(\Z)).
\]
Its image has order \(12\), not \(24\).  If
\(r:\Z/24\to\Z/12\) is ordinary reduction, then
\(r\circ j=[2]\), and there is no group homomorphism
\(\widetilde\chi:SL_2(\Z)\to\Z/24\) with
\(r\circ\widetilde\chi=\chi\).
\end{theorem}

\begin{proof}
Put
\[
 S_0=\begin{pmatrix}0&-1\\1&0\end{pmatrix},\qquad T=C_0.
\]
Then \(C_2'=S_0^3\), \(C_1=T^{-1}S_0\), and the standard integral modular
presentation is
\[
 SL_2(\Z)=
 \langle S_0,T\mid S_0^4=1,\ S_0^2=(S_0T)^3\rangle.
\]
Here is why these relations suffice.  With \(R=S_0T\), the presented
group is \(G=\langle S_0,R:S_0^4=1,S_0^2=R^3\rangle\).
The element \(c=S_0^2=R^3\) is central and has order at most two;
\(G/\langle c\rangle=C_2*C_3\).  The modular normal-form theorem
identifies \(PSL_2(\Z)\) with this same free product on the images of
\(S_0,R\) \cite[Theorem~C.1, pp.~18--19]{ConradSL2Z}.
The matrices \(S_0,T\) generate \(SL_2(\Z)\)
\cite[Theorem~1.1, pp.~1--2]{ConradSL2Z}; hence the map from \(G\)
onto \(SL_2(\Z)\) is surjective and its kernel lies in
\(\langle c\rangle\).  Since \(c\) maps to \(-I_2\ne I_2\), the
kernel is trivial.  This proves the presentation (equivalently
\(C_4*_{C_2}C_6\)) with the stated matrix generators.
After abelianization, with additive generators \(s,t\), its two relations
are
\[
 4s=0,\qquad 2s=3(s+t),
\]
or \(4s=0\) and \(s+3t=0\).  The relation matrix
\(\left(\begin{smallmatrix}4&0\\1&3\end{smallmatrix}\right)\) has Smith
normal form \(\diag(1,12)\), so the abelianization is cyclic of exact order
\(12\), generated by \(t\).  Choosing \(t\mapsto-1\) gives
\[
 s\mapsto3,\qquad
 C_2'=S_0^3\mapsto9=-3,\qquad
 C_1=T^{-1}S_0\mapsto1+3=4,
\]
which proves \eqref{eq:peripheral-mod12-values}.  Direct calculation in
\(\Z/12\) gives
\[
 4^2=4,\qquad9^2=9,\qquad4+9=1,\qquad4\cdot9=0.
\]

Every group homomorphism from \(SL_2(\Z)\) to an abelian group factors through its
abelianization.  An element \(x\in\Z/24\) satisfies \(12x=0\) exactly when
\(x\) is even, proving the unique factorization through \(j\) and
excluding a surjective character.  Also \(rj(a)=2a\pmod {12}\).
Finally a reduction-compatible lift
would send \(C_0\) to either \(11\) or \(23\) modulo \(24\), both of which
have order \(24\); an image of the order-\(12\) abelianized generator
cannot have that order.  This proves the final assertion.
\end{proof}

\begin{proposition}[The primitive cusp quotient singles out mod--3, not
mod--12 or mod--24]
\label[proposition]{prop:cusp-diagonal-modulus}
For \(N\ge1\), let
\[
 d_N:\Z^3\longrightarrow\Z/N,\qquad
 d_N(a_1,a_2,a_3)=a_1+a_2+a_3\pmod N.
\]
There is a homomorphism \(\overline d_N:\Z^2\to\Z/N\) satisfying
\(d_N=\overline d_NP\) if and only if \(N\mid3\).  For \(N=3\), it is
explicitly
\[
 \overline d_3(y_1,y_2)=y_1-y_2\pmod3.
\]
Similarly, \(d_N\) factors through \(Q\) if and only if \(N\mid3\); for
\(N=3\) the exact formula is
\[
 d_3=(1,1)Q.
\]
Thus neither primitive lattice quotient in \cref{thm:lattice-quotient}
produces a diagonal-sum character modulo \(12\) or modulo \(24\).
\end{proposition}

\begin{proof}
Both \(\ker P\) and \(\ker Q\) are generated by
\(k=(1,1,1)^{\mathsf T}\).  A homomorphism out of \(\Z^3\) factors
through either quotient precisely when it kills that kernel.  Since
\(d_N(k)=3\pmod N\), this occurs precisely when \(N\mid3\).  At \(N=3\),
direct multiplication with the matrices in \eqref{eq:PQ} gives
\[
 (1,-1)P=(1,-2,1)\equiv(1,1,1)\pmod3,\qquad
 (1,1)Q=(1,-2,1)\equiv(1,1,1)\pmod3,
\]
which proves both formulas and the claimed boundary.
\end{proof}

\begin{proposition}[Coefficient reduction and determinant winding do not
select a modulus]
\label[proposition]{prop:no-modulus-from-cohomology-or-determinant}
For every \(N\ge1\), coefficient reduction gives a graded-ring
isomorphism
\[
 H^*(\operatorname{Fl}(\Oo);\Z/N)
 \cong H^*(\operatorname{Fl}(\Oo);\Z)\otimes\Z/N,
\]
and its reduced Euler class
\[
 e_N=6u_{\Oo}\bmod N\in H^{24}(\operatorname{Fl}(\Oo);\Z/N)
\]
has additive order \(N/\gcd(N,6)\).  In particular \(e_{12}\) has order
two and \(e_{24}\) has order four.  Ordinary coefficient reduction is
defined for every \(N\); this operation alone makes no choice of
\(N=12\) or \(N=24\).

Likewise, in the lower complex orbit pair \((S^{10},S^8)\), the
determinant phase identifies the fundamental group of the rank-full
complement in \eqref{eq:rank-complement} with \(\Z\).  Reducing its winding
number modulo \(N\) gives a surjection to \(\Z/N\) for every \(N\), so
this plain winding-number reduction makes no choice of \(12\) or \(24\).  This
phase statement is not being transferred to the real octonionic
determinant in \eqref{eq:octonionic-equal-diagonal-determinant}.
\end{proposition}

\begin{proof}
All integral cohomology groups in
\eqref{eq:octonionic-Poincare} are free.  The universal coefficient
sequence therefore has no \(\operatorname{Tor}\) term, and multiplication
is preserved by coefficient change.  Equation
\eqref{eq:octonionic-Euler-six} gives the stated reduced Euler class.
Since the reduction of the primitive \(u_{\Oo}\) generates
\(H^{24}(\operatorname{Fl}(\Oo);\Z/N)\cong\Z/N\), the least positive
integer \(m\) annihilating \(e_N\) is the least \(m\) with \(N\mid6m\),
namely \(N/\gcd(N,6)\).
For the determinant complement,
\eqref{eq:rank-complement} is a deformation retraction onto the phase
circle, and \(\arg\det\) is that retraction coordinate.  Its winding is an
isomorphism on \(\pi_1\), whose reduction modulo any \(N\) is surjective.
\end{proof}

\subsection[The distinct stable mod--12-to-mod--24 injection]
{The distinct stable \(\Z/12\hookrightarrow\Z/24\)}

The injection \(j(a)=2a\) is realized, after compatible generator choices,
by the universal sixfold suspension homomorphism below.  The cusp quotient
supplies one particular element in its domain; it does not supply the
injection itself.  Neither construction comes from the octonionic flag or
from a \(24\)-dimensional smash source.  We state all data needed to type
these assertions.

Let \(v_0\) be the suspension vertex to which the central cusp fibre is
collapsed, write \(W=q_2^{-1}(v_0)\), and choose explicitly
\(x_{\mathrm{base}}\in W\).  The exact cusp collapse in the canonical audit
then gives a based quotient
\[
 q_2:(X,x_{\mathrm{base}})\longrightarrow
 (\mathcal A,v_0),\qquad
 \mathcal A=\operatorname{Susp}(S^1\times\overline T^2).
\]
Let
\[
 K^{(1)}=(S^1\times\{0\})\cup(\{0\}\times S^1)
 \subset\overline T^2,
\]
and let \(c:\overline T^2\to\overline T^2/K^{(1)}\cong S^2\) be the
degree-one cell quotient, with \(*=c(K^{(1)})\).  Write
\[
 \Sigma_r(S^2,*)=
 \frac{[0,\epsilon]\times S^2}
 { (\{0,\epsilon\}\times S^2)\cup([0,\epsilon]\times\{*\})}
\]
for reduced suspension, the entire denominator being collapsed to one
basepoint.  The formula
\[
 \widehat p_2([r,e^{i\phi},u])=[r,c(u)]
 \in\Sigma_r(S^2,*)\cong S^3
\]
defines a based map: it sends both suspension vertices of \(\mathcal A\)
and the entire cylinder over \(K^{(1)}\) to that one basepoint.  Thus
\begin{equation}
 g=\widehat p_2q_2:(X,x_{\mathrm{base}})\longrightarrow(S^3,*).
 \label{eq:one-copy-g}
\end{equation}
After the fixed based identification \(X\cong S^6\), this determines
\([g]\in\pi_6(S^3)\).  The construction and its endpoint continuity are
part of the cited candidate audit \cite{AlpogeS6_2026}.
The exact research-file locators are the reduced-suspension construction in
\begin{center}\small
\texttt{cusp\_sphere\_route\_strengthening.tex}
\end{center}
and the one-copy map in
\begin{center}\small
\texttt{cusp\_s12\_carrier\_homotopy.tex}.
\end{center}

\begin{theorem}[The stable even-subgroup morphism and its exact scope]
\label[theorem]{thm:stable-12-to-24}
Sixfold reduced suspension is an injection
\begin{equation}
 E^6:\pi_6(S^3)\cong\Z/12
 \hookrightarrow
 \pi_{12}(S^9)\cong\Z/24
 \label{eq:stable-even-injection}
\end{equation}
whose image is the even subgroup.  With compatible generators \(\alpha\) and
\(\nu_9\),
\[
 E^6(a\alpha)=2a\nu_9.
\]
Consequently there is a unique \(a\in\Z/12\) for which the particular
map in \eqref{eq:one-copy-g} satisfies
\begin{equation}
 [E^6g]=2a\nu_9.
 \label{eq:one-copy-stable-class}
\end{equation}
The coefficient \(a\) is not determined by point-set surjectivity or by
the ordinary homology map.

For every pair of based maps \(x,y:S^6\to S^3\),
\begin{equation}
 [x\wedge y]=0\in\pi_{12}(S^6).
 \label{eq:external-smash-zero}
\end{equation}
In particular,
\[
 g^{\wedge4}:S^{24}\longrightarrow S^{12}
\]
is based null-homotopic.  Thus the \(\Z/24\) in
\eqref{eq:stable-even-injection} is the ambient target group of the
sixfold-suspended one-copy map \(S^{12}\to S^9\).  The separate
\(24\)-dimensional fourth-smash map is null-homotopic.
\end{theorem}

\begin{proof}
The relevant initial suspension chain is
\[
 \pi_6(S^3)\xrightarrow{E}\pi_7(S^4)
                   \xrightarrow{E}\pi_8(S^5).
\]
Hatcher computes these groups as \(\Z/12\), \(\Z\oplus\Z/12\),
and \(\Z/24\), respectively.  The Hopf fibration identifies the torsion
summand of the middle group with the suspended first group.  In the
quotient defining the second suspension, twice the Hopf generator of
the free summand is identified with a generator of this torsion summand
\cite[Theorem~5.41(b), p.~573, and pp.~577--578]
{HatcherSpectralSequences}.
Choose compatible signs and generators \(\alpha\in\pi_6(S^3)\) and
\(\nu_5=E\nu_4\in\pi_8(S^5)\).  The displayed quotient relation gives
\(E^2\alpha=2\nu_5\).  Multiplication by two from \(\Z/12\) to
\(\Z/24\) is injective: \(24\mid2a\) implies \(12\mid a\).
For \(n\ge5\), Freudenthal gives an isomorphism
\[
 E:\pi_{n+3}(S^n)\longrightarrow\pi_{n+4}(S^{n+1}),
 \qquad n+3\le2n-2.
\]
Applying it for \(n=5,6,7,8\) and putting \(\nu_9=E^4\nu_5\)
gives \(E^6\alpha=2\nu_9\) and
\(\pi_{12}(S^9)\cong\pi_3^{\mathrm S}\cong\Z/24\).
The stable group is independently tabulated in
\cite[Table~1, p.~113]{IsaksenWangXu2023}.  This proves
\eqref{eq:stable-even-injection} and the unique formula
\eqref{eq:one-copy-stable-class}.

For \eqref{eq:external-smash-zero}, the ordered external smash has the
exact factorization
\[
 [x\wedge y]=(E^3[x])\circ(E^6[y])
\]
by the reduced-join formula
\cite[Proposition~3.1]{Toda1962}; its ordered-coordinate sign is
\((-1)^{3\cdot6}=+1\).  The integral low-dimensional table gives
\(\pi_{12}(S^6)\cong\Z/2\)
\cite[opening table, p.~339]{Hatcher2002}.  The class \(E^6[y]\) is even in
\(\pi_{12}(S^9)\cong\Z/24\), and postcomposition by \(E^3[x]\) is a group
homomorphism into the exponent-two group \(\pi_{12}(S^6)\).  Its value on
an even element is therefore zero.  Hence \(g^{\wedge2}\) is
null-homotopic, and smashing its based null-homotopy with
\(g^{\wedge2}\) proves the assertion for \(g^{\wedge4}\).
\end{proof}

\begin{proposition}[The generator-free bridge between the two
\(\Z/12\)-objects]
\label[proposition]{prop:generator-free-bridge}
The marked one-copy class defines a homomorphism
\begin{equation}
 \Theta_g:\Z/12\longrightarrow\pi_6(S^3),\qquad
 \Theta_g([n])=n[g].
 \label{eq:Theta-g}
\end{equation}
Therefore
\[
 \Theta_g\circ\chi:SL_2(\Z)\longrightarrow\pi_6(S^3)
\]
is a well-defined homomorphism from the marked peripheral data to the one-copy
homotopy group.  If \(\alpha\) and \(\nu_9\) are compatible generators as
in \cref{thm:stable-12-to-24} and \([g]=a\alpha\), then
\begin{equation}
 E^6\Theta_g\chi(M)=2a\,\chi(M)\nu_9
 \in\pi_{12}(S^9)\cong\Z/24.
 \label{eq:generator-free-stable-bridge}
\end{equation}
For an integer lift \(\widetilde a\) put \(d=\gcd(\widetilde a,12)\).
Then
\[
 \ker\Theta_g=\langle[12/d]\rangle\subset\Z/12,
 \qquad \im\Theta_g=\langle[g]\rangle=\langle d\alpha\rangle,
 \qquad |\im\Theta_g|=12/d.
\]
The first group has order \(d\).  Since \(\chi\) is surjective and
\(E^6\) injective, both composites have image order \(12/d\), and
\(\ker(E^6\Theta_g\chi)=\chi^{-1}(\langle[12/d]\rangle)\).
This bridge is canonical once the marked character \(\chi\), the based
map \(g\), and the based identification \(S^6\cong X\) are fixed.
It does not prove \(d=1\), equivalently that \([g]\) is a
generator or that \(\Theta_g\) is an isomorphism, and it does not define an
equivariant map of degenerations.
\end{proposition}

\begin{proof}
The source of \(\Theta_g\) is well-defined because
\(\pi_6(S^3)\cong\Z/12\), so \(12[g]=0\).  It is a homomorphism by the
abelian group law on \(\pi_6(S^3)\).  Composition with the already defined
\(\chi\) is therefore a well-defined homomorphism.  Writing
\([g]=a\alpha\) and applying the homomorphism \(E^6\) gives
\[
 E^6\Theta_g\chi(M)
 =\chi(M)E^6(a\alpha)=2a\,\chi(M)\nu_9,
\]
which is \eqref{eq:generator-free-stable-bridge} for every \(M\in SL_2(\Z)\).
The congruence \(n\widetilde a\equiv0\pmod{12}\) is equivalent to
\(12/d\mid n\), giving the kernel.  B\'ezout's identity gives
\(\langle a\alpha\rangle=\langle d\alpha\rangle\), of order \(12/d\).
Surjectivity of \(\chi\) preserves that image, and injectivity of
\(E^6\) preserves its order and the composite kernel.  This group calculation by itself does not compute \(a\).
The explicit framed calculation in \cref{thm:cusp-g-zero} below does:
\(a=0\), so \(d=12\), \(\ker\Theta_g=\Z/12\), and every
image displayed here is zero.
\end{proof}

\subsection{The original framed cusp coefficient is zero}
Let \(X\) be the smooth compact source of the original construction,
\(f:X\to B\) its map, and \(t\) the specified cusp coordinate.
Put \(W=f^{-1}(0)\), choose the original \(x_0\in W\), and retain the
ordered quotient lattice basis induced by
\((\widehat\gamma,\widehat u)\), with kernel generated by
\((\widehat w,\widehat\delta)\). Thus the quotient torus is
\(\overline T^2=\R^2/\Z^2\). The original periods are
\begin{equation}
 Z(s)=sB_0+C(t),\qquad e^{2\pi i s}=t,\qquad
 B_0=\begin{pmatrix}0&1\\-1&0\end{pmatrix},\qquad
 C(t)=\begin{pmatrix}6\mu(t)&h(t)\\b(t)-h(t)&\mu(t)\end{pmatrix}.
 \label{gzero:eq:periods}
\end{equation}
Every entry of \(C\) is holomorphic at \(t=0\).
The dense torus of the original toric cusp has coordinates
\((x_1,x_2,t)\), with
\(x=\exp(2\pi i\zeta)\). The residual lattice action is exactly
\begin{equation}
 \Psi_\lambda(x,t)=
 \bigl(t^{B_0\lambda}\exp(2\pi iC(t)\lambda)x,t\bigr),
 \qquad\lambda\in\Z^2.
 \label{gzero:eq:action}
\end{equation}
It extends freely and properly discontinuously across the original
unimodular toric filling. In particular the quotient projection is a local
biholomorphism at every toric fixed point. These are the reconstructed
data and proved local properties of the original source
\cite[\S4, source pp.~24--29]{AlpogeS6_2026}.

Let
\[
 K^{(1)}=(\R/\Z\times\{0\})\cup(\{0\}\times\R/\Z),\qquad
 c:\overline T^2\longrightarrow \overline T^2/K^{(1)}\cong S^2
\]
be the specified degree-one cell quotient. Its restriction to the open
two-cell is taken with the smooth coordinates used in the original
regular-fibre description. Both radial endpoints and the complete
basepoint cylinder are collapsed in
\[
 \Sigma_r(S^2,*)=
 \frac{[0,\epsilon]\times S^2}
 {(\{0,\epsilon\}\times S^2)\cup([0,\epsilon]\times\{*\})}
 \cong S^3.
\]
The original based map is
\begin{equation}
 g(x)=[|t|,c(\overline\pi_f(x))]\quad(0<|t|<\epsilon),\qquad
 g(W)=g(X\setminus\operatorname{int}N_0)=*.
 \label{gzero:eq:g}
\end{equation}

\begin{theorem}[The original cusp map is null-homotopic]
\label[theorem]{thm:cusp-g-zero}
The based map \(g:(X,x_0)\to(S^3,*)\) in \eqref{gzero:eq:g} is based
null-homotopic. More precisely, a regular inverse image with its actual
ordered normal framing induced by \((r,c_1,c_2)\) is framed nullbordant
inside one original cusp chart times \([0,1]\), disjoint from
\(\{x_0\}\times[0,1]\).
\end{theorem}

\begin{proof}
The proof constructs the complete inverse image, compares its actual
framing with an explicitly extendible framing, and gives the extension.

\paragraph{1. The exact complete inverse image.}
Retain the lift
\[
 u^0=\begin{pmatrix}-1/4\\1/4\end{pmatrix},\qquad
 \overline u^0=(3/4,1/4)\pmod{\Z^2},\qquad
 B_0u^0=\begin{pmatrix}1/4\\1/4\end{pmatrix}.
\]
Thus \(\overline u^0\notin K^{(1)}\), and \(B_0u^0\) lies in the
interior of the original triangle \(\operatorname{conv}(0,e_1,e_2)\).
Use its original unimodular chart
\begin{equation}
 z_0=t/(x_1x_2),\qquad z_1=x_1,\qquad z_2=x_2,
 \qquad t=z_0z_1z_2.
 \label{gzero:eq:chart}
\end{equation}
Set
\[
 \ell=\log|t|<0,\qquad R(t)=-2\pi\operatorname{Im}C(t),\qquad
 A(t)=\ell B_0+R(t).
\]
For a representative \(\zeta=Z(s)u+k\), where \(u,k\in\R^2\),
\(\operatorname{Im}s=-\ell/(2\pi)\) gives the exact identity
\begin{equation}
 \begin{pmatrix}\log|x_1|\\\log|x_2|\end{pmatrix}
 =A(t)u.
 \label{gzero:eq:moduli}
\end{equation}
The quotient value is \(u\bmod\Z^2\).
Since \(B_0\) is orthogonal, for all \(v\in\R^2\),
\begin{equation}
 \|A(t)v\|\geq (|\ell|-\|R(t)\|)\|v\|.
 \label{gzero:eq:invertibility}
\end{equation}
The boundedness of \(R\) therefore makes \(A(t)\) invertible on a
sufficiently small punctured disc.

For any point of the fibre quotient over \(\overline u^0\), subtracting
the appropriate original period, equivalently applying
\(\Psi_{-\lambda}\), gives a representative with the exact lift
\(u=u^0\). This representative is unique: another residual translate
has logarithmic moduli differing by \(A(t)\lambda\), which vanishes
only for \(\lambda=0\). Conversely every dense-torus point satisfying
\(\log|x|=A(t)u^0\) has this quotient value, since the remaining real
parts of \(\zeta\) are exactly the kernel coordinates \(k\).
For each fixed \(t\), varying \(k\bmod\Z^2\) gives all independent
phases of \(x_1,x_2\).

The following change of coordinates uses the full original period units:
\begin{align}
 h_j^0(t)&=\exp\bigl(-2\pi i(C(t)u^0)_j\bigr),\qquad j=1,2,\\
 \mathcal F(z_0,z_1,z_2)
 &=(w_0,w_1,w_2)
 =\left(\frac{z_0}{h_1^0(t)h_2^0(t)},\,
 h_1^0(t)z_1,\,h_2^0(t)z_2\right).
 \label{gzero:eq:straightening}
\end{align}
The functions \(h_j^0\) are nowhere-zero holomorphic units, including
at \(t=0\). The product \(w_0w_1w_2\) is exactly \(t\).
The inverse of \(\mathcal F\) is
\[
 (w_0,w_1,w_2)\longmapsto
 \left(w_0h_1^0(t)h_2^0(t),\,
 \frac{w_1}{h_1^0(t)},\,\frac{w_2}{h_2^0(t)}\right),
 \qquad t=w_0w_1w_2.
\]
Thus this is a biholomorphism of the original chart restricted to the
cusp disc, not a replacement of \(C(t)\).

For \(r_0\in(0,\epsilon)\), the entire inverse image of
\(y_0=[r_0,c(\overline u^0)]\) becomes exactly
\begin{equation}
 T_{r_0}=
 \{|w_0|=\rho_0,\ |w_1|=\rho_1,\ |w_2|=\rho_2\},\qquad
 (\rho_0,\rho_1,\rho_2)=
 (r_0^{1/2},r_0^{1/4},r_0^{1/4}).
 \label{gzero:eq:torus}
\end{equation}
Indeed, \(\log|h^0(t)|=2\pi\operatorname{Im}C(t)u^0=-R(t)u^0\),
so \eqref{gzero:eq:moduli} gives \(\log|w_j|=\ell(B_0u^0)_j\) for
\(j=1,2\); the product identity gives the exponent \(1/2\) for
\(w_0\). Conversely these moduli recover \(u=u^0\) by invertibility
of \(A(t)\). The angular map is explicitly
\[
 (\arg t,\arg w_1,\arg w_2)\longleftrightarrow
 (\arg w_0,\arg w_1,\arg w_2),\qquad
 \arg w_0=\arg t-\arg w_1-\arg w_2\pmod{2\pi},
\]
with inverse \(\arg t=\arg w_0+\arg w_1+\arg w_2\).
Its integral matrix has determinant one. This includes the complete
angular mapping torus and all kernel phases, with their original signs.

Let \(o\) be the origin of \eqref{gzero:eq:chart}. The original free properly
discontinuous action supplies a neighborhood \(V\) of \(o\) on which
the quotient projection is injective. To see the local injectivity
directly, first choose a relatively compact neighborhood; proper
discontinuity leaves only finitely many translates that can intersect it.
None of those nonidentity translates fixes \(o\). Shrinking finitely
many times separates the neighborhood from all of them. The quotient
map is a local biholomorphism, so its restriction to \(V\) is a chart
of the original source.
Choose a polydisc about \(0\) whose closure lies in \(\mathcal F(V)\).
For sufficiently small \(r_0\), both \(T_{r_0}\) and
\begin{equation}
 \{|w_0|\leq\rho_0,\ |w_1|=\rho_1,\ |w_2|=\rho_2\}
 \label{gzero:eq:spatial-bound}
\end{equation}
lie in this polydisc. This proves that the whole inverse image embeds
in one chart and supplies space for the bounding manifold below.

\paragraph{2. The actual normal framing, including its angle dependence.}
Put \(L_j=\log|w_j|\), and use the unique real lift of the quotient
coordinate near \(u^0\). From the exact change of coordinates,
\begin{equation}
 \begin{pmatrix}L_1\\L_2\end{pmatrix}
 =\ell B_0u^0+A(t)(u-u^0),\qquad
 L_0=\ell-L_1-L_2.
 \label{gzero:eq:nearby-moduli}
\end{equation}
Differentiating on \(u=u^0\) annihilates the term
\((dA(t))(u-u^0)\). Hence no derivative of \(C(t)\) has been omitted:
the exact differential there is
\begin{equation}
 \begin{pmatrix}dL_0\\dL_1\\dL_2\end{pmatrix}
 =J(t)\begin{pmatrix}d\ell\\du_1\\du_2\end{pmatrix},\quad
 J(t)=
 \begin{pmatrix}
 1/2&-A_{11}-A_{21}&-A_{12}-A_{22}\\
 1/4&A_{11}&A_{12}\\
 1/4&A_{21}&A_{22}
 \end{pmatrix}.
 \label{gzero:eq:J}
\end{equation}
Adding the second and third rows to the first gives \((1,0,0)\),
and expansion in that row proves \(\det J(t)=\det A(t)\).

Let \(D=Dc_{\overline u^0}\) be the actual two-by-two derivative in
the ordered open-cell target coordinates \((c_1,c_2)\). It is a fixed
invertible matrix; all its entries and signs are retained. Let
\[
 H=\begin{pmatrix}r_0&0&0\\0&D_{11}&D_{12}\\0&D_{21}&D_{22}\end{pmatrix}.
\]
Let \(n_j\) be the real radial vector \(w_j\) in the \(j\)-th complex
coordinate, with other components zero. Then \(dL_i(n_j)=\delta_{ij}\).
Writing \(N=(n_0,n_1,n_2)\) as the matrix with these vector columns, the actual
\(dg\)-dual normal framing is, in the normal quotient bundle,
\begin{equation}
 \mathfrak n(t)=N J(t)H^{-1}.
 \label{gzero:eq:actual-frame}
\end{equation}
Indeed, \(dg(N)=HJ(t)^{-1}\), so multiplication by \(J(t)H^{-1}\)
gives the identity. In the Euclidean metric of this chart the radial
vectors are perpendicular to the product torus; a different source
metric changes their representatives by the canonical projection to its
normal bundle and preserves this normal-quotient framing.

Put \(\ell_0=\log r_0\). Decrease \(r_0\) further, if necessary, so
that \(|\ell_0|>2\sup_{|t|\leq r_0}\|R(t)\|\). For
\(0\leq\sigma\leq1\), define only a homotopy of frames by
\[
 A_\sigma(t)=\ell_0B_0+\sigma R(t),
\]
and let \(J_\sigma(t)\) be \eqref{gzero:eq:J} with \(A_\sigma\) in place
of \(A\). For every \(v\in\R^2\),
\[
 \|A_\sigma(t)v\|\geq
 (|\ell_0|-\|R(t)\|)\|v\|>0\quad(v\ne0).
\]
Thus all matrices are invertible. At \(\sigma=0\),
\begin{equation}
 J_0=
 \begin{pmatrix}
 1/2&\ell_0&-\ell_0\\
 1/4&0&\ell_0\\
 1/4&-\ell_0&0
 \end{pmatrix},\qquad \det J_0=\ell_0^2>0.
 \label{gzero:eq:J0}
\end{equation}
Continuity also proves \(\det J_\sigma(t)>0\) for every \(\sigma,t\).
The homotopy \(N J_\sigma(t)H^{-1}\), run from \(\sigma=1\) to
\(\sigma=0\), compares the actual frame with the exact constant-matrix
frame \(N P\), where
\begin{equation}
 P=J_0H^{-1}.
 \label{gzero:eq:P}
\end{equation}
The original embedded torus and original functions \(C(t)\) have not
changed. A framing homotopy is a framed cylinder: choose a smooth
parameter function constant near its endpoints and use the displayed
frames on \(T_{r_0}\times[0,1]\).

\paragraph{3. A collared framed nullbordism with the exact constant matrix.}
Use a new time coordinate \(v\in[0,1]\). Define
\begin{equation}
 q(v)=
 \begin{cases}
 0,&0\leq v\leq1/4,\\
 \rho_0^2\exp\bigl(4-1/(v-1/4)\bigr),&1/4<v\leq1.
 \end{cases}
 \label{gzero:eq:q}
\end{equation}
This function is smooth: all right derivatives of
\(\exp(-1/(v-1/4))\) at \(1/4\) vanish. It is strictly increasing
for \(v>1/4\), satisfies \(q(1/2)=\rho_0^2\), and
\(q'(1/2)=16\rho_0^2>0\).
In the chart times \([0,1]\), put
\begin{equation}
 \mathcal W=
 \{|w_0|^2+q(v)=\rho_0^2,\quad
       |w_1|^2=\rho_1^2,\quad|w_2|^2=\rho_2^2\}.
 \label{gzero:eq:bordism}
\end{equation}
No point has \(v>1/2\). The three defining differentials are independent:
the latter two have nonzero gradients in their separate complex factors;
the first has nonzero \(w_0\)-gradient whenever \(w_0\ne0\), and at
\(w_0=0\) it has the nonzero \(v\)-derivative \(q'(1/2)\).
Thus \(\mathcal W\) is a smooth compact four-manifold. Its only
boundary is \(T_{r_0}\times\{0\}\), and on \(0\leq v\leq1/4\)
it is the exact product collar \(T_{r_0}\times[0,1/4]\).
The north locus \(w_0=0,v=1/2\) is an interior cap: the implicit
function theorem expresses \(v\) smoothly in terms of
\((\operatorname{Re}w_0,\operatorname{Im}w_0)\) there. The first
factor is a disk obtained by rotating its radial profile with this smooth
cap and one boundary circle. Therefore
\(\mathcal W\cong D^2\times T^2\).
Its spatial coordinates lie in \eqref{gzero:eq:spatial-bound}, so the complete
manifold lies in the injective original chart.

In the Euclidean chart product metric define
\[
 V_0=(w_0,0,0,q'(v)/2),\qquad
 V_1=(0,w_1,0,0),\qquad V_2=(0,0,w_2,0).
\]
These are one half of the gradients of the three defining equations
in \eqref{gzero:eq:bordism}. They are mutually orthogonal and never zero.
For example, if \(F_0=|w_0|^2+q(v)-\rho_0^2\), then
\[
 dF_0(V_0)=2|w_0|^2+\frac12q'(v)^2>0
\]
at every point of \(\mathcal W\). On the entire collar \(q'=0\), so
\((V_0,V_1,V_2)=(N,0)\) exactly. Consequently
\begin{equation}
 (V_0,V_1,V_2)P
 \label{gzero:eq:extended-frame}
\end{equation}
is a framing of the normal three-plane bundle of \(\mathcal W\),
and its restriction to the collar is exactly the framing \(NP\) from
\eqref{gzero:eq:P}, with every original target-coordinate sign retained.
Concatenate the framing cylinder of the preceding paragraph with this
nullbordism, rescaling the two interval lengths. Their framings agree
exactly on the common product collar, so the concatenation is a smooth
framed nullbordism of the actual inverse image.

\paragraph{4. The based map represented by this framed inverse image.}
Choose the injective chart to exclude \(x_0\) if its center is not
\(x_0\). If its center is \(x_0\), the constructed nullbordism still
excludes \(x_0\), since \(|w_1|=\rho_1>0\) and
\(|w_2|=\rho_2>0\) everywhere on it. Compactness therefore gives
tubular neighborhoods for the framed cylinder and cap disjoint from
\(\{x_0\}\times[0,1]\).

The original continuous quotient map is smooth and a submersion near
its entire inverse image of \(y_0\). First make it constant in a
neighborhood of \(x_0\), by a based homotopy that changes nothing near
that inverse image. Explicitly, take a target chart
\(\kappa:O\to\R^3\) with \(\kappa(*)=0\), whose image is a ball,
and \(y_0\notin O\). By continuity choose a neighborhood \(U\) of
\(x_0\) with \(g(U)\subset O\), and a smooth bump function
\(\beta\) supported in \(U\) and equal to one near \(x_0\).
For \(x\in U\), replace \(g(x)\) during this homotopy by
\(\kappa^{-1}((1-a\beta(x))\kappa(g(x)))\), \(0\leq a\leq1\),
and leave \(g\) fixed off \(U\). Convexity of the target chart image
makes the formula well-defined, the vanishing of \(\beta\) near the
boundary of its support proves continuity of the pasted formula, and
\(\kappa(g(x_0))=0\) makes the homotopy based. It never introduces
an inverse image of \(y_0\). Its endpoint is smooth near \(x_0\) and
agrees with \(g\) near the entire old inverse image.

Denote the endpoint of the preceding homotopy by \(g_b\). Choose one
closed tubular neighborhood \(U_1\) of the inverse image, contained in
the smooth submersion region and with the entire inverse image contained
in \(\operatorname{int}U_1\). The compact set
\(X\setminus\operatorname{int}U_1\) has image under \(g_b\) missing
\(y_0\). In a fixed target metric put
\[
 d_0=\min_{x\in X\setminus\operatorname{int}U_1}
       d(g_b(x),y_0)>0.
\]
Relative smooth approximation, fixed on the whole closed set \(U_1\)
and on a neighborhood of \(x_0\), gives a smooth based representative
\(g_{\mathrm{sm}}\) whose error on
\(X\setminus\operatorname{int}U_1\) is less than \(d_0/2\).
It therefore has the same inverse image and differential on \(U_1\),
and \(d(g_{\mathrm{sm}}(x),y_0)>d_0/2\) on its complement. Thus there
is no uncontrolled intervening region and no additional inverse image.
Make the approximation also closer than a fixed target injectivity
radius; its short geodesics give the based homotopy.

The relative Pontryagin construction now applies to the explicit
framed nullbordism. It produces a based homotopy to a map missing
\(y_0\); stereographic contraction of \(S^3\setminus\{y_0\}\cong\R^3\)
to the image of \(*\) is based and contracts that last map.
For completeness, the dependence on the framed inverse image in this
argument is precisely the local product-neighborhood construction and
the interpolation lemma of \cite[\S7, Theorem B and Lemma 4,
pp.~43, 48--50]{Milnor1965}. Their relative versions fix the source
basepoint because all tubular modifications are away from it and
stereographic interpolation between equal values keeps that value
fixed. A framed cylinder realizes a framing homotopy; the empty
Pontryagin manifold realizes a map missing the chosen regular value.
This proves the claimed based null-homotopy of the exact map \(g\).
\end{proof}

\begin{corollary}
Under the original fixed based identification \(X\cong S^6\), the
one-copy coefficient is
\[
 [g]=0\in\pi_6(S^3)\cong\Z/12,\qquad a=0\pmod{12},\qquad
 [E^6g]=0\in\pi_{12}(S^9)\cong\Z/24.
\]
The previously defined homomorphism
\(\Theta_g:\Z/12\to\pi_6(S^3)\), \(\Theta_g([n])=n[g]\), and
its composite with the marked peripheral character are zero.
\end{corollary}
\begin{proof}
The theorem states that the actual based map is null-homotopic.
Precomposition with the fixed based identification preserves a based
null-homotopy. Suspension preserves it as well. Thus the unique
coefficient of \([g]\) in any generator of \(\Z/12\) is zero.
For every \([n]\), \(\Theta_g([n])=n\cdot0=0\), and every
composite with this homomorphism is zero.
\end{proof}

\paragraph{Exact scope.}
The proof computes this sphere-valued postcomposition. It does not
compute the homotopy class of the earlier conical-target map
\(q_2:X\to\operatorname{Susp}(S^1\times\overline T^2)\), and it
does not assert a null-homotopy for that different target. Point-set
surjectivity of \(g\) and all of its stated metric restrictions remain
properties of the original map: a homotopy-class computation does not
alter the defining formula \eqref{gzero:eq:g}.

% BEGIN RETAINED_ANGLE_AND_HOPF_BRIDGE_20260905

\subsection{The retained-angle cusp quotient and its nonzero top-cell class}
\label{ret:section}

\paragraph{The exact object.}
The notation refers to the specified threefold \(X\), its fixed based
smooth identification with \(S^6\), and its cusp filling from
\cite[\S\S3.7,4.1--4.5]{retbib:Canonical}.  The native quotient and its compact
extension are the maps in \cite{retbib:ConeRoutes}; their full formulas are
recalled below.  This calculation uses the original period matrix and
native period lattices.  It computes a new postcomposition of that
particular quotient; it does not use the homotopy class of the
angle-forgetting map \(g:X\to S^3\).

\subsubsection{Original coordinates and the specified sphere-valued map}

The ordered quotient lattice is
\[
 \overline\Lambda=\mathbb Z\langle\overline{\widehat\gamma},
              \overline{\widehat u}\rangle,
 \qquad
 K_\Lambda=\mathbb Z\langle\widehat w,\widehat\delta\rangle,
 \qquad
 \overline T^2=\mathbb R^2/\mathbb Z^2.
\]
In these ordered bases the original period data are
\begin{equation}
 B_0=\begin{pmatrix}0&1\\-1&0\end{pmatrix},\qquad
 Z(s)=sB_0+C(t),\qquad t=e^{2\pi i s},\qquad
 C(t)=\begin{pmatrix}6\mu(t)&h_{\rm per}(t)\\
 b(t)-h_{\rm per}(t)&\mu(t)\end{pmatrix}.
 \label{ret:eq:original-periods}
\end{equation}
Here \(b=\beta+\tau\), \(s=\tau-h_{\rm per}(t)\), and every entry of
\(C\) is holomorphic at \(t=0\), exactly as in the original source.
The subscript in \(h_{\rm per}\) distinguishes that period function
from the sphere-valued map below.  Write
\begin{equation}
 R(t)=-2\pi\operatorname{Im}C(t),\qquad
 \ell=\log|t|,\qquad A(t)=\ell B_0+R(t).
 \label{ret:eq:A}
\end{equation}
Choose a closed cusp disc \(|t|\leq\epsilon\) inside the original
coordinate disc and small enough for the quotient construction.  On a
slightly larger closed disc put \(M=\sup\|R(t)\|<\infty\).  We will
choose a regular value with \(0<r_0<\epsilon\) and
\(|\log r_0|>M\).  This chooses a point of the already specified map;
it does not change that map or the cutoff radius \(\epsilon\).

The original covering model on its dense torus is
\begin{equation}
 x_j=e^{2\pi i\zeta_j},\qquad
 \Psi_n(x,t)=
 \bigl(e^{2\pi i C(t)n}t^{B_0n}x,t\bigr),\qquad n\in\mathbb Z^2.
 \label{ret:eq:deck}
\end{equation}
For a branch of \(s\), writing
\(\zeta=Z(s)u+k\) with \(u,k\in\mathbb R^2\) gives
\begin{equation}
 \log|x|=A(t)u.
 \label{ret:eq:logx}
\end{equation}
Indeed, \(\operatorname{Im}s=-\ell/(2\pi)\), so
\(-2\pi\operatorname{Im}(Z(s)u+k)=\ell B_0u+R(t)u\).
Since \(B_0\) is orthogonal,
\begin{equation}
 \bigl|\bigl(\ell B_0+vR(t)\bigr)\xi\bigr|
 \geq (|\ell|-vM)|\xi|>0
 \quad(0\leq v\leq1,\ \xi\ne0)
 \label{ret:eq:invertibility}
\end{equation}
whenever \(|\ell|>M\).  Thus \(A(t)\) is invertible there.  Equation
\eqref{ret:eq:deck} changes \(A(t)^{-1}\log|x|\) by \(n\), so the native
quotient is exactly
\begin{equation}
 \overline\pi_f([x,t])=[A(t)^{-1}\log|x|]\in\mathbb R^2/\mathbb Z^2.
 \label{ret:eq:native}
\end{equation}
This also follows from the flat marking in \cite{retbib:ConeRoutes}.

Let
\[
 T^3=S^1_\phi\times\overline T^2,
 \qquad S^1_\phi=\mathbb R/(2\pi\mathbb Z),
\]
with coordinate basepoints \(\phi=0\), \(u_1=0\), \(u_2=0\).
Its coordinate two-skeleton is the following specified closed subset:
\begin{equation}
 K^{(2)}=
 \bigl(\{0\}\times\overline T^2\bigr)
 \cup\bigl(S^1_\phi\times\{u_1=0\}\bigr)
 \cup\bigl(S^1_\phi\times\{u_2=0\}\bigr).
 \label{ret:eq:two-skeleton}
\end{equation}
The product CW structure has one open three-cell, with coordinates
\((\phi,u_1,u_2)\in(0,2\pi)\times(0,1)^2\).  Consequently
\[
 c_3:T^3\longrightarrow T^3/K^{(2)}\cong S^3
\]
is the degree-one quotient, where the orientation of \(S^3\) is fixed
by the displayed coordinate order.  One explicit identification of its
open cell with \(\mathbb R^3\) is
\[
 (\phi,u_1,u_2)\longmapsto
 \bigl(-\cot(\phi/2),-\cot(\pi u_1),-\cot(\pi u_2)\bigr).
\]
Each derivative in its own variable is strictly positive, and approach
to any boundary face leaves every compact subset of \(\mathbb R^3\).
Extending by its point at infinity gives the stated based sphere
identification and verifies its continuity on the collapsed skeleton.

The exact conical target and map are
\begin{align}
 \mathcal A&=
 \frac{[0,\epsilon]\times T^3}
 {\{0\}\times T^3\sim v_0,\quad
  \{\epsilon\}\times T^3\sim v_1},
 \label{ret:eq:A-target}\\
 q_2(x)&=[r,e^{i\phi},\overline\pi_f(x)]
       &&(t(f(x))=re^{i\phi},\ 0<r<\epsilon),\notag\\
 q_2(W)&=v_0,\qquad
 q_2\bigl(X\setminus\operatorname{int}N_0\bigr)=v_1.
 \label{ret:eq:q2}
\end{align}
Fix the same source basepoint \(x_{\rm base}\in W\).  Define
\begin{equation}
 p_3:\mathcal A\longrightarrow\Sigma_r(S^3,*),\qquad
 p_3([r,z])=[r,c_3(z)],\qquad
 H=p_3q_2:(X,x_{\rm base})\longrightarrow(S^4,*).
 \label{ret:eq:H}
\end{equation}
The reduced suspension here collapses both endpoint copies of \(S^3\)
and the entire interval over \(*\) to one point.  Thus \(p_3\) sends
both \(v_0\) and \(v_1\) to \(*\) and is well-defined and based.
It is continuous by the defining quotient topology.  An explicit
identification of its open four-cell with \(\mathbb R^4\) is
\begin{equation}
 (r,\phi,u_1,u_2)\longmapsto
 \bigl(-\cot(\pi r/\epsilon),-\cot(\phi/2),
       -\cot(\pi u_1),-\cot(\pi u_2)\bigr).
 \label{ret:eq:four-cell}
\end{equation}
This fixes the target orientation.  Replacing \(r\) locally by
\(\ell=\log r\) preserves this orientation, since \(dr/d\ell=r>0\).
The ordered coordinates in \eqref{ret:eq:four-cell} specify the sphere
identification used throughout this calculation.

\begin{theorem}
\label{ret:main}
For the exact map \eqref{ret:eq:H},
\begin{equation}
 [H]=\eta_4\circ\eta_5\ne0
       \quad\hbox{in}\quad\pi_6(S^4)\cong\mathbb Z/2,
 \label{ret:eq:main-class}
\end{equation}
where \(\eta_n:S^{n+1}\to S^n\) is the \((n-2)\)-fold reduced
suspension of the complex Hopf map.  In particular \(q_2\) is not
based null-homotopic.  Its stabilized top-cell postcomposition is the
element \(\eta^2\in\pi_2^{\rm S}\).
\end{theorem}

\subsubsection{The entire regular inverse image lies in one original chart}

Choose the actual lift
\begin{equation}
 u_0=(-1/4,1/4)^{\mathsf T},\qquad
 p_0=B_0u_0=(1/4,1/4)^{\mathsf T},\qquad \phi_0=\pi.
 \label{ret:eq:chosen-value}
\end{equation}
Modulo \(\mathbb Z^2\), \(u_0\) is \((3/4,1/4)\) in the open cell.  Put
\(t_0=r_0e^{i\phi_0}\) and
\(y_0=[r_0,c_3(e^{i\phi_0},[u_0])]\).  The numbers \(r_0\) and
\(t_0\) will remain nonzero throughout the framing calculation.

The original triangle \(\{0,e_1,e_2\}\) has toric chart
\begin{equation}
 \mathcal U_\sigma\cong\mathbb C^3,\qquad
 (z_0,z_1,z_2)=(t/(x_1x_2),x_1,x_2),\qquad
 t=z_0z_1z_2.
 \label{ret:eq:z-chart}
\end{equation}
These characters are dual to \((0,0,1),(1,0,1),(0,1,1)\), respectively,
as direct pairing verifies.  This is the chart of
\cite[Lemma 4.2(iii)]{retbib:Canonical}, with its original coordinates.

For \(j=1,2\), introduce the nowhere vanishing holomorphic functions
\begin{equation}
 h_j(t)=\exp\bigl(-2\pi i(C(t)u_0)_j\bigr).
 \label{ret:eq:units}
\end{equation}
The following is an exact coordinate map, with the exact inverse
displayed alongside it:
\begin{align}
 F(z)&=(w_0,w_1,w_2)
 =\bigl(z_0/(h_1(t)h_2(t)),\ h_1(t)z_1,\ h_2(t)z_2\bigr),
 \label{ret:eq:F}\\
 F^{-1}(w)&=
 \bigl(w_0h_1(t)h_2(t),\ w_1/h_1(t),\ w_2/h_2(t)\bigr),
 \quad t=w_0w_1w_2.
 \label{ret:eq:Finv}
\end{align}
Both maps are holomorphic wherever \(|t|\) is in the period disc,
both preserve \(t\), and substitution proves that both composites are
the identity.  In particular the maps extend across the origin and
across every coordinate divisor.  No period matrix is replaced by a
different matrix in this coordinate change.

\begin{lemma}
For sufficiently small \(r_0>0\), the whole inverse image
\(M=H^{-1}(y_0)\) is represented injectively inside one contractible
coordinate neighborhood in \(X\).  In its \(w\)-coordinates it is
exactly the embedded two-torus
\begin{equation}
 w(a,b)=
 \bigl(\rho_0e^{i(\phi_0-a-b)},\rho_1e^{ia},\rho_2e^{ib}\bigr),
 \qquad
 \rho_0=r_0^{1/2},\quad \rho_1=\rho_2=r_0^{1/4},
 \quad (a,b)\in\bigl(\mathbb R/(2\pi\mathbb Z)\bigr)^2.
 \label{ret:eq:M}
\end{equation}
\end{lemma}
\begin{proof}
The point \(y_0\) lies off the collapsed subset of the target, so
\eqref{ret:eq:H} and \eqref{ret:eq:q2} imply
\[
 M=\{x:t(f(x))=t_0,\ \overline\pi_f(x)=[u_0]\}.
\]
Any lift of such a point to \((\mathbb C^*)^2\times\{t_0\}\) has
\(A(t_0)^{-1}\log|x|=u_0+n\) for one \(n\in\mathbb Z^2\).
Applying the actual transformation \(\Psi_{-n}\) from
\eqref{ret:eq:deck} produces a lift with
\(\log|x|=A(t_0)u_0\).  If two such lifts represent the same point,
\(\Psi_n\) relates them and changes their lifted quotient coordinate
by \(n\); hence \(n=0\), and the lifts are identical.  This proves
both completeness and injectivity of this chosen representative.

Since
\(\log|h_j(t)|=2\pi\operatorname{Im}(C(t)u_0)_j
=-(R(t)u_0)_j\), its \(w\)-coordinates satisfy
\[
 (\log|w_1|,\log|w_2|)^{\mathsf T}
 =A(t_0)u_0-R(t_0)u_0=\ell_0p_0,
 \qquad \ell_0=\log r_0.
\]
Thus \(|w_1|=|w_2|=r_0^{1/4}\), and the unchanged product
\(w_0w_1w_2=t_0\) forces \(|w_0|=r_0^{1/2}\) and the displayed
phase.  Conversely every pair \((a,b)\) in \eqref{ret:eq:M} gives one
such lift.  The original kernel coordinates \(k_j\) had period one
in \(x_j=e^{2\pi i\zeta_j}\); the phase coordinates \(a,b\) have
period \(2\pi\), with no cover or change of lattice.

The covering \(Y_\epsilon\to N_0\) is injective on a neighborhood of
the origin in \eqref{ret:eq:z-chart}, because the original action is free
and properly discontinuous.  Choose a ball \(V\) centered at that
origin with closure inside such a neighborhood and the period disc.
Choose a ball \(V_w\) about the origin with closure contained in
\(F(V)\).  Every coordinate in \eqref{ret:eq:M} tends uniformly to zero
as \(r_0\downarrow0\).  Therefore for sufficiently small \(r_0\)
the entire torus lies in \(V_w\).  The inverse image of \(V_w\)
under \(F\), followed by the injective covering projection, is the
required contractible chart in \(X\).
\end{proof}

\subsubsection{The full differential and its ordered normal framing}

We calculate in a neighborhood of the compact torus \eqref{ret:eq:M}.
The quotient coordinates admit a unique local lift near \(u_0\).
Set
\[
 v=(\log|w_1|,\log|w_2|)^{\mathsf T},\quad
 \ell=\log|w_0w_1w_2|,\quad
 \phi=\arg(w_0w_1w_2)
\]
using the branch of \(\phi\) near \(\phi_0=\pi\).
Equations \eqref{ret:eq:logx} and \eqref{ret:eq:F} give the exact identity
\begin{equation}
 v=\ell B_0u+R(t)(u-u_0),\qquad
 u=u_0+A(t)^{-1}(v-\ell p_0).
 \label{ret:eq:exact-u}
\end{equation}
Differentiate this identity before restricting to the inverse image:
\begin{equation}
 du=-A(t)^{-1}(dA(t))A(t)^{-1}(v-\ell p_0)
        +A(t)^{-1}(dv-p_0\,d\ell).
 \label{ret:eq:full-derivative}
\end{equation}
On \(M\) the vector \(v-\ell p_0\) is exactly zero.  Therefore
\begin{equation}
 du|_M=A_0^{-1}(dv-p_0\,d\ell),\qquad A_0=A(t_0).
 \label{ret:eq:derivative-on-M}
\end{equation}
This explicitly accounts for all derivatives of the original bounded
period terms; none are omitted or estimated away.

Identify \(T_w\mathbb C^3\) with \(\mathbb C^3\) in the given coordinate chart,
and define real tangent vectors
\begin{equation}
 L_j=(0,\ldots,w_j,\ldots,0),\qquad
 Q_j=(0,\ldots,iw_j,\ldots,0)\qquad(0\leq j\leq2).
 \label{ret:eq:LQ}
\end{equation}
Thus \(d\log|w_i|(L_j)=\delta_{ij}\),
\(d\arg w_i(Q_j)=\delta_{ij}\), and the crossed evaluations vanish.
The ordered tangent vectors of the parametrization \eqref{ret:eq:M} are
\begin{equation}
 T_a=-Q_0+Q_1,\qquad T_b=-Q_0+Q_2.
 \label{ret:eq:tangents}
\end{equation}
They are independent because their \(w_1\) and \(w_2\) components,
respectively, are nonzero.  For the target order
\((\ell,\phi,u_1,u_2)\), define the following four lifts of its
coordinate vectors:
\begin{align}
 N_\ell&=\tfrac12L_0+\tfrac14L_1+\tfrac14L_2,
       &N_\phi&=Q_0,\label{ret:eq:normal-first}\\
 N_i&=-\bigl((A_0)_{1i}+(A_0)_{2i}\bigr)L_0
           +(A_0)_{1i}L_1+(A_0)_{2i}L_2
       &&(i=1,2).\label{ret:eq:normal-u}
\end{align}
Let \(P=(\ell,\phi,u_1,u_2)\).  Evaluation using
\eqref{ret:eq:derivative-on-M} gives the entire differential table
\begin{equation}
\begin{array}{c|rrrrrr}
 &N_\ell&N_\phi&N_1&N_2&T_a&T_b\\\hline
 d\ell&1&0&0&0&0&0\\
 d\phi&0&1&0&0&0&0\\
 du_1&0&0&1&0&0&0\\
 du_2&0&0&0&1&0&0
\end{array}.
 \label{ret:eq:differential-table}
\end{equation}
In particular \(P\), and therefore \(H\) in its open target cell,
is a submersion along \(M\).  The classes of
\((N_\ell,N_\phi,N_1,N_2)\) in \(TX|_M/TM\) give the specified
normal four-framing.  Orthogonal representatives are not required for
this quotient-bundle definition; projecting the displayed lifts along
\(TM\) to any chosen orthogonal complement preserves the quotient
frame.

If the normal frame is specified using the literal four real coordinates
in \eqref{ret:eq:four-cell}, their differential at the chosen value is
\[
 dH_{\rm cell}=K\,dP,\qquad
 K=\operatorname{diag}\!\left(
 \frac{\pi r_0}{\epsilon}\csc^2\frac{\pi r_0}{\epsilon},
 \frac12,2\pi,2\pi\right).
\]
The first factor includes \(dr=r_0\,d\ell\), the second uses
\(\phi_0=\pi\), and the final two use the actual representatives
\((3/4,1/4)\) of \(u_0\). Thus, writing
\(\mathcal N=(N_\ell,N_\phi,N_1,N_2)\), the literal dual frame is
\(\mathcal N K^{-1}\). All four entries of \(K\) are positive
constants on \(M\). The explicitly positive diagonal matrices
\((1-v)K^{-1}+vI_4\), \(0\leq v\leq1\), give a homotopy
from that literal frame to \(\mathcal N\). This retains the original
target coordinates and establishes that either coordinate description
gives the same framed bordism class.

For completeness, the original period matrix admits the following
explicit homotopy of this normal framing.  Put
\begin{equation}
 A_{0,v}=\ell_0B_0+vR(t_0),\qquad 0\leq v\leq1,
 \label{ret:eq:A-homotopy}
\end{equation}
and replace \(A_0\) by \(A_{0,v}\) only in
\eqref{ret:eq:normal-u}.  Equation \eqref{ret:eq:invertibility} proves
invertibility for every \(v\).  Their images under the actual
\((du_1,du_2)\) are the columns of \(A_0^{-1}A_{0,v}\), so these
are normal frames throughout the homotopy.  This is a homotopy of
the computed frame on the same embedded torus, not a replacement of
the original degeneration by a different one.

\begin{lemma}[The exact orientation sign]
Use the complex real-coordinate orientation, whose ordered coordinates are
\[
 (\operatorname{Re}w_0,\operatorname{Im}w_0,
  \operatorname{Re}w_1,\operatorname{Im}w_1,
  \operatorname{Re}w_2,\operatorname{Im}w_2).
\]
Then one has
\begin{equation}
 \det(N_\ell,N_\phi,N_1,N_2,T_a,T_b)
       =-\det(A_0)|w_0w_1w_2|^2<0.
 \label{ret:eq:sign}
\end{equation}
Hence the tangent orientation induced by the ordered normal frame is
the orientation of \((-T_a,T_b)\).
\end{lemma}
\begin{proof}
The ordered frame \((L_0,Q_0,L_1,Q_1,L_2,Q_2)\) has determinant
\(|w_0|^2|w_1|^2|w_2|^2>0\) in real coordinates.  Regroup its rows
as \((L_0,L_1,L_2,Q_0,Q_1,Q_2)\); this permutation has three
inversions and sign \(-1\).  Regroup the columns of the frame in
\eqref{ret:eq:sign} as
\((N_\ell,N_1,N_2,N_\phi,T_a,T_b)\); this has two inversions and
sign \(+1\).  The angular block is
\(\left(\begin{smallmatrix}1&-1&-1\\0&1&0\\0&0&1\end{smallmatrix}\right)\),
of determinant one.  The radial block has first column
\((1/2,1/4,1/4)^{\mathsf T}\) and the other columns given by
\eqref{ret:eq:normal-u}.  Adding its last two rows to its first row gives
first row \((1,0,0)\), with remaining lower-right block \(A_0\).
Its determinant is \(\det A_0\).  This proves the equality.
Equation \eqref{ret:eq:A-homotopy} joins \(A_0\) to \(\ell_0B_0\)
through invertible matrices, so
\(\operatorname{sign}\det A_0=\operatorname{sign}\det(\ell_0B_0)=+1\).
The inequality and the tangent orientation follow.
\end{proof}

\subsubsection{An explicit contraction of the full frame}

Define a complex-linear matrix, also regarded as a real-linear matrix,
\begin{equation}
 D(a,b)=\operatorname{diag}(e^{-i(a+b)},e^{ia},e^{ib})\in\operatorname{SU}(3).
 \label{ret:eq:D}
\end{equation}
By \eqref{ret:eq:M}, every \(L_j\) and \(Q_j\) at \((a,b)\) is
\(D(a,b)\) times its value at \((0,0)\).  The coefficients in
\eqref{ret:eq:normal-first}, \eqref{ret:eq:normal-u}, and
\eqref{ret:eq:tangents} are constant on \(M\), since \(t=t_0\).
It follows that the positively oriented full frame is exactly
\begin{equation}
 (N_\ell,N_\phi,N_1,N_2,-T_a,T_b)_{(a,b)}
       =D(a,b)G_0,
 \label{ret:eq:full-frame}
\end{equation}
where \(G_0\in\operatorname{GL}^+(6,\mathbb R)\) is its single value at \((0,0)\).
Thus neither an unspecified phase nor a position-dependent
real-linear matrix remains.

Here is a periodic contraction of \(D\), with all parameters
specified.  For \(v\in[0,1]\) put
\[
 s_v=\sin(\pi v/2),\qquad c_v=\cos(\pi v/2),\qquad
 Q_v(\theta)=
 \begin{pmatrix}s_ve^{-i\theta}&-c_v\\c_v&s_ve^{i\theta}\end{pmatrix},
 \qquad K_v(\theta)=Q_v(\theta)Q_v(0)^*.
\]
The two columns of \(Q_v\) have norm one and zero Hermitian inner
product, and \(\det Q_v=s_v^2+c_v^2=1\).  Thus
\(Q_v,K_v\in\operatorname{SU}(2)\).  Direct substitution gives
\[
 K_0(\theta)=I,\quad
 K_1(\theta)=\operatorname{diag}(e^{-i\theta},e^{i\theta}),\quad
 K_v(0)=I.
\]
Let \(\iota_{01},\iota_{02}:\operatorname{SU}(2)\to\operatorname{SU}(3)\) act in the
ordered coordinate planes \((0,1)\) and \((0,2)\), respectively,
and as the identity in the unused coordinate.  Then
\begin{equation}
 D_v(a,b)=\iota_{01}(K_v(a))\,\iota_{02}(K_v(b))
 \label{ret:eq:D-homotopy}
\end{equation}
is a well-defined continuous map on
\([0,1]\times(\mathbb R/(2\pi\mathbb Z))^2\), with
\(D_0=I\) and \(D_1=D\).  This proves the needed null-homotopy
directly, including both phase seams.  In particular
\(D_v(a,b)G_0\) contracts the comparison matrix in
\eqref{ret:eq:full-frame} to one constant positive matrix.

\begin{proposition}
The stable tangential framing induced by the actual normal four-frame
of \(H^{-1}(y_0)\) agrees, up to homotopy and the orientation sign
\eqref{ret:eq:sign}, with the Lie framing
\((\partial_a,\partial_b)\) on the native two-torus.
\end{proposition}
\begin{proof}
In the coordinate chart, the stable tangent trivialization is the
isomorphism
\[
 \mathbb R^4\oplus TM\longrightarrow T\mathbb C^3|_M,
 \qquad (v,V)\longmapsto
 v_1N_\ell+v_2N_\phi+v_3N_1+v_4N_2+V.
\]
Using the intrinsic frame \((-\partial_a,\partial_b)\) on \(TM\),
its matrix is precisely \eqref{ret:eq:full-frame}.  The homotopy
\eqref{ret:eq:D-homotopy} therefore makes this comparison constant.
A constant positive matrix changes no homotopy class of stable
framing: Gram--Schmidt deforms it to an element of \(SO(6)\), and
an orthogonal matrix of determinant one can be reduced to the identity
by finitely many oriented plane rotations, each joined to the
identity by varying its angle.  This proves the asserted comparison
within the chart.

It remains to compare the chart with the ambient sphere framing.
Pull back the usual stable tangent trivialization
\(TS^6\oplus\mathbb R\cong\mathbb R^7\) under the fixed based diffeomorphism
\(X\cong S^6\).  On the contractible \(w\)-ball containing \(M\),
it differs from the coordinate trivialization plus one real line by
a continuous map into \(\operatorname{GL}(7,\mathbb R)\).  Contracting the ball to its
center contracts this comparison map to a constant.  Thus its
restriction to \(M\) contributes no varying stable framing class.
If the chosen sphere orientation disagrees with the complex
orientation, the remaining constant has negative determinant and
reverses the resulting framed bordism class.  Both that sign and the
explicit sign \eqref{ret:eq:sign} disappear in the order-two class
computed below.  This argument compares actual ambient
trivializations on a contractible chart; it does not suppose that
\(\mathbb C^3\) globally trivializes \(TX\).
\end{proof}

\subsubsection{Nonzero framed bordism and the exact unstable class}

We give the nonvanishing test for the resulting torus explicitly.
A stable framing gives a spin structure by lifting its oriented frame
to the trivial \(\operatorname{Spin}\)-bundle.  On the Lie-framed
circle this is the disconnected double cover, whereas the spin
structure induced from the boundary of a framed disc is the
connected double cover: as the boundary tangent turns once, its
\(2\pi\)-rotation lifts from \(1\) to \(-1\) in
\(\operatorname{Spin}(2)\).  We call the former nonbounding and the
latter bounding.

For a closed oriented spin surface \(S\), the spin quadratic
refinement
\[
 q_S:H_1(S;\mathbb F_2)\longrightarrow\mathbb F_2
\]
takes a class represented by an embedded circle to zero or one
according to whether the induced spin structure, using its normal
line in the surface, is bounding or nonbounding.  It satisfies
\begin{equation}
 q_S(x+y)=q_S(x)+q_S(y)+(x\mathbin{\cdot}y).
 \label{ret:eq:quadratic}
\end{equation}
This is the usual spin quadratic form: resolving transverse
intersections adds one tangent turn per intersection, and the
\(2\pi\)-spin-lift calculation above gives the last term.  The
well-defined correspondence between spin structures and such
quadratic refinements is given in
\cite[Theorem 3.2 and Remark 3.3]{retbib:DebrayGunningham}; its original
surface proof is \cite{retbib:Johnson}.

For a symplectic basis \((e_i,f_i)_{i=1}^g\), put
\begin{equation}
 \operatorname{Arf}(S)=\sum_{i=1}^gq_S(e_i)q_S(f_i)\in\mathbb F_2.
 \label{ret:eq:Arf}
\end{equation}
The basis independence can be seen without selecting a preferred
basis: factor the finite sum over the symplectic planes to obtain
\[
 \sum_{x\in H_1(S;\mathbb F_2)}(-1)^{q_S(x)}
   =\prod_{i=1}^g
     \sum_{a,b\in\mathbb F_2}(-1)^{a q_S(e_i)+b q_S(f_i)+ab}
   =2^g(-1)^{\sum_i q_S(e_i)q_S(f_i)}.
\]
The left side has no choice of basis, proving the assertion.

Here is also the bordism-invariance argument needed for the present
obstruction.  A spin bordism has a handle decomposition.  An index
zero or three handle creates or removes a spin sphere, whose first
homology and Arf invariant are zero.  An index one handle either
connects two components, in which case the first-homology quadratic
forms add orthogonally, or increases the genus of one component.
In the latter case the new symplectic pair contains the meridian
of the handle.  That meridian bounds the cocore disc with its
induced spin structure, so its quadratic value is zero; the pair
contributes zero to \eqref{ret:eq:Arf}, regardless of the value on its
dual.  Existing cycles can be represented outside the attaching
discs, so their quadratic values are preserved.  An index two
handle is the reverse of an index one handle.  Therefore each
handle preserves the total Arf invariant of the boundary, and a
spin boundary has Arf invariant zero.  This also proves that a
stably framed boundary has Arf invariant zero, because forgetting
the framing gives a spin bordism.  The full spin bordism statement
is recorded in \cite[Theorem 3.6]{retbib:DebrayGunningham}.

For the intrinsic Lie framing on
\(M=(\mathbb R/(2\pi\mathbb Z))^2\), take
\[
 e=[a\mapsto(a,0)],\qquad f=[b\mapsto(0,b)].
\]
Both the tangent and normal vectors in the surface are constant in
the Lie frame along either circle.  Thus both circles have the
nonbounding spin structure, and
\begin{equation}
 e\mathbin{\cdot}f=1,\qquad q_M(e)=q_M(f)=1,
 \qquad q_M(xe+yf)=x+xy+y,
 \qquad \operatorname{Arf}(M)=1.
 \label{ret:eq:actual-Arf}
\end{equation}
Reversing the first Lie tangent vector changes the orientation but
does not change these mod-two values.  In particular the
specifically framed inverse image computed above is not a framed
boundary.

To identify the element, let \(\eta\in\pi_1^{\rm S}\) denote the
stable class of the complex Hopf map.  Its Pontryagin--Thom inverse
image is the circle with its Lie framing.  This can be verified
directly on the Hopf map \(S^3\to\mathbb{CP}^1\): near the fibre
\((z_1,z_2)=(e^{ia},0)\), the target coordinate is \(z_2/z_1\), so
the two horizontal normal vectors are \((0,e^{ia})\) and
\((0,ie^{ia})\).  Add the tangent vector \((ie^{ia},0)\) and the
outward radial vector \((e^{ia},0)\).  In \(\mathbb C^2\) the resulting
full frame makes one complete rotation in each of its two real
coordinate planes.  This loop is contractible explicitly: for the
fixed real-linear involution \(J(z_1,z_2)=(z_1,\overline z_2)\),
\[
 J\operatorname{diag}(e^{ia},e^{ia})J=\operatorname{diag}(e^{ia},e^{-ia})=K_1(-a),
\]
and \(JK_v(-a)J\) contracts the original loop to the identity using
the periodic homotopy already constructed.  Thus the induced stable
tangent frame on the circle is the Lie frame.  This is the basic
framed-circle identification in \cite{retbib:Pontryagin1955}.

The product of two normally framed inverse images represents the
product in stable homotopy: its tubular collapse is the smash of
the two tubular collapse maps, as follows immediately from the
product tubular neighborhood \(\nu_1\times\nu_2\) and its ordered
normal coordinates.  Consequently the product of the two Lie
circles represents \(\eta^2\).  Equation \eqref{ret:eq:actual-Arf}
proves \(\eta^2\ne0\) by an explicit bordism invariant.  The
stable group \(\pi_2^{\rm S}=\mathbb Z/2\) is Pontryagin's computation
\cite{retbib:Pontryagin1950,retbib:Pontryagin1955}; the stable table in
\cite[pp.~384--385]{retbib:Hatcher} gives the same group.  Thus the sign from
\eqref{ret:eq:sign} satisfies \(-\eta^2=\eta^2\).

\begin{proof}[Proof of the theorem]
The globally continuous quotient \(H\) is smooth on a neighborhood
of its compact inverse image \(M\), by
\eqref{ret:eq:exact-u} and \eqref{ret:eq:four-cell}. First make it constant
near the source basepoint, before choosing any approximation bound.
Choose a target chart \(\kappa:O\to\mathbb R^4\) about \(*\) with
\(\kappa(*)=0\), convex ball image, and \(y_0\notin O\).
Choose a source neighborhood \(U\) of \(x_{\rm base}\), disjoint
from \(M\), with \(H(U)\subset O\). Let \(\beta\) be a smooth
function supported in \(U\), equal to one near \(x_{\rm base}\),
with \(0\leq\beta\leq1\). On \(U\) use the based homotopy
\[
 H_s(x)=\kappa^{-1}\bigl((1-s\beta(x))\kappa(H(x))\bigr),
 \qquad 0\leq s\leq1,
\]
and leave \(H\) unchanged off \(U\). Its endpoint \(H_b\) is
smooth and constant near the basepoint; it agrees with \(H\)
near \(M\) and has no new inverse image of \(y_0\).

Choose one closed tubular neighborhood \(U_1\) of \(M\), with
\(M\subset\operatorname{int}U_1\), contained in the unchanged
smooth region and disjoint from the basepoint neighborhood.
Compactness gives the positive bound
\[
 d_0=\min_{x\in X\setminus\operatorname{int}U_1}
             d\bigl(H_b(x),y_0\bigr)>0.
\]
Relative smooth approximation, fixing the same closed set \(U_1\)
and a neighborhood of the basepoint, produces a smooth based map
\(H_{\rm sm}\) whose error from \(H_b\) is less than \(d_0/2\)
on its complement. Choose the approximation closer also than a
fixed target injectivity radius; the short target geodesics give a
based homotopy. Thus the exact inverse image and its differential
are unchanged on \(U_1\), and the complement acquires no new
inverse image. The relative Pontryagin construction applies as in
\cite[\S7, Theorem B and Lemma 4, pp.~43, 48--50]{retbib:Milnor1965}.

Pontryagin--Thom therefore identifies the stabilization of
\([H]\) with the bordism class of \(M\) carrying the four normal
vectors \eqref{ret:eq:normal-first}--\eqref{ret:eq:normal-u}.  The comparison
\eqref{ret:eq:full-frame} and its explicit contraction identify that
class with the Lie torus, with the harmless order-two orientation
sign already computed.  Hence
\[
 E^\infty[H]=\eta^2\ne0.
\]
Freudenthal suspension is an isomorphism
\[
 E:\pi_{n+2}(S^n)\longrightarrow\pi_{n+3}(S^{n+1})
 \qquad(n\geq4),
\]
because \(n+2\leq2n-2\).  Therefore stabilization from
\(\pi_6(S^4)\) is an isomorphism onto \(\pi_2^{\rm S}\).
The composite \(\eta_4\circ\eta_5:S^6\to S^4\) stabilizes to
\(\eta^2\), so injectivity gives the exact equality
\eqref{ret:eq:main-class}.  Finally, a based null-homotopy of \(q_2\)
would postcompose with \(p_3\) to a based null-homotopy of \(H\),
contradicting \eqref{ret:eq:main-class}.
\end{proof}

\paragraph{Exact scope of the comparison with the angle-forgetting map.}
The map here retains the angular coordinate before collapsing the
coordinate two-skeleton.  Its regular inverse image is the specific
normally four-framed \(T^2\) above.  The angle-forgetting
postcomposition has a normally three-framed \(T^3\) inverse image.
The calculation above is independent of its class and supplies an
actual obstruction to null-homotopy of the original \(q_2\), even
when that other postcomposition is null-homotopic.  The exact
morphism carrying this obstruction is
\[
 (p_3)_*:[X,\mathcal A]_*\longrightarrow[X,S^4]_*,\qquad
 [q_2]\longmapsto[H]=\eta_4\eta_5.
\]
No conclusion about an arbitrary earlier quotient follows merely
from the diffeomorphism type of one inverse image; the specified
four-framing is what proves this value.

\begin{corollary}[The exact order-two quotient detected by the original cusp map]
\label{ret:q2-detected-quotient}
Let \(d:(S^6,*)\to(X,x_{\mathrm{base}})\) be the fixed based
diffeomorphism, and put
\[
\alpha=[q_2\circ d]\in\pi_6(\mathcal A,v_0).
\]
The original top-cell postcomposition induces a surjective homomorphism
\[
\varepsilon=(p_3)_*:\pi_6(\mathcal A,v_0)
  \longrightarrow\pi_6(S^4,*)\cong\mathbb Z/2,
\qquad \varepsilon(\alpha)=\eta_4\circ\eta_5.
\]
It gives the exact quotient isomorphism
\[
\pi_6(\mathcal A,v_0)/\ker\varepsilon
 \xrightarrow{\ [\beta]\mapsto\varepsilon(\beta)\ }
\mathbb Z/2.
\]
Every odd multiple of \(\alpha\) is nonzero. Its order is either
infinite or an even positive integer; this calculation does not determine
that order or assert a splitting of the displayed quotient.
\end{corollary}
\begin{proof}
Postcomposition distributes over the pinch sum on the source sphere,
so \((p_3)_*\) is a group homomorphism. Theorem~\ref{ret:main}
gives \((p_3)_*(\alpha)=[p_3q_2d]=\eta_4\eta_5\ne0\).
This is the nonzero element of the cyclic group of order two, proving
surjectivity. The first isomorphism theorem gives the stated quotient
with exactly the displayed map. For every integer \(n\),
\[
\varepsilon(n\alpha)=(n\bmod2)\eta_4\eta_5.
\]
An odd multiple therefore has nonzero image and cannot vanish.
If \(m\alpha=0\), applying \(\varepsilon\) forces \(m\)
to be even. These arguments prove all the assertions without inferring
the exact order of \(\alpha\).
\end{proof}

% Staged proof only. No canonical source has been edited by this worker.
% Uses the already proved retained-angle theorem ret:main.

\subsection{The retained cusp class as a three-dimensional bundle inside the triality carrier}
\label{rettri:section}

Retain the exact based map $H=p_3q_2:X\to S^4$, the fixed based
diffeomorphism $d:S^6\to X$, the original cutoff $\epsilon$, the ordered
native coordinates $(r,\phi,u_1,u_2)$, and every period appearing in
\eqref{ret:eq:original-periods}. Theorem~\ref{ret:main} has proved
$[H d]=\eta_4\eta_5\ne0$ in $\pi_6(S^4)\cong\mathbb Z/2$.
We construct an exact bundle-valued image of this specific map and then
its action on the original three Peirce sectors. A quaternion subalgebra
and its ordered basis are explicit choices in this construction.

\paragraph{The literal Hopf pullback and its transition.}
Fix $\mathbb H=\mathbb R\langle1,i,j,k\rangle$ with
$i^2=j^2=k^2=ijk=-1$. Identify the ordered open four-cell of the target
with $\mathbb H$ by
\begin{equation}
 z=-\cot(\pi r/\epsilon)-\cot(\phi/2)i
       -\cot(\pi u_1)j-\cot(\pi u_2)k.
 \label{rettri:z}
\end{equation}
This is exactly the ordered coordinate map
\eqref{ret:eq:four-cell}, followed by the real linear map carrying
the four standard basis vectors to $(1,i,j,k)$. Each of the four
coordinate derivatives in the original open cell is positive.
The based homeomorphism
\begin{equation}
 \chi:S^4=\mathbb H\cup\{\infty\}\longrightarrow\mathbb H\mathbb P^1,
 \qquad \chi(z)=[1:z],\qquad \chi(*)=[0:1]
 \label{rettri:chi}
\end{equation}
uses right quaternionic lines. The orientation on this presentation of
$\mathbb H\mathbb P^1$ is the one induced by the ordered affine basis
$(1,i,j,k)$. Thus no orientation convention on another presentation of
$\mathbb H\mathbb P^1$ is silently substituted.

Let
\[
 S^7=\{(v_0,v_1)\in\mathbb H^2:|v_0|^2+|v_1|^2=1\},\qquad
 \pi(v_0,v_1)=[v_0:v_1].
\]
The right action is $(v_0,v_1)u=(v_0u,v_1u)$ for
$u\in\operatorname{Sp}(1)=\{u\in\mathbb H:|u|=1\}$. It is free:
at least one $v_j$ is nonzero, and $v_ju=v_j$ implies $u=1$.
Two unit vectors on the same right line differ by a unique unit
quaternion. Hence $\pi$ is the quaternionic Hopf principal bundle.
Its local sections in the affine coordinates $z=v_1v_0^{-1}$ and
$w=v_0v_1^{-1}$ are
\begin{equation}
 s_0(z)=\frac{(1,z)}{\sqrt{1+|z|^2}},\qquad
 s_\infty(w)=\frac{(w,1)}{\sqrt{1+|w|^2}}.
 \label{rettri:sections}
\end{equation}
On $z\ne0$, the coordinate change is $w=z^{-1}$ and
\begin{equation}
 s_\infty(z^{-1})=s_0(z)\,g(z),\qquad
 g(z)=\frac{\overline z}{|z|}.
 \label{rettri:transition}
\end{equation}
Indeed, multiplying the right side gives
$(\overline z/|z|,|z|)/\sqrt{1+|z|^2}$, and the same expression
results from the left side because
$\sqrt{1+|z^{-1}|^2}=\sqrt{1+|z|^2}/|z|$.
This calculation retains both the inverse and the conjugation.

Define
\begin{equation}
 P_H=\{(x,v)\in X\times S^7:\pi(v)=\chi(H(x))\}
       \longrightarrow X,
 \qquad (x,v)u=(x,vu).
 \label{rettri:P}
\end{equation}
The preimages under $\chi H$ of the two affine charts cover $X$.
Equations~\eqref{rettri:sections}--\eqref{rettri:transition}
are local sections and transitions for \eqref{rettri:P} after
composition with $\chi H$. In particular, wherever both charts apply,
the transition is the exact $g(z)$ from \eqref{rettri:transition},
with $z$ given by the original four-coordinate expression
\eqref{rettri:z}. The angle $\phi$ has not been dropped.

\begin{theorem}
\label{rettri:bundleclass}
The principal bundle $P_H$ is nontrivial. Its clutching class is the
nonzero element
\begin{equation}
 c(P_H)=\partial[\chi H d]\in\pi_5(\operatorname{Sp}(1))
                  \cong\mathbb Z/2,
 \label{rettri:clutch}
\end{equation}
where $\partial$ is the boundary of the specified right Hopf fibration.
The associated real three-plane bundle
\begin{equation}
 \mathcal A_H=P_H\times_{\operatorname{Ad}}\operatorname{Im}\mathbb H,
 \qquad \operatorname{Ad}_u(v)=uv\overline u,
 \label{rettri:rankthree}
\end{equation}
is nontrivial, has nonzero order-two clutching class in
$\pi_5(SO(3))$, and admits no nowhere-zero continuous section.
Its oriented orthonormal frame bundle has $P_H$ as its unique spin lift
up to isomorphism over $X$. The maps from the retained homotopy group
to these two clutching groups are isomorphisms:
\begin{equation}
 \pi_6(S^4)\xrightarrow{\chi_*}\pi_6(\mathbb H\mathbb P^1)
 \xrightarrow{\partial}\pi_5(\operatorname{Sp}(1))
 \xrightarrow{\operatorname{Ad}_*}\pi_5(SO(3)).
 \label{rettri:exact-detection}
\end{equation}
\end{theorem}

\begin{proof}
The first map is an isomorphism because $\chi$ is a based
homeomorphism. In the long exact sequence of the displayed Hopf
bundle, the segment is
\[
 \pi_6(S^7)\longrightarrow\pi_6(\mathbb H\mathbb P^1)
 \xrightarrow{\partial}\pi_5(S^3)\longrightarrow\pi_5(S^7).
\]
Both outside groups vanish: cellular approximation homotopes a map
from a sphere of dimension less than seven into the zero-cell of the
CW structure on $S^7$. Exactness therefore makes $\partial$ an
isomorphism. To fix its clutching convention, write the oriented
source $S^6=D^6_+\cup D^6_-$, give the equator the boundary orientation
of $D^6_+$, and choose local lifts $\sigma_\pm$ of the pulled-back
bundle. On the equator define $g_{+-}$ by
$\sigma_- = \sigma_+ g_{+-}$. Choose the values of the two lifts
equal at its fixed basepoint. We use the convention
$\partial[\chi H d]=[g_{+-}]$. Changing disk lifts changes this map
by null-homotopic factors and does not change its homotopy class.
This is the defining boundary construction of the fibration exact
sequence. The already proved nonzero class $[H d]$ therefore gives
the nonzero order-two class in \eqref{rettri:clutch}.

One can also see nontriviality directly. A section of $P_H$ gives a
map $\widetilde h:S^6\to S^7$ with
$\pi\widetilde h=\chi H d$. Since $\pi_6(S^7)=0$, this would make
$\chi H d$ null-homotopic, contrary to Theorem~\ref{ret:main}.
All spaces at issue are simply connected, so based and free
null-homotopy agree here.

Conjugation by a unit quaternion preserves $\operatorname{Im}\mathbb H$
and its Euclidean norm. Its determinant is $+1$ by continuity in
$u$ and its value at $u=1$. The kernel consists of the real unit
quaternions $\{1,-1\}$: a quaternion commuting with $i,j,k$ is real.
For a unit imaginary quaternion $n$ and $v\in\operatorname{Im}\mathbb H$,
write $v=v_\parallel+v_\perp$ relative to the line $\mathbb R n$.
Direct multiplication with $u=\cos\alpha+n\sin\alpha$ gives
\begin{equation}
 u v\overline u=v_\parallel+\cos(2\alpha)v_\perp
                     +\sin(2\alpha)(n\times v_\perp),
 \label{rettri:Rodrigues}
\end{equation}
where $n\times v=(nv-vn)/2$. Every orientation-preserving
orthogonal map of a real three-space fixes an axis and rotates its
orthogonal plane: its real odd-degree characteristic polynomial has
a real eigenvalue, norm preservation gives $\pm1$, and determinant
$+1$ forces an eigenvalue $+1$. Thus
\eqref{rettri:Rodrigues} proves that
$\operatorname{Ad}:S^3\to SO(3)$ is onto with kernel $\{\pm1\}$.
It is the two-sheeted covering and induces an isomorphism on
$\pi_5$. Its value on \eqref{rettri:clutch} is the clutching
class of \eqref{rettri:rankthree}, proving nontriviality and
order two.

The bundle $P_H\to P_H/\{\pm1\}$ is the spin double covering of
the oriented orthonormal frames of $\mathcal A_H$. If that frame
bundle had a section, its inverse image would be a double covering
of the simply connected space $X\cong S^6$. This covering is
trivial, so one of its two sections would give a section of $P_H$,
again a contradiction. Any two lifts of an $SO(3)$ bundle differ
by a principal $\mathbb Z/2$ bundle: locally the ratio between two
lifts of each transition lies in the central kernel $\{\pm1\}$,
and the cocycle identity gives a $\mathbb Z/2$ cocycle. Every such
bundle on the simply connected $X$ is trivial. This proves the
claimed uniqueness of the spin lift.

If $\mathcal A_H$ had a nowhere-zero section, divide that section
by its positive length to obtain a unit section. Orthogonal
complement gives $\mathcal A_H\cong\mathbb R\oplus\mathcal B$
for an oriented two-plane bundle $\mathcal B$. A two-plane
bundle on $S^6$ has clutching map $S^5\to SO(2)=S^1$.
Every such map lifts to $\mathbb R$ and is null-homotopic, so
$\mathcal B$ is trivial. This would trivialize $\mathcal A_H$,
contradicting its nonzero clutching class.
\end{proof}

The fibrewise scalar product, orientation and cross product of
$\operatorname{Im}\mathbb H$ descend to $\mathcal A_H$.
For example,
\[
 (uv\overline u)(uw\overline u)=u(vw)\overline u
\]
by associativity in $\mathbb H$, so conjugation preserves
$v\times w=(vw-wv)/2$. The induced fibrewise Lie bracket is
exactly $[v,w]=2(v\times w)$. No factor of two is discarded.

\paragraph{An explicit subgroup of the three-sector triality group.}
Fix the quaternion subalgebra above inside
$\mathbb O=\mathbb H\oplus\mathbb H\ell$. Its ordered real basis is
\[
 (1,i,j,k,\ell,i\ell,j\ell,k\ell),
\]
and its multiplication is
\begin{equation}
 (a,b)(c,d)=(ac-\overline d b,\;da+b\overline c),\qquad
 K(a,b)=\overline{(a,b)}=(\overline a,-b).
 \label{rettri:CD}
\end{equation}
This is the literal convention of Yokota, \S1.10. In his cyclic
related-triple presentation the group is
\begin{equation}
 \operatorname{RT}(\mathbb O)=
 \{(T_1,T_2,T_3)\in SO(\mathbb O)^3:
       (T_1x)(T_2y)=K T_3 K(xy)\text{ for all }x,y\in\mathbb O\}.
 \label{rettri:RT}
\end{equation}
Yokota, Theorem 1.16.2, identifies this group with
$\operatorname{Spin}(8)$; Theorem 2.7.1 \cite{Yokota2009} gives its diagonal-fixing
Jordan action. In particular the three projection representations
in \eqref{rettri:RT} are the vector and two half-spin sectors
with that fixed marking.

For each $u\in\operatorname{Sp}(1)$ define
\begin{equation}
 \begin{aligned}
 B_u(a,b)&=(u a\overline u,b),\\
 A_u(a,b)&=(u a,b\overline u),\\
 D_u(a,b)&=(a\overline u,b\overline u)=K A_u K(a,b).
 \end{aligned}
 \label{rettri:threeactions}
\end{equation}

\begin{theorem}
\label{rettri:subgroup}
The formula
\begin{equation}
 \rho:\operatorname{Sp}(1)\longrightarrow\operatorname{Spin}(8)
      =\operatorname{RT}(\mathbb O),\qquad
 \rho(u)=(B_u,A_u,D_u)
 \label{rettri:rho}
\end{equation}
is an injective continuous group homomorphism. It satisfies
\begin{equation}
 (B_u x)(A_u y)=A_u(xy),\qquad
 \operatorname{Re}(((B_u x)(A_u y))(D_u z))
           =\operatorname{Re}((xy)z)
 \label{rettri:product}
\end{equation}
for all octonions $x,y,z$ and every unit quaternion $u$.
\end{theorem}

\begin{proof}
All three maps are real linear norm isometries by the quaternionic
norm identity. Their determinants are $+1$: they vary continuously
with $u$ on the connected $S^3$ and equal the identity at $u=1$.
For $u,v\in\operatorname{Sp}(1)$, associativity in $\mathbb H$
and $\overline{uv}=\overline v\,\overline u$ give
\[
 B_uB_v=B_{uv},\qquad A_uA_v=A_{uv},\qquad D_uD_v=D_{uv}.
\]
For example, the second coordinate of $A_uA_v(a,b)$ is
$(b\overline v)\overline u=b\overline{uv}$, retaining the order.

For $x=(a,b)$ and $y=(c,d)$, expansion of
\eqref{rettri:CD} gives
\begin{align*}
 (B_u x)(A_u y)
 &= (u a\overline u,b)(u c,d\overline u)\\
 &=\bigl((u a\overline u)(u c)
          -\overline{d\overline u}\,b,
          (d\overline u)(u a\overline u)
             +b\overline{u c}\bigr)\\
 &=\bigl(u(ac-\overline d b),
                    (da+b\overline c)\overline u\bigr)
 =A_u(xy).
\end{align*}
This uses only quaternionic associativity in the displayed
components. Since $A_u=K D_u K$, it proves membership in
\eqref{rettri:RT}. If $\rho(u)$ is the identity, evaluate
$A_u$ at $(1,0)$ to get $(u,0)=(1,0)$, so $u=1$.
Thus $\rho$ is injective.

The other two cyclic pair identities, with every conjugation
retained, are
\begin{align*}
 (A_u x)(D_u y)
 &=(u a,b\overline u)(c\overline u,d\overline u)\\
 &=\bigl(u(ac-\overline d b)\overline u,\;da+b\overline c\bigr)
 =B_u(xy)=K B_u K(xy),\\
 (D_u x)(B_u y)
 &=(a\overline u,b\overline u)(u c\overline u,d)\\
 &=\bigl((ac-\overline d b)\overline u,
                    (da+b\overline c)\overline u\bigr)
 =D_u(xy)=K A_u K(xy).
\end{align*}

Put $\langle v,w\rangle=\operatorname{Re}(v\overline w)$ on
$\mathbb O$. Then $\operatorname{Re}(vz)=\langle v,Kz\rangle$.
The already proved product relation and $K D_u=A_u K$ imply
\begin{align*}
 \operatorname{Re}(((B_u x)(A_u y))(D_u z))
 &=\langle A_u(xy),K D_u z\rangle\\
 &=\langle A_u(xy),A_u Kz\rangle
 =\langle xy,Kz\rangle
 =\operatorname{Re}((xy)z).
\end{align*}
No reassociation of a non-real octonionic triple product was used.
\end{proof}

\paragraph{The exact original Jordan coordinates and all three bundles.}
Write the original Hermitian matrix with upper entries $(p,q,r)$ as
\begin{equation}
 J(\lambda;p,q,r)=
 \begin{pmatrix}
 \lambda_1&p&q\\
 \overline p&\lambda_2&r\\
 \overline q&\overline r&\lambda_3
 \end{pmatrix},\qquad
 (q_1,q_2,q_3)=(r,\overline q,p).
 \label{rettri:J}
\end{equation}
Thus $q_2$ here is the cyclic Jordan coordinate, not the previously
defined cusp quotient map; only the explicitly displayed ordered
triple in this paragraph uses this coordinate notation.
Under \eqref{rettri:rho}, the cyclic coordinates transform by
\[
 (q_1,q_2,q_3)\longmapsto(B_u q_1,A_u q_2,D_u q_3).
\]
Since $K A_u K=D_u$, the action in the original upper entries is
\begin{equation}
 (r,q,p)\longmapsto(B_u r,D_u q,D_u p),\qquad
 J(\lambda;p,q,r)\longmapsto
       J(\lambda;D_u p,D_u q,B_u r).
 \label{rettri:original-action}
\end{equation}
In particular this is an action on the three specific spaces
$\mathfrak h_{\gamma_1},\mathfrak h_{\gamma_2},
\mathfrak h_{\gamma_3}$ of
\eqref{eq:octonionic-root-spaces}, in their original order.

The determinant and trace form are
\begin{align}
 \det J&=\lambda_1\lambda_2\lambda_3
       -\lambda_1|r|^2-\lambda_2|q|^2-\lambda_3|p|^2
       +2\operatorname{Re}((r\overline q)p),
 \label{rettri:det}\\
 \operatorname{tr}(J^2)&=\lambda_1^2+\lambda_2^2+\lambda_3^2
                       +2(|p|^2+|q|^2+|r|^2).
 \label{rettri:metric}
\end{align}
Theorem~\ref{rettri:subgroup} and norm preservation prove that
\eqref{rettri:original-action} preserves every term, with its
displayed sign and coefficient. Its being a Jordan automorphism
also follows directly from Yokota's diagonal-fixing action in
Theorem 2.7.1, with exactly the identification
\eqref{rettri:J}. Equivalently, the Peirce multiplication rules
$F_i(x)\circ F_i(y)=\langle x,y\rangle(E_{i+1}+E_{i+2})$ and
$F_i(x)\circ F_{i+1}(y)=\tfrac12 F_{i+2}(\overline{xy})$,
where indices are cyclic modulo three, are preserved by the norm
identity and the three cyclic pair identities proved above.
Here $F_1(x)$ has upper entry $x$ at $(2,3)$,
$F_2(x)$ has upper entry $\overline x$ at $(1,3)$, and
$F_3(x)$ has upper entry $x$ at $(1,2)$; all other entries are
zero apart from their Hermitian conjugates. Direct matrix
multiplication gives the displayed rules, with coefficient
$1/2$ from the Jordan product. Finally
$E_i\circ E_j=\delta_{ij}E_i$, $E_i\circ F_i(x)=0$,
and $E_i\circ F_j(x)=F_j(x)/2$ for $i\ne j$, and these rules
are preserved because all $E_i$ are fixed. The displayed
identities specify the Jordan product on every pair by real
bilinearity.

For a right principal bundle our associated-bundle convention is
$(p u,v)\sim(p,\sigma(u)v)$. Let
\[
 \mathcal L_H=P_H\times_{u:a\mapsto u a}\mathbb H.
\]
This is the tautological right quaternionic line pulled back by
$\chi H$: the fibre map $[(x,v),a]\mapsto (x,v a)$ is
well-defined, right quaternionic linear and invertible onto the
line $\chi H(x)$. Its unit sphere bundle is $P_H$, by
$p\mapsto[p,1]$. Consequently it is nontrivial as a quaternionic
line. It is also nontrivial as a real four-plane bundle: if it had
a nowhere-zero real section, dividing by length would give a
section of its unit sphere bundle and hence of $P_H$.

The original ordered Peirce representation produces the real
rank-24 bundle
\begin{equation}
 \mathcal W_H=\mathcal V_{r,H}\oplus\mathcal V_{q,H}
                              \oplus\mathcal V_{p,H},
 \quad
 \mathcal V_{r,H}=P_H\times_B\mathbb O,
 \quad
 \mathcal V_{q,H}=\mathcal V_{p,H}=P_H\times_D\mathbb O.
 \label{rettri:W}
\end{equation}
The equality between the last two names specifies isomorphic
representations with two separately retained labelled copies.
Their explicit decompositions are
\begin{equation}
 \begin{aligned}
 \mathcal V_{r,H}&\cong
       \mathbb R\oplus\mathcal A_H\oplus\mathbb R^4,
       &&(\operatorname{Re}a,\operatorname{Im}a,b),\\
 \mathcal V_{q,H}&\cong\mathcal L_H\oplus\mathcal L_H,
       &&(a,b)\longmapsto(\overline a,\overline b),\\
 \mathcal V_{p,H}&\cong\mathcal L_H\oplus\mathcal L_H,
       &&(a,b)\longmapsto(\overline a,\overline b).
 \end{aligned}
 \label{rettri:splittings}
\end{equation}
For the last two maps,
$\overline{a\overline u}=u\overline a$ proves equivariance,
and similarly for $b$. Each quaternionic conjugation has real
determinant $-1$ in $(1,i,j,k)$, so their ordered pair has
determinant $+1$. The first line keeps the actual coordinate order
$(1,\operatorname{Im}\mathbb H,\mathbb H\ell)$; moving the
one-dimensional first factor across the imaginary three-space
would contribute the sign $(-1)^3=-1$, and no such reorder is
being made. In cyclic coordinates the middle action is $A$;
the equivariant map to two left standard representations is
$(a,b)\mapsto(a,\overline b)$, of determinant $-1$. This is
why the original $\overline q$ is retained explicitly.

All fibrewise maps in \eqref{rettri:product} and
\eqref{rettri:det} descend to the associated bundles: replacing
a local frame $p$ by $p u$ changes inputs by precisely the three
displayed actions, whose product and scalar identities have been
proved. In particular
\[
 \mathbb R^3\oplus\mathcal W_H
       =P_H\times_\rho J_3(\mathbb O)
\]
is a bundle of Jordan algebras with its three fixed diagonal
idempotents, its full product, trace form and determinant. The
three-sector structure on \eqref{rettri:W} includes the
explicit nonzero invariant trilinear map
$(r,q,p)\mapsto2\operatorname{Re}((r\overline q)p)$;
it is nonzero because its value at $r=q=p=1$ is exactly $2$.

\paragraph{What extension to the full triality group forgets.}
Form the principal bundle
\begin{equation}
 Q_H=P_H\times_\rho\operatorname{Spin}(8)\longrightarrow X,
 \quad (p u,s)\sim(p,\rho(u)s).
 \label{rettri:Q}
\end{equation}
It contains the reduction $P_H$ by $p\mapsto[p,1]$. This is
equivariant because $[p u,1]=[p,\rho(u)]$.

\begin{proposition}
\label{rettri:forgetting}
The principal bundle $Q_H$ is trivial. All three rank-eight
bundles in \eqref{rettri:W} are therefore trivial as real
vector bundles, while their specified subbundles
$\mathcal A_H$ and $\mathcal L_H$ remain nontrivial.
After any chosen trivialization of $Q_H$, its reduction determines
a continuous map
\begin{equation}
 f_H:X\longrightarrow\operatorname{Spin}(8)/\rho(\operatorname{Sp}(1))
 \label{rettri:reductionmap}
\end{equation}
whose pullback of the principal bundle
$\operatorname{Spin}(8)\to\operatorname{Spin}(8)/\rho(\operatorname{Sp}(1))$
is $P_H$. This map is non-null-homotopic.
\end{proposition}

\begin{proof}
Let $g:S^5\to\operatorname{Sp}(1)$ be a based clutching
representative for $P_H$ in the convention used in
\eqref{rettri:clutch}. Its homotopy class has order two.
The transition $D_g$ on an eight-dimensional sector is conjugate
by $C(a,b)=(\overline a,\overline b)$ to
$\operatorname{diag}(L_g,L_g)$ on $\mathbb H^2$.
For $0\le t\le\pi/2$ put
\[
 R_t=\begin{pmatrix}\cos t&-\sin t\\
                       \sin t&\cos t\end{pmatrix}
       \quad\text{on }\mathbb H^2.
\]
This is an orthogonal path commuting with quaternionic scalar
right multiplication. Its conjugation carries
$\operatorname{diag}(L_g,I)$ at $t=0$ to
$\operatorname{diag}(I,L_g)$ at $t=\pi/2$. The maps into $SO(8)$
represented by these two matrices are therefore based homotopic.
Pointwise product of maps from a based sphere into a topological
group represents addition in its homotopy group: concatenation
in one loop coordinate and pointwise multiplication have the
same unit and satisfy interchange, yielding the same operation.
Consequently
\[
 [\operatorname{diag}(L_g,L_g)]
 =2[\operatorname{diag}(L_g,I)]
 =\iota_*(2[g])=0\quad\text{in }\pi_5(SO(8)),
\]
where $\iota(u)=\operatorname{diag}(L_u,I)$. Conjugating back by
$C$ gives $[D_g]=0$.

The third projection of \eqref{rettri:RT} is a double covering
$\operatorname{Spin}(8)\to SO(8)$: this is the triality covering
in the third marking, with the same conjugation $K$ as in
\eqref{rettri:RT}. Coverings induce isomorphisms on
$\pi_n$ for $n\ge2$ by unique lifting of based sphere maps and
their homotopies. The projection sends $[\rho g]$ to $[D_g]=0$,
so $[\rho g]=0$. A bundle obtained by gluing two trivial bundles
over disks with a null-homotopic clutching map is trivial: extend
that clutching map across a disk and change its local frame by
the extension. This proves triviality of $Q_H$ directly.
Every associated vector bundle of a trivial principal bundle is
trivial; applying the three projections gives the asserted
rank-eight trivializations. The earlier proofs of nontriviality
of $\mathcal A_H$ and $\mathcal L_H$ still apply to these actual
subbundles, so forgetting their inclusions loses information.

Choose an isomorphism $\tau:Q_H\to X\times G$, where
$G=\operatorname{Spin}(8)$, and put $K_0=\rho(\operatorname{Sp}(1))$.
For $p\in(P_H)_x$ write $\tau([p,1])=(x,a_p)$. Then
$a_{pu}=a_p\rho(u)$, so
$f_H(x)=a_p K_0$ is well-defined, independent of the choice of
$p$, and continuous in local sections. The explicit principal
bundle isomorphism to the pullback is
\[
 P_H\longrightarrow f_H^*(G\to G/K_0),\qquad
 p\longmapsto(x,a_p),
\]
using the identification of $\operatorname{Sp}(1)$ with $K_0$
by the proved injective map $\rho$. The formula is fibrewise
bijective and has inverse in each local trivialization.
If $f_H$ were null-homotopic, its pullback bundle would be
isomorphic to the pullback along a constant map and hence
trivial, contradicting Theorem~\ref{rettri:bundleclass}.
Equivalently, in the fibration exact sequence its clutching
boundary is precisely $c(P_H)\ne0$.
\end{proof}

There is also an explicit map into the space of oriented
three-planes in the first eight-dimensional sector. If
$T_1:G\to SO(\mathbb O)$ is the first projection of
\eqref{rettri:RT}, define
\begin{equation}
 \begin{aligned}
 \Gamma:G/K_0&\longrightarrow\operatorname{Gr}_3^+(\mathbb O),\\
 gK_0&\longmapsto
 \operatorname{span}_{\mathbb R}^{+}
       \bigl(T_1(g)i,T_1(g)j,T_1(g)k\bigr).
 \end{aligned}
 \label{rettri:Grassmann}
\end{equation}
Here the superscript records the orientation of the displayed
ordered basis. This formula is well-defined: multiplication of
$g$ on the right by $\rho(u)$ changes that ordered basis by
the orientation-preserving map $\operatorname{Ad}_u$ of
$\operatorname{Im}\mathbb H$, because $T_1\rho(u)=B_u$.
Continuity follows from local sections of $G\to G/K_0$.
Let $\mathcal T_3$ be the tautological oriented three-plane
bundle on this Grassmannian. The explicit fibre map
\[
 G\times_{K_0}\operatorname{Im}\mathbb H
       \longrightarrow\Gamma^*\mathcal T_3,\qquad
 [g,v]\longmapsto(gK_0,T_1(g)v)
\]
is well-defined since
$T_1(g\rho(u))v=T_1(g)\operatorname{Ad}_u v$.
It is a linear isometry on every fibre, hence a bundle
isomorphism. Pulling back by \eqref{rettri:reductionmap}
therefore gives
\[
 (\Gamma f_H)^*\mathcal T_3\cong\mathcal A_H.
\]
In particular $\Gamma f_H:X\to\operatorname{Gr}_3^+(\mathbb O)$
is non-null-homotopic, since pullback along a null-homotopic
map would make the tautological bundle trivial. This locates
the three-plane in the actual first triality sector, with its
orientation and its nonzero clutching class retained.

Thus the exact bridge is the sequence of explicit constructions
\[
 H\ \longmapsto\ P_H\ \longmapsto\ \mathcal A_H
 \ \longmapsto\ (\mathcal W_H,\rho\text{-reduction},
                   2\operatorname{Re}((r\overline q)p)).
\]
The first two steps retain the entire nonzero order-two class
by \eqref{rettri:exact-detection}. In the last step the
specified reduction, or equivalently the distinguished
three-plane with its spin frame data, retains that class.
The triple of unrestricted rank-eight bundles alone does not.
The space $G/K_0$ parametrizes reductions acting on this very
same rank-24 carrier; no claim that $\dim(G/K_0)=24$ is made.

\paragraph{The precise half-angle sign law and its relation to the cusp.}
For the fixed imaginary unit $i$ define the path
\begin{equation}
 u(\theta)=\cos(\theta/2)+i\sin(\theta/2),
                   \qquad 0\le\theta\le2\pi.
 \label{rettri:halfangle}
\end{equation}
Equation~\eqref{rettri:Rodrigues} gives
\[
 \operatorname{Ad}_{u(\theta)}i=i,\quad
 \operatorname{Ad}_{u(\theta)}j=\cos\theta\,j+\sin\theta\,k,\quad
 \operatorname{Ad}_{u(\theta)}k=-\sin\theta\,j+\cos\theta\,k.
\]
This is one full $2\pi$ rotation in the oriented $(j,k)$ plane.
Its lift starts at $1$ and ends at $-1$. It is therefore the
nontrivial loop in $SO(3)$: a null-homotopy of the loop would
lift over a disk and force its lifted endpoints to agree.
Substitution in \eqref{rettri:threeactions} gives exactly
\begin{equation}
 \rho(-1)=(I,-I,-I).
 \label{rettri:central-sign}
\end{equation}
It is central in \eqref{rettri:RT} because each component
is scalar. Thus a full rotation in the three-dimensional
adjoint sector acts by minus the identity on both associated
spinor sectors while acting by the identity on the vector
sector. This conclusion follows from actual related-triple
maps and their product identity.

The parameter $\theta$ in \eqref{rettri:halfangle} is a
structure-group parameter; it has not been equated with the
original cusp phase $\phi$. The actual dependence on that cusp
phase is already specified, with all four coordinates, by
\eqref{rettri:z} and \eqref{rettri:transition}.
No connection or holonomy is part of the definition of $P_H$.
In particular the order-two class proved here lies in
$\pi_5(\operatorname{Sp}(1))$ through degree-six clutching,
and is not being identified with a degree-one line-bundle
class. Every real line bundle on $X\cong S^6$ is trivial:
its unit sphere bundle is a double covering of a simply
connected space, so it has a section. The exact relation
established here is instead the faithful order-two detection
\eqref{rettri:exact-detection}, the nontrivial real
three-plane \eqref{rettri:rankthree}, its explicit
three-sector triality action, and the structural sign law
\eqref{rettri:central-sign}.

% Primary convention source for integration bibliography:
% I. Yokota, Exceptional Lie Groups, arXiv:0902.0431,
% sections 1.10, 1.14, 1.16 and Theorem 2.7.1.

\subsection{A nonzero Hopf--cusp morphism to the stable mod--24 group}
\label{hbridge:section}
The retained-angle input used below is the complete Theorem~\ref{ret:main} above.

\subsubsection{The exact source maps and conventions}

Retain the original marked source \((X,x_0)\), its fixed based
diffeomorphism \(d:(S^6,*)\to(X,x_0)\), the cusp coordinate \(t\),
the fixed cutoff \(\epsilon\), and the native quotient
\(\overline\pi_f\) of the cusp fibres by
\(\mathbb Z\langle\widehat w,\widehat\delta\rangle\).
The quotient coordinates are ordered by
\((\overline{\widehat\gamma},\overline{\widehat u})\).
Thus the fixed continuous based map is
\[
q_2:X\longrightarrow
\mathcal A=\operatorname{Susp}(S^1\times\overline T^2),
\qquad
q_2(x)=[|t|,t/|t|,\overline\pi_f(x)]
\quad(0<|t|<\epsilon).
\]
It sends the entire central fibre \(W\) to the first suspension
vertex and \(X\setminus\operatorname{int}N_0\) to the second.
The unreduced suspension \(\mathcal A\) is based at its first
vertex. In the original product torus set
\[
K^{(2)}=(\{1\}\times\overline T^2)
\cup(S^1\times\{u_1=0\})
\cup(S^1\times\{u_2=0\}),
\qquad c_3:T^3\longrightarrow T^3/K^{(2)}\cong S^3.
\]
The retained-angle sphere map and its sphere-source representative are
\begin{equation}
p_3([r,z])=[r,c_3(z)],\qquad
H=p_3q_2:X\longrightarrow\Sigma_r(S^3,*)\cong S^4,
\qquad h=H\circ d:S^6\longrightarrow S^4.
\label{hbridge:eq:retained-map}
\end{equation}
Both suspension vertices and the whole cylinder over the basepoint
are identified in \(\Sigma_r(S^3,*)\). The preceding complete
retained-angle framing theorem proves
\begin{equation}
[h]=\eta_4\circ\eta_5\ne0
\quad\hbox{in}\quad\pi_6(S^4)\cong\mathbb Z/2.
\label{hbridge:eq:proved-H}
\end{equation}
Its proof identifies the actual derivative-framed inverse torus
with the Lie-framed torus, with Arf invariant one. We use this
already proved class, retaining the particular map
\eqref{hbridge:eq:retained-map}; see \cite{hbridgebib:HBRetained}.

Fix the complex Hopf map
\[
\eta_2:(S^3,(1,0))\longrightarrow
(\mathbb{CP}^1,[1:0])\cong(S^2,*),\qquad
(z_1,z_2)\longmapsto[z_1:z_2].
\]
Use ordered smash coordinates on spheres. In this note
\(E^k f=f\wedge\operatorname{id}_{S^k}\) denotes right reduced
suspension; the usual ordered sphere homeomorphisms are understood.
Put \(\eta_n=E^{n-2}\eta_2\), and write \(\eta\) for the
stable Hopf class. The order of each factor is retained below.

\subsubsection{The Hopf--cusp homomorphism and its nonzero value}

\begin{theorem}
\label{hbridge:mu-theorem}
The assignment
\begin{equation}
\mu_\eta:\pi_6(S^4)\longrightarrow\pi_9(S^6),\qquad
\mu_\eta([x])=[\eta_2\wedge x]
\label{hbridge:eq:mu}
\end{equation}
is an injective group homomorphism.\par
Its literal representative has
the complete domain and codomain
\[
\eta_2\wedge x:S^3\wedge S^6\cong S^9
\longrightarrow S^2\wedge S^4\cong S^6.
\]
Choose any generator \(\nu_6\) of
\(\pi_9(S^6)\cong\mathbb Z/24\). Then the particular original
cusp map supplies
\begin{equation}
\kappa_H=[\eta_2\wedge h]
=\mu_\eta([h])=12\nu_6\ne0.
\label{hbridge:eq:kappa}
\end{equation}
It has exact order two. Its value is independent of orientation
reversals in the sphere identifications and of the choice of the
generator \(\nu_6\).
\end{theorem}

\begin{proof}
A based homotopy of \(x\) smashes with the fixed based Hopf map,
so \eqref{hbridge:eq:mu} is well-defined. The exact ordered
point-set factorization is
\begin{equation}
\eta_2\wedge x
=(\eta_2\wedge\operatorname{id}_{S^4})
 \circ(\operatorname{id}_{S^3}\wedge x):
S^3\wedge S^6\longrightarrow S^3\wedge S^4
\longrightarrow S^2\wedge S^4.
\label{hbridge:eq:factorization}
\end{equation}
The first applied map has source dimension nine and target dimension
seven; the second has source dimension seven and target dimension six.
The second map is \(E^4\eta_2\). To compare the first with
\(E^3x=x\wedge\operatorname{id}_{S^3}\), interchange its sphere
factors at source and target. The interchange degrees are respectively
\((-1)^{3\cdot6}=+1\) and \((-1)^{3\cdot4}=+1\). A based map
of a sphere of degree one is based homotopic to the identity, so
\begin{equation}
\mu_\eta([x])=(E^4\eta_2)_*(E^3[x]).
\label{hbridge:eq:hom-formula}
\end{equation}
Suspension is a group homomorphism. Postcomposition is also a group
homomorphism on homotopy groups, since it distributes over the
source pinch map defining the sum. Thus
\eqref{hbridge:eq:hom-formula} proves additivity, including zero
and inverses, before passage to stable homotopy.

Freudenthal makes
\[
E:\pi_{n+3}(S^n)\longrightarrow\pi_{n+4}(S^{n+1})
\]
an isomorphism for \(n\ge5\), because \(n+3<2n-1\).
In particular stabilization
\(\pi_9(S^6)\to\pi_3^{\mathrm S}\) is an isomorphism.
Stabilization from \(\pi_6(S^4)\) is likewise an isomorphism
onto \(\pi_2^{\mathrm S}\), because \(n+2<2n-1\) for
\(n\ge4\). The exact factorization
\eqref{hbridge:eq:factorization} shows that its stable class is
the stable composition product of \(\eta\) with \([x]\).
Using the proved original class \eqref{hbridge:eq:proved-H},
\begin{equation}
E^\infty\kappa_H
=\eta\,E^\infty[h]
=\eta\,\eta^2
=\eta^3
=12\nu\ne0\quad\hbox{in}\quad
\pi_3^{\mathrm S}\cong\mathbb Z/24.
\label{hbridge:eq:stable-product}
\end{equation}
The integral relation here is the one explicitly stated in
\cite[p.~385]{hbridgebib:HBHatcher}, not the coefficient \(4\) in the
two-primary quotient. More precisely, the two-primary relation
is \(\eta^3=4\nu\) modulo eight. Bilinearity and \(2\eta=0\)
give \(2\eta^3=0\) integrally. In \(\mathbb Z/24\) the only
elements killed by two are \(0\) and \(12\nu\), and only the
second has nonzero image \(4\nu\) modulo eight. This recovers
the exact integral coefficient and its nonvanishing.

The inverse of the stabilization isomorphism therefore gives
\eqref{hbridge:eq:kappa}. Since the source of \(\mu_\eta\) is
cyclic of order two, generated by \([h]\), its nonzero value
proves injectivity and identifies its image as
\(\{0,12\nu_6\}\).

An orientation reversal at a sphere source or target multiplies
the corresponding stable homotopy class by \(-1\); here
\(-12\nu_6=12\nu_6\). A sign change in the chosen Hopf
representative also leaves the stable product unchanged since
\(-\eta=\eta\). Finally, replacing the generator by
\(u\nu_6\), with \(u\) a unit modulo twenty-four, does not
change the element: every such \(u\) is odd, and
\(12u\equiv12\pmod{24}\). All orientation and generator
assertions consequently hold for the actual order-two element.
\end{proof}

\begin{corollary}[The same nonzero class in the twelve-dimensional receiver]
Put \(\nu_9=E^3\nu_6\). The particular based map
\[
E^3(\eta_2\wedge h):S^{12}\longrightarrow S^9
\]
has class \(12\nu_9\ne0\) in
\(\pi_{12}(S^9)\cong\mathbb Z/24\).
\end{corollary}
\begin{proof}
Each of these three suspensions is an isomorphism in the preceding
stable range. Apply their composite to
\eqref{hbridge:eq:kappa}; it preserves the exact order two.
\end{proof}

\subsubsection{The actual marked peripheral character}

Retain the original three elliptic-kernel matrices, with their
original signs and ordering:
\begin{equation}
C_0=\begin{pmatrix}1&1\\0&1\end{pmatrix},\qquad
C_1=\begin{pmatrix}-1&-1\\1&0\end{pmatrix},\qquad
C_2'=\begin{pmatrix}0&1\\-1&0\end{pmatrix}.
\label{hbridge:eq:matrices}
\end{equation}
The existing peripheral theorem \cite{hbridgebib:HBHigher} defines the
surjective marked abelianization character
\begin{equation}
\chi:SL_2(\mathbb Z)\longrightarrow\mathbb Z/12,
\qquad
(\chi(C_0),\chi(C_1),\chi(C_2'))=(-1,4,-3).
\label{hbridge:eq:chi}
\end{equation}
For an explicit check of this normalization, let
\[
S_0=\begin{pmatrix}0&-1\\1&0\end{pmatrix},\qquad T=C_0.
\]
Then \(C_2'=S_0^3\), \(C_1=T^{-1}S_0\), and direct matrix
multiplication gives \(C_1C_2'=C_0^{-1}\). In the modular
presentation \(S_0^4=1\), \(S_0^2=(S_0T)^3\), the
abelianized generators satisfy \(4s=0\) and \(s+3t=0\).
Thus \(s=-3t\) and \(12t=0\). The assignment
\(t\mapsto-1\), \(s\mapsto3\) satisfies both relations and
maps onto \(\mathbb Z/12\); consequently the order of the
abelianized generator is exactly twelve. It gives, without a
change of peripheral matrices, precisely
\(\chi(C_0)=-1\), \(\chi(C_1)=1+3=4\), and
\(\chi(C_2')=3\cdot3=-3\pmod{12}\).

\begin{theorem}[An exact nonzero peripheral homotopy morphism]
\label{hbridge:peripheral-theorem}
Define
\begin{equation}
\Theta_H:\mathbb Z/12\longrightarrow\pi_6(S^4),
\qquad \Theta_H([n])=n[h],
\qquad
\Psi_H=\mu_\eta\circ\Theta_H\circ\chi:
SL_2(\mathbb Z)\longrightarrow\pi_9(S^6).
\label{hbridge:eq:Psi}
\end{equation}
These are group homomorphisms. The complete coefficient formula is
\begin{equation}
\Psi_H(M)=12\,\widetilde{\chi(M)}\,\nu_6,
\label{hbridge:eq:coefficient}
\end{equation}
where \(\widetilde{\chi(M)}\) is any integer lift of the
marked residue. In the original matrix order of
\eqref{hbridge:eq:matrices}, its exact values are
\begin{equation}
(\Psi_H(C_0),\Psi_H(C_1),\Psi_H(C_2'))
=(12\nu_6,0,12\nu_6).
\label{hbridge:eq:values}
\end{equation}
Its image and kernel are exactly
\begin{equation}
\operatorname{im}\Psi_H=\{0,12\nu_6\},\qquad
\ker\Psi_H=\chi^{-1}(2\mathbb Z/12).
\label{hbridge:eq:kernel}
\end{equation}
In particular its kernel has index two in \(SL_2(\mathbb Z)\).
\end{theorem}

\begin{proof}
The element \([h]\) has exact order two, so replacing \(n\)
by \(n+12k\) leaves \(n[h]\) unchanged. Addition in its
abelian homotopy group proves that \(\Theta_H\) is a group
homomorphism. Moreover
\[
\ker\Theta_H=2\mathbb Z/12,
\qquad \operatorname{im}\Theta_H=\pi_6(S^4).
\]
The composite in \eqref{hbridge:eq:Psi} is consequently a
group homomorphism, and \eqref{hbridge:eq:kappa} gives
\(\mu_\eta(n[h])=12n\nu_6\). Changing the integer lift by
\(12k\) changes its coefficient by \(144k\), which is a
multiple of twenty-four. This proves
\eqref{hbridge:eq:coefficient} with its independence of lifts.
The three coefficients are respectively
\(-12\equiv12\), \(48\equiv0\), and
\(-36\equiv12\pmod{24}\), proving
\eqref{hbridge:eq:values}. Since \(12n\equiv0\pmod{24}\)
holds exactly when \(n\) is even, the asserted kernel is exact.
Surjectivity of \(\chi\) and nonvanishing at \(C_0\) prove
that the image is the full order-two subgroup, so the first
isomorphism theorem gives the index assertion.
\end{proof}

Here is the complete parity factorization. Let
\(\operatorname{red}_2:\mathbb Z/12\to\mathbb Z/2\) be
ordinary parity reduction, let
\(\iota_2([e])=12e\pmod{24}\), and let
\(c_{\nu_6}:\pi_9(S^6)\to\mathbb Z/24\) send \(\nu_6\)
to one. Every arrow and its codomain in the diagram are fixed:
\begin{equation}
\begin{CD}
SL_2(\mathbb Z) @>{\Psi_H}>> \pi_9(S^6)\\
@V{\operatorname{red}_2\circ\chi}VV @VV{c_{\nu_6}}V\\
\mathbb Z/2 @>{\iota_2}>> \mathbb Z/24.
\end{CD}
\label{hbridge:eq:parity-square}
\end{equation}
Commutativity is exactly the coefficient calculation above; the
bottom arrow is injective because \(12e\) is zero modulo
twenty-four exactly when \(e\) is zero modulo two.

\subsubsection{Comparison with the earlier maps, with their exact types}

Write \(\psi_H=c_{\nu_6}\circ\Psi_H\). Ordinary reduction
\(r:\mathbb Z/24\to\mathbb Z/12\) gives
\begin{equation}
r\circ\psi_H=0,
\qquad (r\circ\psi_H)(C_0)=0\ne-1=\chi(C_0).
\label{hbridge:eq:not-lift}
\end{equation}
Thus this nonzero character is not a reduction-compatible lift of
\(\chi\). Its natural commuting square is
\eqref{hbridge:eq:parity-square}. The existing nonexistence
theorem for a reduction-compatible lift remains valid: every
homomorphism from \(SL_2(\mathbb Z)\) to an abelian group
factors through its abelianization \(\mathbb Z/12\), whereas
a lift of \(\chi(C_0)=-1\) would require \(C_0\) to map
to \(11\) or \(23\) modulo twenty-four, each of order
twenty-four. The abelianized generator has order twelve and
cannot have either image.

The universal even-subgroup injection is the separate map
\[
j:\mathbb Z/12\hookrightarrow\mathbb Z/24,
\qquad j([a])=[2a].
\]
The new coefficient morphism factors through it by the exact
identity
\begin{equation}
\psi_H=j\circ[6]\circ\chi,
\qquad [6]:\mathbb Z/12\longrightarrow\mathbb Z/12,
\quad [a]\longmapsto[6a].
\label{hbridge:eq:j-factor}
\end{equation}
Indeed \(2(6a)=12a\pmod{24}\). The map \(j\) is
injective with image of order twelve; the composite
\eqref{hbridge:eq:j-factor} has image of order two because
\([6]\) retains precisely parity. The map \(j\) is realized
by the already proved universal suspension
\(E^6:\pi_6(S^3)\to\pi_{12}(S^9)\) after compatible
generator choices. Equation \eqref{hbridge:eq:j-factor} is a
coefficient identity; the actual new geometric representative is
\(\eta_2\wedge h\), followed, when desired, by three
suspensions as above.

For the original angle-forgetting map
\(g=\widehat p_2q_2:X\to S^3\), the separate framed
calculation has proved \([g\circ d]=0\). Therefore the
previous morphism
\[
\Theta_g:\mathbb Z/12\longrightarrow\pi_6(S^3),
\qquad [n]\longmapsto n[g\circ d]
\]
and \(E^6\Theta_g\chi\) are zero. In contrast,
\(E^3\Psi_H(C_0)=12\nu_9\ne0\) in this same stable
receiver \(\pi_{12}(S^9)\). These statements agree because
the exact postcompositions from \(\mathcal A\) are different:
\(p_3\) retains the angular coordinate in its top-cell map,
while \(\widehat p_2\) forgets it. The new homomorphism
uses the already proved value of \(p_3q_2\) and the specified
Hopf map, with all maps displayed in
\eqref{hbridge:eq:retained-map}--\eqref{hbridge:eq:Psi}.

The constructed representative has source \(S^3\wedge X\cong
S^9\); its threefold suspension has source \(S^{12}\).
The number twenty-four is the order of the receiving stable
homotopy group. This calculation does not identify either source
with a twenty-four-dimensional carrier, does not supply new
periods or analytic overlap maps, and does not complete the
global higher torus family. It supplies the exact nonzero
homotopy and peripheral morphisms just proved.

\paragraph{Source checks.}
The three matrices and the marked values were checked directly
against equations \emph{elliptic-peripheral-matrices} and
\emph{peripheral-mod12-values} in the existing higher-rung draft
\cite{hbridgebib:HBHigher}. The calculation above reproduces their signs
using the same \(S_0,T\). The retained-angle input is the
complete derivative-framing theorem in \cite{hbridgebib:HBRetained}.
For the stable facts, the author-hosted text
\cite{hbridgebib:HBHatcher} was read at Corollary~4.24, printed p.~360,
and at printed pp.~384--385. Page~385 explicitly distinguishes
the two-primary coefficient four from the integral coefficient
twelve. Proposition~4.56 on that page, with its proof on
p.~388, gives the stable composition product and its smash
description. The fixed-factor proof
\eqref{hbridge:eq:factorization}--\eqref{hbridge:eq:hom-formula}
establishes the needed homomorphism here directly.

\subsection{Boundary of the three tested numerical mechanisms}


% END RETAINED_ANGLE_AND_HOPF_BRIDGE_20260905

\begin{theorem}[Boundary of the three tested numerical mechanisms]
\label[theorem]{thm:exact-numerical-boundary}
For the chamber indexing, peripheral abelianization, stable suspension,
and the three additional mechanisms tested below, the verified relations
are the following.
\begin{enumerate}[label=\textup{(\roman*)}]
\item The six Kuhn zero-strata and the six octonionic fixed flags have
      the marked \(S_3\)-equivariant finite-set bijection
      \eqref{eq:Kuhn-flag-bijection}.  Their common labels also index
      the flag cells, and both Euler characteristics are \(6\).
\item The elliptic-kernel peripheral matrices produce the surjective
      character \(\chi:SL_2(\Z)\to\Z/12\), with no
      reduction-compatible or surjective mod--\(24\) lift.
\item The separate stable suspension map embeds \(\Z/12\) as
      \(2(\Z/24)\), and \cref{prop:generator-free-bridge} gives the exact
      composite determined by \(g\). The actual coefficient is \(a=0\)
      by \cref{thm:cusp-g-zero}; the composite is zero, and the associated \(S^{24}\)
      fourth-smash map is null-homotopic.
\item The ambient \(F_4\) root system has \(h=12\), and
      \(W(F_4)/W(D_4)\cong S_3\).  The three specific tests give:
      diagonal-sum factorization through \(P\) or \(Q\) exactly for
      \(N\mid3\); ordinary cohomology reduction for every \(N\), with
      \(e_N\) of order \(N/\gcd(N,6)\); and plain determinant-winding
      reduction onto \(\Z/N\) for every \(N\).  None of these specified
      operations selects \(N=12\) or \(N=24\) from the dimension or
      Coxeter number alone.
\end{enumerate}
No arrow between these four lines is defined merely by equality of their
integers.
\end{theorem}

\begin{proof}
Item (i) is \cref{thm:Kuhn-octonionic-comparison}; item (ii) is
\cref{thm:peripheral-moduli}; item (iii) is
\cref{thm:stable-12-to-24,prop:generator-free-bridge}; and the Coxeter
assertion in (iv) is \cref{prop:two-Coxeter-systems}.  The three stated
tests for an additional arithmetic arrow are evaluated by
\cref{prop:cusp-diagonal-modulus,prop:no-modulus-from-cohomology-or-determinant}:
the primitive lattice
diagonal-sum factorization allows only \(N\mid3\), cohomology reduces for every
modulus, and determinant winding reduces onto every modulus.  These
computations prove the stated boundary for those mechanisms without
asserting that arbitrary future structures cannot be added.
\end{proof}

% BEGIN ES-NIEMEIER HOST-ONLY INPUT: exclude from generated h24: fragment.
% Proof-preserving conversion of ORBIT_CARTAN_BRIDGE.md, Sections 1--7.
% Staged fragment only; uses the host's original Albert coordinates.
\section{The full Albert orbit quotient and the Cartan family}
\label{ocb:sec:full-family}

The fixed-spectrum comparison of
\cref{thm:octonionic-triality,att:sec:ordered-tensor} extends to the
whole Albert algebra with all three eigenvalue variables retained.
We give the invariant quotient, its smooth regular part, the exact
normal framing, and the two focal strata. The diagonalization and flag
geometry are established in \cite[Proposition~2.8.2 and
Theorems~2.7.4, 2.8.3]{Yokota2009} and
\cite[Definition~2.2 and Sections~2.3--2.6]{MareWillems2013}; the
calculations below specify their maps in the retained coordinates.

\subsection{Original space, invariants, and direct diagonalization}
\label{ocb:sec:diagonalization}

Let $\mathbb O$ be the real octonion division algebra with its retained
conjugation, norm $N$, and inner product
$h(a,b)=\operatorname{Re}(a\bar b)$. Set
$J=\mathfrak h_3(\mathbb O)$, with Jordan product
$X\circ Y=(XY+YX)/2$, and retain the cyclic coordinates of
\eqref{att:eq:Q}:
\begin{equation}
 X=\begin{pmatrix}\xi_1&q_3&\bar q_2\\
 \bar q_3&\xi_2&q_1\\q_2&\bar q_1&\xi_3\end{pmatrix}
 =\sum_{i=1}^3\xi_i e_i+\sum_{i=1}^3E_i(q_i).
 \label{ocb:eq:coordinates}\tag{OCB1}
\end{equation}
Here $e_i$ are the original diagonal idempotents, $E_i$ are the original
off-diagonal inclusions, and $I=e_1+e_2+e_3$. Matrix products in the
Jordan product involve only products of two octonions. Powers of a single
Jordan element are unambiguous by the Jordan identity; all subsequent
polynomial formulas use these powers.

The positive trace metric and determinant are
\begin{equation}
 \begin{split}
 G(X,Y)&=\operatorname{tr}(X\circ Y),\\
 G(X,X)&=\xi_1^2+\xi_2^2+\xi_3^2+2\sum_iN(q_i),\\
 d(X)&=\xi_1\xi_2\xi_3-\sum_i\xi_iN(q_i)
              +2\operatorname{Re}((q_1q_2)q_3).
 \end{split}
 \label{ocb:eq:metric-determinant}\tag{OCB2}
\end{equation}
Thus $\dim_{\mathbb R}J=27$, with every diagonal coordinate, metric
factor, and mixed cubic retained. Define
\begin{equation}
 t(X)=\operatorname{tr}X,\qquad u(X)=G(X,X),\qquad
 e_2(X)=\frac{t(X)^2-u(X)}2,\qquad
 \Phi(X)=(t(X),u(X),d(X)).
 \label{ocb:eq:invariants}\tag{OCB3}
\end{equation}
Let $F_4=\operatorname{Aut}_{\mathbb R}(J,\circ)$. An automorphism
fixes $I$, preserves trace and $G$, and therefore preserves the
determinant by the trace identity
\begin{equation}
 d(X)=\tfrac13\operatorname{tr}(X^3)
 -\tfrac12t(X)u(X)+\tfrac16t(X)^3.
 \label{ocb:eq:trace-determinant}\tag{OCB4}
\end{equation}
For completeness, trace preservation follows without choosing a
preferred matrix realization: the Peirce table gives
$\operatorname{Tr}_{\operatorname{End}(J)}(L_X)=9\operatorname{tr}X$;
conjugation of $L_X$ by an automorphism proves the assertion. Product
preservation is a closed polynomial condition, so $F_4$ is closed in the
orthogonal group of $(J,G)$ and hence compact.

\begin{proposition}[Diagonalization and the exact invariant fibres]
\label{ocb:prop:invariant-fibres}
Every $X\in J$ is $F_4$-conjugate to a real diagonal matrix. The fibres
of $\Phi$ are exactly the $F_4$ orbits, including repeated-eigenvalue
and scalar orbits. The map $\Phi:J\to\Phi(J)$ is proper.
\end{proposition}

\begin{proof}
On the compact orbit of $X$, maximize
$\xi_1^2+\xi_2^2+\xi_3^2$. If $q_i\ne0$ at a maximizing point,
put $a=q_i/|q_i|$, let $(i,j,k)$ be cyclic, and use the already proved
derivation
\begin{equation}
 D_i(a)=4[L_{E_i(a)},L_{e_j}].
 \label{ocb:eq:derivation}\tag{OCB5}
\end{equation}
Its action satisfies
$D_i(a)(e_j-e_k)=2E_i(a)$ and
$D_i(a)E_i(a)=-2(e_j-e_k)$. The remaining off-diagonal slots do not
contribute diagonal entries under this action. Write the relevant
component as $v(e_j-e_k)+wE_i(a)$, with
$v=(\xi_j-\xi_k)/2$ and $w=|q_i|>0$. Exponentiation rotates $(v,w)$
through angle $2s$. Some rotation gives squared diagonal-difference
coordinate $v^2+w^2$, increasing the diagonal-square sum by $2w^2$.
This contradicts maximality. Hence every off-diagonal coordinate
vanishes at a maximum, proving diagonalization with the diagonal mean
retained.

If $X=g\sum_i\lambda_i e_i$, then
\eqref{ocb:eq:invariants}--\eqref{ocb:eq:trace-determinant} show that
the multiset of real eigenvalues is exactly the multiset of roots of
\begin{equation}
 p_X(z)=z^3-t(X)z^2+e_2(X)z-d(X).
 \label{ocb:eq:characteristic-polynomial}\tag{OCB6}
\end{equation}
It is independent of the diagonalizing automorphism. Every permutation
of a real diagonal is realized by conjugation with a real orthogonal
permutation matrix. Such conjugation preserves the Jordan product;
a commuting sign adjustment can be used if determinant $+1$ is
desired. Equal invariant triples thus give the same eigenvalue multiset
and the same orbit. Conversely, the invariants are constant on orbits,
as proved above. This includes repeated roots and scalar matrices.
Finally, a compact set of invariant triples bounds $u=G(X,X)$, and
its inverse image is closed and bounded in $J$, hence compact.
\end{proof}

This is the compact-orbit diagonalization argument of
\cite[Proposition~2.8.2, pp.~56--57]{Yokota2009}, reproduced in the
retained derivation coordinates. It also proves the polynomial identities
used below by transporting their diagonal versions.

\subsection{Exact projectors, smooth orbit inverse, and regular family}
\label{ocb:sec:regular-family}

Fix labelled, pairwise distinct real numbers
$\lambda_1,\lambda_2,\lambda_3$, with arbitrary sum. Put
$d_\lambda=\sum_i\lambda_i e_i$ and
$\mathcal O_\lambda=F_4d_\lambda$. The ordered frame space is
\[
 \mathcal F=\{(p_1,p_2,p_3):p_i^2=p_i,\
 \operatorname{tr}p_i=1,\ p_i\circ p_j=0\ (i\ne j),\
 \textstyle\sum_i p_i=I\}.
\]
The correspondence $(p_1,p_2,p_3)\mapsto(p_1,p_2)$ and its inverse
$(p_1,p_2)\mapsto(p_1,p_2,I-p_1-p_2)$ identify this space with the
flag space of \eqref{eq:octonionic-flag-definition}; in particular
$\mathcal F=F_4/\operatorname{Spin}(8)$ by
\cref{thm:octonionic-triality}. Define
\begin{equation}
 \Xi_\lambda:\mathcal F\longrightarrow\mathcal O_\lambda,
 \qquad (p_i)\longmapsto\sum_i\lambda_i p_i,
 \label{ocb:eq:orbit-map}\tag{OCB7}
\end{equation}
and, for $\{i,j,k\}=\{1,2,3\}$, define
\begin{equation}
 P_i^\lambda(X)=
 \frac{(X-\lambda_jI)\circ(X-\lambda_kI)}
      {(\lambda_i-\lambda_j)(\lambda_i-\lambda_k)}.
 \label{ocb:eq:projectors}\tag{OCB8}
\end{equation}

\begin{proposition}[The ordered spectral diffeomorphism]
\label{ocb:prop:spectral-diffeomorphism}
The map $\Xi_\lambda$ is an $F_4$-equivariant diffeomorphism, with
inverse $(P_1^\lambda,P_2^\lambda,P_3^\lambda)$. In particular
$\dim\mathcal O_\lambda=24$. If
$\mathcal C=\{\lambda_1>\lambda_2>\lambda_3\}$ and
$J_{\mathrm{reg}}$ is the set of elements with distinct eigenvalues,
then $\mathcal F\times\mathcal C\to J_{\mathrm{reg}}$, given by the
same spectral sum, is a diffeomorphism.
\end{proposition}

\begin{proof}
For $X=\sum_l\lambda_l p_l$, orthogonality and idempotency give
\[
 \begin{aligned}
 (X-\lambda_jI)\circ(X-\lambda_kI)
 &=\sum_l(\lambda_l-\lambda_j)(\lambda_l-\lambda_k)p_l\\
 &=(\lambda_i-\lambda_j)(\lambda_i-\lambda_k)p_i.
 \end{aligned}
\]
Thus \eqref{ocb:eq:projectors} recovers every ordered projector and
proves both inverse compositions. It also proves that the stabilizer
of $d_\lambda$ is the pointwise diagonal-frame stabilizer
$\operatorname{Spin}(8)$. The quotient orbit map is an injective
immersion of the compact homogeneous space
$F_4/\operatorname{Spin}(8)$ into the Hausdorff vector space $J$,
hence an embedding. Its inverse is polynomial in $X$ divided by
nonzero fixed real constants, proving smoothness of the inverse and
the asserted equivariant diffeomorphism.

The differential at fixed $\lambda$ is explicitly
\begin{equation}
 dP_i^\lambda|_X(H)=
 \frac{2X\circ H-(\lambda_j+\lambda_k)H}
      {(\lambda_i-\lambda_j)(\lambda_i-\lambda_k)}.
 \label{ocb:eq:projector-differential}\tag{OCB9}
\end{equation}
At $d_\lambda$, the sector $E_i(\mathbb O)$ joins eigenspaces $j,k$.
On this sector $dP_j(H)=H/(\lambda_j-\lambda_k)$,
$dP_k(H)=-H/(\lambda_j-\lambda_k)$, and $dP_i(H)=0$.
Substitution into $\sum_l\lambda_l dP_l(H)$ returns $H$, verifying
the differential inverse with its signs.

The original derivations satisfy
\begin{equation}
 \begin{gathered}
 D_i(a)d_\lambda=(\lambda_j-\lambda_k)E_i(a)=\Delta_i E_i(a),\\
 (\Delta_1,\Delta_2,\Delta_3)=
 (\lambda_2-\lambda_3,\lambda_3-\lambda_1,\lambda_1-\lambda_2),
 \qquad(i,j,k)\text{ cyclic}.
 \end{gathered}
 \label{ocb:eq:tangent-multipliers}\tag{OCB10}
\end{equation}
All three multipliers are nonzero, so the three eight-dimensional
off-diagonal sectors constitute the tangent space. Their transported
Peirce subbundles at every orbit point have metric $2\Delta_i^2h$
in the homogeneous tangent coordinates of
\eqref{att:eq:bundle-comparison}.

For the full map
\begin{equation}
 \mathcal F\times\mathcal C\longrightarrow J_{\mathrm{reg}},
 \qquad ((p_i),(\lambda_i))\longmapsto\sum_i\lambda_i p_i,
 \label{ocb:eq:full-regular-map}\tag{OCB11}
\end{equation}
surjectivity follows from diagonalization and sorting. Injectivity
follows from the unique ordered roots and \eqref{ocb:eq:projectors}.
At a simple root,
$\partial p_X/\partial z=\prod_{j\ne i}(\lambda_i-\lambda_j)\ne0$;
the implicit function theorem therefore gives a smooth local root
function. Ordering glues these functions globally.
Equation~\eqref{ocb:eq:projectors}, with the smoothly varying roots,
provides the smooth inverse of \eqref{ocb:eq:full-regular-map}.
\end{proof}

\subsection{Retained mean, radius, and angular parameter}
\label{ocb:sec:mean-radius-angle}

For every original $X$, put
\begin{equation}
 \mu=\frac{t(X)}3,\qquad Y=X-\mu I,\qquad
 s=G(Y,Y)=u(X)-\frac{t(X)^2}3,\qquad d_0=d(Y).
 \label{ocb:eq:mean-radius}\tag{OCB12}
\end{equation}
Write $J_0=\mathfrak h_3^0(\mathbb O)=\{Y\in J:\operatorname{tr}Y=0\}$.
The displayed change of variables has inverse $X=\mu I+Y$;
$\mu$ and $s$ retain their actual values. The eigenvalues
$a_i=\lambda_i-\mu$ of $Y$ satisfy $\sum_i a_i=0$, so expansion gives
\begin{equation}
 d(X)=\mu^3-\frac{\mu s}{2}+d_0.
 \label{ocb:eq:mean-determinant}\tag{OCB13}
\end{equation}
For $s=r^2>0$, the ordered eigenvalues have the exact parametrization
\begin{equation}
 \begin{aligned}
 \lambda_1&=\mu+\sqrt{2/3}\,r\cos\theta,\\
 \lambda_2&=\mu+\sqrt{2/3}\,r\cos(\theta-2\pi/3),\\
 \lambda_3&=\mu+\sqrt{2/3}\,r\cos(\theta+2\pi/3),
 \end{aligned}
 \qquad 0\le\theta\le\pi/3.
 \label{ocb:eq:angular-eigenvalues}\tag{OCB14}
\end{equation}
The cosines sum to zero, their squared sum is $3/2$, and their
product is $\tfrac14\cos3\theta$. Consequently
\begin{equation}
 \begin{aligned}
 t(X)&=3\mu,&u(X)&=3\mu^2+r^2,\\
 d_0&=\frac{r^3}{3\sqrt6}\cos3\theta,&
 d(X)&=\mu^3-\frac{\mu r^2}{2}
                    +\frac{r^3}{3\sqrt6}\cos3\theta.
 \end{aligned}
 \label{ocb:eq:angular-invariants}\tag{OCB15}
\end{equation}
For $0<\theta<\pi/3$ the ordered inequalities are strict. The inverse is
\begin{equation}
 \begin{gathered}
 \mu=t/3,\qquad r=\sqrt{u-t^2/3},\\
 \theta=\tfrac13\arccos\left(
 \frac{3\sqrt6\,[d-\mu^3+\mu r^2/2]}{r^3}\right).
 \end{gathered}
 \label{ocb:eq:angular-inverse}\tag{OCB16}
\end{equation}
To prove that every ordered triple is listed, put $a_3=-a_1-a_2$
and expand the discriminant:
\begin{equation}
 \mathcal D=\prod_{i<j}(a_i-a_j)^2
            =\frac{s^3}{2}-27d_0^2\ge0.
 \label{ocb:eq:discriminant}\tag{OCB17}
\end{equation}
Thus the arccos argument lies in $[-1,1]$. The triple in
\eqref{ocb:eq:angular-eigenvalues} has the same characteristic
polynomial and the correct ordering, hence equals the original
triple. Equality in \eqref{ocb:eq:discriminant} means repeated
eigenvalues. If $r=0$, positivity of $G$ gives $Y=0$ and $X=\mu I$;
this scalar stratum requires no angle.

The full invariant image is therefore exactly
\begin{equation}
 \left\{(t,u,d):s=u-t^2/3\ge0,\quad
 27(d-\mu^3+\mu s/2)^2\le s^3/2,\quad\mu=t/3\right\}.
 \label{ocb:eq:invariant-image}\tag{OCB18}
\end{equation}
Every triple in this set is attained by
\eqref{ocb:eq:angular-eigenvalues}, including the scalar case just
described. This gives the whole three-dimensional orbit quotient;
specified mean and radius select its corresponding fibre.

The full regular metric also records these parameters. Let $h_i$ be
the metric $h$ on the $i$th associated bundle, expressed by the
homogeneous tangent coordinate $K_i[g,a]$ of \eqref{att:eq:Kmap}.
The pullback of $G$ by
\eqref{ocb:eq:full-regular-map}--\eqref{ocb:eq:angular-eigenvalues} is
\begin{equation}
 3\,d\mu^2+dr^2+r^2d\theta^2
             +\sum_i2\Delta_i(\mu,r,\theta)^2h_i.
 \label{ocb:eq:full-metric}\tag{OCB19}
\end{equation}
Indeed the diagonal differential has coefficient vectors
$\mathbf1$, $a/r$, and $\partial_\theta a$. Their pairwise inner
products vanish, and their squared norms are $3,1,r^2$. Every
diagonal differential is orthogonal to every off-diagonal Peirce
sector. The sector terms follow from
\eqref{ocb:eq:tangent-multipliers} and $E_i^*G=2h$, with every
eigenvalue difference and radial factor retained.

\subsection{Exact codimension and global normal framing}
\label{ocb:sec:normal-framing}

The same $24$-dimensional regular orbit has codimension three in $J$,
two in the affine fixed-trace space $\mu I+J_0$, and one in the
fixed-trace, fixed-radius sphere $\mu I+S_r^{25}$. The change of
variables \eqref{ocb:eq:mean-radius} and the radius-level inclusion
give the actual maps between these ambient descriptions.

For a fixed distinct triple, its normal space in $J$ is
\begin{equation}
 \nu_X\mathcal O_\lambda
 =\operatorname{span}_{\mathbb R}
     \{P_1^\lambda(X),P_2^\lambda(X),P_3^\lambda(X)\}.
 \label{ocb:eq:projector-normal-space}\tag{OCB20}
\end{equation}
At a diagonal point this is the orthogonal complement of the three
off-diagonal tangent sectors; transport by $F_4$ proves the equality
everywhere. The projectors are pairwise $G$-orthogonal with squared
norm one.

The differential of the invariant map has gradients
\begin{equation}
 \nabla t=I,\qquad\nabla u=2X,\qquad
 \nabla d=X^\#=X^2-tX+e_2I.
 \label{ocb:eq:invariant-gradients}\tag{OCB21}
\end{equation}
To check the determinant gradient, diagonalize $X$. By
\eqref{ocb:eq:metric-determinant}, its derivative on a diagonal
variation is $\sum_i\lambda_j\lambda_k H_{ii}$, where
$\{i,j,k\}=\{1,2,3\}$, and its derivative on an off-diagonal
variation is zero. The right side of
\eqref{ocb:eq:invariant-gradients} is the diagonal vector of these
coefficients; automorphism invariance proves the formula for every
$X,H$. In the projector basis the coefficient determinant is
\begin{equation}
 \det\begin{pmatrix}
 1&1&1\\2\lambda_1&2\lambda_2&2\lambda_3\\
 \lambda_2\lambda_3&\lambda_1\lambda_3&\lambda_1\lambda_2
 \end{pmatrix}
 =-2(\lambda_1-\lambda_2)(\lambda_1-\lambda_3)
                                      (\lambda_2-\lambda_3).
 \label{ocb:eq:gradient-determinant}\tag{OCB22}
\end{equation}
Hence $\Phi$ has rank three exactly on the regular stratum, where
its tangent kernel is the orbit tangent space.

In the fixed-trace hyperplane define the two normal vectors
\begin{equation}
 Y=X-\mu I,\qquad Z=Y^2-\frac{s}{3}I.
 \label{ocb:eq:two-normal-fields}\tag{OCB23}
\end{equation}
They have trace zero and are normal by
\eqref{ocb:eq:projector-normal-space}. The trace-zero eigenvalue
identities give
\begin{equation}
 \operatorname{tr}Y^3=3d_0,\qquad
 \operatorname{tr}Y^4=s^2/2,\qquad
 \operatorname{Gram}(Y,Z)=
 \begin{pmatrix}s&3d_0\\3d_0&s^2/6\end{pmatrix}.
 \label{ocb:eq:normal-gram}\tag{OCB24}
\end{equation}
Specifically, $\sum_i a_i^3=3a_1a_2a_3$ follows from
$\sum_i a_i=0$. Squaring $\sum_{i<j}a_ia_j=-s/2$ gives
$\sum_{i<j}a_i^2a_j^2=s^2/4$, since the remaining cross term is
$2a_1a_2a_3\sum_i a_i=0$. Expanding $s^2$ now gives
$\sum_i a_i^4=s^2/2$. These identities and
$Z=Y^2-sI/3$ prove every Gram entry. Therefore
\begin{equation}
 \det\operatorname{Gram}(Y,Z)=s^3/6-9d_0^2
                             =\mathcal D/3>0
 \label{ocb:eq:normal-gram-determinant}\tag{OCB25}
\end{equation}
for distinct eigenvalues. Thus $(Y,Z)$ is a smooth ordered global
two-framing of the orbit normal bundle in $\mu I+J_0$.

The corresponding explicitly specified orthonormal frame is
\begin{equation}
 n_1=\frac{Y}{\sqrt s},\qquad
 n_2=\frac{Z-(3d_0/s)Y}{\sqrt{s^2/6-9d_0^2/s}}.
 \label{ocb:eq:orthonormal-frame}\tag{OCB26}
\end{equation}
Every square root here is the positive real root; its argument is
positive by \eqref{ocb:eq:normal-gram-determinant}. In $J$, prepend
$n_0=I/\sqrt3$ to obtain $(n_0,n_1,n_2)$. In the radius sphere the
normal direction is $n_2$. These formulas identify the three normal
bundles explicitly.

These normal fields are parallel for the normal connection along the
fixed orbit. Each is a linear combination of the $P_i^\lambda$ with
coefficients constant on that orbit. Along $g(s)X$, its ambient
derivative is $g'(s)$ applied to the corresponding diagonal normal
vector, and lies in the transported off-diagonal tangent sectors by
\eqref{ocb:eq:tangent-multipliers}. Its normal component is zero.
Thus the normal bundle is flat with the displayed global parallel
orthonormal frame.

For any transported normal $\eta=\sum_i\nu_iP_i^\lambda$ with fixed
real $\nu_i$, use the convention $A_\eta H=-(D_H\eta)^T$. Then
\begin{equation}
 A_\eta\big|_{E_i(\mathbb O)}
 =-\frac{\nu_j-\nu_k}{\lambda_j-\lambda_k}\,\operatorname{id},
 \qquad(i,j,k)\text{ cyclic},
 \label{ocb:eq:shape-operator}\tag{OCB27}
\end{equation}
with the sectors transported by $g$ at other points. Indeed a path
generated by $D_i(a)$ has tangent
$(\lambda_j-\lambda_k)E_i(a)$, whereas the derivative of its
normal field is $(\nu_j-\nu_k)E_i(a)$. This proves the formula
and its minus sign. The shape operator for $n_0$ is zero, and the
one for the radial normal $n_1$ is $-\operatorname{id}/r$.

\subsection{Cartan equations and spherical principal curvatures}
\label{ocb:sec:cartan-equations}

On the whole trace-zero space $J_0$, with its original trace metric,
define the explicitly scaled cubic
\begin{equation}
 C(Y)=3\sqrt6\,d(Y).
 \label{ocb:eq:cartan-polynomial}\tag{OCB28}
\end{equation}
The inverse relation is $d(Y)=C(Y)/(3\sqrt6)$, and
\eqref{ocb:eq:mean-determinant} recovers the full determinant with
its mean and radius. Orthogonal projection of
\eqref{ocb:eq:invariant-gradients} to $J_0$ gives
$\nabla_{J_0}d=Y^2-sI/3=Z$. Hence
\eqref{ocb:eq:normal-gram} gives
\begin{equation}
 \|\nabla C\|_G^2=54\|Z\|_G^2=9s^2,
 \qquad\Delta_{J_0}C=0.
 \label{ocb:eq:cartan-euclidean}\tag{OCB29}
\end{equation}
For the second equality, the full $J$ Laplacian of
\eqref{ocb:eq:metric-determinant} is $-8\operatorname{tr}X$.
Each eight-dimensional octonion slot has metric $2h$ and contributes
$-8\xi_i$; each pure diagonal second derivative is zero, and the
mixed cubic is separately linear in each real coordinate. Restricting
to $J_0$ subtracts the second derivative in the unit scalar direction
$I/\sqrt3$. The identity
\[
 d(Y+zI/\sqrt3)=d(Y)-\frac{s z}{2\sqrt3}
                              +\frac{z^3}{3\sqrt3}
\]
shows that this second derivative is zero at $z=0$. The full
Laplacian also vanishes at trace zero, proving
\eqref{ocb:eq:cartan-euclidean} with the original metric factor two.

On each specified radius sphere $S_r^{25}\subset J_0$, $r>0$,
set $f_r=C/r^3$. Then \eqref{ocb:eq:angular-invariants} gives
$f_r=\cos3\theta$, and homogeneity yields
\begin{equation}
 \|\nabla_{S_r}f_r\|^2=\frac9{r^2}(1-f_r^2),\qquad
 \Delta_{S_r}f_r=-\frac{81}{r^2}f_r.
 \label{ocb:eq:cartan-spherical}\tag{OCB30}
\end{equation}
Indeed the radial component of $\nabla C$ is $(3C/r)n_1$.
The polar Laplacian on the $26$-dimensional Euclidean space,
applied to the degree-three homogeneous harmonic function, is
$r[3(3+24)f_1+\Delta_{S_1}f_1]$. These calculations prove both
constants in \eqref{ocb:eq:cartan-spherical}. Its regular level
hypersurfaces are therefore isoparametric: both the squared gradient
and the Laplacian depend only on the function value. By
\eqref{ocb:eq:angular-invariants}--\eqref{ocb:eq:invariant-image},
the regular levels are precisely the fixed-mean, fixed-radius regular
$F_4/\operatorname{Spin}(8)$ orbits.

The principal curvatures admit a direct calculation. On the diagonal
section let $N_\theta=(1/r)\partial_\theta Y$, a unit vector
perpendicular to $Y$, and transport it by the projectors. The exact
eigenvalue differences in \eqref{ocb:eq:angular-eigenvalues} are
\begin{equation}
 \Delta_1=\sqrt2r\sin\theta,\qquad
 \Delta_2=-\sqrt2r\sin(\theta+\pi/3),\qquad
 \Delta_3=-\sqrt2r\sin(\theta-\pi/3).
 \label{ocb:eq:angular-differences}\tag{OCB31}
\end{equation}
Apply \eqref{ocb:eq:shape-operator} with
$\nu_i=(\partial_\theta\lambda_i)/r$. For the declared normal
$N_\theta$, the three principal curvatures are
\begin{equation}
 -\frac1r\cot\theta,\qquad
 -\frac1r\cot(\theta+\pi/3),\qquad
 -\frac1r\cot(\theta-\pi/3),
 \label{ocb:eq:principal-curvatures}\tag{OCB32}
\end{equation}
each on its original eight-dimensional Peirce sector. They are
distinct for $0<\theta<\pi/3$. In the invariant frame
\eqref{ocb:eq:orthonormal-frame}, one has $n_2=-N_\theta$:
differentiating $C=r^3\cos3\theta$ gives
$\nabla_{S_r}C=-3r^2\sin3\theta\,N_\theta$, whereas $n_2$
is the positive-length version of the projected determinant
gradient. The principal curvatures for $n_2$ are consequently the
opposites of \eqref{ocb:eq:principal-curvatures}. This proves the
three-curvature, multiplicity-eight Cartan family with both normal
orientations, mean, and radius specified.

\subsection{Focal projective planes and their endpoint maps}
\label{ocb:sec:focal-strata}

At $\theta=0$ one has
$\lambda_2=\lambda_3=\mu-r/\sqrt6$ and
$\lambda_1=\mu+2r/\sqrt6$; at $\theta=\pi/3$ one has
$\lambda_1=\lambda_2=\mu+r/\sqrt6$ and
$\lambda_3=\mu-2r/\sqrt6$. The two singular orbits in
$\mu I+S_r^{25}$ are exactly
\begin{equation}
 \begin{aligned}
 \mathcal P^+_{\mu,r}
 &=\{\mu I+\tfrac r{\sqrt6}(3p-I):p^2=p,\
                                       \operatorname{tr}p=1\},\\
 \mathcal P^-_{\mu,r}
 &=\{\mu I-\tfrac r{\sqrt6}(3p-I):p^2=p,\
                                       \operatorname{tr}p=1\}.
 \end{aligned}
 \label{ocb:eq:focal-planes}\tag{OCB33}
\end{equation}
Their inverse maps are, respectively,
\begin{equation}
 p=\tfrac13\left(I+\tfrac{\sqrt6}{r}(X-\mu I)\right),\qquad
 p=\tfrac13\left(I-\tfrac{\sqrt6}{r}(X-\mu I)\right).
 \label{ocb:eq:focal-inverses}\tag{OCB34}
\end{equation}
These affine maps are smooth, have the displayed two-sided inverses
on the specified sets, and are $F_4$-equivariant. The primitive
idempotent space is the sixteen-dimensional Cayley projective plane
$\mathbb O\mathbb P^2=F_4/\operatorname{Spin}(9)$
\cite[Theorems~2.7.4 and 2.8.3]{Yokota2009}
\cite[Section~2.3]{MareWillems2013}. Hence the formulas give embedded
projective-plane strata.

The endpoint maps of the regular family are the explicit projections
\begin{equation}
 \begin{aligned}
 (p_1,p_2,p_3)&\longmapsto
             \mu I+\tfrac r{\sqrt6}(3p_1-I),\\
 (p_1,p_2,p_3)&\longmapsto
             \mu I-\tfrac r{\sqrt6}(3p_3-I).
 \end{aligned}
 \label{ocb:eq:endpoint-maps}\tag{OCB35}
\end{equation}
Each fibre is $\operatorname{Spin}(9)/\operatorname{Spin}(8)\cong S^8$:
fixing one primitive idempotent leaves two ordered complementary
idempotents, the octonionic projective line. These frame bundles
are identified in \cite[Sections~2.3--2.6]{MareWillems2013}.
They are also the focal maps. Indeed the diagonal path $Y(\theta)$
is a great circle because $\partial_\theta^2Y=-Y$ and
$|\partial_\theta Y|=r$. Transport by $F_4$ therefore gives the
normal geodesics to the regular hypersurfaces. At each endpoint one
of the tangent multipliers in \eqref{ocb:eq:angular-differences}
vanishes with its eight-dimensional sector, giving exactly the
collapse in \eqref{ocb:eq:endpoint-maps}. Each focal stratum has
codimension nine in the ambient radius sphere and a normal sphere
of dimension eight.

For $r=0$, the construction gives the scalar line $X=\mu I$,
each point fixed by $F_4$. For $r>0$ and endpoint angle, the full
repeated-eigenvalue stratum retains the two parameters $(\mu,r)$
in addition to its projective-plane coordinate, and has dimension
eighteen in $J$. These descriptions retain every scalar, radial,
and mean stratum.

\subsection{Exact map from the equal-diagonal triality slice}
\label{ocb:sec:triality-slice}

For the original matrix
$Q_{\mathbb O}(\lambda;q)=\lambda I+E(q)$, with
$E(q)=\sum_iE_i(q_i)$ and
$\tau(q)=\operatorname{Re}((q_1q_2)q_3)$, the maps above give
\begin{equation}
 \mu=\lambda,\qquad Y=E(q),\qquad
 r^2=2\sum_iN(q_i),\qquad d_0=2\tau(q),\qquad
 C(Y)=6\sqrt6\,\tau(q).
 \label{ocb:eq:triality-parameters}\tag{OCB36}
\end{equation}
Thus its invariant triple is
\begin{equation}
 \begin{aligned}
 \Phi(Q_{\mathbb O}(\lambda;q))
 =\bigl(&3\lambda,\ 3\lambda^2+2\sum_iN(q_i),\\
        &\lambda^3-\lambda\sum_iN(q_i)+2\tau(q)\bigr).
 \end{aligned}
 \label{ocb:eq:triality-invariant-map}\tag{OCB37}
\end{equation}
This is the explicit invariant morphism from the original triality
coordinates to the full orbit quotient. At every simple spectrum,
its ordered-projector lift is \eqref{ocb:eq:projectors} applied to
$Q_{\mathbb O}(\lambda;q)$. The positive boundary graph of
\cref{tri:thm:spectral-graph} chooses the original scalar
$\lambda=-\lambda_{\min}(E(q))$. In the full orbit coordinates,
this says that the least eigenvalue is zero and the other two are
nonnegative. Equations~\eqref{ocb:eq:triality-parameters}--
\eqref{ocb:eq:triality-invariant-map} prove the relation to the whole
Albert orbit family with all coefficients and this scalar choice
retained.

\paragraph{Attribution and scope.}
Yokota's diagonalization is Proposition~2.8.2, pp.~56--57;
his primitive-idempotent stabilizer and Cayley-plane identification
are Theorems~2.7.4 and 2.8.3, pp.~54 and 57
\cite{Yokota2009}. Mare and Willems give the primitive idempotents,
full real diagonal space, frame stabilizer, projective plane,
$S^8$ frame projections, and the polar isotropy representation
in Definition~2.2 and Sections~2.3--2.6, pp.~488--493
\cite{MareWillems2013}. Their description identifies $J_0$ with
the isotropy representation of $E_{6(-26)}/F_4$ and displays its
diagonal two-plane and three eight-dimensional root spaces. The
Cartan-family equations and focal maps used here are proved in
\eqref{ocb:eq:cartan-euclidean}--\eqref{ocb:eq:endpoint-maps}.
The orbit and flag geometry is established literature; these
calculations assert neither complex integrability nor physical
dynamics.


% Proof-preserving conversion of NIEMEIER_CARTAN_SPECTRAL_BRIDGE.md.
% Staged fragment only; depends on the common-lattice and orbit fragments.
\section{The exact spectral map of the common marked lattice}
\label{ncs:sec:spectral-bridge}

The literal common-lattice embedding of
\cref{wcl:thm:literal-intersection} has a matrix-valued and
frame-valued lift into the Albert orbit family of
\cref{ocb:sec:full-family}. We calculate this map with the three
octonionic slots separately labelled, the trace metric, the arbitrary
scalar diagonal, and the original determinant coefficient two.

\subsection{Domain and the original matrix}
\label{ncs:sec:domain}

Use the ordered quaternion subalgebra
\[
 \mathbb H=\mathbb R\langle1,i,j,k\rangle\subset\mathbb O,
 \qquad i^2=j^2=k^2=ijk=-1.
\]
For $a=(a_0,a_1,a_2,a_3)$, write
\begin{equation}
 \begin{gathered}
 v=a_1i+a_2j+a_3k,\qquad
 t=\sqrt{a_1^2+a_2^2+a_3^2},\qquad a=a_0+v,\\
 Q=a_0^2+t^2,\qquad P=a_0(a_0^2-3t^2).
 \end{gathered}
 \label{ncs:eq:quaternion-coordinates}\tag{NCS1}
\end{equation}
Every square root in this section means the nonnegative real root.
Here $t$ denotes the length of the imaginary quaternion, while the
trace invariant will always be written $\operatorname{tr}X$.
The integral domain and its retained embedding are
\begin{equation}
 D_4=\{a\in\mathbb Z^4:a_0+a_1+a_2+a_3\in2\mathbb Z\},
 \qquad j(a)=\sqrt2(a,a,a)\in\mathbb O^3.
 \label{ncs:eq:lattice-embedding}\tag{NCS2}
\end{equation}
The literal intersection theorem
\cref{wcl:thm:literal-intersection} identifies $j(D_4)$ with the
common lattice of the two specified marked Niemeier/Leech embeddings
and the specified Wilson embedding. Its original Niemeier coordinates
are
\begin{equation}
 x_{i,m}=(-1)^m a_{i-1}\quad(1\le i\le4,\ 0\le m\le5),
 \qquad d=0.
 \label{ncs:eq:niemeier-coordinates}\tag{NCS3}
\end{equation}
The six coordinates in each $A_5$ block sum to zero. Pairing with
the retained Weyl vector gives $3\sum_i a_i$, which is divisible
by six for $a\in D_4$; hence these vectors belong to the retained
neighbor sublattice $K$. The full Wilson membership certificate and
both intersection equalities are proved in
\cref{wcl:thm:literal-intersection}.

For arbitrary $\mu\in\mathbb R$, define
\begin{equation}
 Y(a)=\begin{pmatrix}
 0&\sqrt2a&\sqrt2\bar a\\
 \sqrt2\bar a&0&\sqrt2a\\
 \sqrt2a&\sqrt2\bar a&0
 \end{pmatrix},\qquad
 X(\mu,a)=\mu I+Y(a).
 \label{ncs:eq:original-matrix}\tag{NCS4}
\end{equation}
This is the original cyclic convention
$(q_1,q_2,q_3)=(\sqrt2a,\sqrt2a,\sqrt2a)$, including the
conjugated upper-right entry. On $J=\mathfrak h_3(\mathbb O)$,
the product is $A\circ B=(AB+BA)/2$ and the metric is
$G(A,B)=\operatorname{tr}(A\circ B)$. Consequently
\begin{equation}
 \|j(a)\|^2=6Q,\qquad r^2:=G(Y,Y)=12Q.
 \label{ncs:eq:two-metrics}\tag{NCS5}
\end{equation}
This retains the factor two between the two metrics. The map
$(\mu,a)\mapsto X(\mu,a)$ from $\mathbb R\times D_4$ to $J$
is injective: its inverse on its image is
$\mu=\operatorname{tr}X/3$ and $a=X_{12}/\sqrt2$.

\subsection{Complete determinant and invariant map}
\label{ncs:sec:invariant-map}

The quaternion identity $v^2=-t^2$ and associativity give
\[
 \operatorname{Re}(a^3)=a_0^3-3a_0t^2=P.
\]
Thus the triality cubic and determinant of $Y$ are exactly
\begin{equation}
 \begin{aligned}
 \tau(j(a))
 &=\operatorname{Re}((\sqrt2a\sqrt2a)\sqrt2a)=2\sqrt2P,\\
 \det Y&=2\tau(j(a))=4\sqrt2P.
 \end{aligned}
 \label{ncs:eq:cubic-determinant}\tag{NCS6}
\end{equation}
The complete invariant map is therefore
\begin{equation}
 \boxed{
 (\operatorname{tr}X,\operatorname{tr}X^2,\det X)
   =\bigl(3\mu,\ 3\mu^2+12Q,\
                     \mu^3-6\mu Q+4\sqrt2P\bigr).
 }
 \label{ncs:eq:full-invariants}\tag{NCS7}
\end{equation}
Indeed each of the three squared off-diagonal norms is $2Q$.
The negative terms in the original determinant formula thus sum
to $-6\mu Q$, while the mixed term is precisely
\eqref{ncs:eq:cubic-determinant}. This proves the displayed
invariant map with every diagonal variable and original lattice
scale retained. The characteristic polynomial of $Y$ is
\begin{equation}
 p_Y(z)=z^3-6Qz-4\sqrt2P.
 \label{ncs:eq:characteristic-polynomial}\tag{NCS8}
\end{equation}

\subsection{Exact eigenvalues, projectors, and directional inverse}
\label{ncs:sec:projectors}

The three labelled eigenvalues of $Y$ are
\begin{equation}
 \nu_0=2\sqrt2a_0,\qquad
 \nu_1=-\sqrt2a_0-\sqrt6t,\qquad
 \nu_2=-\sqrt2a_0+\sqrt6t.
 \label{ncs:eq:labelled-eigenvalues}\tag{NCS9}
\end{equation}
Their sum is zero, their squared sum is $12Q$, and their product
is $4\sqrt2P$. Direct multiplication of
$\prod_{j=0}^2(z-\nu_j)$ gives
\eqref{ncs:eq:characteristic-polynomial}, proving that these are
the eigenvalues. The eigenvalues of $X$ are $\mu+\nu_j$, with
$\mu$ arbitrary.

For $t>0$, let $n=v/t$, so $n^2=-1$ and $v=tn$.
The plane $\mathbb C_n=\mathbb R\langle1,n\rangle$ is a
commutative associative subalgebra. Set
\begin{equation}
 \begin{gathered}
 T=\begin{pmatrix}0&1&0\\0&0&1\\1&0&0\end{pmatrix},
 \qquad T^3=I,\\
 \omega_n=-\frac12+\frac{\sqrt3}{2}n,\qquad
 \omega_n^3=1,\qquad1+\omega_n+\omega_n^2=0.
 \end{gathered}
 \label{ncs:eq:cyclic-matrix}\tag{NCS10}
\end{equation}
Every entry of \eqref{ncs:eq:original-matrix} lies in
$\mathbb C_n$, so associative matrix multiplication is legitimate
throughout this calculation. Define
\begin{equation}
 \Pi_j(n)=\frac13\left(I+\omega_n^{-j}T
                               +\omega_n^{-2j}T^2\right),
 \qquad j=0,1,2.
 \label{ncs:eq:fourier-projectors}\tag{NCS11}
\end{equation}
These matrices are Hermitian: $T^*=T^2$,
$\bar\omega_n=\omega_n^{-1}$, and the two coefficients are
conjugates. Each trace is one. To verify all projector identities,
use
\[
 \sum_{k=0}^2\omega_n^{mk}=
 \begin{cases}3,&3\mid m,\\0,&3\nmid m.\end{cases}
\]
Expanding \eqref{ncs:eq:fourier-projectors} gives
\begin{equation}
 \Pi_i\Pi_j=\delta_{ij}\Pi_j,\qquad
 \sum_{j=0}^2\Pi_j=I,\qquad T\Pi_j=\omega_n^j\Pi_j.
 \label{ncs:eq:projector-identities}\tag{NCS12}
\end{equation}
Furthermore $Y=\sqrt2(aT+\bar aT^2)$, whence
\begin{equation}
 \begin{aligned}
 Y\Pi_j&=\sqrt2(a\omega_n^j+\bar a\omega_n^{2j})\Pi_j
            =\nu_j\Pi_j,\\
 X&=\sum_{j=0}^2(\mu+\nu_j)\Pi_j.
 \end{aligned}
 \label{ncs:eq:spectral-decomposition}\tag{NCS13}
\end{equation}
Multiplication in $\mathbb C_n$ gives exactly the three signs in
\eqref{ncs:eq:labelled-eigenvalues}. Thus
\eqref{ncs:eq:fourier-projectors} gives an actual map into the
ordered Jordan-frame space $F_4/\operatorname{Spin}(8)$.

When the eigenvalues are pairwise distinct, these are also the
intrinsic spectral projectors:
\begin{equation}
 \Pi_j=
 \frac{(Y-\nu_kI)\circ(Y-\nu_lI)}
      {(\nu_j-\nu_k)(\nu_j-\nu_l)},
 \qquad\{j,k,l\}=\{0,1,2\}.
 \label{ncs:eq:intrinsic-projectors}\tag{NCS14}
\end{equation}
Substitution of \eqref{ncs:eq:spectral-decomposition}, together
with \eqref{ncs:eq:projector-identities}, proves this formula.
Its denominators are retained; their exceptional loci are
computed below.

The labelled frame retains the direction $n$. Specifically,
\begin{equation}
 (\Pi_1)_{12}=-\frac16-\frac{\sqrt3}{6}n,
 \qquad n=-2\sqrt3\,\operatorname{Im}(\Pi_1)_{12}.
 \label{ncs:eq:directional-inverse}\tag{NCS15}
\end{equation}
Together with
$a_0=\nu_0/(2\sqrt2)$,
$t=(\nu_2-\nu_1)/(2\sqrt6)$, and $v=tn$, this recovers the
original quaternion from the labelled spectrum and frame. The
scalar parameter is recovered from $\operatorname{tr}X/3$.
Thus the matrix and frame lift has an exact inverse on its
labelled image. The invariant triple
\eqref{ncs:eq:full-invariants} alone loses $n$; the lift
\eqref{ncs:eq:fourier-projectors} records it explicitly.

At $t=0$, one has $Y=\sqrt2a_0(T+T^2)$.
The rank-one projector $\Pi_0=(I+T+T^2)/3$ remains fixed, and
$I-\Pi_0$ is the rank-two projector of the repeated eigenvalue
$-\sqrt2a_0$ when $a_0\ne0$. No direction $n$ is assigned to a
zero imaginary part. At $a=0$, the matrix is the scalar $\mu I$.

\subsection{Singular loci and the exact Cartan angle}
\label{ncs:sec:singular-angle}

The discriminant of the original polynomial
\eqref{ncs:eq:characteristic-polynomial}, with its full scalar
factor, is
\begin{equation}
 \boxed{
 \prod_{i<j}(\nu_i-\nu_j)^2
       =864(Q^3-P^2)
       =864t^2(3a_0^2-t^2)^2.
 }
 \label{ncs:eq:discriminant}\tag{NCS16}
\end{equation}
Indeed expansion of
$(a_0^2+t^2)^3-a_0^2(a_0^2-3t^2)^2$ gives
$t^2(3a_0^2-t^2)^2$. The first equality also follows by
multiplying the three differences in
\eqref{ncs:eq:labelled-eigenvalues}. The spectrum is simple
exactly when $t>0$ and $t^2\ne3a_0^2$.

For nonzero $a$, use the Cartan polynomial of
\eqref{ocb:eq:cartan-polynomial} in the original trace metric.
Equations~\eqref{ncs:eq:two-metrics}--\eqref{ncs:eq:cubic-determinant}
give
\begin{equation}
 \begin{gathered}
 C(Y)=3\sqrt6\det Y=24\sqrt3P,\qquad r^2=12Q,\\
 \boxed{C(Y)=r^3\cos3\theta,\qquad
 \theta=\frac13\arccos\left(\frac{P}{Q^{3/2}}\right)
                              \in[0,\pi/3].}
 \end{gathered}
 \label{ncs:eq:cartan-angle}\tag{NCS17}
\end{equation}
The discriminant identity proves that the arccos argument belongs
to $[-1,1]$. This is an identity for $C$ with the whole amplitude
retained: $Y$ still has radius $r=\sqrt{12Q}$, and its scalar-shifted
determinant is still \eqref{ncs:eq:full-invariants}.

To link this to the quaternionic argument, define
$\varphi\in[0,\pi]$ by
\[
 a_0=\sqrt Q\cos\varphi,\qquad t=\sqrt Q\sin\varphi.
\]
When $t>0$, the inverse is
$a=\sqrt Q(\cos\varphi+n\sin\varphi)$, retaining
$Q,\varphi,n$. The identity
$4\cos^3\varphi-3\cos\varphi=\cos3\varphi$ gives
\begin{equation}
 \cos3\theta=\cos3\varphi.
 \label{ncs:eq:quaternion-angle}\tag{NCS18}
\end{equation}
The ordered Cartan chamber thus folds the quaternionic angular
variable by this explicit triple-angle map. The labelled frame,
\eqref{ncs:eq:labelled-eigenvalues}, and
\eqref{ncs:eq:directional-inverse} retain the inverse information.
This quaternionic angle is defined by
\eqref{ncs:eq:quaternion-coordinates}; its definition makes no
identification with a cusp angular coordinate.

For $t>0$, the decreasing orders of the labelled spectrum are
\begin{equation}
 \begin{array}{c|c}
 a_0>t/\sqrt3&\nu_0>\nu_2>\nu_1\\
 -t/\sqrt3<a_0<t/\sqrt3&\nu_2>\nu_0>\nu_1\\
 a_0<-t/\sqrt3&\nu_2>\nu_1>\nu_0.
 \end{array}
 \label{ncs:eq:ordered-chambers}\tag{NCS19}
\end{equation}
The differences
$\nu_0-\nu_2=\sqrt2(3a_0-\sqrt3t)$ and
$\nu_0-\nu_1=\sqrt2(3a_0+\sqrt3t)$ prove each row.
At $a_0=t/\sqrt3$, precisely $\nu_0=\nu_2$; at
$a_0=-t/\sqrt3$, precisely $\nu_0=\nu_1$. Thus these
rows specify how to order the projectors before applying
\eqref{ocb:eq:orbit-map}. At the two double-eigenvalue loci,
the formulas \eqref{ncs:eq:fourier-projectors} remain defined,
and \eqref{ncs:eq:projector-identities}--
\eqref{ncs:eq:spectral-decomposition} give repeated-eigenspace
projectors $\Pi_0+\Pi_2$ and $\Pi_0+\Pi_1$, respectively.
Their complementary simple-eigenvalue projectors are $\Pi_1$
and $\Pi_2$. This follows directly by grouping the equal
coefficients in \eqref{ncs:eq:spectral-decomposition}; no
vanishing denominator in \eqref{ncs:eq:intrinsic-projectors}
is used.

\subsection{Actual integral witnesses in the orbit family}
\label{ncs:sec:integral-witnesses}

The image proof of \cref{d4i:thm:complete-image} supplies three
original-coordinate lattice lines:
\begin{equation}
 \begin{gathered}
 L_\varepsilon(h)=(-3h-\varepsilon,h,h,h+\varepsilon),
       \qquad\varepsilon=\pm1,\\
 L_0(h)=(-3h,h+1,h-1,h),\qquad h\in\mathbb Z.
 \end{gathered}
 \label{ncs:eq:integral-lines}\tag{NCS20}
\end{equation}
Their coordinate sums are zero. Direct expansion gives
\begin{equation}
 \begin{array}{c|c|c}
 a&P(a)&Q(a)\\\hline
 L_\varepsilon(h)&6h+2\varepsilon&12h^2+8\varepsilon h+2\\
 L_0(h)&18h&12h^2+2.
 \end{array}
 \label{ncs:eq:line-values}\tag{NCS21}
\end{equation}
The eigenvalues, projectors, full radius, and scalar-shifted
determinant of each actual common-lattice vector are therefore
given by \eqref{ncs:eq:full-invariants}--
\eqref{ncs:eq:cartan-angle}, evaluated at
\eqref{ncs:eq:line-values} and its original imaginary vector $v$.
The Cartan angles satisfy
\begin{equation}
 \begin{aligned}
 \cos3\theta_\varepsilon(h)
    &=\frac{6h+2\varepsilon}
                  {(12h^2+8\varepsilon h+2)^{3/2}},\\
 \cos3\theta_0(h)&=\frac{18h}{(12h^2+2)^{3/2}}.
 \end{aligned}
 \label{ncs:eq:line-angles}\tag{NCS22}
\end{equation}
As $|h|\to\infty$, both ratios tend to zero: the numerators
have degree one, while their positive denominators grow as
$|h|^3$. Hence the angles tend to $\pi/6$, and simultaneously
$r^2=12Q\to\infty$. The original matrices and their radii
remain those in \eqref{ncs:eq:original-matrix}--
\eqref{ncs:eq:two-metrics}.

On the shared direction $a=h(-3,1,1,1)$ one has $P=0$ and
$Q=12h^2$. For $h\ne0$, the decreasing spectrum is exactly
\begin{equation}
 (6\sqrt2|h|,\ 0,\ -6\sqrt2|h|),\qquad
 r=12|h|,\qquad\theta=\pi/6.
 \label{ncs:eq:zero-cubic-direction}\tag{NCS23}
\end{equation}
The discriminant \eqref{ncs:eq:discriminant} is strictly positive,
so this zero-cubic direction lies on a regular Cartan orbit.
The spectrum of $X(\mu,a)$ is obtained by adding $\mu$ to each
displayed eigenvalue.

\paragraph{Sources and scope of the bridge.}
The literal-intersection theorem \cref{wcl:thm:literal-intersection}
supplies the marked Niemeier/Leech and Wilson membership proofs for
\eqref{ncs:eq:lattice-embedding}--\eqref{ncs:eq:niemeier-coordinates};
the mixed cubic is \eqref{wcl:eq:common-cubic}. Wilson's construction
is \cite{esn:Wilson2009}. The full invariant quotient, retained
mean and radius, ordered spectral diffeomorphism, and focal maps are
proved in \cref{ocb:sec:full-family}, using the flag geometry of
\cite[Sections~2.3--2.6]{MareWillems2013} and
\cite[Proposition~2.8.2 and Theorems~2.7.4, 2.8.3]{Yokota2009}.
The complete image theorem \cref{d4i:thm:complete-image} proves
\[
 P(D_4)=18\mathbb Z\cup(6\mathbb Z+2)\cup(6\mathbb Z-2)
\]
using exactly \eqref{ncs:eq:integral-lines}. This section provides
its matrix-valued and frame-valued lift with an inverse on the
labelled image; it makes no assertion that every Cartan orbit meets
this rank-four lattice. The construction is independent of the
assertion that the original completed complex threefold is
diffeomorphic to $S^6$. No interacting quantum Hamiltonian is
specified here. The retained cusp bundle has its separate explicit
triality-subgroup map; its source angular coordinate and the
quaternionic angular coordinate above retain their respective
definitions.


% Shared mathematical module. Input exactly once in each host document.
% No document class, package loads, macros, or theorem environments are defined.
% Requirements and bibliography entries: INTEGRATION_MAP.md.
% Source locators and snapshot provenance: ES_NIEMEIER_PROVENANCE.md.
% All labels and citation keys introduced here use the esn: namespace.

\section{The modulo-840 lattice, triality, and exact shell maps}
\label{esn:sec:shared}

This section constructs the original modulo-$840$ lattice $N$ with
root system $A_5^4D_4$, its full integral isometry group, and the exact
maps from denominator shells and finite prime samples to the retained
residue coordinates. The original data and source labels are recorded in
\cite{esn:Provenance}. The shared continuation proves the index-six Leech
neighbour $\Lambda$ and the marked isometry $\mathcal L$ into the three
octonion sectors. In the specified Cayley frame, put
$\tau(q)=\operatorname{Re}((q_1q_2)q_3)$ and $F=\tau\circ\mathcal L$;
the Albert determinant is $\lambda^3-\lambda(y,y)+2F(y)$.
The value-generated groups have quotient
$\mathcal V_F(\Lambda)/\mathcal V_F(N)\cong\mathbb Z/12$
(Theorem~\ref{nca:thm:value-modules}). The three residue-cycled neighbours
meet in $J$, with determinant $81$ and index $9$ in each; the fixed
fourth extension has exactly $72$ roots of type $A_2^{12}$
(Theorems~\ref{esrf:thm:common-J} and~\ref{esrf:thm:fourth-extension}).
The common lattice and this fourth extension have the exact same
value-generated group
$\frac1{27}\mathbb Z\oplus\frac{\sqrt2}{216}\mathbb Z
\oplus\frac{\sqrt3}{18}\mathbb Z$.
The inclusions from $\mathcal V_F(N)$ and into $\mathcal V_F(\Lambda)$
have matrices $\operatorname{diag}(1,4,1)$ and
$\operatorname{diag}(1,1,3)$ and quotients $\mathbb Z/4$, $\mathbb Z/3$
(Theorems~\ref{jmv:thm:JM-values} and~\ref{jmv:thm:chain}).
Wilson's construction \cite{esn:Wilson2009} is realized with both metrics
and the direction $c_{\rm W}=c^{-1}$ retained. Its fixed Gram matrix is
$3G_{E_8}$, its complement is explicitly isometric to $A_2\otimes E_8$,
and the index is $3^8$, with quotient $\mathbb F_3^8$. Its literal
intersections with $\mathcal L(N)$ and $\mathcal L(\Lambda)$ are the
same rank-four lattice of Gram matrix $6G_{D_4}$, minimum squared norm
$12$ and value-generated subgroup $4\sqrt2\mathbb Z$ for $\tau$
($8\sqrt2\mathbb Z$ for $2\tau$). This is not the actual image:
$12\sqrt2$ is not attained by $\tau$ on that intersection.
All coordinates, multiplication conventions, domain lattices and maps
are specified in the complete proofs below.

\subsection{The finite multiplicative carrier and its four fibres}
\label{esn:subsec:residues}

For an odd prime \(q\) and a unit \(r\) modulo \(q\), write
\(\left(\frac rq\right)=1\) when \(r\) is a square modulo \(q\),
and \(-1\) otherwise.  Let
\begin{equation}
 \begin{split}
 V_{840}&=\{r\in(\mathbb Z/840\mathbb Z)^\times:
                                      r\equiv1\pmod {24}\},\\
 \beta(r)&=\left(\frac{1-\left(\frac r5\right)}2,
                  \frac{1-\left(\frac r7\right)}2\right)
                          \in\mathbb F_2^2,\\
 S_{840}&=\ker\beta,\qquad F_\eta=\beta^{-1}(\eta)
                                      \quad(\eta\in\mathbb F_2^2).
 \end{split}
 \label{esn:eq:residue-carrier}
\end{equation}
Multiplication is used in \(V_{840}\) and addition in \(\mathbb F_2^2\).

\begin{theorem}[The marked cyclic fibre]
\label{esn:thm:cyclic-fibre}
The map \(\beta\) is an onto group homomorphism, each \(F_\eta\) has
six elements, and
\begin{equation}
 \begin{split}
 S_{840}&=\{1,121,169,289,361,529\}=\langle289\rangle,\\
 (289^j)_{j=0}^5&=(1,289,361,169,121,529)\pmod {840}.
 \end{split}
 \label{esn:eq:ordered-fibre}
\end{equation}
Each fibre is a torsor under multiplication by \(S_{840}\).  In
particular the marking
\begin{equation}
 \kappa:S_{840}\xrightarrow{\ \cong\ }\mathbb Z/6\mathbb Z,
 \qquad r=289^{\kappa(r)}\pmod {840},
 \label{esn:eq:kappa}
\end{equation}
is an isomorphism of groups.
\end{theorem}

\begin{proof}
The Chinese remainder map modulo \(24,5,7\), whose moduli are pairwise
coprime, identifies \(V_{840}\) with
\((\mathbb Z/5)^\times\times(\mathbb Z/7)^\times\).  The powers of
\(2\) modulo \(5\) are \(1,2,4,3\), and the powers of \(3\) modulo
\(7\) are \(1,3,2,6,4,5\); these prove the respective cyclic orders
four and six.  In a cyclic group of either even order, the squares are
exactly the even powers.  Therefore each Legendre character is the
homomorphism recording exponent parity and is onto.  Their product is
onto \(\mathbb F_2^2\).  Its kernel has
\((4/2)(6/2)=6\) elements and every fibre has that same cardinality.

Successive multiplication of the six displayed residues by \(289\)
gives the next entry and then returns from \(529\) to \(1\).
They are distinct.  Moreover \(289\equiv1\pmod {24}\),
\(289\equiv4\pmod5\), and \(289\equiv2\pmod7\), so it belongs to
the simultaneous square kernel.  Its six powers therefore exhaust that
kernel.  If \(r,s\in F_\eta\), then \(sr^{-1}\in S_{840}\) and
\(s=(sr^{-1})r\), uniquely.  This proves the torsor assertion and all
claims about \(\kappa\).
\end{proof}

Fix the coordinate order
\[
 (\eta_1,\eta_2,\eta_3,\eta_4)=(00,10,01,11)
\]
and choose one origin \(f_i\in F_{\eta_i}\) in each fibre, with
\(f_1=1\).  The six ordered elements of its fibre are
\(f_i289^j\), \(0\le j\le5\).  These origins are part of the
coordinate identifications used below; no selection of them is inferred
from an individual denominator certificate.

\subsection{The root lattice and its exact discriminant form}
\label{esn:subsec:root-and-discriminant}

For a finite set \(F\), give the free abelian group
\(\mathbb Z[F]=\bigoplus_{r\in F}\mathbb Ze_r\) the metric
\((e_r,e_s)=1\) for \(r=s\) and zero otherwise.  Put
\[
 A(F)=\left\{\sum_{r\in F}n_re_r:\sum_{r\in F}n_r=0\right\}.
\]
The ordered fibre coordinates identify each \(A(F_{\eta_i})\)
isometrically with
\[
 A=\{(n_0,\ldots,n_5)\in\mathbb Z^6:\textstyle\sum_j n_j=0\}.
\]
In a separate Euclidean space with orthonormal basis
\(e_1,e_2,e_3,e_4\), define
\begin{equation}
 D=\{x\in\mathbb Z^4:x_1+x_2+x_3+x_4\equiv0\pmod2\},
 \qquad R_{\mathrm{ES}}=
       \left(\bigoplus_{i=1}^4 A(F_{\eta_i})\right)\mathbin{\perp}D.
 \label{esn:eq:root-lattice}
\end{equation}
All five summands are orthogonal with precisely their displayed metrics.
The symbol \(R_{\mathrm{ES}}\) denotes this lattice; the unadorned
\(R\) later used for an integer shell modulus has a different declared
domain.

For a full integral lattice \(L\) in its real span, let
\[
 L^\vee=\{x:(x,L)\subset\mathbb Z\},\qquad D_L=L^\vee/L.
\]
If \(L\) is even integral, its discriminant quadratic form is
\(q_L(x+L)=(x,x)\pmod {2\mathbb Z}\).  This is well-defined because
replacing \(x\) by \(x+\ell\) changes its squared norm by
\(2(x,\ell)+(\ell,\ell)\in2\mathbb Z\).

\begin{theorem}[Root metric and marked discriminant]
\label{esn:thm:root-discriminant}
The lattice \(R_{\mathrm{ES}}\) is even integral of rank 24 and has
Gram matrix
\[
 G_{R_{\mathrm{ES}}}=G_A^{\oplus4}\oplus G_D,
\]
where
\[
 G_A=\begin{pmatrix}
 2&-1&0&0&0\\-1&2&-1&0&0\\0&-1&2&-1&0\\
 0&0&-1&2&-1\\0&0&0&-1&2
 \end{pmatrix},\qquad
 G_D=\begin{pmatrix}
 2&-1&0&0\\-1&2&-1&-1\\0&-1&2&0\\0&-1&0&2
 \end{pmatrix}.
\]
Their determinants are 6 and 4, respectively, and
\(\det R_{\mathrm{ES}}=6^4\cdot4=5184\).

In each identified \(A\) summand use
\(\lambda=(5,-1,-1,-1,-1,-1)/6\) as the discriminant generator.
For \(D\), use the representatives
\begin{equation}
 v=(1,0,0,0),\qquad s=\tfrac12(1,1,1,1),\qquad
 c=\tfrac12(-1,1,1,1)
 \label{esn:eq:d4-representatives}
\end{equation}
and the exact marking \(10\mapsto v+D\), \(01\mapsto s+D\),
\(11\mapsto c+D\).  Then
\begin{equation}
 \begin{split}
 D_{R_{\mathrm{ES}}}&=(\mathbb Z/6\mathbb Z)^4\oplus\mathbb F_2^2,\\
 q_{R_{\mathrm{ES}}}(x_1,x_2,x_3,x_4;y)
   &=\frac56\sum_{i=1}^4x_i^2+\mathbf1_{y\ne0}
                                              \pmod {2\mathbb Z}.
 \end{split}
 \label{esn:eq:full-discriminant}
\end{equation}
\end{theorem}

\begin{proof}
For a six-element ordered fibre the five differences
\(e_j-e_{j+1}\), \(0\le j\le4\), form a basis of its augmentation
lattice: the coefficients of a zero-sum vector in this basis are its
first five successive partial sums.  Dot products give \(G_A\).
The determinant \(d_n\) of the analogous tridiagonal \(n\)-by-\(n\)
matrix satisfies \(d_n=2d_{n-1}-d_{n-2}\), \(d_0=1,d_1=2\), by
first-row expansion; induction gives \(d_n=n+1\).
The four vectors
\(e_1-e_2,e_2-e_3,e_3-e_4,e_3+e_4\) lie in \(D\) and their
coordinate determinant has absolute value 2.  Their integer span
therefore has index 2 in \(\mathbb Z^4\).  The parity homomorphism
\(\mathbb Z^4\to\mathbb F_2\), \(x\mapsto\sum x_i\pmod2\), is
onto and has kernel \(D\), so this span equals \(D\).
Their dot products give \(G_D\), whose determinant is the square of
the coordinate determinant, namely 4.  Integer coordinates make all
pairings integral.  In either lattice the squared norm is congruent
modulo two to the coordinate sum, hence is even.  Rank and determinant
now follow from the orthogonal direct sum.

The vector \(\lambda\) has sum zero and integral pairing with each
augmentation root.  Its norm is \(5/6\), and \(k\lambda\in A\)
holds exactly when \(6\) divides \(k\).  Thus it has exact order six
in \(D_A\).  In an integral basis of an integral lattice with Gram
matrix \(G\), the dual coordinates are \(G^{-1}\mathbb Z^n\);
therefore \(|L^\vee/L|=|\det G|\).  Since \(|D_A|=6\), this
element generates the full discriminant group and gives its stated form.

A vector pairing integrally with all \(e_i-e_j\) has common
fractional parts in its four coordinates.  Pairing also with
\(e_i+e_j\) forces this common fractional part to be either zero or
one half.  Conversely every entirely integral or entirely half-integral
vector pairs integrally with every root of \(D\).  There are two
integer classes according to sum parity, and two half-integer classes
according to the parity of their difference from \(s\).  The four
representatives \(0,v,s,c\) are therefore all the classes, and their
squared norms are \(0,1,1,1\).  Their addition is as marked because
\(v+s-c=(2,0,0,0)\in D\).  The dual and quadratic form of an
orthogonal sum are the direct sums of those of its summands, which proves
\eqref{esn:eq:full-discriminant}.
\end{proof}

\subsection{The index-72 glue and all glue-class minima}
\label{esn:subsec:glue}

Define the binary and ternary codes
\begin{equation}
 \begin{split}
 C_2&=\operatorname{span}_{\mathbb F_2}
       \{(1111;00),(1100;10),(1010;01)\},\\
 C_3&=\operatorname{span}_{\mathbb F_3}
       \{z_1=(1,1,1,0),\ z_2=(1,2,0,1)\}.
 \end{split}
 \label{esn:eq:codes}
\end{equation}
The semicolon distinguishes the four binary A5 coordinates from the
two D4 discriminant coordinates.  Let
\begin{equation}
 \begin{split}
 \mathcal H&=\{(3u+2z;y):(u;y)\in C_2, z\in C_3\}
                                  \subset D_{R_{\mathrm{ES}}},\\
 N&=\{x\in R_{\mathrm{ES}}^\vee:
                                x+R_{\mathrm{ES}}\in\mathcal H\}.
 \end{split}
 \label{esn:eq:glue-and-overlattice}
\end{equation}
Here \(3u+2z\) is computed coordinatewise modulo six using the
integer representatives of \(u\) and \(z\).

\begin{theorem}[Even unimodular completion with exactly 144 roots]
\label{esn:thm:niemeier}
The subgroup \(\mathcal H\) is isotropic of order 72.  Its inverse
image \(N\) is a positive-definite even unimodular lattice of rank 24,
with \([N:R_{\mathrm{ES}}]=72\).  Among the 71 nonzero glue classes,
46 have minimum squared norm 4 and 25 have minimum squared norm 6.
The roots of \(N\) are exactly those of \(R_{\mathrm{ES}}\), with
root system \(A_5^4D_4\): 30 roots in each A5 summand and 24 in D4.
\end{theorem}

\begin{proof}
The binary generators are independent, hence give eight words.  Listing
their sums shows that the seven nonzero words comprise one word
\((1111;00)\) and six words whose first part has weight two and whose
second part is nonzero.  A ternary word is
\[
 az_1+bz_2=(a+b,a+2b,a,b).
\]
Its last two coordinates prove injectivity of the parameter pair
\((a,b)\), so there are nine words.  If \(a=0\ne b\) the third
coordinate alone vanishes; if \(b=0\ne a\) the fourth alone vanishes;
and if both are nonzero exactly one of \(a+b,a+2b\) vanishes.
Thus every nonzero word has weight three, and each of the four
coordinates is the unique vanishing coordinate for two words.

The map \((u_i,z_i)\mapsto3u_i+2z_i\pmod6\) is bijective:
reduction modulo two recovers \(u_i\), and reduction modulo three
recovers \(2z_i\), from which multiplication by 2 recovers \(z_i\).
It is a homomorphism of additive groups, since changing binary or
ternary representatives by their moduli changes \(3u_i+2z_i\) by a
multiple of six.  Its product with the retained binary \(y\)-coordinate
therefore maps \(C_2\times C_3\) injectively onto the subgroup
\(\mathcal H\).  It follows that \(|\mathcal H|=8\cdot9=72\).  For
\(u_i\in\{0,1\}\) and \(z_i\in\{0,1,2\}\),
\[
 \frac56(3u_i+2z_i)^2
 =\frac{15}{2}u_i+10u_iz_i+\frac{10}{3}z_i^2
 \equiv\frac32u_i+\frac43z_i^2\pmod {2\mathbb Z}.
\]
The cross-term is even; a nonzero \(z_i\) has square 1 or 4, whose
contributions in the last expression differ by 4.  The binary
contribution including D4 is therefore zero, 6, or \(3+1=4\), and
the ternary contribution is zero or \((4/3)3=4\) modulo
\(2\mathbb Z\).  This proves isotropy.

The inverse image is a lattice because it is a subgroup between the
full lattices \(R_{\mathrm{ES}}\) and \(R_{\mathrm{ES}}^\vee\).
Isotropy gives \((x,x)\in2\mathbb Z\) for every \(x\in N\), and
\[
 (x,y)=\frac{(x+y,x+y)-(x,x)-(y,y)}2\in\mathbb Z
                                    \qquad(x,y\in N).
\]
Thus it is even integral.  Its index is \(|\mathcal H|=72\).
For a finite-index inclusion of equal-rank lattices, the covolume
changes by that index, and the Gram determinant is the square of the
covolume in orthonormal coordinates.  Hence
\[
 \det N=\frac{6^4\cdot4}{72^2}=1.
\]
The dual-index formula proved above implies \(N=N^\vee\).
Its rank and positive definiteness are inherited from the unchanged real
span and metric.

For \(0\le k\le5\), every vector in the A5 class \(k\lambda+A\)
has the form
\[
 m-\frac{k}{6}(1,1,1,1,1,1),\qquad
 m\in\mathbb Z^6,\quad\sum m_i=k.
\]
Its squared norm is \(\sum m_i^2-k^2/6\).  Since
\(m_i^2\ge m_i\) for every integer \(m_i\), the minimum is at least
\(k-k^2/6\).  It is attained by choosing exactly \(k\) of the six
entries equal to one and the others zero.  Thus
\begin{equation}
 \mu_A(k)=\frac{k(6-k)}6.
 \label{esn:eq:a5-coset-minimum}
\end{equation}
Every nonzero integral D4 class has minimum one, attained by \(v\);
every half-integral vector has norm at least
\(4(1/2)^2=1\), attained by \(s,c\) in their respective classes.
The four A5 minima and the D4 minimum add and are simultaneously attained
because the metric is an orthogonal sum.

For a coordinate with \(k_i=3u_i+2z_i\pmod6\), its minimum
contribution is 0 for \((u_i,z_i)=(0,0)\), \(3/2\) for
\((1,0)\), \(4/3\) for \((0,z_i\ne0)\), and \(5/6\) for
\((1,z_i\ne0)\).  The complete nonzero case table is
\begin{equation}
\begin{array}{c|c|c|r|l}
 \operatorname{wt}(u)&\operatorname{wt}(z)&
 |\operatorname{supp}(u)\cap\operatorname{supp}(z)|&
 \text{classes}&\text{minimum squared norm}\\\hline
 2&0&0&6&2(3/2)+1=4\\
 4&0&0&1&4(3/2)=6\\
 0&3&0&8&3(4/3)=4\\
 2&3&1&24&(3/2)+2(4/3)+(5/6)+1=6\\
 2&3&2&24&(4/3)+2(5/6)+1=4\\
 4&3&3&8&(3/2)+3(5/6)=4.
\end{array}
 \label{esn:eq:glue-minimum-table}
\end{equation}
For each of the six weight-two binary words, four nonzero ternary
words omit a coordinate inside its support, and four omit a coordinate
outside.  This proves the two counts \(6\cdot4=24\); all other counts
follow from the enumerated binary and ternary words.  Thus the total
number at minimum 4 is \(6+8+24+8=46\), and the total at minimum 6 is
\(1+24=25\).  Their sum is 71, leaving exactly the zero glue class.

No nonzero glue class can contain a norm-two vector.  In the zero
class, an A5 norm-two vector has exactly one coordinate 1 and another
\(-1\), so is \(e_i-e_j\), with \(6\cdot5=30\) possibilities.
A D4 norm-two vector has exactly two coordinates of absolute value one,
so is \(\pm e_i\pm e_j\), with \(4\binom42=24\) possibilities.
An orthogonal sum of two nonzero root-lattice vectors has norm at least
4 because each summand is positive even.  These are therefore all
\(4\cdot30+24=144\) roots of \(N\).
\end{proof}

A Niemeier lattice means a positive-definite even unimodular lattice of
rank 24.  The theorem constructs one with root system \(A_5^4D_4\)
directly.  Its root type appears with lambency 6 in the primary umbral
table \cite[Table~3, p.~21]{esn:CDH2014}.  This table reference is not
needed to prove the inverse-image construction or any minimum above.

For comparison with the earlier labeled glue presentation in the source
\cite{esn:ESMain2026,esn:Provenance}, put
\[
 h_1=(1,3,1,1;00),\qquad h_2=(1,4,2,3;11),\qquad
 h_3=(3,0,3,0;01),
\]
and let \(\mathcal H_{\mathrm{old}}\) be their subgroup.  The exact
transport is
\begin{equation}
 P(x_1,x_2,x_3,x_4;y)=(x_3,x_4,x_1,x_2;y),\qquad
 P(\mathcal H_{\mathrm{old}})=\mathcal H.
 \label{esn:eq:old-code-transporter}
\end{equation}
Indeed, if \(rh_1+th_2+wh_3=0\) with \(r,t\in\mathbb Z/6\)
and \(w\in\mathbb Z/2\), the final two coordinates imply that \(t\)
is even and \(w=0\).  The first and third coordinates then give
\(r+t=0\) and \(r+2t=0\) modulo six, whence \(t=r=0\).
Thus \(|\mathcal H_{\mathrm{old}}|=6\cdot6\cdot2=72\).  Directly,
\[
\begin{split}
 P(h_1)&=3(1111;00)+2(2,2,2,0;00),\\
 P(h_2)&=3(0110;11)+2(1,0,2,2;00),\\
 P(h_3)&=3(1010;01).
\end{split}
\]
The displayed ternary parts are \(2z_1\) and \(2z_1+2z_2\).
Consequently \(P(\mathcal H_{\mathrm{old}})\subset\mathcal H\),
and the equal orders prove equality.  This transporter is realized by
the corresponding isometry permuting the full identified A5 summands.

\subsection{The common D4 root lattice, shortest weights, and triality}
\label{esn:subsec:d4-flag}

Use orthonormal coordinates \(\varepsilon_1,\ldots,\varepsilon_4\)
for the ordinary F4 root realization.  Its long roots and the three
parts of its short roots are
\begin{equation}
\begin{split}
 \Phi_D&=\{\pm\varepsilon_i\pm\varepsilon_j:1\le i<j\le4\},\\
 \mathcal V&=\{\pm\varepsilon_i:1\le i\le4\},\\
 \mathcal S_\pm&=
 \left\{\frac12\sum_{i=1}^4\delta_i\varepsilon_i:
          \delta_i\in\{1,-1\},\ \prod_i\delta_i=\pm1\right\},\\
 \Phi_F&=\Phi_D\cup\mathcal V\cup\mathcal S_+\cup\mathcal S_-.
\end{split}
 \label{esn:eq:root-weight-sets}
\end{equation}
For a nonzero vector \(a\), reflection means the exact map
\(s_a(x)=x-2(x,a)a/(a,a)\).  Write \(W_D\) and \(W_F\) for the
groups generated by the reflections in \(\Phi_D\) and \(\Phi_F\),
respectively.

\begin{theorem}[Integral D4 and the marked triality quotient]
\label{esn:thm:d4-f4-triality}
The map
\begin{equation}
 \iota_D:D\longrightarrow\mathbb Z\langle\Phi_D\rangle,
 \qquad(x_1,x_2,x_3,x_4)\longmapsto
                         \sum_{i=1}^4x_i\varepsilon_i
 \label{esn:eq:d4-isometry}
\end{equation}
is an integral isometry onto the long-root lattice.  It extends to duals
and discriminant forms, sending the shortest-vector sets of the marked
classes \(v,s,c\) respectively to
\(\mathcal V,\mathcal S_+,\mathcal S_-\).

The two matrices
\begin{equation}
 R_D=\operatorname{diag}(-1,1,1,1),\qquad
 H_D=\frac12\begin{pmatrix}
 1&1&1&1\\1&1&-1&-1\\1&-1&1&-1\\1&-1&-1&1
 \end{pmatrix}
 \label{esn:eq:d4-triality-matrices}
\end{equation}
are the short-root reflections \(s_{\varepsilon_1}\) and \(s_c\),
where \(c=(-1,1,1,1)/2\) is read in the \(\varepsilon\) coordinates.
They induce \((s\ c)\) and \((v\ s)\) on the three nonzero
discriminant classes.  There is a split exact sequence
\begin{equation}
 1\longrightarrow W_D\longrightarrow W_F
   \longrightarrow S_3\longrightarrow1,
 \qquad W_D\cong\{\text{even signed permutations of four coordinates}\},
 \label{esn:eq:weyl-triality-sequence}
\end{equation}
whose displayed section is \(\langle R_D,H_D\rangle\cong S_3\).
In particular \(|W_D|=192\) and \(|W_F|=1152\).
\end{theorem}

\begin{proof}
The four D4 basis vectors in Theorem~\ref{esn:thm:root-discriminant}
are among \(\Phi_D\), and every vector of \(\Phi_D\) has even
coordinate sum.  Thus the displayed identity of coordinates has exactly
the stated image, preserves the metric, and extends to duals by the
integer-pairing definition.  In a nonzero integral dual class, norm one
forces one coordinate \(\pm1\) and all others zero.  In a
half-integral class, norm one forces all four coordinates to be
\(\pm1/2\).  The difference from \(s\) has even coordinate sum
exactly when the number of negative signs is even.  This proves the
three shortest-vector identifications with the stated marking.

The reflection identity for \(R_D\) is immediate, and
\(H_D=I-2cc^T\) follows by multiplying the four coordinates of \(c\).
Both are orthogonal involutions.  The first preserves the parity of an
integer coordinate sum.  For \(x\in D\), the four coordinates of
\(H_Dx\) are integers because each numerator is congruent to
\(\sum x_i\) modulo two, and their sum is \(2x_1\).  Thus both
preserve \(D\).  Their effects on the representatives show that
\(R_D\) fixes the class \(v\) and exchanges \(s,c\), whereas
\(H_Dv=s\), \(H_Ds=v\), and \(H_Dc=-c\equiv c\pmod D\).
They therefore preserve all of \(\Phi_F\) by the shortest-vector
description.  On the plane with ordered basis \(\varepsilon_1,c\),
their product has matrix
\[
 \begin{pmatrix}0&-1\\1&-1\end{pmatrix},
\]
whose cube is the identity, while both reflections fix the orthogonal
complement.  Consequently \(R_D^2=H_D^2=(R_DH_D)^3=1\).
Every word in two such involutions reduces to
\((R_DH_D)^k\) or \((R_DH_D)^kR_D\), \(0\le k\le2\).
Their discriminant permutations are the six elements of \(S_3\),
so these six words are distinct and give the stated section.

Reflection in \(\varepsilon_i-\varepsilon_j\) swaps coordinates
\(i,j\).  Reflection in \(\varepsilon_i+\varepsilon_j\) swaps them
and changes both signs.  Hence the long-root reflections generate every
permutation and every even sign change, and each long-root reflection
is of this form.  This proves the assertion about \(W_D\) and its
order \(2^3\cdot4!=192\).  Such a reflection acts trivially on
\(D^\vee/D\): for a root \(a\) of squared norm two,
\(s_a(x)-x=-(x,a)a\in D\) for all \(x\in D^\vee\).

Let \(K=\langle W_D,R_D,H_D\rangle\).  Each of its generators
preserves the long roots, so \(W_D\) is normal in \(K\) by
\(g s_a g^{-1}=s_{ga}\).  The even signed permutations act
transitively on each of \(\mathcal V,\mathcal S_+,\mathcal S_-\),
and the displayed section permutes these three sets transitively.
Thus every short root is \(g\varepsilon_1\) for some \(g\in K\),
and its reflection \(gR_Dg^{-1}\) belongs to \(K\).  All generators
of \(W_F\) therefore belong to \(K\); the reverse inclusion follows
from their definitions, proving \(W_F=K\).
An element in the kernel of its permutation action on the three short
sets preserves \(\mathcal V\), so it is a signed permutation matrix.
Preservation of \(\mathcal S_+\) forces an even number of sign
changes, so the kernel is \(W_D\).  This proves exactness, splitting,
and the order \(6\cdot192=1152\).
\end{proof}

\paragraph{The exact geometric datum and weight marking.}
Let \(\mathbb O\) denote the real normed octonion algebra and let
\(J=\mathfrak h_3(\mathbb O)\) be the real vector space of Hermitian
three-by-three matrices, with Jordan product
\(a\circ b=(ab+ba)/2\).  Define
\[
 \mathbb O\mathbb P^2=\{a\in J:a^2=a,\ \operatorname{tr}a=1\},
 \qquad
 \operatorname{Fl}(\mathbb O)=
 \{(a,b)\in(\mathbb O\mathbb P^2)^2:
                                 \operatorname{Re}\operatorname{tr}(ab)=0\}.
\]
For the compact group \(F_4=\operatorname{Aut}_{\mathbb R}(J,\circ)\),
Mare and Willems prove the homogeneous identification
\(\operatorname{Fl}(\mathbb O)=F_4/\operatorname{Spin}(8)\) and
identify its tangent representation as
\(V_8\oplus S_8^+\oplus S_8^-\)
\cite[Definition~2.2, Proposition~2.4, Proposition~2.7, and
pp.~491--492]{esn:MareWillems2013}.  These are the prior geometric
results used here.  With their torus coordinate forms \(L_i\), their
equations (5.4)--(5.6) assign the vector, positive half-spin, and negative
half-spin weights respectively the patterns
\(\pm L_i\), even half-sign sums, and odd half-sign sums
\cite[pp.~510--511]{esn:MareWillems2013}.
The coordinate map \(L_i\mapsto\varepsilon_i\) retains that sign
marking and identifies these three weight sets with those of
\eqref{esn:eq:root-weight-sets}.

For completeness the half-spin weight patterns admit the following
direct calculation.  In the complex exterior algebra
\(\Lambda^*\mathbb C^4\), let \(a_i\) be exterior multiplication
by the \(i\)-th basis vector and \(b_i\) contraction by its dual.
On a basis wedge, insertion and deletion give
\[
 a_ia_j+a_ja_i=b_ib_j+b_jb_i=0,\qquad
 a_ib_j+b_ja_i=\delta_{ij}I.
\]
Thus \(C_i=a_i+b_i\) and \(D_i=\mathrm i(b_i-a_i)\) square to
one and mutually anticommute for distinct Clifford generators.
Writing \(N_i=a_ib_i\), direct multiplication gives
\[
 \tfrac12 C_iD_i=\mathrm i(N_i-\tfrac12 I).
\]
The four commuting infinitesimal torus operators therefore have weights
\(\sum_i(\mathbf1_{i\in I}-1/2)\varepsilon_i\) on the basis
wedge indexed by \(I\subset\{1,2,3,4\}\).  The even exterior
part has an even number of negative signs, and the odd part has an odd
number.  There are eight one-dimensional weight spaces in each part.
For the vector representation the four real coordinate planes
complexify into their two rotation eigenspaces, with weights
\(\varepsilon_i,-\varepsilon_i\).  This proves the displayed
weight formulas directly in the stated marking.

The map proved here is consequently the actual chain
\[
 D\xrightarrow{\ \iota_D\ }\mathbb Z\langle\Phi_D\rangle,
 \qquad
 \{v,s,c\}\longleftrightarrow
       \{\mathcal V,\mathcal S_+,\mathcal S_-\}
       =\{\operatorname{wt}(V_8),
          \operatorname{wt}(S_8^+),\operatorname{wt}(S_8^-)\}.
\]
The same \(S_3\) acts through the explicitly proved permutations of
these three sets.  The companion higher-rung treatment records the
geometric datum and root-coordinate comparison at the respective labels
\begin{center}
\texttt{h24:thm:octonionic-triality}\\
\texttt{h24:prop:two-Coxeter-systems}.
\end{center}
The exact source pointers and passages are included in
\cite{esn:Provenance}.  The integral metric map
in this section is the map on D4 and its dual.  The definitions make no
map from the whole rank-24 lattice \(N\) into the flag tangent space.

\subsection{The full indicated glue stabilizer and its triality quotient}
\label{esn:subsec:glue-gl2-triality}

All vector spaces over finite fields in this subsection retain their displayed
bases.  Write
\[
 j:\mathbb F_3^2\longrightarrow\mathbb F_3^4,
 \qquad j(a,b)=(a+b,a+2b,a,b),
 \qquad C_3=j(\mathbb F_3^2).
\]
The four row forms defining this injection are
\[
 \ell_1=(1,1),\quad \ell_2=(1,2),\quad
 \ell_3=(1,0),\quad \ell_4=(0,1).
\]
They represent all four one-dimensional subspaces of the dual of
\(\mathbb F_3^2\).  Indeed each nonzero row has either first coordinate zero,
in which case its line is represented by \(\ell_4\), or a unique scalar
multiple with first coordinate one, namely one of
\((1,0),(1,1),(1,2)\).

Let \(E\subset\mathbb F_2^4\) be the even-weight subspace and define
\[
 \tau:E\longrightarrow\mathbb F_2^2,
 \qquad \tau(u_1,u_2,u_3,u_4)
       =(u_2+u_4,u_3+u_4).
\]
The binary code in the index-72 glue is exactly
\[
 C_2=\{(u,\tau(u)):u\in E\}.
\]
To verify this, the independent vectors
\(1111,1100,1010\) form a basis of the three-dimensional space \(E\),
and their images under \(\tau\) are respectively
\(00,10,01\).  These are precisely the three stated binary code generators.
In particular \(\tau\) is onto and
\[
 \ker\tau=\langle1111\rangle.
\]
Explicitly, \(\tau(u)=0\) gives \(u_2=u_3=u_4\), and the even-weight
equation then gives \(u_1=u_4\).

For a permutation \(p\) of \(\{1,2,3,4\}\), use the coordinate convention
\[
 (P_pu)_i=u_{p(i)}.
\]
The permutation preserves \(E\) and \(\langle1111\rangle\), so it induces
a unique invertible map \(B_p\) satisfying
\begin{equation}
 B_p\tau(u)=\tau(P_pu)\qquad(u\in E).
 \label{esn:eq:binary-permutation-map}
\end{equation}
For an entirely explicit formula, the two columns of \(B_p\) are
\(\tau(P_p(1100))\) and \(\tau(P_p(1010))\).
The convention implies \(P_pP_q=P_{q\circ p}\).  Accordingly
\(B_pB_q=B_{q\circ p}\); no reversal is suppressed in what follows.

The subgroup of \eqref{esn:eq:glue-and-overlattice} can therefore be written
\[
 \mathcal H=
 \{(3u+2z;\tau(u)):u\in E,\ z\in C_3\}
 \subset(\mathbb Z/6)^4\oplus\mathbb F_2^2.
\]
Let \(G\) be its stabilizer among the transformations
\[
 g=(\varepsilon,p,B),\qquad
 (x;y)\longmapsto
 \bigl((\varepsilon_i x_{p(i)})_{i=1}^4;By\bigr),
\]
where \(\varepsilon_i\in\{1,-1\}\), \(p\in S_4\), and
\(B\in GL_2(\mathbb F_2)\).  Denote the signed permutation matrix on the
first four coordinates by \(M_g\).

\begin{theorem}[An explicit isomorphism to \(GL_2(\mathbb F_3)\)]
\label{esn:thm:glue-gl2}
Restriction to the ternary code gives an isomorphism
\[
 \Theta:G\xrightarrow{\ \cong\ }GL_2(\mathbb F_3),
 \qquad M_gj=j\Theta(g).
\]
For every \(A\in GL_2(\mathbb F_3)\), its inverse image is specified
uniquely by
\begin{equation}
 \ell_iA=\varepsilon_i\ell_{p(i)}\quad(1\le i\le4),
 \qquad B=B_p.
 \label{esn:eq:inverse-gl2-glue-map}
\end{equation}
In particular \(|G|=48\).
\end{theorem}

\begin{proof}
The elements of \(\mathcal H\) of order dividing two are precisely
\((3u;\tau(u))\), and those of order dividing three are precisely
\((2z;0)\).  Every indicated ambient transformation preserves these primary
parts, so an element of \(G\) satisfies \(M_gC_3=C_3\).  The injection
\(j\) therefore yields one and only one \(\Theta(g)\in GL_2(\mathbb F_3)\).
The binary part is preserved exactly when
\(B\tau(u)=\tau(P_pu)\), because all signs become one modulo two.
Surjectivity of \(\tau\) forces \(B=B_p\).

If \(\Theta(g)=I\), the equality \(M_gj=j\) implies
\(\varepsilon_i\ell_{p(i)}=\ell_i\) for every \(i\).
The four row lines are distinct, so \(p(i)=i\).  Since each \(\ell_i\)
is nonzero and \(1\ne-1\) in \(\mathbb F_3\), also
\(\varepsilon_i=1\) for every \(i\).  It follows that \(B_p=I\) and
\(g=1\).  Thus \(\Theta\) is injective.

Conversely fix \(A\in GL_2(\mathbb F_3)\).  The rows \(\ell_iA\) are
nonzero and belong to four distinct row lines.  Since the \(\ell_i\)
represent every row line, each \(\ell_iA\) has a unique expression
\(\varepsilon_i\ell_{p(i)}\), where the nonzero scalar in \(\mathbb F_3\)
is represented by the exact integer sign \(\varepsilon_i=1\) or \(-1\).
Distinctness of the row lines makes \(p\) a permutation.
The resulting signed permutation satisfies \(M_gj=jA\), hence preserves
\(C_3\).  Choosing \(B=B_p\) preserves \(C_2\) as proved above, hence
preserves \(\mathcal H\).  This constructs an element of \(G\) mapping
to \(A\), proving surjectivity and the stated inverse formula.

Finally
\[
 M_{gh}j=M_gM_hj=M_gj\Theta(h)
 =j\Theta(g)\Theta(h),
\]
so \(\Theta(gh)=\Theta(g)\Theta(h)\).
An invertible \(2\)-by-\(2\) matrix over \(\mathbb F_3\) has eight
choices for its first nonzero column and six choices for its second column
outside the span of the first.  Thus its group order is \(8\cdot6=48\).
\end{proof}

The map to the permutation of the four component labels is more precisely
\(g\mapsto p^{-1}\), because a vector supported on component \(k\) is
sent to component \(p^{-1}(k)\).  Its kernel under \(\Theta\) is
\(\{I,-I\}\).  In fact a matrix preserving all four row lines is diagonal
by preservation of \(\ell_3\) and \(\ell_4\), and its diagonal entries
are equal by preservation of \(\ell_1\).  Conversely scalar matrices
preserve all four lines.  Therefore the component permutation image has
order \(48/2=24\), so is all \(S_4\), with
\[
 G/\Theta^{-1}(\{\pm I\})
       \xrightarrow{\ \overline\Theta\ }PGL_2(\mathbb F_3)\cong S_4.
\]
The center of \(GL_2(\mathbb F_3)\) is exactly \(\{\pm I\}\): a
central matrix commutes with \(\operatorname{diag}(1,-1)\), forcing its
off-diagonal entries to vanish, and commuting with
\(\left(\begin{smallmatrix}0&1\\1&0\end{smallmatrix}\right)\) forces
the diagonal entries to agree.

\begin{theorem}[The exact triality quotient and its quaternion kernel]
\label{esn:thm:triality-q8}
The projection \(\pi:G\to GL_2(\mathbb F_2)\),
\((\varepsilon,p,B)\mapsto B\), is onto, and under \(\Theta\) its
kernel is the quaternion subgroup generated by
\[
 I_q=\begin{pmatrix}0&1\\2&0\end{pmatrix},
 \qquad J_q=\begin{pmatrix}1&1\\1&2\end{pmatrix}.
\]
There is a split exact sequence
\[
 1\longrightarrow Q_8\longrightarrow GL_2(\mathbb F_3)
 \xrightarrow{\ \pi\Theta^{-1}\ }GL_2(\mathbb F_2)
 \longrightarrow1.
\]
The splitting is the subgroup generated by
\[
 A_R=\begin{pmatrix}1&0\\0&2\end{pmatrix},
 \qquad A_H=\begin{pmatrix}1&2\\0&2\end{pmatrix}.
\]
It identifies \(GL_2(\mathbb F_3)\) with \(Q_8\rtimes S_3\) for the
conjugation action of this displayed subgroup on this displayed \(Q_8\).
\end{theorem}

\begin{proof}
The nonzero elements of \(E/\langle1111\rangle\) are represented by
the complementary pairs of two-element subsets
\[
 \{12\mid34\},\qquad\{13\mid24\},\qquad\{14\mid23\},
\]
marked respectively by \(10,01,11\) via \(\tau\).
Thus \(B_p\) is exactly the permutation action on these three
partitions, expressed in the displayed binary marking.

Using \eqref{esn:eq:inverse-gl2-glue-map} gives
\[
 M(A_R)x=(x_2,x_1,x_3,-x_4),\quad
 B_R=\begin{pmatrix}1&1\\0&1\end{pmatrix},
\]
\[
 M(A_H)x=(x_1,x_3,x_2,-x_4),\quad
 B_H=\begin{pmatrix}0&1\\1&0\end{pmatrix}.
\]
These binary matrices act respectively as the transpositions
\((01\ 11)\) and \((10\ 01)\), so they generate all permutations of
the three nonzero binary vectors.  There are
\((4-1)(4-2)=6\) invertible binary matrices; a matrix fixing all nonzero
vectors fixes the basis.  Hence this permutation action identifies
\(GL_2(\mathbb F_2)\) with \(S_3\) and proves surjectivity of \(\pi\).
Its kernel therefore has order eight.

For \(I_q\) the row rule gives
\[
 p=(2,1,4,3),\qquad\varepsilon=(-1,1,1,-1),
\]
and for \(J_q\) it gives
\[
 p=(3,4,1,2),\qquad\varepsilon=(-1,-1,1,1).
\]
Both permutations fix each of the three displayed complementary
partitions, so both matrices lie in the kernel.  Direct matrix
multiplication over \(\mathbb F_3\) gives
\[
 I_q^2=J_q^2=-I,\qquad I_qJ_q=-J_qI_q,
 \qquad I_qJ_q=\begin{pmatrix}1&2\\2&2\end{pmatrix}.
\]
The eight matrices
\(I,-I,\pm I_q,\pm J_q,\pm I_qJ_q\) are pairwise distinct by their
displayed entries.  The relations show this set is a group: powers reduce
using \(I_q^2=J_q^2=-I\), and factors can be reordered using
\(J_qI_q=-I_qJ_q\).  These relations identify it with the quaternion
group \(Q_8\).  Since the kernel has order eight, this is the whole kernel.

The matrices \(A_R,A_H\) satisfy
\[
 A_R^2=A_H^2=I,\qquad
 A_RA_H=\begin{pmatrix}1&2\\0&1\end{pmatrix},\qquad
 (A_RA_H)^3=I.
\]
Every word in the two involutions reduces to one of
\((A_RA_H)^k\) or \((A_RA_H)^kA_R\), \(k=0,1,2\), by alternation
and the cubic relation.  Their images under \(\pi\) are the six distinct
permutations generated by \(B_R,B_H\).  Hence the subgroup has exactly
six elements and maps isomorphically onto \(GL_2(\mathbb F_2)\).
This gives the splitting.  For any \(g\), the unique splitting element
with image \(\pi(g)\) leaves \(g\) times its inverse in the kernel;
the intersection of kernel and splitting is trivial.  Multiplication
therefore realizes the asserted semidirect product with the indicated
conjugation action.
\end{proof}

The matrices in the kernel have an additional marked interpretation:
their projective component permutations are the four permutations
\(1,(12)(34),(13)(24),(14)(23)\).  Each fixes all three partitions, and
these are exactly the kernel of the onto action \(S_4\to S_3\), whose
kernel has order \(24/6=4\).  Thus the preceding exact sequence is the
precise map from the order-48 ternary-code symmetry to the three
\(D_4\) discriminant classes.

\begin{theorem}[A faithful action on the actual index-72 lattice]
\label{esn:thm:gl2-full-lattice-lift}
Fix the four identifications of the \(A_5\) coordinate lattices used in
the glue presentation, and the discriminant marking
\(10\mapsto v,01\mapsto s,11\mapsto c\) of \(D_4\).
The group \(G\cong GL_2(\mathbb F_3)\) has a faithful action by integral
isometries of the index-72 overlattice.  Write
\(\varphi_i:A\longrightarrow A(F_{\eta_i})\) for the restriction of
the fixed free-coordinate isometry \(e_j\mapsto e_{f_i289^j}\).
For \(X_i\in A(F_{\eta_i})\), the \(i\)-th output component is exactly
\(\varepsilon_i\varphi_i\varphi_{p(i)}^{-1}X_{p(i)}\).
In the identified \(A\) coordinates this is the action
\[
 L_g(X_1,X_2,X_3,X_4;d)
 =\bigl(\varepsilon_1X_{p(1)},\ldots,
        \varepsilon_4X_{p(4)};\delta(B)d\bigr),
\]
where \(\delta:GL_2(\mathbb F_2)\to O(D_4)\) is the section given by
\[
 \delta(B_R)=R_D=\operatorname{diag}(-1,1,1,1),\qquad
 \delta(B_H)=H_D=\frac12
 \begin{pmatrix}
 1&1&1&1\\1&1&-1&-1\\1&-1&1&-1\\1&-1&-1&1
 \end{pmatrix}.
\]
Under \(\Theta\), the two displayed full triality lifts are exactly
\(A_R\) and \(A_H\).
\end{theorem}

\begin{proof}
The four identified \(A_5\) summands carry the same Euclidean metric,
so the displayed signed permutation is a root-lattice isometry.  The
matrices \(R_D,H_D\) preserve \(D_4\): \(R_D\) negates one integral
coordinate and preserves parity of the sum; for \(x\in D_4\), each
coordinate of \(H_Dx\) is integral because its numerator has the same
parity as \(\sum x_i\), and the coordinate sum of \(H_Dx\) is \(2x_1\).
Both matrices are orthogonal involutions by multiplication.
Furthermore \(R_D=s_{e_1}\), \(H_D=s_c\), with
\(c=(-1,1,1,1)/2\), and \((e_1,c)=-1/2\).
On the plane with basis \(e_1,c\), their product has matrix
\(\left(\begin{smallmatrix}0&-1\\1&-1\end{smallmatrix}\right)\)
and has cube the identity; both reflections fix the orthogonal
complement pointwise.  The discriminant actions are the two distinct
transpositions \(B_R,B_H\).  The same word reduction as in the preceding
proof therefore shows that \(\langle R_D,H_D\rangle\) has exactly six
elements and that its discriminant action is an isomorphism onto
\(GL_2(\mathbb F_2)\).  Its inverse is the homomorphism \(\delta\).

Thus \(g\mapsto L_g\) is a homomorphism of root-lattice isometries.
It extends to the dual lattice because an isometry preserves the
integer-pairing condition defining the dual.  Its discriminant action is
the specified \((\varepsilon,p,B)\), which preserves \(\mathcal H\).
Consequently it preserves the inverse image
\[
 N=\{x\in R_{\mathrm{ES}}^\vee:
                  x+R_{\mathrm{ES}}\in\mathcal H\}.
\]
It is therefore an integral isometry of \(N\).
If \(L_g\) is the identity on \(N\), it is the identity on the full
real span, because \(N\) contains the full-rank lattice
\(R_{\mathrm{ES}}\).  It is therefore also the identity on
\(R_{\mathrm{ES}}^\vee/R_{\mathrm{ES}}\), in particular on the ternary
code.  Hence \(\Theta(g)=I\), and injectivity of \(\Theta\) gives \(g=1\).
This proves faithfulness.
The exact formulas for \(A_R,A_H\) were established above, so their
lifts are the stated full triality transformations.
\end{proof}

This constructs an explicit finite subgroup of \(O(N)\), including its
action on all five root summands and on their glue, and its exact
\(D_4\)-triality quotient.  The group and its projective quotient agree
with the \(A_5^4D_4\) row of \cite[Table~3, p.~21]{esn:CDH2014};
the isomorphisms here have been proved by their actual code actions.

For clarity, the two full triality elements, denoted
\(\widehat R=L_{\Theta^{-1}(A_R)}\) and
\(\widehat H=L_{\Theta^{-1}(A_H)}\), act in the fixed coordinates by
\[
\begin{split}
 \widehat R(X_1,X_2,X_3,X_4;d)
     &=(X_2,X_1,X_3,-X_4;R_Dd),\\
 \widehat H(X_1,X_2,X_3,X_4;d)
     &=(X_1,X_3,X_2,-X_4;H_Dd).
\end{split}
\]
Their squares and the cube of their product are the identity on the
entire lattice, by the proved faithful group homomorphism.
The product on the first three A5 summands is the indicated
three-cycle; its two negations on the fourth summand cancel.

\subsection{The exact reduced-shell divisor bijection}
\label{esn:subsec:shells}

Let \(p\) be a positive prime with \(p\equiv1\pmod {24}\), and
let the positive integer shell \(R\) satisfy \(R\equiv3\pmod4\).
Retain
\begin{equation}
 a=a_R(p)=\frac{p+R}{4},\qquad
 N_{\mathrm{sh}}=pa,\qquad g=\gcd(R,N_{\mathrm{sh}}),\qquad
 R_0=\frac Rg,\quad N_0=\frac{N_{\mathrm{sh}}}{g}.
 \label{esn:eq:shell-data}
\end{equation}
The integer \(N_{\mathrm{sh}}\) is the source's scalar \(N=pa\);
the subscript distinguishes it from the lattice \(N\) without changing
either object.  In particular \(a,R_0,N_0\) are positive integers,
\(\gcd(R_0,N_0)=1\), and the prescribed first denominator equation is
\begin{equation}
 \frac4p-\frac1a=\frac{4a-p}{pa}
       =\frac{R}{N_{\mathrm{sh}}}=\frac{R_0}{N_0}
       =\frac1b+\frac1c.
 \label{esn:eq:fixed-shell-equation}
\end{equation}
Here \(b,c\) are positive integer denominators, and their order is
initially retained.  Define the finite set
\[
 \mathcal D_{R,p}=
 \{d\in\mathbb Z_{>0}:d\mid N_0^2,
                             \ d\equiv-N_0\pmod {R_0}\},
 \qquad A_R(p)=|\mathcal D_{R,p}|.
\]

\begin{theorem}[Certificates, ordered pairs, and complementary pairs]
\label{esn:thm:shell-bijection}
There is an exact bijection
\begin{equation}
 \begin{split}
 \mathcal D_{R,p}&\longrightarrow
  \{(b,c)\in\mathbb Z_{>0}^2:
                    4/p=1/a+1/b+1/c\},\\
 d&\longmapsto
 \left(\frac{N_0+d}{R_0},
       \frac{N_0+N_0^2/d}{R_0}\right),
 \end{split}
 \label{esn:eq:shell-bijection}
\end{equation}
with inverse \(d=R_0b-N_0\).  The involution \(d\mapsto N_0^2/d\)
corresponds exactly to \((b,c)\mapsto(c,b)\).  Consequently the
number of unordered pairs of remaining denominators is
\begin{equation}
 U_R(p)=\frac{A_R(p)+\delta_{R,p}}2,
 \qquad
 \delta_{R,p}=\begin{cases}1,&R_0\mid2,\\0,&R_0\nmid2.\end{cases}
 \label{esn:eq:unordered-count}
\end{equation}
The unique possible fixed certificate is \(d=N_0\), giving
\(b=c=2N_0/R_0\).
\end{theorem}

\begin{proof}
For a positive solution of \eqref{esn:eq:fixed-shell-equation},
\(R_0b/N_0=1+b/c>1\), so \(d=R_0b-N_0>0\); likewise
\(e=R_0c-N_0>0\).  Multiplying the fraction equation by
\(N_0bc\) gives \(R_0bc=N_0(b+c)\), and therefore
\[
 (R_0b-N_0)(R_0c-N_0)=N_0^2.
\]
Thus \(d\mid N_0^2\), \(d\equiv-N_0\pmod {R_0}\), and
\(e=N_0^2/d\).

Conversely, let \(d\in\mathcal D_{R,p}\) and \(e=N_0^2/d\).
For \(R_0>1\), the coprimality \(\gcd(R_0,N_0)=1\) makes
\(d\equiv-N_0\) a unit, and hence
\[
 e\equiv N_0^2(-N_0)^{-1}\equiv-N_0\pmod {R_0}.
\]
For \(R_0=1\), this congruence is vacuous and holds automatically.
Thus both expressions in \eqref{esn:eq:shell-bijection} are positive
integers.  Substitution into \((R_0b-N_0)(R_0c-N_0)=N_0^2\)
and expansion give the original fraction equation.  Each construction
recovers exactly \(d=R_0b-N_0\) and \(e=R_0c-N_0\), so they
are inverse, proving the bijection and the involution statement.

For a positive divisor, \(d=N_0^2/d\) holds exactly when
\(d=N_0\).  This divisor is a certificate exactly when
\(N_0\equiv-N_0\pmod {R_0}\), that is, \(R_0\mid2N_0\).
Coprimality makes this equivalent to \(R_0\mid2\).
Every other orbit contains two certificates, so counting the fixed
orbits and the two-element orbits gives
\(U_R(p)=\delta_{R,p}+(A_R(p)-\delta_{R,p})/2\), exactly
\eqref{esn:eq:unordered-count}.
\end{proof}

Factor the exact integer \(N_0\) as
\(N_0=\prod_{j=1}^{m}q_j^{e_j}\), with distinct positive primes
\(q_j\) and positive exponents \(e_j\), and define
\begin{equation}
 \begin{split}
 B(N_0)&=\prod_{j=1}^{m}\{-e_j,-e_j+1,\ldots,e_j\},\\
 \lambda_{R_0,N_0}:B(N_0)&\longrightarrow(\mathbb Z/R_0\mathbb Z)^\times,
 \qquad
 \lambda_{R_0,N_0}(\beta)=\prod_{j=1}^{m}q_j^{\beta_j}\pmod {R_0}.
 \end{split}
 \label{esn:eq:exponent-box}
\end{equation}
The group for modulus one is taken to be the one-element group, so that
the same notation also expresses the vacuous congruence in that case.
For \(R_0>1\), every \(q_j\) is a unit by coprimality, so all
negative powers in this formula are defined.

\begin{theorem}[The exact exponent-box map and the two raw channels]
\label{esn:thm:raw-channels}
The map
\begin{equation}
 \beta\longmapsto d=\prod_{j=1}^{m}q_j^{e_j+\beta_j}
 \label{esn:eq:exponent-to-divisor}
\end{equation}
bijects \(\lambda_{R_0,N_0}^{-1}(-1)\) with
\(\mathcal D_{R,p}\).  It retains each valuation and both original
integers \(R_0,N_0\).

In addition, suppose \(\gcd(R,pa)=1\) and \(p\nmid a\).
Write
\[
 a=\prod_{j=1}^{s}q_j^{e_j},\qquad
 B(a)=\prod_{j=1}^{s}\{-e_j,\ldots,e_j\},
\]
and define, for each unit
\(t\) modulo \(R\),
\[
 c_R(a;t)=\#\left\{\beta\in B(a):
                                 \prod_j q_j^{\beta_j}\equiv t\pmod R\right\}.
\]
Then the number of raw divisor certificates, equivalently ordered
denominator pairs, is
\begin{equation}
 A_R(p)=c_R(a;-1)+2c_R(a;-p^{-1}).
 \label{esn:eq:two-raw-channels}
\end{equation}
Its complementary-divisor quotient is the count in
\eqref{esn:eq:unordered-count}, with \(R_0=R\) under these additional
coprimality hypotheses.
\end{theorem}

\begin{proof}
A divisor of \(N_0^2\) has one and only one exponent vector
\(\alpha_j\in\{0,\ldots,2e_j\}\).  The exact change
\(\beta_j=\alpha_j-e_j\) gives the entire displayed box and
\(dN_0^{-1}\equiv\prod_jq_j^{\beta_j}\pmod {R_0}\).
Multiplication of \(d\equiv-N_0\) by the unit \(N_0^{-1}\)
gives precisely the target \(-1\); modulo one both statements are
vacuous.  This proves the first bijection and its inverse.

Under the additional hypotheses, \(R_0=R\), \(N_0=pa\), and
the exponent of \(p\) in \(pa\) is exactly one.  Every positive integer
divisor \(d\) of \(p^2a^2\) therefore has exactly one valuation-index
representation
\[
 d=p^{t+1}\prod_jq_j^{e_j+\beta_j},\qquad
 (t,\beta)\in\{-1,0,1\}\times B(a).
\]
After division by \(pa\), its unit-group residue is
\[
 d(pa)^{-1}\equiv p^t\prod_jq_j^{\beta_j}\pmod R.
\]
Distinct valuation-index pairs may have the same unit-group residue;
they remain distinct integer divisors in the count.
The three conditions for this to equal \(-1\) are
\[
\begin{array}{c|c}
 t&\prod_jq_j^{\beta_j}\pmod R\\\hline
 -1&-p\\0&-1\\1&-p^{-1}.
\end{array}
\]
The negation \(\beta\mapsto-\beta\) is a bijection of \(B(a)\)
and inverts the unit-group value, so it bijects the fibres at
\(-p\) and \(-p^{-1}\).  Summing the sizes of these three disjoint
raw exponent classes proves \eqref{esn:eq:two-raw-channels}.  The
complementary divisor sends \((t,\beta)\) to \((-t,-\beta)\),
so identifying it changes the raw count as proved in
Theorem~\ref{esn:thm:shell-bijection}; it does not leave the factor
two in \eqref{esn:eq:two-raw-channels} unchanged.
\end{proof}

This counting convention corrects the wording attached to the source's
two-channel formula: the displayed sum is a raw count, whereas a count
after identifying complementary divisors is \eqref{esn:eq:unordered-count}.
The exact source locator and formula are retained in
\cite{esn:Residual2026,esn:Provenance}.

\subsection{A fixed-shell count does not descend to the residue fibre}
\label{esn:subsec:non-descent}

\begin{theorem}[Two primes in the same residue with different shell counts]
\label{esn:thm:non-descent}
The two primes \(2521\) and \(3361\) both reduce to \(1\in S_{840}\),
but for the same shell \(R=3\) their raw and unordered counts are
\begin{equation}
 A_3(2521)=0,\qquad A_3(3361)=6,
 \qquad U_3(2521)=0,\quad U_3(3361)=3.
 \label{esn:eq:non-descent-counts}
\end{equation}
One actual denominator certificate in the second case is
\begin{equation}
 \frac4{3361}=\frac1{841}+\frac1{974690}+\frac1{28266010}.
 \label{esn:eq:explicit-denominator-example}
\end{equation}
Thus the raw count, its unordered quotient, and the indicator of
completion of this fixed shell each fail to factor through prime
reduction modulo 840.
\end{theorem}

\begin{proof}
The residue equations are the literal equalities
\(2521=3\cdot840+1\) and \(3361=4\cdot840+1\).
For an explicit finite primality verification, the primes at most 57,
in increasing order, are
\[
 \mathcal P=(2,3,5,7,11,13,17,19,23,29,31,37,41,43,47,53).
\]
The remainders against its first 15 entries for \(2521\), all 16
entries for \(3361\), and its first nine entries for \(631\), are
respectively
\[
\begin{split}
 &(1,1,1,1,2,12,5,13,14,27,10,5,20,27,30),\\
 &(1,1,1,1,6,7,12,17,3,26,13,31,40,7,24,22),\\
 &(1,1,1,1,4,7,2,4,10).
\end{split}
\]
All are nonzero, and the respective integer square roots are 50, 57,
and 25.  The prime 29 has remainders \(1,2,4\) against \(2,3,5\),
the primes at most its integer square root 5.
Every composite positive integer greater than one has a prime divisor
at most its square root: in a factorization into two factors greater
than one, the smaller factor is at most the square root and has a prime
divisor no larger than itself.  These nonzero remainder lists therefore
prove each stated primality.

For \(R=3\),
\[
 a_3(2521)=631,\qquad a_3(3361)=841=29^2.
\]
Both primes are 1 modulo three; also \(631\equiv1\pmod3\) and
\(841\equiv1\pmod3\).  Thus \(\gcd(3,pa)=1\) in both cases.
Both \(a\) values are smaller than their respective primes, so
\(p\nmid a\).  Every additional hypothesis of
Theorem~\ref{esn:thm:raw-channels} is now verified.
The two targets \(-1\) and \(-p^{-1}\) coincide modulo three.

For \(a=631\), all exponents \(-1,0,1\) give residue one, so
the target \(-1\) has count zero.  For \(a=29^2\), among the five
exponents \(-2,-1,0,1,2\), exactly \(-1,1\) give \(-1\) modulo
three because \(29\equiv-1\pmod3\).  Both target counts are
therefore 2 and the raw count is \(2+2\cdot2=6\).  Here
\(R_0=3\nmid2\), so both unordered counts are half the corresponding
raw counts.  This proves \eqref{esn:eq:non-descent-counts}.

For the second prime, retain
\[
 N_{\mathrm{sh}}=3361\cdot841=2826601,
 \qquad d=3361\cdot29=97469,
 \qquad \frac{N_{\mathrm{sh}}^2}{d}=3361\cdot29^3=81971429.
\]
Then
\[
 \frac{N_{\mathrm{sh}}+d}{3}=974690,
 \qquad
 \frac{N_{\mathrm{sh}}+N_{\mathrm{sh}}^2/d}{3}=28266010.
\]
The bijection of Theorem~\ref{esn:thm:shell-bijection} now proves
\eqref{esn:eq:explicit-denominator-example}.  The first prime has
no completion with this prescribed first denominator by its zero
certificate count.  Two arguments with the same residue and unequal
function values cannot be the pullback of a single residue function:
if \(f(p)=h(p\bmod840)\), equality of the residues forces equality
of the values.  Applying this observation to the three displayed
functions proves the claimed failures of factorization.
\end{proof}

\subsection{An actual residue cocycle and a retained-sample Fourier map}
\label{esn:subsec:cocycle-fourier}

For \(r\in S_{840}\), multiplication by \(r\) permutes each fibre.
Let \(\rho(r)\) act by
\begin{equation}
 \rho(r)e_s=e_{rs}\quad(s\in V_{840}),\qquad
 \rho(r)|_D=I,
 \qquad b(r)=e_r-e_1\in A(S_{840})\subset N.
 \label{esn:eq:residue-action-cocycle}
\end{equation}
The two free coordinates in \(b(r)\) belong to the principal fibre.

\begin{theorem}[The full-lattice residue action and its exact cocycle]
\label{esn:thm:residue-cocycle}
The permutations \(\rho(r)\) extend to a faithful homomorphism
\(\rho:S_{840}\to O(N)\).  The map \(b:S_{840}\to N\) satisfies
\begin{equation}
 b(rs)=b(r)+\rho(r)b(s)\qquad(r,s\in S_{840}),
 \qquad b(1)=0,\qquad (b(r),b(r))=2\quad(r\ne1).
 \label{esn:eq:cocycle-identity}
\end{equation}
Thus reduction of any prime into \(S_{840}\), followed by \(b\),
is an explicit lattice-valued map on that prime domain.
\end{theorem}

\begin{proof}
Multiplication permutes the orthonormal free coordinates, preserves
their augmentation kernels, and obeys \(\rho(r)\rho(s)=\rho(rs)\).
On the \(i\)-th fibre its dual discriminant generator is
\[
 \lambda_i=e_{f_i}-\frac16\sum_{t\in F_{\eta_i}}e_t.
\]
Its image differs from it by
\(e_{rf_i}-e_{f_i}\in A(F_{\eta_i})\).  Hence \(\rho(r)\)
acts trivially on the entire A5 discriminant group, since
\(\lambda_i\) generates it.  It fixes D and therefore the whole
discriminant group, in particular \(\mathcal H\).
It follows from the inverse-image definition that it preserves \(N\).
If \(\rho(r)\) were the identity, applying it to
\(e_1-e_{289}\) would give
\(e_r-e_{289r}=e_1-e_{289}\).  Since \(289\ne1\), the
positive coordinate on the left is exactly \(r\), so comparison of
free-coordinate coefficients forces \(r=1\).  This proves faithfulness.

Finally
\[
 b(r)+\rho(r)b(s)
   =(e_r-e_1)+(e_{rs}-e_r)=e_{rs}-e_1=b(rs).
\]
For \(r\ne1\), the two basis vectors in the difference are
orthogonal unit vectors, giving squared norm two.  The stated prime map
is the composition \(p\mapsto p\bmod840\mapsto b(p\bmod840)\),
whose domains have all been specified.
\end{proof}

Fix a positive shell \(R\equiv3\pmod4\), and let \(P\) be any
finite set of positive primes whose reduction lies in \(S_{840}\).
For every \(p\in P\), use its actual count \(A_R(p)\) defined by
the reduced divisor set, including the full data of
\eqref{esn:eq:shell-data}.  Define
\begin{equation}
 \begin{split}
 w_{R,P}(r)&=\sum_{\substack{p\in P\\p\equiv r\ (840)}}A_R(p),\\
 \widehat w_{R,P}(k)&=
       \sum_{r\in S_{840}}w_{R,P}(r)
                     \exp\left(-\frac{2\pi\mathrm i k\kappa(r)}6\right),
             \qquad k\in\mathbb Z/6\mathbb Z.
 \end{split}
 \label{esn:eq:sample-pushforward}
\end{equation}
This is the pushforward along prime reduction, followed by the finite
Fourier transform with the already fixed generator 289.

\begin{theorem}[Inversion and an integral augmentation-valued count map]
\label{esn:thm:sample-fourier}
The function \(w_{R,P}\) is recovered by
\begin{equation}
 w_{R,P}(r)=\frac16\sum_{k=0}^5\widehat w_{R,P}(k)
                      \exp\left(\frac{2\pi\mathrm i k\kappa(r)}6\right).
 \label{esn:eq:fourier-inverse}
\end{equation}
Let \(u=\sum_{r\in S_{840}}e_r\),
\(W_{R,P}=\sum_rw_{R,P}(r)e_r\), and
\(M_{R,P}=\sum_rw_{R,P}(r)=\widehat w_{R,P}(0)\).
Then
\begin{equation}
 Q_{R,P}=6W_{R,P}-M_{R,P}u\in A(S_{840})\subset N
 \label{esn:eq:integral-count-augmentation}
\end{equation}
is a well-defined integral lattice-valued count map.  Put
\(\zeta_6=e^{2\pi\mathrm i/6}\) and
\[
 f_k=\sum_{j=0}^5\zeta_6^{kj}e_{289^j}\qquad(0\le k\le5).
\]
The five vectors \(f_1,\ldots,f_5\) form a basis of
\(A(S_{840})\otimes\mathbb C\), with
\begin{equation}
 \rho(289)f_k=\zeta_6^{-k}f_k,\qquad
 Q_{R,P}=\sum_{k=1}^5\widehat w_{R,P}(k)f_k.
 \label{esn:eq:fourier-augmentation-eigenvectors}
\end{equation}
Both the prime sample \(P\) and the shell \(R\) are retained in the
domain of these maps.
\end{theorem}

\begin{proof}
For an integer \(m\), the sum
\(\sum_{k=0}^5\zeta_6^{km}\) equals 6 if \(6\mid m\).
Otherwise it is the geometric sum
\((1-\zeta_6^{6m})/(1-\zeta_6^m)=0\).
Substituting \eqref{esn:eq:sample-pushforward} into the right side
of \eqref{esn:eq:fourier-inverse} therefore gives
\[
 \sum_{s\in S_{840}}w_{R,P}(s)
       \frac16\sum_{k=0}^5
              \zeta_6^{k(\kappa(r)-\kappa(s))}=w_{R,P}(r),
\]
because \(\kappa\) is bijective.  This proves inversion.

The coefficient sum of \(6W_{R,P}-M_{R,P}u\) is
\(6M_{R,P}-6M_{R,P}=0\), and its coefficients are integers, so
it belongs to the exact augmentation lattice.  The same geometric sum
shows that \(f_k\) has coefficient sum zero for \(1\le k\le5\)
and \(f_0=u\).  Moreover the six \(f_k\) are independent: the
coefficient Fourier transform of \(f_k\) is six in coordinate \(k\)
and zero in all other coordinates, by that same identity.  Since the
augmentation space has dimension five, its five nonzero-frequency
vectors form a basis.  Replacing \(j\) by \(j+1\) in
\(\rho(289)f_k\) gives the multiplier \(\zeta_6^{-k}\), including
its displayed sign.

Inversion in each free coordinate gives
\(W_{R,P}=(1/6)\sum_{k=0}^5\widehat w_{R,P}(k)f_k\).
Multiplication by 6 and subtraction of
\(M_{R,P}u=\widehat w_{R,P}(0)f_0\) prove the second identity
in \eqref{esn:eq:fourier-augmentation-eigenvectors}.
All sums were taken over the specified finite set \(P\) and its
residue image.  The transform therefore recovers these aggregate counts;
it makes no identification between the separate counts of distinct
primes with the same residue.  Their possible inequality is proved in
Theorem~\ref{esn:thm:non-descent}.
\end{proof}

% Input after esn:thm:gl2-full-lattice-lift and before the shell subsection.
% This fragment introduces no packages, macros, or theorem environments.
% New labels have prefix esn:fa:.  Its residue corollary uses the formula
% defining rho, which is restated here so placement before that formula is safe.

\subsection{The entire lattice automorphism group and its Weyl quotient}
\label{esn:fa:subsec:full-automorphisms}

Cheng, Duncan and Harvey define the umbral group as the quotient of
the Niemeier lattice automorphism group by its root-reflection subgroup;
their Table~3 identifies this quotient as \(GL_2(3)\) for
\(A_5^4D_4\) \cite[Section~2.4 and Table~3]{esn:CDH2014}.
The argument here proves that identification in the present residue,
metric, discriminant and glue coordinates, and locates the arithmetic
six-cycle in the full group. It does not claim the abstract group as new.
The displayed complement is the particular lift already constructed here;
no equality of its representation with a simple-root-preserving section
is assumed.

Retain the exact lattice \(N\), the root lattice \(R_{\mathrm{ES}}\),
the four fibre identifications \(\varphi_i\), the discriminant marking,
the glue subgroup \(\mathcal H\), and the lift \(g\mapsto L_g\)
constructed above.  For a lattice \(K\) in a Euclidean space, \(O(K)\)
means the group of real linear isometries of its real span that map
\(K\) onto itself.  For the present root lattice define
\[
 W(R_{\mathrm{ES}})
 =\langle s_\alpha:\alpha\in R_{\mathrm{ES}},\ (\alpha,\alpha)=2\rangle,
 \qquad s_\alpha(x)=x-(x,\alpha)\alpha.
\]
All actions below retain the previously fixed six-coordinate order in
each \(A_5\) summand and the marking \(10\mapsto v\),
\(01\mapsto s\), \(11\mapsto c\) in \(D^\vee/D\).

\begin{theorem}[All isometries of the displayed \(A_5\) lattice]
\label{esn:fa:thm:a5-automorphisms}
Let \(A=\{x\in\mathbb Z^6:\sum_{j=0}^5x_j=0\}\) with its displayed
Euclidean metric, and for a permutation \(\sigma\) of
\(\{0,1,2,3,4,5\}\) let \(Q_\sigma e_j=e_{\sigma(j)}\).
Every element of \(O(A)\) has exactly one expression
\(\epsilon Q_\sigma|_{\mathbb R A}\), where
\(\epsilon\in\{1,-1\}\).  Thus
\[
 O(A)=\langle Q_\sigma|_{\mathbb R A}:\sigma\in S_6\rangle
                  \times\{I,-I\},
 \qquad W(A)=\langle Q_\sigma|_{\mathbb R A}:\sigma\in S_6\rangle
                  \cong S_6.
\]
On the marked discriminant group
\(D_A=(\mathbb Z/6)\lambda\),
\(\lambda=e_0-\frac16\sum_{j=0}^5e_j\), this isometry acts by
\(k\lambda+A\mapsto\epsilon k\lambda+A\).  The kernel of this
discriminant action is exactly \(W(A)\).
\end{theorem}

\begin{proof}
The norm-two vectors are precisely
\(\alpha_{ij}=e_i-e_j\) for \(i\ne j\).  For two distinct such roots,
\[
 (\alpha_{ij},\alpha_{kl})
       =\delta_{ik}-\delta_{il}-\delta_{jk}+\delta_{jl}.
\]
Their inner product is one exactly when they have the same first index
or the same second index.  Indeed equality of first indices, with
distinct second indices, makes the displayed expression one; equality
of second indices, with distinct first indices, does the same.  If both
positive Kronecker terms vanish the expression is nonpositive, and if
both are one then the roots are equal, the excluded case.

Put
\[
 C_i=\{\alpha_{ij}:j\ne i\},\qquad
 -C_i=\{\alpha_{ji}:j\ne i\}\qquad(0\le i\le5).
\]
These twelve sets are exactly the maximal sets of roots in which every
two distinct members have inner product one.  To prove this, a singleton
can be enlarged by keeping its first index and choosing a third index,
so a maximal such set contains a pair.  If the pair is
\(\alpha_{ij},\alpha_{ik}\), with \(j\ne k\), a further root
\(\alpha_{lm}\) has inner product one with both only if
\[
 (l=i\ \text{or}\ m=j)\quad\text{and}\quad
 (l=i\ \text{or}\ m=k).
\]
Since \(j\ne k\), these conditions force \(l=i\).  Thus the whole
set is contained in \(C_i\), and maximality makes it equal to \(C_i\).
For a pair with common second index the same argument, after negating
both roots and every possible further member, gives \(-C_i\).
Each \(C_i\) and \(-C_i\) has five roots and is maximal by this argument.

Form a graph whose vertices are these twelve maximal sets, with an edge
when the two sets have nonempty intersection.  Its intersections are
exactly
\[
 C_i\cap C_j=(-C_i)\cap(-C_j)=\varnothing\quad(i\ne j),
 \qquad
 C_i\cap(-C_j)=
 \begin{cases}\{\alpha_{ij}\},&i\ne j,\\
               \varnothing,&i=j.
 \end{cases}
\]
Consequently it is the bipartite graph on the two displayed six-element
families with every cross-edge except the six pairs \((C_i,-C_i)\).
It is connected: two distinct vertices on the same side have a common
neighbor indexed by any index different from both; vertices on opposite
sides with different indices are adjacent; and \(C_i\) and \(-C_i\)
are joined by the path
\(C_i,-C_j,C_k,-C_i\) for distinct \(i,j,k\).
A connected bipartite graph has only two choices for its ordered
bipartition, because the side of a fixed vertex determines the side of
each other vertex along any connecting path.  Hence an isometry of
\(A\), which permutes roots and preserves their inner products and set
intersections, either preserves both displayed families or exchanges them.

In the first case write
\(C_i\mapsto C_{\sigma(i)}\) and
\(-C_i\mapsto-C_{\tau(i)}\).
The unique opposite-family vertex not adjacent to \(C_i\) is \(-C_i\).
Preservation of that missing edge forces \(\tau(i)=\sigma(i)\)
for each \(i\).  The singleton intersection formula then gives
\(\alpha_{ij}\mapsto\alpha_{\sigma(i)\sigma(j)}\).
In the case exchanging the families, composition with \(-I\) reduces
to the first case, so the root action is
\(\alpha_{ij}\mapsto-\alpha_{\sigma(i)\sigma(j)}\).
The roots \(e_j-e_5\), \(0\le j\le4\), span \(\mathbb R A\),
so the root action determines the whole linear isometry.  These cases
give precisely the asserted expressions.

The permutation \(\sigma\) is determined by the permutation of the
six vertices \(C_i\), after the sign is determined by whether the two
families are preserved or exchanged.  Thus the expression is unique.
Every \(Q_\sigma\) and \(-I\) preserves \(A\), and \(-I\)
commutes with every linear map; uniqueness proves the direct product.
Reflection in \(e_i-e_j\) interchanges free coordinates \(i,j\)
and fixes the others, as follows on substituting this root in
\(x\mapsto x-(x,e_i-e_j)(e_i-e_j)\).
Every permutation is a product of transpositions: each cycle
\((i_1\ i_2\ \cdots\ i_m)\) is
\((i_1\ i_m)(i_1\ i_{m-1})\cdots(i_1\ i_2)\).
These reflections therefore generate precisely the displayed \(S_6\).
Finally
\[
 Q_\sigma\lambda-\lambda=e_{\sigma(0)}-e_0\in A.
\]
Thus the discriminant action is the stated sign.  Since \(\lambda+A\)
has exact order six, its negative is distinct from itself; the action
is trivial exactly when \(\epsilon=1\).
\end{proof}

\begin{theorem}[All isometries of the displayed \(D_4\) lattice]
\label{esn:fa:thm:d4-automorphisms}
For the exact lattice
\(D=\{x\in\mathbb Z^4:\sum_i x_i\equiv0\pmod2\}\),
the action on its marked discriminant group gives the split exact sequence
\[
 1\longrightarrow W_D\longrightarrow O(D)
   \xrightarrow{\ d_D\ }GL_2(\mathbb F_2)\longrightarrow1,
 \qquad
 O(D)=W_D\rtimes\delta(GL_2(\mathbb F_2)).
\]
Here \(W_D\) consists of all signed coordinate permutations with an even
number of minus signs, and \(\delta\) is exactly the section generated
by the earlier matrices \(R_D,H_D\).  In particular \(O(D)=W_F\)
in the root and weight coordinates of
Theorem~\ref{esn:thm:d4-f4-triality}.
\end{theorem}

\begin{proof}
Every isometry of \(D\) preserves its dual: if \(T\in O(D)\),
\(x\in D^\vee\), and \(a\in D\), then
\((Tx,a)=(x,T^{-1}a)\in\mathbb Z\).
It therefore induces an additive automorphism of
\(D^\vee/D\cong\mathbb F_2^2\), giving the homomorphism \(d_D\).
The three nonzero classes have the exact norm-one representative sets
\[
 \begin{aligned}
 \mathcal V&=\{\pm e_i:1\le i\le4\},\\
 \mathcal S_+&=\{\tfrac12\textstyle\sum_i\delta_i e_i:
                                      \prod_i\delta_i=1\},\\
 \mathcal S_-&=\{\tfrac12\textstyle\sum_i\delta_i e_i:
                                      \prod_i\delta_i=-1\},
 \qquad \delta_i\in\{1,-1\}.
 \end{aligned}
\]
in the respective marked classes \(v,s,c\).
For the integral nonzero class, norm one forces exactly one coordinate
\(\pm1\).  A half-integral vector has norm at least
\(4(1/2)^2=1\), with equality exactly when every coordinate is
\(\pm1/2\).  The difference of such a vector from
\(s=(1,1,1,1)/2\) has one coordinate \(-1\) per minus sign,
so it belongs to \(D\) exactly when their number is even.
This proves the three sets and their marking directly.

If \(d_D(T)=I\), then \(T\) preserves \(\mathcal V\).
The vectors \(Te_i\) are orthonormal members of \(\mathcal V\),
so \(Te_i=\delta_i e_{\sigma(i)}\) for a permutation \(\sigma\)
of four indices and signs \(\delta_i\).  Preservation of
\(\mathcal S_+\) forces \(Ts\in\mathcal S_+\), hence
\(\prod_i\delta_i=1\).  Thus the kernel consists only of even
signed permutations.

Conversely, reflection in \(e_i-e_j\) interchanges coordinates
\(i,j\), and reflection in \(e_i+e_j\) interchanges them and
negates both.  Their composition produces a sign change in exactly
those two coordinates.  Any even set of sign changes is a product of
such disjoint pairs, and transpositions generate all permutations.
The root reflections therefore generate every even signed permutation,
and each root reflection is itself such a matrix.  If
\(\alpha\in D\) has norm two and \(x\in D^\vee\), then
\(s_\alpha x-x=-(x,\alpha)\alpha\in D\).  Every root reflection
acts trivially on the discriminant group.  This proves
\(\ker d_D=W_D\) in both directions.

For completeness the section has the original exact matrices
\[
 R_D=\operatorname{diag}(-1,1,1,1),\qquad
 H_D=\frac12
 \begin{pmatrix}
 1&1&1&1\\1&1&-1&-1\\1&-1&1&-1\\1&-1&-1&1
 \end{pmatrix}.
\]
They preserve \(D\), are orthogonal involutions, and act on
\((v,s,c)\) by \((s\ c)\) and \((v\ s)\), respectively,
as calculated in Theorem~\ref{esn:thm:d4-f4-triality}.
Their binary matrices in the retained marking are
\[
 B_R=\begin{pmatrix}1&1\\0&1\end{pmatrix},\qquad
 B_H=\begin{pmatrix}0&1\\1&0\end{pmatrix}.
\]
The same earlier coordinate calculation gives
\(R_D^2=H_D^2=(R_DH_D)^3=I\); thus their subgroup has at most six
elements, by reducing words to \((R_DH_D)^k\) or
\((R_DH_D)^kR_D\), \(k=0,1,2\).
The induced two transpositions generate the six permutations of the
three nonzero binary vectors.  There are exactly six invertible binary
matrices, by the \(3\cdot2\) choices of their independent columns,
and their action on the three nonzero vectors is faithful.  Hence the
six words have distinct discriminant actions, and the restriction of
\(d_D\) to \(\langle R_D,H_D\rangle\) is an isomorphism onto
\(GL_2(\mathbb F_2)\).  Its inverse is the specified \(\delta\).
For every \(T\in O(D)\),
\(T\delta(d_D(T))^{-1}\in\ker d_D=W_D\).
The two factors intersect trivially because \(d_D\delta=I\).
This proves the asserted splitting and exhausts \(O(D)\).
The earlier theorem proved
\(W_F=\langle W_D,R_D,H_D\rangle\), so it also proves the last equality.
\end{proof}

\begin{theorem}[The exact full automorphism quotient of \(N\)]
\label{esn:fa:thm:full-quotient}
Every isometry of \(N\) has a unique factorization
\[
 T=wL_g,\qquad w\in W(R_{\mathrm{ES}}),\quad g\in G.
\]
The Weyl group is normal, and the induced discriminant action gives
the split exact sequence and the indicated isomorphisms
\begin{equation}
 \begin{split}
 1&\longrightarrow W(R_{\mathrm{ES}})
   \longrightarrow O(N)\xrightarrow{\ d_{\mathrm{ES}}\ }G
   \longrightarrow1,\\
 O(N)&=W(R_{\mathrm{ES}})\rtimes L(G),\\
 L(G)&\xrightarrow{\ \cong\ }O(N)/W(R_{\mathrm{ES}})
               \cong G\xrightarrow{\ \Theta\ }GL_2(\mathbb F_3),\\
 W(R_{\mathrm{ES}})&\cong S_6^4\times W_D,
 \qquad |O(N)|=48\cdot(6!)^4\cdot(2^3\cdot4!).
 \end{split}
 \label{esn:fa:eq:full-group}
\end{equation}
In identified coordinates, if
\(w=(Q_{\sigma_1},Q_{\sigma_2},Q_{\sigma_3},Q_{\sigma_4};w_D)\)
and \(g=(\varepsilon,p,B)\), the exact semidirect-product action is
\begin{equation}
 L_gwL_g^{-1}
   =\bigl(Q_{\sigma_{p(1)}},Q_{\sigma_{p(2)}},
          Q_{\sigma_{p(3)}},Q_{\sigma_{p(4)}};
          \delta(B)w_D\delta(B)^{-1}\bigr).
 \label{esn:fa:eq:weyl-conjugation}
\end{equation}
In the original fibre coordinates the \(i\)-th permutation factor
on the right means
\(\varphi_iQ_{\sigma_{p(i)}}\varphi_i^{-1}\).
\end{theorem}

\begin{proof}
Theorem~\ref{esn:thm:niemeier} proves that the norm-two vectors of \(N\)
are exactly the 144 roots of its five original orthogonal summands.
Join two distinct roots by an edge when their inner product is nonzero.
The five summands are exactly the connected components of this graph.
There are no edges between different summands by orthogonality.
Within an \(A_5\) summand, the simple roots
\(e_0-e_1,e_1-e_2,e_2-e_3,e_3-e_4,e_4-e_5\)
have consecutive inner product \(-1\), forming a connected graph,
and span its real space.  Every root not already among them has nonzero
inner product with at least one of them; otherwise it is orthogonal
to its entire real space and has norm zero, contrary to norm two.
Thus its whole root graph is connected.  In \(D\), the four roots
\[
 e_1-e_2,\quad e_2-e_3,\quad e_3-e_4,\quad e_3+e_4
\]
span \(\mathbb R D\), and the second has inner product \(-1\)
with each of the other three.  The same spanning argument proves
connectedness of the \(D_4\) root graph.

An element \(T\in O(N)\) permutes the norm-two vectors and preserves
these edges, hence permutes their connected components.  Their
cardinalities are \(30,30,30,30,24\), so \(T\) preserves the
\(D_4\) component and permutes the four \(A_5\) components.
The roots generate each original lattice summand: the five and four
displayed simple roots are their integral bases by
Theorem~\ref{esn:thm:root-discriminant}.
Consequently \(T\) preserves their integer span
\(R_{\mathrm{ES}}\).  The same component argument applies to every
element of \(O(R_{\mathrm{ES}})\), since its norm-two vectors are
the same 144 roots.  It follows that every such element has, in the
fixed identified coordinates, the form
\begin{equation}
 T(X_1,X_2,X_3,X_4;d)
    =\bigl(\varepsilon_iQ_{\sigma_i}X_{p(i)}\bigr)_{i=1}^4
                                  \mathbin{;}T_Dd,
 \label{esn:fa:eq:all-root-isometries}
\end{equation}
where \(p\in S_4\), \(\varepsilon_i\in\{1,-1\}\),
\(\sigma_i\in S_6\), and \(T_D\in O(D)\).
This follows by transporting each isometry from its input \(A_5\)
summand to its output summand using the fixed \(\varphi_i\), then
applying Theorem~\ref{esn:fa:thm:a5-automorphisms}.
All choices in \eqref{esn:fa:eq:all-root-isometries} are unique.
The convention is precisely the earlier convention: a vector supported
on component \(k\) is sent to component \(p^{-1}(k)\).

Let \(\mathcal A_{\mathrm{ES}}\) be the group of transformations
of the exact discriminant group
\((\mathbb Z/6)^4\oplus\mathbb F_2^2\) given by
\[
 (x;y)\longmapsto
        \bigl((\varepsilon_i x_{p(i)})_{i=1}^4;By\bigr),
 \qquad \varepsilon_i\in\{1,-1\},\quad p\in S_4,
                 \quad B\in GL_2(\mathbb F_2).
\]
Each tuple gives a distinct transformation: the image of the order-six
generator in each individual summand determines its output summand
and its sign, and the images of the two binary basis vectors determine
\(B\).  The discriminant action of
\eqref{esn:fa:eq:all-root-isometries}, calculated by the two preceding
theorems, is exactly this tuple with \(B=d_D(T_D)\).
This defines a homomorphism
\[
 d_{\mathrm{ES}}:O(R_{\mathrm{ES}})\longrightarrow
                                      \mathcal A_{\mathrm{ES}}.
\]
It is a homomorphism because extending an isometry to the dual and then
passing to the quotient commutes with composition.  In particular the
component permutation rule remains
\(P_pP_q=P_{q\circ p}\); there is no change of the earlier convention.

This homomorphism is onto.  Given any tuple
\(a=(\varepsilon,p,B)\in\mathcal A_{\mathrm{ES}}\), define
\[
 \widetilde L_a(X_1,X_2,X_3,X_4;d)
   =\bigl(\varepsilon_iX_{p(i)}\bigr)_{i=1}^4
                                     \mathbin{;}\delta(B)d.
\]
It is a root-lattice isometry and has discriminant action \(a\).
The map \(a\mapsto\widetilde L_a\) is a homomorphism: if
\(a=(\varepsilon,p,B)\) and \(a'=(\varepsilon',q,B')\),
its composite has \(i\)-th output
\(\varepsilon_i\varepsilon'_{p(i)}X_{q(p(i))}\)
and \(D\)-component \(\delta(B)\delta(B')d=\delta(BB')d\),
which are exactly the tuple product in the stated convention.

The kernel of \(d_{\mathrm{ES}}\) consists exactly of the tuples
in \eqref{esn:fa:eq:all-root-isometries} with
\(p=1\), every \(\varepsilon_i=1\), and \(T_D\in W_D\).
It is therefore the direct product \(S_6^4\times W_D\).
Every root reflection acts on its own orthogonal summand and fixes
the other four pointwise, and the preceding two theorems identify the
groups generated inside those summands.  Hence this direct product
is exactly \(W(R_{\mathrm{ES}})\).

Every root reflection also preserves \(N\).  Indeed \(\alpha\in N\)
and integrality of \(N\) give
\(s_\alpha x=x-(x,\alpha)\alpha\in N\) for every \(x\in N\).
The reflection is an involution, so inclusion is equality.
Thus \(W(R_{\mathrm{ES}})\subset O(N)\).
It is normal there, since for \(T\in O(N)\)
\[
 Ts_\alpha T^{-1}(x)
   =x-(x,T\alpha)T\alpha=s_{T\alpha}(x),
\]
and \(T\alpha\) is again a norm-two root.

It remains to impose the actual glue, which completely determines the
remaining isometries.  For \(T\in O(R_{\mathrm{ES}})\) and
\(x\in R_{\mathrm{ES}}^\vee\), the inverse-image definition gives
\[
 Tx\in N
 \quad\Longleftrightarrow\quad
 d_{\mathrm{ES}}(T)(x+R_{\mathrm{ES}})\in\mathcal H.
\]
Consequently \(T(N)=N\) if and only if
\(d_{\mathrm{ES}}(T)(\mathcal H)=\mathcal H\).
For the forward implication, represent each class of \(\mathcal H\)
by a vector of \(N\) and apply both \(T\) and \(T^{-1}\).
For the reverse implication, the displayed equivalence applies to
\(T\) and \(T^{-1}\) because the action on \(\mathcal H\) is
bijective.  The group preserving \(\mathcal H\) inside precisely
\(\mathcal A_{\mathrm{ES}}\) is the group \(G\) of
Theorem~\ref{esn:thm:glue-gl2}, proved there to be exactly
\(GL_2(\mathbb F_3)\) through the code map \(\Theta\).
We have therefore proved the equality of subgroups
\[
 O(N)=d_{\mathrm{ES}}^{-1}(G)
                          \quad\text{inside }O(R_{\mathrm{ES}}).
\]
The section \(\widetilde L\) restricts on \(G\) to the already
specified lift \(L\).  This proves surjectivity of the restricted
map \(d_{\mathrm{ES}}:O(N)\to G\), with exactly the asserted kernel.
For \(T\in O(N)\), set \(g=d_{\mathrm{ES}}(T)\).
Then \(w=TL_g^{-1}\) lies in the kernel and \(T=wL_g\).
If \(wL_g=w'L_{g'}\), applying \(d_{\mathrm{ES}}\) gives
\(g=g'\), then cancellation gives \(w=w'\).
In particular \(L(G)\cap W(R_{\mathrm{ES}})=\{I\}\).
Normality and this unique factorization prove the semidirect product
and the quotient isomorphism, including its surjectivity.
The four \(S_6\) factors have order \((6!)^4\);
\(W_D\) has \(2^3\) sign choices and \(4!\) permutation choices,
and \(|G|=48\).  Unique factorization proves the displayed order.

Finally put \(Z=L_g^{-1}X\).  From the \(i\)-th output equation
\(X_i=\varepsilon_iZ_{p(i)}\) follows
\(Z_{p(i)}=\varepsilon_iX_i\), with each vector read in the fixed
identified \(A\) coordinates.  Therefore the \(i\)-th output of
\(L_gwL_g^{-1}X\) is
\[
 \varepsilon_iQ_{\sigma_{p(i)}}Z_{p(i)}
   =\varepsilon_iQ_{\sigma_{p(i)}}\varepsilon_iX_i
   =Q_{\sigma_{p(i)}}X_i.
\]
The \(D\)-component is
\(\delta(B)w_D\delta(B)^{-1}d\).
It belongs to \(W_D\) by normality of \(\ker d_D\).
These are exactly the components of
\eqref{esn:fa:eq:weyl-conjugation}, proving the asserted action with
all signs and component indices retained.
\end{proof}

\begin{theorem}[The residue six-cycle in the full automorphism group]
\label{esn:fa:thm:residue-weyl}
For \(r\in S_{840}=\langle289\rangle\), let
\(\rho(r)e_s=e_{rs}\) on each of the four fibres and
\(\rho(r)|_D=I\), with the exact fibre action used in
\eqref{esn:eq:residue-action-cocycle}.
Then
\begin{equation}
 \rho(S_{840})\subset W(R_{\mathrm{ES}}),\qquad
 d_{\mathrm{ES}}(\rho(r))=1,\qquad
 L_g\rho(r)L_g^{-1}=\rho(r)
       \quad(g\in G,\ r\in S_{840}).
 \label{esn:fa:eq:residue-position}
\end{equation}
In particular the map
\[
 S_{840}\times G\longrightarrow O(N),\qquad(r,g)\longmapsto\rho(r)L_g
\]
is an injective homomorphism with image of order \(6\cdot48=288\),
isomorphic to \(C_6\times GL_2(\mathbb F_3)\).
\end{theorem}

\begin{proof}
Put \(C e_j=e_{j+1\bmod6}\) on the free six-coordinate space and
restrict it to \(\mathbb R A\).  If
\(r=289^m\), with \(m\in\{0,1,2,3,4,5\}\), the exact fibre
identification gives
\[
 \rho(r)e_{f_i289^j}=e_{f_i289^{j+m}},\qquad
 \rho(r)|_{A(F_{\eta_i})}=\varphi_iC^m\varphi_i^{-1}.
\]
Thus in all four identified summands \(\rho(r)\) is the same
coordinate permutation \(C^m\), and it is the identity on \(D\).
Theorem~\ref{esn:fa:thm:a5-automorphisms} puts it in the corresponding
product of Weyl groups, giving the first two assertions.
Each sign \(\varepsilon_i I\) commutes with \(C^m\), and the
same operator \(C^m\) occurs in every component.  Directly the
\(i\)-th outputs of \(L_g\rho(r)\) and \(\rho(r)L_g\) are
respectively
\(\varepsilon_iC^mX_{p(i)}\) and
\(C^m\varepsilon_iX_{p(i)}\), which agree; the \(D\)-outputs
also agree since \(\rho(r)|_D=I\).
This proves the conjugation identity without changing the chosen origins.

The six restrictions \(C^m|_{\mathbb R A}\) are distinct: if
\(C^m(e_0-e_1)=e_0-e_1\), comparison of the unique coordinate with
coefficient \(+1\) gives \(m=0\pmod6\).
Thus \(\rho\) is faithful.  Commutation makes the displayed product
map a homomorphism.  If \(\rho(r)L_g=I\), then
\(L_g=\rho(r)^{-1}\) lies in
\(L(G)\cap W(R_{\mathrm{ES}})=\{I\}\).
Faithfulness of both factors gives \(r=1\) and \(g=1\).
This proves injectivity and the exact image order.
\end{proof}

% Exact proof module, 6 September 2026. Input after es_niemeier_shared.tex.
% All original coordinates and Euclidean root lengths are retained.
\section{The exact Leech neighbor of the retained index-72 lattice}
\label{esnl:sec:neighbor}

Retain the lattice, metric, discriminant marking, and glue literally:
\[
 \begin{split}
 A&=\{(x_0,\ldots,x_5)\in\mathbb Z^6:\sum x_j=0\},\\
 D&=\{(d_1,\ldots,d_4)\in\mathbb Z^4:\sum d_i\equiv0\pmod2\},
 \qquad R=A^{\perp4}\perp D,\\
 C_2&=\langle b_0=(1111;00),b_1=(1100;10),b_2=(1010;01)\rangle_{\mathbb F_2},\\
 C_3&=\langle z_1=(1,1,1,0),z_2=(1,2,0,1)\rangle_{\mathbb F_3},\\
 H&=\{(3u+2z;y):(u;y)\in C_2,\ z\in C_3\},\\
 N&=\{x\in R^\vee:x+R\in H\}.
 \end{split}
\]
The coordinates in each $A$ remain the six ordered elements $f_i289^j$
of its original residue fibre. The $A$ discriminant generator is
$\lambda=(5,-1,-1,-1,-1,-1)/6$. The $D$ markings $10,01,11$
remain represented by $v_D=(1,0,0,0)$,
$s_D=(1,1,1,1)/2$, $c_D=(-1,1,1,1)/2$.
The preceding proof establishes that $N$ is even unimodular,
$[N:R]=72$, $\det R=5184$, and its roots are exactly those of
$A_5^4D_4$. Every root has squared norm two. No metric rescaling
occurs anywhere below.

\subsection{The Weyl vector and the exact neighbor}

Use the original simple roots $a_j=e_j-e_{j+1}$, $0\le j\le4$,
in each $A$ component and
\[
 d_1=e_1-e_2,\quad d_2=e_2-e_3,\quad
 d_3=e_3-e_4,\quad d_4=e_3+e_4
\]
in $D$. Their highest roots are
$\theta_A=e_0-e_5=\sum_ja_j$ and
$\theta_D=e_1+e_2=d_1+2d_2+d_3+d_4$.
Both highest-root heights are five, so the common Coxeter number
is six. Define
\[
 \rho_A=(5,3,1,-1,-3,-5)/2,\qquad
 \rho_D=(3,2,1,0),\qquad \rho=\rho_A^{\oplus4}\perp\rho_D.
 \label{esnl:eq:rho}
\]
Each simple root pairs to one with $\rho$.
Moreover $\rho_A-3\lambda=(0,2,1,0,-1,-2)\in A$ and
$\rho_D\in D$. Thus the exact glue class of $\rho$ is
$(3,3,3,3;00)$, the image of $b_0$, and $\rho\in N$.
Directly
\[
 \rho_A^2=35/2,\quad \rho_D^2=14,\quad
 \rho^2=4(35/2)+14=84=2\cdot6\cdot7.
\]
Consequently $(\rho,N)\subset\mathbb Z$.
Fix $\alpha=a_0$ in the first $A$ component and put
\[
 \ell_\rho:N\to\mathbb Z/6,\quad x\mapsto(\rho,x)\bmod6,\qquad
 K=\ker\ell_\rho,\quad R_0=R\cap K,\quad
 v=\rho/6-\alpha,\quad \Lambda=K+\mathbb Z v.
 \tag{NL1}
\]
The pairing $(\rho,\alpha)=1$ proves
$[N:K]=[R:R_0]=6$. For $x\in K$, the number
$(v,x)=(\rho,x)/6-(\alpha,x)$ is integral, and
\[
 v^2=84/36-2/6+2=4,\qquad
 6v=\rho-6\alpha\in K,\qquad(\rho,6v)=78.
\]
Hence $\Lambda$ is positive definite even integral of rank 24.
If $mv\in N$, its pairing after adding $m\alpha$ with $\alpha$
is $m/6\in\mathbb Z$, so $6\mid m$.
It follows that
\[
 N\cap\Lambda=K,\qquad [\Lambda:K]=6,\qquad
 \det K=36,\qquad \det\Lambda=36/6^2=1.
 \tag{NL2}
\]
Changing $\alpha$ to another root of $\rho$-pairing one changes
$v$ by an element of $K$, and so does not change $\Lambda$.
We shall use the equivalent exact congruence description
\[
 \Lambda=\{x+k\rho/6:x\in N,\ k\in\mathbb Z,\
                                (\rho,x)+k\equiv0\pmod6\}.
 \tag{NL3}
\]
Indeed write $x+k\rho/6=(x+k\alpha)+kv$; the first summand
is in $K$ exactly under this congruence. Changing $k$ by six
can be absorbed in $x$ using $\rho\in N$, and preserves the
congruence since $\rho^2=84$.

\subsection{A complete rootlessness proof}

\begin{theorem}[The Leech neighbor in the retained coordinates]
\label{esnl:thm:leech-neighbor}
The lattice $\Lambda=K+\mathbb Z(\rho/6-\alpha)$ is positive
definite even unimodular of rank 24 and has minimum squared norm
exactly four. It is the holy-construction Leech lattice of the
retained $A_5^4D_4$ glue, and
$N\cap\Lambda=K$ with $[N:K]=[\Lambda:K]=6$.
\end{theorem}

\begin{proof}
The rank, parity, unimodularity and indices were proved in
(NL1)--(NL2). The holy-presentation identification is proved in
(NL8)--(NL11) below. We now prove the required minimum directly.
For $k=0$ in (NL3), a possible root is an original root of $N$.
An $A$ root pairs with $\rho$ to one of $\pm1,\ldots,\pm5$.
For $D$, all roots $\pm e_i\pm e_j$ have nonzero $\rho_D$-pairing
of absolute value at most five. Thus none belongs to $K$.
Negation and the reduction of $k$ modulo six leave $k=1,2,3$.

Define the exact shifted class minima by
\[
 m_A(k,t)=\min_{x\in t\lambda+A}|x+k\rho_A/6|^2,\qquad
 m_D(k,y)=\min_{x+D=y}|x+k\rho_D/6|^2.
\]
Their values are
\[
 \begin{array}{c|rrrrrr}
 k\backslash t&0&1&2&3&4&5\\\hline
 1&35/72&35/72&35/72&35/72&35/72&35/72\\
 2&11/18&4/9&11/18&4/9&11/18&4/9\\
 3&3/8&17/24&17/24&3/8&17/24&17/24
 \end{array}
 \quad
 \begin{array}{c|rrrr}
 k\backslash y&00&10&01&11\\\hline
 1&7/18&7/18&7/18&7/18\\
 2&2/9&5/9&5/9&5/9\\
 3&1/2&1/2&1/2&1/2
 \end{array}.
 \tag{NL4}
\]
We prove the table and its needed equality information.
The shifted $A$ coordinates are precisely
\[
 \left(n_j-\frac t6+\frac{k(5-2j)}{12}\right)_{j=0}^5,
               \qquad n_j\in\mathbb Z,\quad\sum n_j=t.
 \tag{NL5}
\]
At $k=1$, coordinatewise nearest representatives are the six
numbers $\pm1/12,\pm3/12,\pm5/12$ in an order depending on $t$.
They sum to zero and their squares sum to $35/72$.
At $k=2$ and even $t$, the nearest coordinates consist of
two copies of each sign of $1/6$ and one copy of each sign of
$1/2$; the half-integer ties admit exactly the needed zero-sum
choice and the squared sum is $11/18$.
At odd $t$, there are two each of $0,1/3,-1/3$, giving $4/9$.
For $t=1,3,5$ their actual ordered shifted vectors are
\[
 (-1,1,0,-1,1,0)/3,\quad
 (1,0,-1,1,0,-1)/3,\quad
 (0,-1,1,0,-1,1)/3.
 \tag{NL6}
\]
Subtracting $\rho_A/3$ gives original squared norms
$17/6,3/2,17/6$, respectively.
At $k=3,t=0,3$, the nearest coordinates alternate $\pm1/4$,
giving $3/8$ with zero sum.
At each other $t$, the coordinatewise nearest representatives
have three absolute values $1/12$ and three $5/12$, square sum
$13/24$, and coordinate sum $1$ or $-1$.
Meeting the zero-sum constraint requires an integer change.
Every nonzero coordinate change increases the squared sum by
at least $(7/12)^2-(5/12)^2=1/6$.
Changing one suitable $5/12$ coordinate attains this increase
and restores the zero sum. Thus the constrained minimum is
$17/24$. At $t=0,3$, the original minimizing vectors are
\[
 (-1,-1,0,0,1,1),\qquad (-3,-1,-1,1,1,3)/2,
 \tag{NL7}
\]
of squared norms $4,11/2$.

For $D$ and $k=1$ the shift is $(1/2,1/3,1/6,0)$.
The nearest integer coordinate distances are
$(1/2,1/3,1/6,0)$; the first-coordinate tie meets either integer
parity. The nearest half-integer distances are
$(0,1/6,1/3,1/2)$; the last-coordinate tie meets either
half-integer class. Both squared sums are $7/18$.
At $k=2$ the shift is $(1,2/3,1/3,0)$.
The unique nearest integer vector before shifting is
$(-1,-1,0,0)\in D$, giving $2/9$.
Changing parity costs at least $1/3$, attained by changing the
second or third coordinate, giving $5/9$.
The half-integer coordinate distances are
$(1/2,1/6,1/6,1/2)$, giving $5/9$; ties meet either marking.
At $k=3$ the shift is $(3/2,1,1/2,0)$.
Integer-class minimum distances are $(1/2,0,1/2,0)$,
and half-integer-class distances $(0,1/2,0,1/2)$.
Ties meet every marking and all square sums are $1/2$.
In class $00$, the two original minimizing vectors are
\[
 (-1,-1,0,0),\qquad(-2,-1,-1,0),
\]
of squared norms $2,6$. In class $10$, they are
\[
 (-1,-1,-1,0),\qquad(-2,-1,0,0),
\]
of squared norms $3,5$.
In either half-integer class they have the shapes
$(-3/2,-1/2,-1/2,\pm1/2)$ and
$(-3/2,-3/2,-1/2,\pm1/2)$.
Exactly one last-coordinate sign in each shape meets each
original marking; their squared norms are $3,5$.
This proves every entry and the relevant equality cases in (NL4).

For $k=1$, every full shifted vector has squared norm at least
$4(35/72)+7/18=7/3>2$.
For $k=2$, the lower bound from (NL4) is two.
Equality forces all $t_i$ odd and $y=00$.
In the actual glue this forces $u=1111$.
If $z=0$, all $t_i=3$ and (NL6) gives
$x^2=4(3/2)+2=8$.
If $z\ne0$, its weight is three, so exactly one $t_i=3$
and three $t_i\in\{1,5\}$; then
$x^2=3/2+3(17/6)+2=12$.
The norm-two equation is
\[
 2=x^2+\tfrac23(\rho,x)+\tfrac{28}{3},
 \qquad(\rho,x)=-11-\tfrac32x^2.
\]
Its pairing is $-23$ or $-29$, always $1$ modulo six.
Condition (NL3) requires residue $4$, so all nine equality
candidates, one per ternary word, are excluded.

For $k=3$, the lower bound two requires $z=0$:
every nonzero ternary word raises the minimum by
$3(17/24-3/8)=1$.
If $u=0$, the possible original squared norms are $18,22$.
If $u=1111$, they are $24,28$.
If $u$ has weight two, its $D$ class is nonzero and they are
$2(4)+2(11/2)+3=22$ or $2(4)+2(11/2)+5=24$.
These are all sixteen equality candidates.
But
\[
 2=x^2+(\rho,x)+21,\qquad(\rho,x)=-19-x^2,
\]
and (NL3) would require $x^2\equiv2\pmod6$.
None of $18,22,24,28$ has that residue.
Thus $\Lambda$ has no roots. Together with (NL2), this proves
that it is positive definite even unimodular of rank 24,
with minimum squared norm exactly four, attained by $v$.
\end{proof}

\subsection{The exact holy-construction presentation}

For $0\le t\le5$ set
\[
 \omega_t=(\underbrace{1,\ldots,1}_{t},
             \underbrace{0,\ldots,0}_{6-t})-\tfrac t6(1,\ldots,1).
\]
Then $\omega_t+A=t\lambda+A$ and
$(\rho_A,\omega_t)=t(6-t)/2=3\omega_t^2$.
For the $D$ markings set
\[
 \omega_{00}=0,\quad\omega_{10}=e_1,\quad
 \omega_{01}=(1,1,1,1)/2,\quad
 \omega_{11}=(1,1,1,-1)/2.
\]
The last differs from the original $c_D$ by $(1,0,0,-1)\in D$,
so retains its exact marking.
For $a=(t_1,t_2,t_3,t_4;y)\in H$ define
\[
 u_a=(\omega_{t_1},\omega_{t_2},\omega_{t_3},\omega_{t_4};\omega_y),
 \qquad g_a=u_a-\rho/6.
 \tag{NL8}
\]
These give all 72 representatives of $N/R$. Each nonzero $D$
representative has norm one and pairing three; hence
\[
 (\rho,u_a)=3u_a^2\in6\mathbb Z,\quad
 u_a\in K,\quad g_a^2=u_a^2-(\rho,u_a)/3+84/36=7/3.
 \tag{NL9}
\]
Let $\mathcal F$ be the 29 vectors consisting, in each original
component, of the negative original simple roots and the positive
highest root. Their pairings with $\rho$ are $-1$ or $5$.
There are positive integral zero relations of coefficient sum six:
\[
 \theta_A+\sum_{j=0}^4(-a_j)=0,\qquad
 \theta_D+(-d_1)+2(-d_2)+(-d_3)+(-d_4)=0.
 \tag{NL10}
\]
Every component of $g_a$ pairs to $1/6$ with the corresponding
vertices except one tip, where the pairing is $-5/6$.
For $A$ this follows because $\omega_t$ has simple-root pairing
one only at $a_{t-1}$ when $t\ne0$, and has
$(\omega_t,\theta_A)=1$ for $t\ne0$.
For $D$ the three nonzero representatives have simple-root
pairing one only at $d_1,d_4,d_3$, respectively; all have
highest-root pairing one. The zero representative points to
the highest root. These computations fix all orientations in
the holy generator diagram.

The following coefficient presentations are exact equalities:
\[
 \begin{split}
 N&=\left\{\sum_fa_ff+\sum_ab_ag_a:
                    a_f,b_a\in\mathbb Z,\ \sum_ab_a=0\right\},\\
 \Lambda&=\left\{\sum_fa_ff+\sum_ab_ag_a:
            a_f,b_a\in\mathbb Z,\ \sum_fa_f+\sum_ab_a=0\right\}.
 \end{split}
 \tag{NL11}
\]
For the first, the $f$ span $R$ and $g_a-g_0=u_a$ generate its
given extension $N$; the zero coefficient sum cancels $\rho/6$.
For the second, write $A=\sum a_f$, $B=\sum b_a$,
$r=\sum a_ff$, and $x=r+\sum b_au_a$.
The vector is $x-B\rho/6$, and by (NL9)
$(\rho,x)\equiv-A\equiv B$, exactly (NL3) with $k=-B$.
Conversely, for any $x+k\rho/6$ in (NL3), write
$x=r+\sum b_au_a$. The coefficient of $u_0=0$ can be set so
$B=-k$. Write $r$ with the negative simple-root vertices.
The congruence gives $A+B\equiv0\pmod6$.
Adding a suitable multiple of a zero relation (NL10) makes
$A+B=0$ exactly. Thus both directions of (NL11) are proved.

These are exactly the holy coefficient rules of Conway--Sloane
\cite[pp.~277--280]{esn:ConwaySloane1982}.
Borcherds gives the uniform proof of the holy-construction
theorem \cite[Section~7, especially the coefficient rules
(i)--(ii) and Lemma~7.3]{esn:Borcherds1985}.
Thus (NL11) is an explicit realization of that Leech
construction from the actual 72 marked glue words; all its
metric, determinant, integrality and rootlessness assertions
have also been independently proved here.

\subsection{The indices, discriminant morphism, and cyclic order twelve}

Reduction modulo $R$ gives an isomorphism $K/R_0\cong H$:
its kernel is $R_0$, and (NL9) supplies representatives in $K$
for every $H$ class. Consequently
\[
 [K:R_0]=72,\quad [N:R_0]=[\Lambda:R_0]=432,\quad
 \det R_0=6^2\cdot5184=186624.
 \tag{NL12}
\]
Thus the original and new inclusions, with every index, are
\[
 R_0\subset R\subset N,\quad [R:R_0]=6,\ [N:R]=72;
 \qquad
 R_0\subset K\subset\Lambda,\quad[K:R_0]=72,\ [\Lambda:K]=6.
\]
In particular $R\cap\Lambda=R_0$, an exact common sublattice.

\begin{theorem}[The exact order-twelve quotient generator]
\label{esnl:thm:cyclic-twelve}
The quotient of $\Lambda$ by the original height-zero root
sublattice $R_0$ is
$\Lambda/R_0\cong\mathbb Z/12\oplus(\mathbb Z/6)^2$.
The class of $v=\rho/6-\alpha$ has exact order twelve.
\end{theorem}

\begin{proof}
Let $\bar b_i$ denote the image in $H$ of the binary generator
$b_i$, namely its first four coordinates multiplied by three
and its last two coordinates retained.
Since $6v=\rho-6\alpha$ and $6\alpha\in R_0$,
\[
 6[v]=\bar b_0\quad\text{in }\Lambda/R_0.
 \tag{NL13}
\]
This is a nonzero order-two element. The classes
$\bar b_0,\bar b_1,\bar b_2,2z_1,2z_2$ form the displayed
independent binary and ternary generators of $H$.
Eliminate $\bar b_0$ using (NL13). The resulting surjective
presentation is
\[
 \Lambda/R_0\cong
 \mathbb Z/12\oplus(\mathbb Z/2)^2\oplus(\mathbb Z/3)^2
 \cong\mathbb Z/12\oplus(\mathbb Z/6)^2.
 \tag{NL14}
\]
There are no further relations: this presentation has order
$12\cdot4\cdot9=432$, equal to (NL12).
Thus $[v]$ has exact order twelve. This particular occurrence
of twelve is proved from the Coxeter denominator six and the
actual order-two Weyl-vector glue class.
\end{proof}

Set $M=N+\Lambda=N+\mathbb Z(\rho/6)$.
Then
\[
 M=K^\vee,\qquad K^\vee/K=
    \langle[\alpha],[v]\rangle\cong(\mathbb Z/6)^2.
 \tag{NL15}
\]
Indeed $M\subset K^\vee$, its index over $N$ is six by pairing
with $\alpha$, and $[K^\vee:N]=[N:K]=6$.
The cyclic images of $N/K$ and $\Lambda/K$ have order six and
trivial intersection, proving the stated generators.
Their exact discriminant data are
\[
 q_K([\alpha])=q_K([v])=0,\quad
 b_K([\alpha],[v])=1/6\pmod{\mathbb Z},\quad
 q_K(a[\alpha]+b[v])=ab/3\pmod{2\mathbb Z}.
 \tag{NL16}
\]
These follow from $\alpha^2=2$, $v^2=4$ and
$(\alpha,v)=1/6-2=-11/6$.
The two original lattices are therefore the inverse images
of the two indicated order-six isotropic subgroups in this
specified discriminant form.

\subsection{An exact Lorentzian quotient and symmetry transport}

Take $\mathcal I=N\perp U$ in coordinates $(x,m,n)$ with
\[
 \langle(x,m,n),(x',m',n')\rangle
 =(x,x')-mn'-m'n .
\]
Thus $U$ has Gram matrix
$\left(\begin{smallmatrix}0&-1\\-1&0\end{smallmatrix}\right)$.
The vectors $z=(0,0,1)$ and $w=(\rho,6,7)$ are primitive
isotropic: $w^2=84-2\cdot6\cdot7=0$, and six and seven are
coprime. There are integral quotient isometries
\[
 \begin{split}
 z^\perp/\mathbb Zz&\longrightarrow N,\quad[(x,0,n)]\mapsto x,\\
 w^\perp/\mathbb Zw&\longrightarrow\Lambda,\quad
                                  [(x,m,n)]\mapsto x-m\rho/6.
 \end{split}
 \tag{NL17}
\]
For the second map, orthogonality means $(\rho,x)=7m+6n$,
giving (NL3) with $k=-m$. Conversely take $m=-k$ and
$n=((\rho,x)+7k)/6$ for a vector in (NL3).
If the image vanishes, $x=m\rho/6$; pairing with $\alpha$
forces $m=6j$, and then $n=7j$, so the kernel is exactly
$\mathbb Zw$. Finally
\[
 |x-m\rho/6|^2
 =x^2-\tfrac m3(\rho,x)+\tfrac73m^2=x^2-2mn.
\]
Polarization gives the full pairing, proving the isometry.

For every $T\in O(N)$ the exact transport laws are
\[
 T(K_\rho)=K_{T\rho},\qquad
 T(\Lambda_\rho)=\Lambda_{T\rho},\qquad
 (x,m,n)\mapsto(Tx,m,n):w_\rho\mapsto w_{T\rho}.
 \tag{NL18}
\]
They follow by substitution in (NL3).
In particular the stabilizer of $\rho$ acts on the fixed
$\Lambda_\rho$. General isometries act on this specified
family of neighbors; no assertion that an arbitrary chosen
automorphism complement fixes this Weyl vector is required.

\subsection{The exact $F_4$ Coxeter action of order twelve}

The order twelve in (NL14) is a lattice quotient invariant of
this explicitly marked construction. There is also the ambient
$F_4$ Coxeter action on the same original $D_4$ real root space.
Use the precise $F_4$ simple roots
\[
 \beta_1=e_2-e_3,\quad\beta_2=e_3-e_4,\quad
 \beta_3=e_4,\quad\beta_4=(e_1-e_2-e_3-e_4)/2,
 \qquad C=s_{\beta_1}s_{\beta_2}s_{\beta_3}s_{\beta_4},
\]
where $s_\beta(x)=x-2(x,\beta)\beta/\beta^2$ and the rightmost
factor acts first. The first two roots have squared norm two,
and the last two squared norm one. No rescaling of the $D_4$
roots has occurred.

\begin{theorem}[Coxeter twelve and its marked triality quotient]
\label{esnl:thm:f4-coxeter-twelve}
In the retained four-coordinate metric,
\[
 C=\frac12
 \begin{pmatrix}
 1&1&1&1\\-1&1&1&-1\\1&1&-1&-1\\1&-1&1&-1
 \end{pmatrix},\qquad
 C^3=
 \begin{pmatrix}
 0&1&0&0\\-1&0&0&0\\0&0&0&-1\\0&0&1&0
 \end{pmatrix}.
 \tag{NL19}
\]
It satisfies $C^4-C^2+I=0$, $C^6=-I$ and has exact order
twelve. Its marked discriminant action is
$v_D+D\mapsto c_D+D\mapsto s_D+D\mapsto v_D+D$.
The restriction of the original triality quotient is the exact
sequence
\[
 1\longrightarrow\langle C^3\rangle\cong\mathbb Z/4
 \longrightarrow\langle C\rangle\cong\mathbb Z/12
 \longrightarrow\langle(v_D\,c_D\,s_D)\rangle\cong\mathbb Z/3
 \longrightarrow1,
 \tag{NL20}
\]
with $\langle C^3\rangle\subset W(D_4)$.
\end{theorem}

\begin{proof}
The first reflection exchanges coordinates two and three; the
second exchanges coordinates three and four; the third changes
the sign of coordinate four; the last is
$I-2\beta_4\beta_4^T$. Their product gives the first matrix in
(NL19). Multiplying it three times gives the second, whose
square is $-I$. Multiplying the first matrix also gives
$C^4-C^2+I=0$, and hence $C^{12}=I$.
Its first column gives
$Cv_D=c_D+(1,-1,0,0)$, with the displayed correction in $D$.
Directly $Cs_D=v_D$ and $Cc_D=(0,0,-1,-1)+s_D$.
The displayed integer correction lies in $D$ and proves the
stated three-cycle with the original $v_D,s_D,c_D$ markings.
The first two reflections preserve $D$ as root reflections.
The last two preserve $D$ because a sign change preserves
integer sum parity, and reflection in $\beta_4$ sends an
integer even-sum vector to an integer vector of even sum.
The inverse reflections have the same property, so these
maps are isometries of $D$.
Consequently their induced action on $D^\vee/D$ is defined
and has the exact order-three image calculated above.
The matrix $C^3$ is an even signed permutation of exact order
four, and thus lies in the already identified group $W(D_4)$.
If $C^m=I$, its discriminant image forces $3\mid m$,
and the exact order of $C^3$ then forces $4\mid m/3$.
Thus $C$ has exact order twelve and the discriminant kernel
on its cyclic subgroup is exactly $\langle C^3\rangle$.
This proves every map and group order in (NL20).
\end{proof}

In particular $C^2$ still has nontrivial order-three
discriminant action, while $C^3$ is in the long-root Weyl
group. The Coxeter denominator six in (NL1), the cyclic order
twelve in (NL14), and the $F_4$ Coxeter twelve in (NL20)
therefore have their stated, distinct domains and explicit
relationships through the original $D_4$ marking.

% Staged only. No canonical file has been edited.
% Requires amsmath, amssymb, amsthm, and the existing esn theorem labels.
% All new labels have namespace n24:.

\subsection{An exact cyclic-triality isometry of the entire Niemeier space}
\label{n24:sec:cyclic-isometry}

Retain the ordered fibres
\(f_i289^j\), \(1\leq i\leq4\), \(0\leq j\leq5\),
the Euclidean augmentation coordinates of \(A(F_{\eta_i})\), the
ordered orthonormal coordinates of \(D\subset\mathbb R^4\), and the
72-element glue \(\mathcal H\) of
\eqref{esn:eq:glue-and-overlattice}.  Thus the actual real span is
\[
 V_N=N\otimes_{\mathbb Z}\mathbb R
 =\left\{(x_{i,j};d_k)\in\mathbb R^{24}\oplus\mathbb R^4:
               \sum_{j=0}^5x_{i,j}=0\ (1\leq i\leq4)\right\},
\]
with squared norm \(\sum_{i,j}x_{i,j}^2+\sum_kd_k^2\).
The four displayed linear constraints, rather than an implicit change of
metric, give its dimension 24.  Choose an ordered orthonormal octonion
basis \((u_0,\ldots,u_7)\) with \(u_0=1\).  Every octonion product
below is the product of the chosen real octonion algebra, without a
change of multiplication convention.

For \(p\in\{0,1\}\) and \(j\in\{0,1,2\}\), set
\begin{align}
 m_{i,p}&=\frac13\sum_{j=0}^2x_{i,p+2j},&
 t_{i,p,j}&=x_{i,p+2j}-m_{i,p},\label{n24:eq:triplets}\\
 \alpha_i&=\frac{\sum_{j=0}^2x_{i,2j}
                    -\sum_{j=0}^2x_{i,1+2j}}{\sqrt6},&
 (b_0,\ldots,b_7)&=(\alpha_1,\alpha_2,\alpha_3,\alpha_4,
                           d_1,d_2,d_3,d_4).
 \label{n24:eq:fixed-coordinates}
\end{align}
In particular \(m_{i,0}=\alpha_i/\sqrt6\),
\(m_{i,1}=-\alpha_i/\sqrt6\), and
\(\sum_jt_{i,p,j}=0\).  The ordered map is
\begin{equation}
 \begin{split}
 \mathcal L:V_N&\longrightarrow\mathbb O^3,\\
 (x;d)&\longmapsto(q_1,q_2,q_3),\\
 q_{j+1}&=\sum_{i=1}^4\sum_{p=0}^1
       \left(t_{i,p,j}+\frac{b_{2(i-1)+p}}{\sqrt3}\right)
                                          u_{2(i-1)+p}.
 \end{split}
 \label{n24:eq:full-isometry}
\end{equation}

\begin{theorem}[Actual rank-24 map, inverse, pairing, and residue action]
\label{n24:thm:full-isometry}
The map \(\mathcal L\) is a real linear isometry onto
\(\mathbb O^3\) with its product norm.  Write
\(q_{j+1}=\sum_{a=0}^7q_{j,a}u_a\), and define
\begin{equation}
 b_a=\frac1{\sqrt3}\sum_{j=0}^2q_{j,a},\qquad
 t_{i,p,j}=q_{j,2(i-1)+p}
                    -\frac13\sum_{h=0}^2q_{h,2(i-1)+p}.
 \label{n24:eq:inverse-first}
\end{equation}
Its inverse is exactly
\begin{equation}
 x_{i,p+2j}=t_{i,p,j}+
                 \frac{(-1)^p b_{i-1}}{\sqrt6},
 \qquad d_k=b_{k+3}.
 \label{n24:eq:inverse-second}
\end{equation}
If \(r=\rho(289^2)=\rho(361)\), with the residue action
of Theorem~\ref{esn:fa:thm:residue-weyl}, and
\(c(q_1,q_2,q_3)=(q_3,q_1,q_2)\), then
\begin{equation}
 \mathcal Lr=c\mathcal L,\qquad r^3=c^3=I,
 \qquad rN=N.
 \label{n24:eq:cyclic-intertwining}
\end{equation}
Thus \(\mathcal L|_N\) embeds the original even unimodular lattice,
with every pairing and glue class retained, into the actual three
octonion sectors and intertwines an actual order-three lattice
isometry with cyclic triality.
\end{theorem}

\begin{proof}
The two means sum to zero because each six-coordinate fibre has total
sum zero.  The equations for \(m_{i,p}\) follow from
\(\alpha_i=6m_{i,0}/\sqrt6\).  Orthogonality to constant triplets gives
\[
 \sum_{j=0}^5x_{i,j}^2
 =\sum_{p=0}^1\sum_{j=0}^2t_{i,p,j}^2
                    +3m_{i,0}^2+3m_{i,1}^2
 =\sum_{p,j}t_{i,p,j}^2+\alpha_i^2.
\]
For each \(a=2(i-1)+p\), the same zero-sum identity gives
\[
 \sum_{j=0}^2\left(t_{i,p,j}+\frac{b_a}{\sqrt3}\right)^2
                      =\sum_{j=0}^2t_{i,p,j}^2+b_a^2.
\]
Summing over all eight values of \(a\) proves
\(\sum_{h=1}^3|q_h|^2=\sum_{i,j}x_{i,j}^2+\sum_kd_k^2\).
Polarization proves equality of all pairings.  Equations
\eqref{n24:eq:inverse-first} recover the means and residuals of each
output triplet.  Equation~\eqref{n24:eq:inverse-second} then recovers
the original means, the original fibre coordinates, and \(d\).
Conversely those formulas applied to arbitrary \(q\) yield
\(\sum_jx_{i,j}=0\) and return the prescribed \(q\).  This proves
bijectivity without a dimension-only argument.

The residue action sends the coordinate vector \(e_{f_i289^h}\)
to \(e_{f_i289^{h+2}}\), so the coefficient in output position
\(p+2j\) is the old coefficient in position \(p+2(j-1)\), with
\(j\) read modulo three.  Both means, all \(\alpha_i\), and all
\(d_k\) are fixed.  Hence the three output octonions become
\((q_3,q_1,q_2)\), proving the intertwining formula.  Each coordinate
permutation acts trivially on \(A^\vee/A\): for
\(\lambda=e_0-\frac16\sum_he_h\), its difference from \(\lambda\)
is an augmentation root.  The \(D\)-coordinates are fixed.  Thus
\(r\) acts trivially on the full discriminant group, preserves
\(\mathcal H\), and preserves its exact inverse image \(N\).
Its order is three, as seen on \(e_0-e_1\) in any fibre.
\end{proof}

\paragraph{The exact octonion-coordinate glue test.}
This embedding has the following fully explicit image, denoted
\(\Lambda_{\mathbb O}=\mathcal L(N)\).  Given \(q\), first recover
\((x;d)\) by \eqref{n24:eq:inverse-first}--\eqref{n24:eq:inverse-second}.
For each fibre require
\(x_{i,h}-x_{i,k}\in\mathbb Z\) for every \(h,k\), and put
\[
 k_i=-6x_{i,0}\pmod6\in\mathbb Z/6\mathbb Z.
\]
This is defined because
\(6x_{i,0}=\sum_h(x_{i,0}-x_{i,h})\in\mathbb Z\).
Require also that the four \(d_k\) are all integral or all
half-integral.  Let \(y(d)\in\mathbb F_2^2\) be the unique class with
\[
 d+D=0+D,\ v+D,\ s+D,\ c+D
 \quad\hbox{marked respectively by }00,10,01,11.
\]
Then the final and exact test is
\begin{equation}
 (k_1,k_2,k_3,k_4;y(d))
   =\bigl(3u+2(a+b,a+2b,a,b);\tau(u)\bigr)
 \label{n24:eq:image-glue}
\end{equation}
for \(u\in\mathbb F_2^4\) of even weight,
\((a,b)\in\mathbb F_3^2\), with the first four coordinates read
modulo six, and \(\tau(u)=(u_2+u_4,u_3+u_4)\).
Indeed the integral-difference conditions are precisely membership in
\(A^\vee\), and the common fractional part equals
\(-k_i/6\pmod{\mathbb Z}\), which is the fractional part of
\(k_i\lambda\).  The half-integral test is precisely membership in
\(D^\vee\), and \eqref{n24:eq:image-glue} is the original
\(\mathcal H\).  This proves both directions of the image
description.  No new lattice or rescaled glue is substituted.

\paragraph{The determinant and the factor two in the Albert metric.}
The matrix convention, also specifying this construction independently
of its host document, is
\[
 Q_{\mathbb O}(\lambda;q)=
 \begin{pmatrix}
 \lambda&q_3&\overline q_2\\
 \overline q_3&\lambda&q_1\\
 q_2&\overline q_1&\lambda
 \end{pmatrix}\in\mathfrak h_3(\mathbb O),\qquad
 X\circ Y=(XY+YX)/2.
\]
Put
\[
 E(q)=Q_{\mathbb O}(0;q_1,q_2,q_3),\qquad
 \mathfrak t(q_1,q_2,q_3)=\operatorname{Re}((q_1q_2)q_3),
 \qquad \mathfrak t_N(X)=\mathfrak t(\mathcal LX).
\]
The original ordered Albert determinant becomes the exact polynomial
\begin{equation}
 \det_{\mathbb O}\bigl(\lambda I+E(\mathcal LX)\bigr)
   =\lambda^3-\lambda(X,X)+2\mathfrak t_N(X).
 \label{n24:eq:niemeier-albert-cubic}
\end{equation}
Moreover
\begin{equation}
 \operatorname{tr}\bigl(E(\mathcal LX)\circ E(\mathcal LY)\bigr)
                  =2(X,Y).
 \label{n24:eq:albert-metric-two}
\end{equation}
For completeness, direct binary multiplication gives
$\operatorname{tr}E=0$, $\operatorname{tr}(E\circ E)=2\sum_i|q_i|^2$,
and $\operatorname{tr}(E\circ(E\circ E))=6\mathfrak t(q)$.
For the last equality, the three cyclic off-diagonal coordinates of
$E\circ E$ are $\overline{q_2q_3},\overline{q_3q_1},
\overline{q_1q_2}$; their trace pairings with the corresponding
$q_i$ each give $2\mathfrak t(q)$.
The Jordan trace formula
$\det X=\operatorname{tr}(X\circ(X\circ X))/3
-\operatorname{tr}X\operatorname{tr}(X\circ X)/2
+(\operatorname{tr}X)^3/6$
therefore gives the first identity after expanding $X=\lambda I+E$
and substituting the proved norm equality. For the second,
each off-diagonal octonion
occurs in both transposed matrix entries, giving exactly
\(2\sum_i\operatorname{Re}(q_i\overline{p_i})\).
Real associativity and cyclicity of octonion real parts give
\(\mathfrak t(q_3,q_1,q_2)=\mathfrak t(q_1,q_2,q_3)\):
\[
 \operatorname{Re}((q_3q_1)q_2)
 =\operatorname{Re}(q_3(q_1q_2))
 =\operatorname{Re}((q_1q_2)q_3).
\]
Consequently \(\mathfrak t_N(rX)=\mathfrak t_N(X)\), and the
whole determinant in \eqref{n24:eq:niemeier-albert-cubic} is
\(r\)-invariant with \(\lambda\) fixed.  Equivalently, for the real
permutation matrix \(P e_i=e_{i+1}\), indices modulo three,
\[
 P\bigl(\lambda I+E(q)\bigr)P^T
             =\lambda I+E(q_3,q_1,q_2).
\]
This is also a Jordan automorphism: multiplication by a real
permutation matrix merely reindexes the two-factor matrix products,
so it preserves \((AB+BA)/2\).  It therefore carries the positive
cone and its determinant boundary to themselves.  The cyclic bridge
is simultaneously linear, isometric on \(\mathbb O^3\), integral
on the stated lattice image, and compatible with the actual Albert
cubic and positive cone.

\paragraph{Dependence and the four-dimensional lattice summand.}
The map uses precisely the four specified origins \(f_i\), the
generator \(289\), the even/odd triplet order, the displayed order
of \(b\), the ordered \(D\)-coordinates, and the chosen orthonormal
octonion basis.  Changing one of these changes the displayed map;
the construction makes no claim of independence from these data.
Its restriction to the actual \(D\)-summand is
\begin{equation}
 d\longmapsto\frac1{\sqrt3}(z(d),z(d),z(d)),\qquad
 z(d)=d_1u_4+d_2u_5+d_3u_6+d_4u_7.
 \label{n24:eq:d-diagonal-embedding}
\end{equation}
This is an isometry on \(D\), and on \(D^\vee\) in its real span.
In particular its 24 shortest nonzero dual vectors become 24
specified points in this four-dimensional diagonal subspace.
The separate Cartan-weight map \(d\mapsto\sum_kd_k\varepsilon_k\)
is related to it by the exact isometry
\(\sum_kd_k\varepsilon_k\mapsto(z(d),z(d),z(d))/\sqrt3\).

\subsection{The Leech neighbour in the same three octonion coordinates}
\label{n24:sec:leech-coordinates}

The preceding map also transports the entire neighbour construction,
not only the real span of $N$. Put $N_{\mathbb O}=\mathcal L(N)$ and
retain the actual $\rho,\alpha,K,\Lambda$ of
Theorem~\ref{esnl:thm:leech-neighbor}. In the already fixed basis let
\[
 U=\sum_{a=0}^7u_a,\qquad
 V=\frac{u_0+u_1+u_2+u_3}{\sqrt2}
       +\sqrt3u_4+\frac{2u_5+u_6}{\sqrt3}.
\]
Substitution in the original isometry gives the explicit vectors
\begin{equation}
 \begin{aligned}
 \varrho=\mathcal L\rho&=(V+2U,V,V-2U),\\
 \delta=\mathcal L\alpha
 &=\left(\frac{(2+\sqrt2)u_0-2u_1}{3},
          \frac{(\sqrt2-1)u_0+u_1}{3},
          \frac{(\sqrt2-1)u_0+u_1}{3}\right).
 \end{aligned}
 \label{n24:eq:weyl-in-octonions}
\end{equation}
Indeed each even triplet of $\rho_A$ has mean $1/2$, each odd
triplet has mean $-1/2$, and both residual triplets are $(2,0,-2)$.
The four mean-difference coordinates are $\sqrt6/2$, while the
four $D$ coordinates are $(3,2,1,0)$. For $\alpha$, the two residual
triplets in the first fibre are $(2,-1,-1)/3$ and $(-2,1,1)/3$;
its only nonzero mean-difference coordinate is $2/\sqrt6$.
These facts prove every coefficient in
\eqref{n24:eq:weyl-in-octonions}. In particular
$|\varrho|^2=84$, $|\delta|^2=2$, and
$\langle\varrho,\delta\rangle=1$ in the product octonion metric.

Define
\begin{equation}
 \begin{aligned}
 K_{\mathbb O}&=\{q\in N_{\mathbb O}:
                          \langle\varrho,q\rangle\equiv0\pmod6\},\\
 \Lambda_{\mathrm{Leech},\mathbb O}
       &=K_{\mathbb O}+\mathbb Z(\varrho/6-\delta).
 \end{aligned}
 \label{n24:eq:leech-in-octonions}
\end{equation}
The isometry and its inverse prove, in both directions,
$\mathcal L(K)=K_{\mathbb O}$ and
$\mathcal L(\Lambda)=\Lambda_{\mathrm{Leech},\mathbb O}$.
Thus \eqref{n24:eq:leech-in-octonions} is a full coordinate
description of the Leech lattice inside the very same three sectors.
The glue test above, followed by its displayed height congruence,
is an exact membership test for $K_{\mathbb O}$; adjoining the six
cosets $k(\varrho/6-\delta)$, $0\le k<6$, gives all of the Leech
lattice. The common-sublattice indices, discriminant form, and
order-twelve quotient are transported without alteration.

The cyclic action also acts on the marked family of neighbours.
Retain $r=\rho(361)$ and its corresponding map
$c(q_1,q_2,q_3)=(q_3,q_1,q_2)$.
For any transported Weyl vector $\rho'=T\rho$, $T\in O(N)$,
the neighbour is defined using $T\alpha$, so that
$\langle\rho',T\alpha\rangle=1$. Then
\[
 c\,\mathcal L(K_{\rho'})=\mathcal L(K_{r\rho'}),\qquad
 c\,\mathcal L(\Lambda_{\rho'})=\mathcal L(\Lambda_{r\rho'}).
\]
This follows directly from $\mathcal Lr=c\mathcal L$ and
(NL18). The inverse maps use $r^{-1}$ and $c^{-1}$.
It specifies exactly which marked neighbour is the target; it does
not silently assume that the chosen $\rho$ is fixed by $r$.
Finally the common cubic evaluates on every one of these lattice
points as
\[
 \det Q_{\mathbb O}(\lambda;\mathcal LX)
    =\lambda^3-\lambda(X,X)
                   +2\operatorname{Re}(((\mathcal LX)_1
                                         (\mathcal LX)_2)(\mathcal LX)_3).
\]
This identity holds for all $X\in N_{\mathbb R}$, hence also
on the explicitly embedded $N$, $K$ and $\Lambda$.

\subsection{Exactly which root isometries preserve this glue}
\label{n24:sec:glue-symmetry}

Retain the coordinate convention \((P_px)_i=x_{p(i)}\).
An arbitrary root-lattice isometry has the original form
\[
 T(X_1,X_2,X_3,X_4;d)
   =\bigl((\varepsilon_iQ_{\sigma_i}X_{p(i)})_{i=1}^4;\ T_Dd\bigr),
\]
where \(\varepsilon_i\in\{\pm1\}\), \(\sigma_i\in S_6\),
\(p\in S_4\), and \(T_D\in O(D)\).  If \(B=d_D(T_D)\),
the necessary and sufficient condition for \(T(N)=N\) is
\begin{equation}
 B\tau(u)=\tau(P_pu)\ (u\in E),\qquad
 \bigl(\varepsilon_i\ell_{p(i)}\bigr)_{i=1}^4
                  =\bigl(\ell_iA\bigr)_{i=1}^4
 \quad\hbox{for some }A\in GL_2(\mathbb F_3),
 \label{n24:eq:all-glue-preservers}
\end{equation}
where
\(\ell_1=(1,1),\ell_2=(1,2),\ell_3=(1,0),\ell_4=(0,1)\).
The \(A\), and hence \(B=B_p\), are unique.  Indeed the
two-primary part of the glue is the graph of \(\tau\), and signs
are trivial modulo two.  Its three-primary part is the image of the
injective map \(j(a,b)=(a+b,a+2b,a,b)\), so preserving it is
equivalent to \(M_\varepsilon P_pj=jA\).  These are exactly
the two conditions in \eqref{n24:eq:all-glue-preservers}; the
primary decomposition proves their sufficiency as well as necessity.
Since the four \(\ell_i\) represent the four distinct projective
row lines over \(\mathbb F_3\), each \(A\) gives exactly one
signed permutation by this equation.  This recovers the precise
48-element glue quotient, not merely its action on three labels.

The canonical full-automorphism theorem proves, with these same
coordinates and section, that
\[
 O(N)=\bigl(S_6^4\times W(D_4)\bigr)\rtimes L(G),
 \qquad G\cong GL_2(\mathbb F_3).
\]
Conjugation by the actual map \(\mathcal L\) now gives the entire
orthogonal stabilizer of the embedded lattice, with explicit matrices
obtained from \eqref{n24:eq:full-isometry} and its inverse:
\begin{equation}
 O(\Lambda_{\mathbb O})
       =\mathcal L O(N)\mathcal L^{-1},\qquad
 \Gamma_G(g)=\mathcal LL_g\mathcal L^{-1},\qquad
 \mathcal LL_g=\Gamma_G(g)\mathcal L.
 \label{n24:eq:full-transported-action}
\end{equation}
Both inclusions in the first equality follow by conjugating an
isometry preserving the specified lattice.  This is a faithful
action on the whole 24-dimensional space. Each \(\Gamma_G(g)\)
commutes with \(c\),
because the canonical equation
\(L_g\rho(289^2)=\rho(289^2)L_g\) conjugates to that equation.

For completeness, the previously displayed \(D_4\)-triality
section and geometric frame permutations have a precise comparison
even though that particular proposed linear conjugacy fails.  On
\(V_N\), each of \(\widehat R,\widehat H\) has trace
\(5-5+2=2\): the transposed \(A_5\) pair contributes zero,
the other two components contribute \(5,-5\), and the
\(D_4\) reflection contributes 2.  Their three-cycle has trace
\(5+1=6\).  A geometric Albert-frame transposition acts on
\(\mathbb O^3\) by swapping two sectors with conjugation and by
conjugating the fixed sector; its trace is
\(1-7=-6\).  Its three-cycle has trace zero.  Therefore an
invertible map cannot conjugate these two particular \(S_3\)
representations, since trace is invariant under conjugation.
The complete real character decompositions are, respectively,
\[
 V_N|_{\langle\widehat R,\widehat H\rangle}
       \cong7\mathbf1\oplus5\operatorname{sgn}\oplus6\operatorname{Std},
 \qquad
 \mathbb O^3|_{S_3^{\rm frame}}
       \cong\mathbf1\oplus7\operatorname{sgn}\oplus8\operatorname{Std}.
\]
These multiplicities follow by taking inner products with the
characters \((1,1,1),(1,-1,1),(2,0,-1)\), on classes of sizes
\(1,3,2\).  The exact whole-space relation for the original section
is \eqref{n24:eq:full-transported-action}; the exact
determinant-compatible cyclic relation is
\eqref{n24:eq:cyclic-intertwining}.  In the latter relation the
residue element has trivial discriminant action while its octonion
action cycles the sectors.  Those are two explicitly calculated
actions of the same operator on two explicitly different constructions;
no identification between their quotient homomorphisms is assumed.

Finally, the connected group \(\operatorname{Spin}(8)\) acts on
the ambient triality module, and its subgroup preserving
\(\Lambda_{\mathbb O}\) is exactly its intersection with
\eqref{n24:eq:full-transported-action}.  It is finite: the columns
of a lattice isometry in a fixed lattice basis have prescribed
bounded lengths, so each column has finitely many possible lattice
values.  A connected group preserving a discrete full lattice
pointwise as a set acts trivially, since each orbit of a lattice
vector is connected and discrete and the lattice spans.  The
nontrivial connected triality action therefore does not preserve
the whole lattice.  The proved cyclic subgroup \(\langle c\rangle\)
does preserve both the lattice and its Albert cubic.  Equations
\eqref{n24:eq:all-glue-preservers} and
\eqref{n24:eq:full-transported-action} specify the full glue
stabilizer; they do not assert that every element of that finite
orthogonal group also preserves the cubic.

% Proof-bearing staged module. Namespace nca:. No canonical edits.
\section{The arithmetic of the retained Albert cubic on both lattices}
\label{nca:sec:cubic-arithmetic}

\subsection{Original data and the exact polynomial}

Retain the ordered coordinates \(X=(x_{i,j};d_k)\),
\(1\leq i\leq4\), \(0\leq j\leq5\), \(1\leq k\leq4\), with
\(\sum_{j=0}^5x_{i,j}=0\), and their original Euclidean metric.
Let
\[
 A_5=\{x\in\mathbb Z^6:\sum_jx_j=0\},\qquad
 D_4=\{d\in\mathbb Z^4:\sum_kd_k\equiv0\pmod2\},\qquad
 R=A_5^{\perp4}\perp D_4.
\]
Write \(\lambda=(5,-1,-1,-1,-1,-1)/6\),
\(v_D=(1,0,0,0)\), and \(s_D=(1,1,1,1)/2\).
The original discriminant glue is generated by
\[
 \begin{split}
 &(3,3,3,3;00),\quad(3,3,0,0;10),\quad(3,0,3,0;01),\\
 &(2,2,2,0;00),\quad(2,4,0,2;00)
 \end{split}
\]
in \((\mathbb Z/6)^4\oplus\mathbb F_2^2\), with the same marking
\(10=v_D+D_4\), \(01=s_D+D_4\). Its order is 72: the first
three generators have independent two-primary reductions, and the
last two independent three-primary reductions.  The lattice \(N\)
is its inverse image in \(R^\vee\).

Retain
\[
 \rho_A=(5,3,1,-1,-3,-5)/2,\qquad
 \rho_D=(3,2,1,0),\qquad \rho=\rho_A^{\oplus4}\perp\rho_D,
 \qquad \alpha=(e_0-e_1,0,0,0;0),
\]
and the exact neighbour
\[
 K=\{X\in N:(\rho,X)\equiv0\pmod6\},\qquad
 v=\rho/6-\alpha,\qquad \Lambda=K+\mathbb Zv.
\]
Here \(\rho\in N\), \(\rho^2=84\), and \((\rho,\alpha)=1\),
as follows immediately from these coordinates and the first glue
generator.  The name \(\Lambda\) always denotes this neighbour.

Fix an ordered orthonormal octonion basis \(u_0=1,u_1,\ldots,u_7\).
For the first part of the calculation no multiplication table is
imposed beyond the actual multiplication of these basis elements.
Define, in the original coordinates,
\begin{align}
 s_i&=\sum_{j=0}^5(-1)^jx_{i,j},\label{nca:eq:original-s}\\
 A_{2(i-1)+p}
  &=x_{i,p}-\frac{x_{i,p}+x_{i,p+2}+x_{i,p+4}}3,
       &&p=0,1,\label{nca:eq:original-A}\\
 B_{2(i-1)+p}
  &=x_{i,p+2}-\frac{x_{i,p}+x_{i,p+2}+x_{i,p+4}}3,
       &&p=0,1.\label{nca:eq:original-B}
\end{align}
Write \(A=A_0+\mathbf a\), \(B=B_0+\mathbf b\), where
\(\mathbf a=\sum_{h=1}^7A_hu_h\) and
\(\mathbf b=\sum_{h=1}^7B_hu_h\).  The common mean octonion is
\begin{equation}
 M=\frac{\sqrt2}{6}\sum_{a=0}^3s_{a+1}u_a
       +\frac{\sqrt3}{3}\sum_{k=1}^4d_ku_{k+3}.
 \label{nca:eq:mean}
\end{equation}
The retained isometry therefore has precisely the ordered form
\[
 \mathcal LX=(A+M,B+M,-A-B+M).
\]
This is an identity with its original three triplets: the residuals
sum to zero, the first four common mean coefficients are
\(s_i/\sqrt{18}=\sqrt2s_i/6\), and the last four are
\(d_k/\sqrt3=\sqrt3d_k/3\).
For completeness the original squared norm and the output norm are both
\[
 |A|^2+|B|^2+|A+B|^2+\frac16\sum_i s_i^2+\sum_kd_k^2.
\]
In the original fibres this follows by separating the means
\(s_i/6,-s_i/6\) from their zero-sum triplets; in the output it
follows from \(A+B+(-A-B)=0\) and
\(3|M|^2=\sum_i s_i^2/6+\sum_kd_k^2\). The inverse recovers
\(M=(q_1+q_2+q_3)/3\), then \(s_i=3\sqrt2M_{i-1}\),
\(d_k=\sqrt3M_{k+3}\), and \(A=q_1-M,B=q_2-M\).
Adding \((-1)^ps_i/6\) to the \(p\)-th residual triplet in the
\(i\)-th fibre recovers each original \(x_{i,p+2j}\). Thus the
map is onto and preserves the original pairing by polarization.

For imaginary octonions put
\[
 \mathbf a\times\mathbf b=\frac{\mathbf a\mathbf b-
                                    \mathbf b\mathbf a}{2},\qquad
 \phi(\mathbf a,\mathbf b,\mathbf c)
                     =\langle\mathbf a\times\mathbf b,\mathbf c\rangle.
\]
Let \(\phi_{ijk}=\phi(u_i,u_j,u_k)\).  Thus \(\phi\) is the
actual alternating octonion three-form in the chosen frame, including
all signs. Define
\begin{align}
 D&=A_0^2+A_0B_0+B_0^2-\langle\mathbf a,\mathbf a\rangle
             -\langle\mathbf a,\mathbf b\rangle
             -\langle\mathbf b,\mathbf b\rangle,\label{nca:eq:D}\\
 Z&=(2A_0+B_0)\mathbf a+(A_0+2B_0)\mathbf b
                               -3\mathbf a\times\mathbf b
       =\sum_{h=1}^7Z_hu_h.\label{nca:eq:Z}
\end{align}
Equivalently, the component in the final formula is the exact
quadratic polynomial
\[
 Z_h=(2A_0+B_0)A_h+(A_0+2B_0)B_h
                  -3\sum_{i,j=1}^7\phi_{ijh}A_iB_j.
\]

\begin{theorem}[Exact pullback with every radical and coefficient]
\label{nca:thm:polynomial}
The original cubic
\(F(X)=\operatorname{Re}(((\mathcal LX)_1(\mathcal LX)_2)
                                           (\mathcal LX)_3)\)
has the complete polynomial expression
\begin{equation}
 F(X)=P_0(X)+\sqrt2P_2(X)+\sqrt3P_3(X),
 \label{nca:eq:polynomial}
\end{equation}
where every variable is the displayed linear function of the original
\(x_{i,j},d_k\), and
\begin{align}
 P_0={}&-A_0B_0(A_0+B_0)+B_0\langle\mathbf a,\mathbf a\rangle
              +A_0\langle\mathbf b,\mathbf b\rangle
              +2(A_0+B_0)\langle\mathbf a,\mathbf b\rangle,
                \label{nca:eq:P0}\\
 P_2={}&\frac{-s_1D+s_2Z_1+s_3Z_2+s_4Z_3}{6}
       +\frac{s_1(4s_1^2-3\sum_{i=1}^4s_i^2)}{108}
       -\frac{s_1\sum_{k=1}^4d_k^2}{6},\label{nca:eq:P2}\\
 P_3={}&\frac{d_1Z_4+d_2Z_5+d_3Z_6+d_4Z_7}{3}.
                  \label{nca:eq:P3}
\end{align}
These identities hold throughout the original real vector space, in
particular on \(N\), \(K\), and \(\Lambda\). They retain the ordered
Albert determinant as
\[
 \det Q_{\mathbb O}(\lambda;\mathcal LX)
                =\lambda^3-\lambda(X,X)+2F(X).
\]
\end{theorem}

\begin{proof}
For \(a=a_0+\mathbf a\), \(b=b_0+\mathbf b\),
\(c=c_0+\mathbf c\), binary octonion multiplication gives
\[
 \operatorname{Re}((ab)c)
  =a_0b_0c_0-a_0\langle\mathbf b,\mathbf c\rangle
   -b_0\langle\mathbf a,\mathbf c\rangle
   -c_0\langle\mathbf a,\mathbf b\rangle
   -\phi(\mathbf a,\mathbf b,\mathbf c).
\]
Call the first four, symmetric, terms \(S(a,b,c)\).  Substitute
\((a,b,c)=(A+M,B+M,-A-B+M)\).  Terms with exactly two copies
of \(M\) in \(S\) sum to \(S(M,M,A+B-A-B)=0\). Terms
with exactly one copy sum to
\(-S(M,A,A)-S(M,A,B)-S(M,B,B)\).  The term with no copy is
\(S(A,B,-A-B)=P_0\), by multiplying its four displayed terms.
The alternating term with no copy vanishes; terms with two or
three copies vanish. Its three terms with one copy all equal
\(\phi(\mathbf a,\mathbf b,\operatorname{Im}M)\), since
\[
 \phi(\mathbf a,\operatorname{Im}M,-\mathbf a-\mathbf b)
 =\phi(\mathbf a,\mathbf b,\operatorname{Im}M),\quad
 \phi(\operatorname{Im}M,\mathbf b,-\mathbf a-\mathbf b)
 =\phi(\mathbf a,\mathbf b,\operatorname{Im}M).
\]
Thus all one-copy terms together are
\(-M_0D+\langle\operatorname{Im}M,Z\rangle\).
Finally the term with three copies is
\[
 \operatorname{Re}(M^3)=4M_0^3-3M_0|M|^2,
 \qquad |M|^2=\frac{\sum_i s_i^2}{18}
                              +\frac{\sum_kd_k^2}{3},
 \qquad M_0=\frac{\sqrt2s_1}{6}.
\]
Substitution proves \eqref{nca:eq:P0}--\eqref{nca:eq:P3} with
their displayed factors. In particular it proves the absence of a
\(\sqrt6\) term without discarding any mixed term. The determinant
formula uses the original isometry and its unchanged squared norm.
\end{proof}

\subsection{Finite generators of both original lattices}

Order the 24 original simple roots as
\[
 g_{5(i-1)+j+1}=a_{i,j}=e_{i,j}-e_{i,j+1}
 \quad(1\leq i\leq4,\ 0\leq j\leq4),
\]
followed by
\[
 g_{21}=e_1-e_2,\quad g_{22}=e_2-e_3,\quad
 g_{23}=e_3-e_4,\quad g_{24}=e_3+e_4
\]
in the original \(D_4\) coordinates.  All unspecified components
are zero.  Adjoin the five actual glue lifts
\begin{align*}
 g_{25}&=(3\lambda,3\lambda,3\lambda,3\lambda;0),&
 g_{26}&=(3\lambda,3\lambda,0,0;v_D),\\
 g_{27}&=(3\lambda,0,3\lambda,0;s_D),&
 g_{28}&=(2\lambda,2\lambda,2\lambda,0;0),\\
 g_{29}&=(2\lambda,4\lambda,0,2\lambda;0).
\end{align*}
Then
\begin{equation}
 N=\sum_{i=1}^{29}\mathbb Zg_i,\qquad
 H_i=(\rho,g_i)=
 \begin{cases}
 1,&1\leq i\leq24,\\
 30,18,18,15,20,&i=25,26,27,28,29\text{ respectively}.
 \end{cases}
 \label{nca:eq:N-generators}
\end{equation}
Indeed the first 24 vectors span \(R\), the last five generate
the original glue, and subtraction of the corresponding glue lift
leaves every element of \(N\) in \(R\).  The height calculation
uses \((\rho_A,\lambda)=5/2\),
\((\rho_D,v_D)=(\rho_D,s_D)=3\), and the simple-root heights one.

Define the following 30, fully specified neighbour generators:
\begin{equation}
 h_1=6\alpha,\qquad h_i=g_i-H_i\alpha\quad(2\leq i\leq29),
 \qquad h_{30}=v=\rho/6-\alpha.
 \label{nca:eq:Lambda-generators}
\end{equation}
Then \(K=\sum_{i=1}^{29}\mathbb Zh_i\) and
\(\Lambda=\sum_{i=1}^{30}\mathbb Zh_i\).
For if \(X=\sum_i n_ig_i\in K\), then
\[
 X=\sum_{i=2}^{29}n_i(g_i-H_i\alpha)
                         +\left(\sum_{i=1}^{29}n_iH_i\right)\alpha,
\]
and the last coefficient is divisible by six. Conversely every
displayed generator before \(h_{30}\) has height divisible by six.
Adjoining \(h_{30}\) is exactly the original neighbour definition.

\subsection{A complete finite calculation of the cubic value modules}

For a lattice \(L\) and a homogeneous real cubic \(F\), define
its additive value module by
\[
 \mathcal V_F(L)=\sum_{X\in L}\mathbb ZF(X)\subset\mathbb R.
\]
This definition retains every value and takes its generated additive
group. It does not claim that every element of the group is one
individual value of \(F\).

\begin{theorem}[Finite cubic-value generators, including negative integers]
\label{nca:thm:finite-values}
If \(L=\sum_{i=1}^m\mathbb Zb_i\), allowing dependent generators,
then \(\mathcal V_F(L)\) is generated by the following finite list:
\begin{align}
 A_i={}&F(b_i),&&1\leq i\leq m,\label{nca:eq:single}\\
 D^+_{ij}={}&F(b_i+b_j)-F(b_i)-F(b_j),&&i<j,\label{nca:eq:Dplus}\\
 D^-_{ij}={}&F(b_i-b_j)-F(b_i)+F(b_j),&&i<j,\label{nca:eq:Dminus}\\
 D_{ijk}={}&F(b_i+b_j+b_k)-F(b_i+b_j)-F(b_i+b_k)\notag\\
          &{}-F(b_j+b_k)+F(b_i)+F(b_j)+F(b_k),&&i<j<k.
 \label{nca:eq:Dthree}
\end{align}
Consequently \eqref{nca:eq:N-generators} gives exactly 4495
specified generators of \(\mathcal V_F(N)\), and
\eqref{nca:eq:Lambda-generators} gives exactly 4960 specified
generators of \(\mathcal V_F(\Lambda)\), for every fixed
orthonormal octonion frame and its actual \(\phi_{ijk}\).
\end{theorem}

\begin{proof}
Let \(B_F\) be the symmetric trilinear polarization, with
\(B_F(X,X,X)=F(X)\); explicitly its value is
\(1/6\) times the seven-term expression in
\eqref{nca:eq:Dthree}, with its three inputs replacing the three
generators. For every integer tuple \(n\), expansion gives
\[
 F\left(\sum_i n_ib_i\right)
 =\sum_i n_i^3A_i
 +\sum_{i<j}\bigl(n_i^2n_j C_{ij}+n_in_j^2 E_{ij}\bigr)
 +\sum_{i<j<k}n_in_jn_kD_{ijk},
\]
where \(C_{ij}=3B_F(b_i,b_i,b_j)\) and
\(E_{ij}=3B_F(b_i,b_j,b_j)\).  The two pair generators are
\(D^+_{ij}=C_{ij}+E_{ij}\) and
\(D^-_{ij}=-C_{ij}+E_{ij}\). Therefore their contribution is
\[
 \frac{n_in_j(n_i+n_j)}2D^+_{ij}
       +\frac{n_in_j(n_j-n_i)}2D^-_{ij}.
\]
Both coefficients are integers, for all positive, zero, and negative
integers \(n_i,n_j\): if either integer is even the claim is immediate,
and if both are odd their sum and difference are even. This proves
that the finite list generates all values. Conversely each entry
of the list is an integer linear combination of actual values on
\(L\). The proof does not use linear independence of the generators.
Finally \(29+2\binom{29}{2}+\binom{29}{3}=4495\) and
\(30+2\binom{30}{2}+\binom{30}{3}=4960\).
\end{proof}

For numerical arithmetic now specify the additional octonion frame
choice completely.  Use \(\mathbb O=\mathbb H\oplus\mathbb H\ell\)
in ordered coordinates \((a,b)\), with
\begin{equation}
 \begin{gathered}
 (a,b)(c,d)=(ac-\overline d\,b,\ da+b\overline c),\\
 u_0=(1,0),\quad u_1=(i,0),\quad u_2=(j,0),\quad u_3=(k,0),\\
 u_4=(0,1),\quad u_5=(0,i),\quad u_6=(0,j),\quad u_7=(0,k).
 \end{gathered}
 \label{nca:eq:Cayley-frame}
\end{equation}
Here \(ij=k\), \(jk=i\), \(ki=j\), and the reversed products
are negative.  The positive oriented three-form triples are exactly
\begin{equation}
 (123),\ (145),\ (176),\ (246),\ (257),\ (347),\ (365).
 \label{nca:eq:seven-triples}
\end{equation}
Every \(\phi_{ijk}\) is obtained from these by alternation, or is
zero.  This is a specified frame for the retained, previously
frame-dependent isometry, without a change to any lattice coordinates
or to the Albert matrix convention.

\begin{theorem}[Exact value modules and their quotient of order twelve]
\label{nca:thm:value-modules}
For the frame \eqref{nca:eq:Cayley-frame}, the exact additive modules are
\begin{align}
 \mathcal V_F(N)
    &=\frac1{27}\mathbb Z\ \oplus\
       \frac{\sqrt2}{54}\mathbb Z\ \oplus\
       \frac{\sqrt3}{18}\mathbb Z,\label{nca:eq:N-values}\\
 \mathcal V_F(\Lambda)
    &=\frac1{27}\mathbb Z\ \oplus\
       \frac{\sqrt2}{216}\mathbb Z\ \oplus\
       \frac{\sqrt3}{54}\mathbb Z.\label{nca:eq:Lambda-values}
\end{align}
In particular the inclusion of these two additive subgroups of
\(\mathbb R\) has the exact quotient
\begin{equation}
 \mathcal V_F(\Lambda)/\mathcal V_F(N)
                  \cong\mathbb Z/4\oplus\mathbb Z/3
                  \cong\mathbb Z/12.
 \label{nca:eq:value-quotient}
\end{equation}
This is an inclusion of cubic value modules; it does not assert the
false lattice inclusion \(N\subset\Lambda\).
\end{theorem}

\begin{proof}
The polynomial coefficients in
\eqref{nca:eq:P0}--\eqref{nca:eq:P3} are rational on both lattices
for this exact multiplication table.  The numbers
\(1,\sqrt2,\sqrt3\) are linearly independent over \(\mathbb Q\).
For example, isolating one radical in a rational relation and
squaring either gives a rational multiple of another irrational
radical, or a rational square equal to two, three, or \(3/2\),
which is impossible by parity of prime exponents. Thus a value has
a unique rational coefficient triple \((P_0,P_2,P_3)\).

Apply the complete finite list in
Theorem~\ref{nca:thm:finite-values} to the displayed original
generators. Exact substitution gives the following least common
multiples of coefficient denominators. The denominators in each
entry are ordered as \((P_0,P_2,P_3)\), and a zero coefficient has
denominator one:
\begin{equation}
\begin{array}{c|c|c|c|c}
 &A_i&D^+_{ij}&D^-_{ij}&D_{ijk}\\\hline
 N&(27,27,1)&(9,18,18)&(9,18,18)&(9,18,18)\\
 \Lambda&(27,216,27)&(27,108,54)&(27,108,54)&(9,36,54)
\end{array}.
\label{nca:eq:denominator-certificate}
\end{equation}
This table is a finite rational calculation, with no bound on
lattice points or heuristic sampling.  The complete arithmetic
certificate is the accompanying \texttt{verify\_cubic.py}: its
\texttt{residual\_data} and \texttt{cubic} functions are exactly
\eqref{nca:eq:original-s}--\eqref{nca:eq:P3}; it forms every
\(A_i,D^+_{ij},D^-_{ij},D_{ijk}\) for the 29 or 30 explicitly
defined vectors, using rational arithmetic.  The \texttt{cross}
function executes the three cyclic positive and three negative
entries for each of the seven triples
\eqref{nca:eq:seven-triples}.  Hence the full finite calculation
is specified by these equations, the finite index bounds in
Theorem~\ref{nca:thm:finite-values}, and integer arithmetic.
Its 4495 and 4960 entries verify the upper inclusions in
\eqref{nca:eq:N-values}--\eqref{nca:eq:Lambda-values}.

Here are short certificates for the reverse inclusions. Coordinates
in the next two tables are relative to the respective three
proposed generators on the right sides of
\eqref{nca:eq:N-values} and \eqref{nca:eq:Lambda-values}.
In the first table all pair symbols use the \(g_i\):
\begin{equation}
\begin{array}{c|r r r}
 F(g_1)&-4&4&0\\
 F(g_2)&-5&-4&0\\
 F(g_7)&0&18&0\\
 F(g_{28})&-80&52&0\\
 D^+_{21,28}&0&-36&16\\
 D^+_{2,27}&-27&27&3\\
 D^+_{2,25}&-81&45&0
\end{array}.
\label{nca:eq:N-lower-certificate}
\end{equation}
The determinants of the three columns obtained by transposing
rows \((1,4,5)\), \((2,3,6)\), and \((2,7,6)\) are,
respectively, \(1792,-270,-1647\). Their greatest common divisor
is one. In the second table all pair symbols use the \(h_i\):
\begin{equation}
\begin{array}{c|r r r}
 F(h_2)&14&-272&0\\
 F(h_4)&5&-128&0\\
 F(h_{28})&13780&-74672&0\\
 F(h_{30})&14&-63&-2\\
 D^+_{2,27}&7425&-60624&-9
\end{array}.
\label{nca:eq:Lambda-lower-certificate}
\end{equation}
Rows \((1,3,4)\), \((1,4,5)\), and \((2,4,5)\) give
determinants \(-5405504,2315394,1281267\), whose greatest
common divisor is also one.

For clarity, the elementary integer argument using these
determinants does not require a normal-form theorem.  Let
\(W\subset\mathbb Z^3\) be the group generated by one displayed
table.  For any three of its columns forming a matrix \(C\),
the adjugate identity \(C\operatorname{adj}(C)=\det(C)I\)
implies \(\det(C)\mathbb Z^3\subset W\).  Bezout's identity
for the three displayed determinants, whose gcd is one, therefore
gives \(\mathbb Z^3\subset W\). The opposite inclusion is
immediate from the integer entries. Every table entry is an
integer combination of actual values by its definition, proving
the two reverse inclusions.  Finally the coordinates of the
inclusion of the first module in the second are
\((a,b,c)\mapsto(a,4b,3c)\). Its cokernel is the displayed
product of cyclic groups, with order twelve. This also proves
the asserted exact quotient map.
\end{proof}

\subsection{Frame dependence and exact symmetry transport}

The finite-list theorem is valid for the actual coefficients of
every orthonormal frame. The numerical modules
\eqref{nca:eq:N-values}--\eqref{nca:eq:Lambda-values} refer to
the frame \eqref{nca:eq:Cayley-frame}.  This dependence has an
exact comparison map. If \(u_a'=Tu_a\) for a real isometry
\(T\) fixing \(1\), then
\[
 \mathcal L_{u'}=(T,T,T)\mathcal L_u.
\]
The symmetric part \(S\) of the trilinear tensor is unchanged,
because \(T\) preserves the scalar coordinate and inner product.
Writing \(\mathbf m=\operatorname{Im}M\), its exact change is
\begin{equation}
 F_{u'}(X)-F_u(X)
  =-3\bigl(\phi(T\mathbf a,T\mathbf b,T\mathbf m)
                  -\phi(\mathbf a,\mathbf b,\mathbf m)\bigr).
 \label{nca:eq:frame-transport}
\end{equation}
This follows from the complete expansion in the polynomial proof.
If \(T\) is an octonion automorphism it preserves \(\phi\), so
the cubic and its value modules are unchanged.

The dependence on a general orthonormal frame is substantive. Take
the actual lattice vector with the only nonzero fibre coordinates
\[
 x_{1,1}=1,\ x_{1,5}=-1,\quad
 x_{2,2}=1,\ x_{2,4}=-1,\qquad d=(1,1,0,0).
\]
It lies in \(R\subset N\), has \(A=u_1\), \(B=u_2\), and
\(M=(u_4+u_5)/\sqrt3\). In the Cayley frame its value is zero.
Replace only
\[
 u'_3=(2u_3-\sqrt5u_4)/3,\qquad
 u'_4=(\sqrt5u_3+2u_4)/3.
\]
This is an exact orthonormal frame change; the same original
lattice vector now has value \(-\sqrt{15}/3\), because
\(\phi(u_1,u_2,u'_4)=\sqrt5/3\) and
\(\phi(u_1,u_2,u_5)=0\). This value is outside
\(\mathbb Q+\sqrt2\mathbb Q+\sqrt3\mathbb Q\). Indeed the
independent sign changes of \(\sqrt2,\sqrt3\) are automorphisms
of \(\mathbb Q(\sqrt2,\sqrt3)\). If an element \(z\) in that
field had \(z^2=15\), each sign change would send it to
\(z\) or \(-z\). In the basis
\(1,\sqrt2,\sqrt3,\sqrt6\) this makes \(z\) a rational
multiple of a single basis vector. Its square would force one
of \(15,15/2,5,5/2\) to be a rational square, impossible by
parity of its prime exponents.
Thus the numerical specialization is not silently extended to
an unspecified orthonormal frame.

Return to the explicit Cayley frame for the following numerical
witnesses. Let \(\sigma=\rho(289)\) act simultaneously by
\(e_{i,j}\mapsto e_{i,j+1\bmod6}\), and trivially on \(D_4\).
The exact isometry gives
\(\mathcal L\sigma^2=c\mathcal L\) for
\(c(q_1,q_2,q_3)=(q_3,q_1,q_2)\). Cyclic real associativity
therefore proves \(F(\sigma^{2k}X)=F(X)\) for every real \(X\).
These maps preserve \(N\) by their original discriminant action.
The stabilizer of this cubic within the six-element residue group
is exactly its three-element subgroup: direct substitution at
\(\alpha=g_1\) gives
\begin{equation}
 F(\alpha)=\frac{-4+2\sqrt2}{27},\qquad
 F(\sigma\alpha)=\frac{-5-2\sqrt2}{27},
 \label{nca:eq:residue-nonpreserver}
\end{equation}
which differ. All odd powers have the same effect on the cubic
as \(\sigma\), by its proved invariance under even powers.

The original reflection exchanging \(e_{1,0},e_{1,1}\) sends
\(\alpha\) to \(-\alpha\), hence changes its nonzero cubic
value to its negative. The original full glue-preserving lifts
\[
 \begin{aligned}
 \widehat R(X_1,X_2,X_3,X_4;d)
       &=(X_2,X_1,X_3,-X_4;R_Dd),\\
 \widehat H(X_1,X_2,X_3,X_4;d)
       &=(X_1,X_3,X_2,-X_4;H_Dd).
 \end{aligned}
\]
also have explicit witnesses:
\[
 F(\widehat R\alpha)=0,\qquad
 F(g_7)=\frac{\sqrt2}{3},\qquad F(\widehat Hg_7)=0.
\]
Their original matrices remain
\[
 R_D=\operatorname{diag}(-1,1,1,1),\qquad
 H_D=\frac12\begin{pmatrix}1&1&1&1\\1&1&-1&-1\\
                           1&-1&1&-1\\1&-1&-1&1\end{pmatrix}.
\]
These maps preserve the original lattice, as can be checked on its
displayed generators.  Both \(R_D\) and \(H_D\) preserve \(D_4\):
the former changes one sign, while the latter has integral coordinates
on every even-sum integer vector and has output sum twice its first
input coordinate. They are orthogonal involutions. Their marked
discriminant actions fix \(v_D\) and exchange \(s_D,c_D\), and
exchange \(v_D,s_D\) while fixing \(c_D\), respectively.
On the binary glue generators, \(\widehat R\) fixes the first
two and sends the third to their second-plus-third sum;
\(\widehat H\) fixes the first and exchanges the second and third.
On the ternary code generators \(z_1=(1,1,1,0)\),
\(z_2=(1,2,0,1)\), the two signed component permutations act as
\((z_1,z_2)\mapsto(z_1,2z_2)\) and
\((z_1,z_2)\mapsto(z_1,2z_1+2z_2)\), respectively.
Thus both preserve the full original glue and its inverse image.
Every within-fibre coordinate permutation also preserves that inverse
image, since its action on \(\lambda+A_5\) is trivial:
\(Q\lambda-\lambda=Qe_0-e_0\in A_5\).
These are concrete failures to preserve this cubic, without
discarding their actual lattice actions or glue.

Negation preserves each lattice \(N\) and \(\Lambda\), and its
exact polynomial law is \(F(-X)=-F(X)\). In the neighbour the
nonzero witness is its displayed norm-four generator:
\begin{equation}
 F(\rho)=15\sqrt2+16\sqrt3,\qquad
 F(v)=\frac{14}{27}-\frac{7\sqrt2}{24}-\frac{\sqrt3}{27}.
 \label{nca:eq:special-values}
\end{equation}
Both equalities follow by substituting the original coordinates
of \(\rho\) and \(v\) into the proved polynomial.

Every original lattice isometry \(T\in O(N)\), whether or not
it preserves the cubic, has the exact replacement transport
\begin{equation}
 \Gamma_T=\mathcal LT\mathcal L^{-1},\qquad
 F^T(Y)=F(T^{-1}Y),\qquad F^T(TX)=F(X),
 \label{nca:eq:all-transport}
\end{equation}
and
\[
 T\Lambda_\rho=\Lambda_{T\rho},\qquad
 \mathcal V_{F^T}(\Lambda_{T\rho})=\mathcal V_F(\Lambda_\rho),
 \qquad \mathcal V_{F^T}(N)=\mathcal V_F(N).
\]
For the neighbour transport, write its exact elements as
\(x+k\rho/6\), where \(x\in N\) and
\((\rho,x)+k\equiv0\pmod6\); application of \(T\) gives
the identical congruence with \((T\rho,Tx)\). The inverse is
\(T^{-1}\). The value-module equalities then follow by the
bijections and the displayed identity \(F^T(TX)=F(X)\).
In particular for \(T=\sigma^{2k}\), one has \(F^T=F\), so
the same retained cubic has exactly the same value module on
every neighbour \(\Lambda_{\sigma^{2k}\rho}\) in this cyclic
orbit. This statement names its target neighbour and does not
assume the chosen \(\rho\) is fixed.

% Exact follow-on module. Input after the original Niemeier and Leech proofs.
% Original Euclidean coordinates, residue action, scales and glue are retained.
\section{The residue cubic action and the common integral Leech correspondence}
\label{esrf:sec:fixed-neighbor}

Retain the exact lattices $R=A^{\perp4}\perp D\subset N$ and the
binary and ternary codes from the index-72 construction.
The four $A$ coordinates remain the ordered six-element residue
fibres. To avoid confusing a vector with a representation, write
$\varrho$ for the Weyl vector
\[
 \varrho_A=(5,3,1,-1,-3,-5)/2,\qquad
 \varrho_D=(3,2,1,0),\qquad
 \varrho=\varrho_A^{\oplus4}\perp\varrho_D.
\]
Thus $\varrho\in N$ and $\varrho^2=84$.
The original residue representation is still denoted by $\rho$.
The specific element studied here is
\[
 r=\rho(361)=\rho(289)^2,\qquad
 r|_{A_i}=C^2,\quad Ce_j=e_{j+1\bmod6},\qquad r|_D=I.
 \tag{RF1}
\]
In particular $r^3=I$. This is the actual residue coordinate
permutation, with no conjugation or replacement.

For each $\sigma\in\Omega=O(N)\varrho$, choose an original root
$\alpha_\sigma$ with $(\sigma,\alpha_\sigma)=1$ and set
\[
 \begin{split}
 \ell_\sigma(x)&=(\sigma,x)\bmod6,\qquad
 K_\sigma=\ker(\ell_\sigma:N\to\mathbb Z/6),\\
 \Lambda_\sigma&=K_\sigma+
                      \mathbb Z(\sigma/6-\alpha_\sigma)\\
 &=\{x+k\sigma/6:x\in N,\ k\in\mathbb Z,\
                                    (\sigma,x)+k\equiv0\pmod6\}.
 \end{split}
 \tag{RF2}
\]
The preceding neighbor proof shows that the definition is independent
of $\alpha_\sigma$, every $\Lambda_\sigma$ is the specified rootless
even unimodular lattice, and
$N\cap\Lambda_\sigma=K_\sigma$ with both indices six.
The family considered in this section is precisely
$\mathscr F=\{\Lambda_\sigma:\sigma\in\Omega\}$.

\subsection{A necessary and sufficient equality criterion}

\begin{theorem}[Equality of the marked Weyl neighbors]
\label{esrf:thm:equality}
For $\sigma,\tau\in\Omega$, the following three statements are equivalent:
\[
 \Lambda_\sigma=\Lambda_\tau,\qquad
 K_\sigma=K_\tau,\qquad
 \tau-\epsilon\sigma\in6N
                 \text{ for some }\epsilon\in\{1,-1\}.
 \tag{RF3}
\]
\end{theorem}

\begin{proof}
Equality of the two neighbors gives equality of the kernels by
intersection with $N$.
Each $\ell_\sigma$ is onto because
$(\sigma,\alpha_\sigma)=1$.
If the kernels are equal, their induced isomorphisms
$N/K_\sigma\to\mathbb Z/6$ differ by an automorphism of
$\mathbb Z/6$. Its two units are $1,-1$.
Therefore $(\tau-\epsilon\sigma,N)\subset6\mathbb Z$.
Unimodularity $N^\vee=N$ gives
$(\tau-\epsilon\sigma)/6\in N$.
Conversely this congruence makes the two kernels equal.

It remains to show that this kernel criterion actually gives
equality of the chosen unimodular extensions.
The exact formula (RF2), after replacing $k$ by $-k$, gives
$\Lambda_{-\sigma}=\Lambda_\sigma$.
We may consequently suppose $\tau=\sigma+6t$ with $t\in N$.
Both Weyl vectors have squared norm 84, so
\[
 0=\tau^2-\sigma^2=12(\sigma,t)+36t^2,\qquad
                       (\sigma,t)=-3t^2\in6\mathbb Z.
 \tag{RF4}
\]
The last inclusion uses the actual evenness of $N$.
For $x+k\tau/6$ in (RF2), put $x'=x+kt\in N$.
Then $x+k\tau/6=x'+k\sigma/6$, and
\[
 (\sigma,x')+k
 \equiv(\sigma,x)+k
 \equiv(\tau,x)+k\equiv0\pmod6.
\]
This proves $\Lambda_\tau\subset\Lambda_\sigma$.
Interchanging $\sigma,\tau$ proves the reverse inclusion.
Thus no additional unverified choice of an extension is being
identified merely from equality of its common kernel.
\end{proof}

It follows in particular that the map
\[
 \Omega/{\sim}\ \longrightarrow\ \mathscr F,\qquad
 [\sigma]\longmapsto\Lambda_\sigma,\qquad
 \sigma\sim\tau\ \Longleftrightarrow\
                          \tau\equiv\pm\sigma\pmod{6N},
 \tag{RF5}
\]
is a bijection. It is $O(N)$-equivariant: the original formula
$T\Lambda_\sigma=\Lambda_{T\sigma}$ follows by substitution in
(RF2), and $T$ preserves the congruence relation since $T(N)=N$.

\subsection{The exact tetracode obstruction to a fixed neighbor}

\begin{theorem}[No fixed member of this specified family]
\label{esrf:thm:no-fixed-neighbor}
The actual residue element $r=\rho(361)$ fixes no lattice of
$\mathscr F$. Every orbit of its action on $\mathscr F$ has
exactly three members.
\end{theorem}

\begin{proof}
Suppose first that $r\sigma\equiv-\sigma\pmod{6N}$.
Applying $r$ successively, using $r(N)=N$ and $r^3=I$, gives
$\sigma\equiv-\sigma\pmod{6N}$.
Then $\sigma/3\in N$, contradicting its pairing
$1/3$ with the root $\alpha_\sigma\in N$.
Thus (RF3) permits fixedness only if
$t=(r\sigma-\sigma)/6\in N$.

Every isometry of $N$ preserves the original root lattice, permutes
its four $A_5$ components, and acts in each of them by a signed
coordinate permutation, as proved in the full automorphism theorem.
Negation permutes the six entries of $\varrho_A$.
Consequently each $A$ block $s=(s_0,\ldots,s_5)$ of every
$\sigma\in\Omega$ is a permutation of the six distinct numbers
$5/2,3/2,1/2,-1/2,-3/2,-5/2$.
In this block put $\Delta_j=s_{j-2\bmod6}-s_j$.
If the corresponding block $\Delta/6$ of $t$ belongs to $A^\vee$,
its coordinate differences are integral, so
\[
 \Delta_j\equiv c\pmod6\qquad\text{for every }j
 \tag{RF6}
\]
for one integer $c$ modulo six.
On either three-cycle of the index permutation, the three
$\Delta_j$ sum to zero. Hence $3c\equiv0\pmod6$ and
$c\in\{0,2,4\}$.
Each $\Delta_j$ is an integer between $-5$ and $5$.
If $c=0$, (RF6) forces all $\Delta_j=0$, contrary to the
distinctness of the entries on each three-cycle.
Therefore $c=2$ or $4$.

In the original discriminant marking, a vector of the $A$ class
$t_i\lambda+A$ has common fractional part $-t_i/6$.
The corresponding class of $\Delta/6$ is therefore
$t_i=-c\pmod6$, also in $\{2,4\}$.
The difference on the $D$ block is exactly zero because $r|_D=I$.
Thus membership of the full vector $(r\sigma-\sigma)/6$ in $N$
would require a glue word
\[
 (t_1,t_2,t_3,t_4;00)\in H,\qquad t_i\in\{2,4\}
                                                \text{ for every }i.
 \tag{RF7}
\]
In $H=\{(3u+2z;y)\}$ this forces $u=0$ and all four $z_i$
nonzero. But the exact ternary code is
\[
 C_3=\{(a+b,a+2b,a,b):a,b\in\mathbb F_3\}.
 \tag{RF8}
\]
If exactly one parameter is zero, exactly one coordinate is zero.
If both are nonzero, precisely one of $a+b,a+2b$ is zero.
Thus every nonzero word has weight three, and the zero word
has weight zero. No word in (RF7) belongs to $H$.
This contradicts $t\in N$.

Neither sign in (RF3) can occur. Since $r^3=I$, every orbit
therefore has cardinality three.
\end{proof}

The conclusion concerns exactly $\mathscr F$ and the action (RF1).
The integral correspondence for these three-member orbits is
constructed next; no assertion about other embeddings or
constructions of a Leech lattice follows from this obstruction.

\subsection{The canonical orbit and its common integral lattice}

Put $\sigma_j=r^j\varrho$, $\Lambda_j=\Lambda_{\sigma_j}$ and
$K_j=K_{\sigma_j}$, with $j$ read modulo three.
Keep the same original root
$\alpha=e_0-e_1$ in the first $A$ block.
Directly $(\sigma_j,\alpha)=1$ for all three $j$.
Define
\[
 v_j=\sigma_j/6-\alpha,\quad v=v_0,\qquad
 q_1=(\sigma_1-\sigma_0)/6,\quad
 q_2=(\sigma_2-\sigma_0)/6.
 \tag{RF9}
\]
The complete original coordinates are
\[
 \begin{split}
 q_1&=\bigl((-2,-2,1,1,1,1)/3\bigr)^{\oplus4}\perp0,\\
 q_2&=\bigl((-1,-1,-1,-1,2,2)/3\bigr)^{\oplus4}\perp0.
 \end{split}
 \tag{RF10}
\]
Their classes in $R^\vee/R$ are respectively
$(4,4,4,4;00)$ and $(2,2,2,2;00)$.
Neither is in $H$, while $3q_1,3q_2\in R$.
In particular each $q_i+N$ has exact order three.
Direct pairings give
\[
 q_1^2=q_2^2=16/3,\quad(q_1,q_2)=8/3,\quad
                    (\varrho,q_1)=(\varrho,q_2)=-16,
 \qquad q_1+q_2\in R.
 \tag{RF11}
\]

For $x\in N$, write its original glue class as $(3u+2z;y)$,
with $z=(a+b,a+2b,a,b)\in C_3$, and define the homomorphism
\[
 \beta:N\longrightarrow\mathbb F_3,\qquad\beta(x)=b=z_4.
 \tag{RF12}
\]
This map is defined through the original quotient $N/R$.
This definition uses the retained residue marking without changing it.
The sum of the four $z$ coordinates is also $b$ modulo three.
Write $\ell_j(x)=(\sigma_j,x)\bmod6$.
Then the exact three character formulas are
\[
 \ell_1(x)=\ell_0(x)+4\beta(x),\qquad
 \ell_2(x)=\ell_0(x)+2\beta(x)\quad\text{in }\mathbb Z/6.
 \tag{RF13}
\]
The maps $b\mapsto2b,4b$ from $\mathbb F_3$ to $\mathbb Z/6$
are well-defined homomorphisms.
To prove (RF13), the discriminant bilinear form gives
\[
 (q_1,x)\equiv\tfrac56\,4\sum t_i
       \equiv\tfrac13\sum t_i
       \equiv\tfrac23\beta(x)\pmod{\mathbb Z},
 \qquad t_i=3u_i+2z_i\pmod6.
\]
For $q_2$ the same calculation gives
$(q_2,x)\equiv\beta(x)/3\pmod{\mathbb Z}$.
Multiplication by six and (RF9) prove both formulas.

Using the original dominant representatives
$\omega_t=(1^t,0^{6-t})-\frac t6(1,\ldots,1)$, put
\[
 u_*=(\omega_2,\omega_4,0,\omega_2;0)\in N.
 \tag{RF14}
\]
Its ternary class is exactly the original generator
$z_2=(1,2,0,1)$.
It has $u_*^2=4$, $(\varrho,u_*)=12$ and $\beta(u_*)=1$.
The map
\[
 (\ell_0,\beta):N\longrightarrow\mathbb Z/6\oplus\mathbb F_3
 \tag{RF15}
\]
is onto: $\alpha$ maps to $(1,0)$ and $u_*$ to $(0,1)$.
Consequently
\[
 \begin{gathered}
 I=\ker(\ell_0,\beta)=K_0\cap K_1=K_1\cap K_2
     =K_2\cap K_0=\bigcap_jK_j,\\
 [N:I]=18,\qquad[K_j:I]=3,\qquad\det I=324.
 \end{gathered}
 \tag{RF16}
\]
The kernel equalities follow from (RF13), because either
$2b=0$ or $4b=0$ in $\mathbb Z/6$ forces $b=0$ in
$\mathbb F_3$.
The action $r$ permutes the $K_j$, so it preserves $I$.
Also $3q_i\in I$: these are in $R$, hence have $\beta=0$,
and $(\varrho,3q_i)=-48$ is divisible by six.

Define the full-rank lattice
\[
 J=I+\mathbb Z(3v).
 \tag{RF17}
\]
It is contained in $\Lambda_0$.
The element $6v=\varrho-6\alpha$ lies in $I$: its
$\ell_0$ value is zero, and its glue has zero ternary part.
The element $3v$ is not in $N$, because
$(3v,\alpha)=-11/2$ is not an integer.
Thus
\[
 N\cap J=I,\qquad[J:I]=2,\qquad
                    \det J=324/4=81,\qquad[\Lambda_0:J]=9.
 \tag{RF18}
\]
This lattice is positive definite, even, and rootless, because
it is a sublattice of $\Lambda_0$.

All exact root corrections in the cyclic action are retained.
Put $\kappa_j=\alpha-r^j\alpha\in R$.
The vectors $r\alpha=e_2-e_3$ and $r^2\alpha=e_4-e_5$
are original simple roots with $\varrho$-pairing one.
Hence $\kappa_j\in I$.
The precise formulas, rather than literal identification of
different generators, are
\[
 r^jv=v+q_j+\kappa_j\quad(j=1,2),\qquad
 r(3v)=3v+3q_1+3\kappa_1.
 \tag{RF19}
\]
Since $rI=I$, $3q_1\in I$ and $\kappa_1\in I$, this proves
$rJ\subset J$. Applying $r^3=I$ gives equality $rJ=J$.

\begin{theorem}[All three pairwise Leech intersections]
\label{esrf:thm:common-J}
The actual three-member orbit satisfies
\[
 \Lambda_0\cap\Lambda_1
 =\Lambda_1\cap\Lambda_2
 =\Lambda_2\cap\Lambda_0
 =\bigcap_j\Lambda_j=J,\qquad[\Lambda_j:J]=9.
 \tag{RF20}
\]
The element $r$ is an integral isometry of $J$, and induces
isometries $\Lambda_j\to\Lambda_{j+1}$ with this common restriction.
\end{theorem}

\begin{proof}
First $J\subset\Lambda_0\cap\Lambda_1$: $I$ lies in both,
and $3v=3v_1-3q_1$ with $3q_1\in I$.
Suppose $y\in\Lambda_0\cap\Lambda_1$.
Using the same $\alpha$, write
$y=x_0+k_0v_0=x_1+k_1v_1$, $x_j\in K_j$.
The fractional part of $(v_j,\alpha)$ is $1/6$ for both
$j$, so pairing with $\alpha\in N$ gives
$k_0\equiv k_1\pmod6$.
Absorbing multiples of $6v_j\in K_j$, take the same
$k\in\{0,\ldots,5\}$ in both expressions.
Then $x_0-x_1=kq_1\in N$.
The exact order three of $q_1+N$ forces $k=0$ or $3$.
For these values $kq_1\in I\subset K_1$, so
$x_0=x_1+kq_1\in K_0\cap K_1=I$.
Thus $y=x_0+kv\in J$, proving the first equality.
Apply $r$ and $r^2$, using $rJ=J$, to prove the other pairwise
equalities. Their common index is nine by (RF18) and the
isometries induced by $r$.
\end{proof}

For comparison with the original index-72 glue, let
$R_0=R\cap K_0$. Since $\beta|_R=0$, (RF13) gives
$R\cap K_j=R_0$ for every $j$, and $R\cap I=R_0$.
The reduction map $K_0/R_0\cong H$ identifies $I/R_0$
with the kernel of the original ternary coefficient $b$.
Thus
\[
 [I:R_0]=24,\qquad[J:R_0]=48,\qquad
 [K_0:R_0]=72,\qquad[\Lambda_j:R_0]=432.
 \tag{RF21}
\]
All these are inclusions of the displayed original lattices.

\subsection{The exact discriminant form and its cyclic three planes}

\begin{theorem}[The integral correspondence over $\mathbb F_3^4$]
\label{esrf:thm:discriminant-correspondence}
One has
\[
 J^\vee=\Lambda_i+\Lambda_j\quad(i\ne j),\qquad
 J^\vee/J\cong\mathbb F_3^4.
 \tag{RF22}
\]
In the ordered basis
\[
 e_1=u_*+J,\quad e_2=v+J,\quad
 e_3=ru_*+J,\quad e_4=rv+J,
 \tag{RF23}
\]
the discriminant quadratic form is
\[
 q_J(a,b,c,d)=\frac23(2ad+bc+bd)\pmod{2\mathbb Z},
 \tag{RF24}
\]
and three times its discriminant bilinear form has matrix
\[
 B=
 \begin{pmatrix}
 0&0&0&2\\0&0&1&1\\0&1&0&0\\2&1&0&0
 \end{pmatrix}\quad\text{over }\mathbb F_3.
 \tag{RF25}
\]
The three original Leech extensions are exactly the inverse
images in $J^\vee$ of the isotropic planes
\[
 \begin{split}
 P_0&=\{(a,b,0,0):a,b\in\mathbb F_3\},\\
 P_1&=\{(0,0,c,d):c,d\in\mathbb F_3\},\\
 P_2&=\{(a,b,a-b,b):a,b\in\mathbb F_3\}.
 \end{split}
 \tag{RF26}
\]
The induced action of the actual $r$ is the matrix
\[
 \mathcal R=
 \begin{pmatrix}
 0&0&2&1\\
 0&0&0&2\\
 1&0&2&2\\
 0&1&0&2
 \end{pmatrix},\qquad
 \mathcal R^3=I,\quad \mathcal R^TB\mathcal R=B,
 \quad\mathcal R P_j=P_{j+1}.
 \tag{RF27}
\]
All entries in (RF23)--(RF27) use the original glue generator
$z_2$, root $\alpha$, pairing, and residue permutation.
\end{theorem}

\begin{proof}
Since $J\subset\Lambda_i=\Lambda_i^\vee$, each $\Lambda_i$
lies in $J^\vee$.
By (RF20), the two subgroups $\Lambda_i/J$ and $\Lambda_j/J$
have order nine and intersect trivially for $i\ne j$.
Their sum has order 81, equal to
$|J^\vee/J|=\det J=81$, proving the first equality in (RF22).

The homomorphism $\beta:K_0/I\to\mathbb F_3$ is an isomorphism,
with generator $u_*+I$. Hence $\Lambda_0/J$ is generated by
$u_*+J,v+J$.
Both have order dividing three because $3u_*\in I$ and
$3v\in J$.
They are independent: if $a u_*+b v\in J$, reduce first modulo
$N$; the subgroup $J+N$ has the two cosets represented by
$0,3v$, whereas $v+N$ has order six.
For $a,b\in\{0,1,2\}$ this forces $b=0$.
Then membership of $a u_*$ in $I=N\cap J$ forces $a=0$.
Thus $\Lambda_0/J\cong\mathbb F_3^2$ with the first two
generators in (RF23).
Applying $r$ proves the corresponding statement for
$\Lambda_1/J$ and the last two generators.
Their trivial intersection and total order 81 prove that
(RF23) is indeed a basis of the full discriminant group.

Here are the full Euclidean Gram entries before any passage
to a quotient:
\[
 \operatorname{Gram}(u_*,v,ru_*,rv)=
 \begin{pmatrix}
 4&2&-2&-4/3\\
 2&4&-2/3&-2/3\\
 -2&-2/3&4&2\\
 -4/3&-2/3&2&4
 \end{pmatrix}.
 \tag{RF28}
\]
For an explicit verification, each of the three nonzero
$A$ blocks of $u_*$ pairs to $-2/3$ with its shifted block,
giving $(u_*,ru_*)=-2$.
One has
$(u_*,\alpha)=(u_*,r\alpha)=(ru_*,\alpha)=0$,
$(\varrho,u_*)=12$,
$(\sigma_1,u_*)=-8$, and
$(\sigma_2,u_*)=-4$, directly from (RF10) and (RF14).
These give respectively
$(u_*,v)=2$, $(u_*,rv)=-4/3$, and
$(v,ru_*)=-2/3$.
Also $(\varrho,\sigma_1)=-12$,
$(\varrho,r\alpha)=(\sigma_1,\alpha)=1$,
and $(\alpha,r\alpha)=0$. Thus
\[
 (v,rv)=-12/36-1/6-1/6=-2/3.
\]
The four diagonal entries equal four, and
$(ru_*,rv)=(u_*,v)=2$ by the exact isometry $r$.
This proves (RF28).
Reducing its entries modulo $\mathbb Z$ gives (RF25).
The diagonal and integral cross terms in the squared norm
are even. Reducing the remaining cross terms modulo
$2\mathbb Z$ gives exactly (RF24).
Changing any coordinate by three changes (RF24) by an even
integer, so it is well-defined on the asserted field basis.

It remains to calculate the action without identifying a
family permutation with a fixed Leech automorphism.
By definition $r e_1=e_3$, $r e_2=e_4$.
The vector $u_*+ru_*+r^2u_*$ lies in $I$: its $\beta$ value
is $1+1+1=0$, because $r$ acts trivially on $N/R$,
and its $\ell_0$ value is $12-4-8=0$.
Consequently
\[
 r^2u_*\equiv-u_*-ru_*\pmod J.
 \tag{RF29}
\]
By (RF19),
\[
 v+rv+r^2v
       =3v+q_1+q_2+\kappa_1+\kappa_2.
\]
The $\kappa_j$ belong to $I$, and $3v\in J$.
Furthermore $q_1+q_2-u_*+ru_*\in I$.
Indeed $q_1+q_2\in R$, so its $\beta$ value is zero;
subtracting $u_*$ and adding $ru_*$ keeps that value zero.
Its pairing with $\varrho$ is
$-32-12-4=-48$, divisible by six.
Thus
\[
 r^2v\equiv u_*-v-ru_*-rv\pmod J.
 \tag{RF30}
\]
The four images in (RF29)--(RF30) give the four columns of
$\mathcal R$ in (RF27).
The matrix has cube the identity and preserves $B$ either
by multiplication, or by the already proved facts $r^3=I$
and $rJ=J$ with its original isometric pairing.

The first two planes in (RF26) are respectively
$\Lambda_0/J$ and $\Lambda_1/J$ by their proved bases.
Their next image is the span of
$(-1,0,-1,0)$ and $(1,-1,-1,-1)$, from (RF29)--(RF30).
Solving their two linear parameters gives exactly
$(a,b,a-b,b)$, proving the stated $P_2=\Lambda_2/J$.
Each plane is isotropic: the first two vanish in (RF24),
and the third has
$2ab+b(a-b)+b^2=3ab=0$ in $\mathbb F_3$.
They have nine elements and pairwise intersection zero,
also following directly from (RF26).
Finally each $\Lambda_j$ is exactly the inverse image of
its quotient plane because it contains $J$.
This proves the claimed integral correspondence and all
its maps, indices, pairings, and cyclic actions.
\end{proof}

\subsection{The fixed fourth plane and its exact \texorpdfstring{$A_2^{12}$}{A2 to the twelfth power} extension}

The same finite discriminant form has the additional plane
\[
 P_\infty=\{(a,b,-a+b,-b):a,b\in\mathbb F_3\}.
 \tag{RF31}
\]
Put
\[
 p=u_*-ru_*,\qquad q=v+ru_*-rv,\qquad
                       M=J+\mathbb Zp+\mathbb Zq.
 \tag{RF32}
\]
These are literal vectors in the original Euclidean space.

\begin{theorem}[The fixed fourth extension]
\label{esrf:thm:fourth-extension}
The plane $P_\infty$ is isotropic and is fixed by the induced
residue action. Its inverse image $M$ is a positive definite
even unimodular lattice with exactly 72 roots, of root system
$A_2^{12}$ and minimum squared norm two.
The original residue element $r$ is an integral isometry of $M$.
It fixes four of these $A_2$ components pointwise and acts as
a coordinate three-cycle on each of the other eight.

The exact intersection and index relations are
\[
 \begin{gathered}
 M\cap\Lambda_j=J,\qquad [M:J]=[\Lambda_j:J]=9,\\
 M\cap N=L_0:=\{x\in N:\beta(x)=0,\
                                (\varrho,x)\equiv0\pmod2\},\\
 [M:L_0]=[N:L_0]=6.
 \end{gathered}
 \tag{RF33}
\]
If $R_M$ denotes the lattice generated by all its roots, then
\[
 R_M\cong A_2^{\perp12},\qquad
 \det R_M=3^{12}=531441,\qquad[M:R_M]=3^6=729.
 \tag{RF34}
\]
\end{theorem}

\begin{proof}
Substitution of (RF31) in (RF24) gives the numerator
$2a(-b)+b(-a+b)+b(-b)=-3ab=0$ in $\mathbb F_3$.
Thus $P_\infty$ is isotropic.
The matrix (RF27) maps its vector with parameters $(a,b)$ to
\[
 (a+b,b,-a,-b),
 \tag{RF35}
\]
which is again in $P_\infty$, now with parameters $(a+b,b)$.
The images of $p,q$ in (RF23) are
$(1,0,-1,0)$ and $(0,1,1,-1)$, a basis of $P_\infty$.
Consequently (RF32) is its exact inverse image under
$J^\vee\to J^\vee/J$, and $[M:J]=9$.
Isotropy gives even norms on $M$; polarization gives integral
pairings. Its determinant is $\det J/9^2=1$, and its
positive definiteness and rank are inherited from the unchanged
ambient real space. Equation (RF35), together with $rJ=J$,
proves $rM=M$.
Each of $P_0,P_1,P_2$ intersects $P_\infty$ in zero, as follows
by equating their coordinates in (RF26) and (RF31).
Taking inverse images proves $M\cap\Lambda_j=J$.

For its relation with $N$, (RF19) gives the exact expression
\[
 q=ru_*-q_1-\kappa_1 .
 \tag{RF36}
\]
Therefore
\[
 N+M=N+\mathbb Z(\varrho/2)+\mathbb Zq_1.
 \tag{RF37}
\]
The two indicated cosets modulo $N$ have orders two and three.
For the first, pair with $\alpha$; for the second, use (RF10).
Subgroups of coprime orders intersect trivially, so the quotient
in (RF37) has order six.
Also $p\in N$ has $\beta(p)=0$ and
$(\varrho,p)=12-(-4)=16$, hence image $(4,0)$ in (RF15).
Thus $p+I$ has order three.
The vector $3q$ lies in $J$ because $q+J$ has order three,
and lies in $N$ by (RF36) and $3q_1\in N$.
It is therefore in $N\cap J=I$.
We also have $6v\in I$.
Write an element of $M$ using generators $I,3v,p,q$ and reduce
the coefficient of $3v$ modulo two and of $q$ modulo three.
If this element is in $N$, their two coprime-order cosets in
(RF37) force both reduced coefficients to vanish.
Hence
\[
 M\cap N=I+\mathbb Zp.
 \tag{RF38}
\]
Its image in (RF15) is precisely $2(\mathbb Z/6)\oplus0$,
which proves the displayed formula for $L_0$ and
$[N:L_0]=6$. Since $\det N=\det M=1$, the equal covolumes
give $[M:L_0]=6$, proving (RF33).

We next prove that the following root count is exhaustive.
Every vector in $N+M$ has one of the six forms
\[
 x+a\varrho/2-bq_1,\qquad
                   x\in N,\ a\in\{0,1\},\ b\in\{0,1,2\}.
 \tag{RF39}
\]
When $a=0$, the vector is in $R^\vee$.
Every norm-two vector of this particular $R^\vee$ is already
a root of $R$.
Indeed, the $A$ class minima are
$0,5/6,4/3,3/2,4/3,5/6$, and the three nonzero $D$ classes
have minimum one. In a nonzero class of $R^\vee/R$ with
minimum at most two, the possible minima are exactly among
\[
 \{5/6,1,4/3,3/2,5/3,11/6\}.
 \tag{RF40}
\]
This exhausts them because at most two $A$ classes can be
nonzero under the bound; if two are nonzero, each has minimum
$5/6$; with a nonzero $D$ class there is at most one nonzero
$A$ class and its minimum must be $5/6$.
No number in (RF40) is congruent to two modulo $2\mathbb Z$.
All norms within a fixed discriminant class differ by an even
integer, so no such class has a norm-two vector.
This proves the claim about $R^\vee$.

All original roots have $\beta=0$, so (RF33) says that those
lying in $M$ are exactly the ones with even $\varrho$ height.
In each original $A$ block they are
$e_i-e_j$ with $i\ne j$ and $i\equiv j\pmod2$.
These give two orthogonal $A_2$ systems per block, on the
coordinate triples $\{0,2,4\}$ and $\{1,3,5\}$, and 48 roots
in total. In $D$ they are the eight roots
$\pm e_1\pm e_3,\ \pm e_2\pm e_4$.
They form four pairwise orthogonal $A_1$ systems.
This accounts for exactly 56 roots with $a=0$ in (RF39).

For $a=1$, absorb the $A$ dual vector $-bq_1$ into the
marked $A$ discriminant classes.
If the original class of $x$ is $(t_1,t_2,t_3,t_4;y)\in H$,
the shifted class in each $A$ is $t_i-4b$.
The exact half-Weyl-shift minima proved in the neighbor
calculation are
\[
 \min_{w\in t\lambda+A}|w+\varrho_A/2|^2
 =\begin{cases}3/8& t\equiv0\pmod3,\\17/24&t\not\equiv0\pmod3,\end{cases}
 \qquad
 \min_{w+D=y}|w+\varrho_D/2|^2=1/2.
 \tag{RF41}
\]
For completeness the first line follows by choosing centered
quarter coordinates when $t=0,3$; for other $t$, the nearest
twelfth coordinates have square sum $13/24$ and coordinate
sum $\pm1$, and restoring zero sum costs $1/6$.
The second line follows from the integer coordinate distances
$(1/2,0,1/2,0)$ and half-integer distances
$(0,1/2,0,1/2)$, with their ties meeting every $D$ marking.
Thus the total minimum in (RF39) with $a=1$ is at least two.
Equality requires $t_i-4b\equiv0\pmod3$ in every component.
Since $t_i\equiv2z_i\pmod3$, this says
$z_i=2b$ for all four coordinates.
For $b=1,2$ this is excluded by the exact code (RF8).
For $b=0$ it requires $z=0$.
The eight binary codewords and the two $D$ ties then give
exactly sixteen equality vectors $y=x+\varrho/2$.
Every one belongs to $M$: the norm equation gives
\[
 2=x^2+(\varrho,x)+21,\qquad
                         (\varrho,x)=-19-x^2,
 \tag{RF42}
\]
which is odd because $N$ is even.
Also $\beta(x)=0$ because $z=0$.
Therefore $x+3\alpha\in L_0$ and
$y=(x+3\alpha)+3v\in M$.
This proves exhaustively that $M$ has exactly $56+16=72$ roots.

Here are the full root coordinates and the resulting exact
orthogonal decomposition.
Let $\zeta_i$ be supported in the $i$th original $A$ block
with entries
\[
 \zeta=(1,-1,1,-1,1,-1)/4,\qquad \zeta^2=3/8.
\]
In the following table every $Z_y$ is an actual sum of those
four mutually orthogonal vectors, and every $\delta_y$ is
an original $D$ root:
\[
 \begin{array}{c|l|l}
 y&Z_y&\delta_y\\\hline
 00&\zeta_1+\zeta_2+\zeta_3+\zeta_4&e_1+e_3\\
 10&-\zeta_1-\zeta_2+\zeta_3+\zeta_4&e_1-e_3\\
 01&-\zeta_1+\zeta_2-\zeta_3+\zeta_4&e_2+e_4\\
 11&\zeta_1-\zeta_2-\zeta_3+\zeta_4&e_2-e_4
 \end{array}.
 \tag{RF43}
\]
The sixteen new roots are exactly
\[
 \{\epsilon Z_y+\eta\delta_y/2:
                   y\in\mathbb F_2^2,\ \epsilon,\eta\in\{1,-1\}\}.
 \tag{RF44}
\]
To verify the signs from the original glue, the $A$ half-shift
minimizer for binary coordinate $u_i=0$ is $\zeta_i$ and
for $u_i=1$ it is $-\zeta_i$.
For $y=00$ the two codewords are $0000,1111$; for $10$ they
are $1100,0011$; for $01$ they are $1010,0101$; and for $11$
they are $0110,1001$.
These give respectively the four opposite pairs $\pm Z_y$
in (RF43).
The two shifted $D$ minimizers in these exact original
markings are respectively
$\pm(e_1+e_3)/2,\pm(e_1-e_3)/2,
 \pm(e_2+e_4)/2,\pm(e_2-e_4)/2$.
This proves (RF44) with every sign and marking retained.

The four $Z_y$ are mutually orthogonal and have squared norm
$3/2$, as follows by dotting their four displayed sign rows.
The $\delta_y$ are mutually orthogonal of squared norm two.
The $Z_y$ are orthogonal to the entire original $D$ space,
and to the eight old $A_2$ systems because their coordinates
are constant on each parity triple.
For each $y$, the six roots
\[
 \{\pm\delta_y\}
       \ \cup\ \{\epsilon Z_y+\eta\delta_y/2:\epsilon,\eta=\pm1\}
 \tag{RF45}
\]
are an $A_2$ system: take simple roots
$\delta_y$ and $Z_y-\delta_y/2$, whose Gram matrix is
$\left(\begin{smallmatrix}2&-1\\-1&2\end{smallmatrix}\right)$;
their sum is $Z_y+\delta_y/2$.
The four such systems are mutually orthogonal and orthogonal
to the eight old systems. They exhaust the 72 proved roots.
Thus their root lattice is $A_2^{\perp12}$ of determinant
$3^{12}$, and unimodularity of $M$ gives the index $3^6$.
This proves (RF34) and the Coxeter number three of this
rooted Niemeier lattice directly from its twelve components.

Finally $r$ fixes each $\zeta_i$ because it permutes coordinates
within their parity triples; it fixes each $\delta_y$ because
it is the identity on $D$. Hence it fixes the four systems
(RF45) pointwise. On each of the eight old $A_2$ systems it
cyclically permutes its three coordinate positions, giving its
exact order-three Coxeter action.
This completes the asserted classification of the fixed
fourth extension and its precise relation to the three Leech
lattices through $J$.
\end{proof}

% New J/M extension; prerequisite definitions are in es_niemeier_cubic_arithmetic.tex.
\section{The cubic on the common lattice and the fourth Niemeier extension}
\label{jmv:sec:JM-extension}

Retain the original coordinates, Cayley multiplication, $F$, the generators $g_i$
and their heights $H_i$ from Section~\ref{nca:sec:cubic-arithmetic}, and the
exact common lattice and fourth extension of Section~\ref{esrf:sec:fixed-neighbor}.
The symbol $u$ below is the unchanged vector $u_*$ in (RF32).
\subsection{Integral generators of the common lattice and the fixed extension}

Define the unchanged residue action by
\[
 (rX)_{i,j}=x_{i,j-2\pmod6},\qquad (rX)_{D,k}=d_k.
\]
The two primary parts of the original glue have the exact descriptions
\begin{align*}
 H_2={}&\{(\varepsilon_1+\varepsilon_2+\varepsilon_3,
 \varepsilon_1+\varepsilon_2,\varepsilon_1+\varepsilon_3,
 \varepsilon_1;\varepsilon_2,\varepsilon_3):\varepsilon_i\in\mathbb F_2\},\\
 H_3={}&\{(t+s,t+2s,t,s;00):t,s\in\mathbb F_3\}.
\end{align*}
In the original \((\mathbb Z/6)^4\) coordinates this means
\(k=3w+2z\), with \(w\) the four binary entries of \(H_2\)
and \(z\) the four ternary entries of \(H_3\). The binary
\(D_4\) marking is retained. Chinese remainders recover \(w,z\)
uniquely, and the displayed parameterizations are injective.
Consequently
\[
 \beta:N\longrightarrow\mathbb F_3,\qquad \beta(X)=z_4=s
\]
is a well-defined additive homomorphism. The actual lifts
\eqref{nca:eq:N-generators} satisfy
\[
 \beta(g_i)=0\quad(1\leq i\leq28),\qquad \beta(g_{29})=1.
\]
Indeed roots have zero glue, the first three lifts are two-primary,
and the last two have \((t,s)=(1,0),(0,1)\), respectively.

Set, with no change of representatives,
\begin{align}
 I&=\{X\in N:(\rho,X)\equiv0\pmod6,\ \beta(X)=0\},
       &J&=I+\mathbb Z(3v),\label{jmv:eq:I-J}\\
 \omega_t&=(\underbrace{1,\ldots,1}_{t},0,\ldots,0)
                    -\frac t6(1,1,1,1,1,1),
       &u&=(\omega_2,\omega_4,0,\omega_2;0),\notag\\
 p&=u-r(u),&q&=v+r(u)-r(v),\notag\\
 M&=J+\mathbb Zp+\mathbb Zq.&&\label{jmv:eq:M}
\end{align}
Thus \(J\) and \(M\) are precisely the common-lattice and
residue-fixed extension presentations being evaluated here; the
following value-group proof needs no root count, lattice determinant,
or classification label for \(M\). For explicitness, the two extra
vectors in the original five blocks are
\begin{align*}
 p={}&((1,1,-1,-1,0,0),(1,1,0,0,-1,-1),0,
                         (1,1,-1,-1,0,0);0),\\
 q={}&\tfrac13((-2,4,4,-2,-2,-2),0,(2,2,-1,-1,-1,-1),
                         (1,1,1,1,-2,-2);0).
\end{align*}
The zero blocks have their original sizes six and four. These
coordinates follow by substituting \(v=\rho/6-\alpha\) into
\eqref{jmv:eq:M}; in particular the shifted \(\alpha\) term in
\(r(v)\) is retained. Also \(u\in N\),
\(\beta(u)=1\), \((\rho,u)=12\), and \((u,u)=4\):
its class is \((2,4,0,2;00)\),
\((\rho_A,\omega_t)=t(6-t)/2\), and
\((\omega_t,\omega_t)=t(6-t)/6\).

\begin{theorem}[Integral generating lists with both kernel conditions]
\label{jmv:thm:IJM-generators}
Define the following 32 ordered vectors:
\begin{align}
 b_1&=6\alpha,&
 b_i&=g_i-H_i\alpha\quad(2\leq i\leq28),\notag\\
 b_{29}&=3g_{29}-60\alpha,&
 b_{30}&=3v=\rho/2-3\alpha,\notag\\
 b_{31}&=p,&b_{32}&=q.\label{jmv:eq:IJM-generators}
\end{align}
Then the integral, not merely rational, equalities are
\[
 I=\sum_{i=1}^{29}\mathbb Zb_i,\qquad
 J=\sum_{i=1}^{30}\mathbb Zb_i,\qquad
 M=\sum_{i=1}^{32}\mathbb Zb_i.
\]
\end{theorem}
\begin{proof}
Every one of the first 29 vectors has \(\beta=0\) and height
divisible by six, by the exact heights and characters already
calculated. Conversely take any integral expression
\(X=\sum_{i=1}^{29}n_ig_i\) for \(X\in I\). Applying
\(\beta\) shows \(n_{29}=3a\) for an integer \(a\),
regardless of which expression was chosen. The exact equality
\[
 X=\sum_{i=2}^{28}n_i(g_i-H_i\alpha)
        +a(3g_{29}-60\alpha)
        +\left(n_1+\sum_{i=2}^{28}n_iH_i+60a\right)\alpha
\]
has last coefficient \((\rho,X)\), which is divisible by six.
It therefore expresses \(X\) as an integer combination of the
first 29 \(b_i\). The reverse containment was proved for each
generator. The last two equalities follow by adjoining exactly
the vectors in \eqref{jmv:eq:I-J} and \eqref{jmv:eq:M}.
No assertion of linear independence is used.
\end{proof}

\subsection{Complete finite certificates for the two new value groups}

\begin{theorem}[Equality of the common-lattice and fixed-extension value-generated groups]
\label{jmv:thm:JM-values}
For the retained Cayley multiplication and the original ordered
coordinates, the two additive cubic-value groups are exactly
\begin{equation}
 \mathcal V_F(J)=\mathcal V_F(M)
 =\frac1{27}\mathbb Z\ \oplus\frac{\sqrt2}{216}\mathbb Z
                         \ \oplus\frac{\sqrt3}{18}\mathbb Z.
 \label{jmv:eq:JM-values}
\end{equation}
The equality concerns generated additive groups, not equality
of individual-value sets \(F(J)\) and \(F(M)\).
\end{theorem}
\begin{proof}
Apply Theorem~\ref{nca:thm:finite-values} to the first 30 and
all 32 vectors in \eqref{jmv:eq:IJM-generators}. The resulting
list lengths are exactly
\[
 30+2\binom{30}{2}+\binom{30}{3}=4960,\qquad
 32+2\binom{32}{2}+\binom{32}{3}=5984.
\]
Write all pair and triple symbols in the next tables using these
\(b_i\), not the \(g_i\) or \(h_i\). Substitution in the full
polynomial gives the complete denominator certificate
\begin{equation}
\begin{array}{c|c|c|c|c}
 &A_i&D^+_{ij}&D^-_{ij}&D_{ijk}\\\hline
 J&(27,216,1)&(9,36,18)&(9,36,18)&(9,36,18)\\
 M&(27,216,1)&(9,36,18)&(9,36,18)&(9,36,18)
\end{array}.
\label{jmv:eq:JM-denominators}
\end{equation}
As before each triple is the least common multiple of all reduced
rational denominators in its full category, in the order
\((P_0,P_2,P_3)\); a zero has denominator one. The finite
arithmetic is completely specified by
\eqref{nca:eq:original-s}--\eqref{nca:eq:P3},
\eqref{jmv:eq:IJM-generators}, and all indices in
\eqref{nca:eq:single}--\eqref{nca:eq:Dthree}.
The accompanying standard-library checker
\texttt{check\_jm\_values.py}, using the included
\texttt{cubic\_engine.py}, executes every entry with exact
fractions and asserts this table. It reconstructs all 72 marked
glue classes, checks all 29 \(\beta(g_i)\), and checks both
kernel conditions on all the displayed \(I\) generators.
These finite computations establish the upper containment in
\eqref{jmv:eq:JM-values}, including all mixed polarization terms
involving \(p,q\), not just their individual cubic values.

Here is a short certificate for the reverse containment, using
only \(J\). Coordinates of each row are taken in the ordered
proposed value-group basis
\((1/27,\sqrt2/216,\sqrt3/18)\):
\begin{equation}
\begin{array}{c|r r r}
 F(b_2)&14&-272&0\\
 F(b_4)&5&-128&0\\
 F(b_{27})&23220&-111672&-90\\
 F(b_{28})&13780&-74672&0\\
 F(b_{30})&378&-1701&-18\\
 D^+_{2,27}&7425&-60624&-3\\
 D^+_{11,21}&36&48&-16
\end{array}.
\label{jmv:eq:JM-lower-certificate}
\end{equation}
Let \(C_1,C_2,C_3\) have, respectively, columns obtained by
transposing rows \((2,3,5)\), \((2,5,6)\), and \((1,4,7)\)
of this table. Direct three-by-three determinants are
\begin{align*}
 d_1&=-39859290,&d_2&=11531403,&d_3&=-43244032.
\end{align*}
The following explicit integer identity certifies their gcd one:
\begin{equation}
 (-6882196252)d_1+(-23788908901)d_2+(-907)d_3=1.
 \label{jmv:eq:JM-bezout}
\end{equation}
Put \((c_1,c_2,c_3)=(-6882196252,-23788908901,-907)\).
For each standard basis column \(e_t\in\mathbb Z^3\) the exact
integer reconstruction is
\begin{equation}
 e_t=\sum_{j=1}^3 C_j\bigl(c_j\operatorname{adj}(C_j)e_t\bigr).
 \label{jmv:eq:JM-axis-reconstruction}
\end{equation}
All columns of every \(C_j\) represent integer combinations of
actual cubic values on \(J\); their integer coefficient vectors
in \eqref{jmv:eq:JM-axis-reconstruction} therefore give each
of the three claimed generators as such a combination. This is
an explicit reverse-containment morphism, not an index inference.
The checker also expands every adjugate and verifies all three
axis reconstructions, saving their integer coefficients in
\texttt{EXACT\_RESULTS.json}. Thus \(\mathcal V_F(J)\) contains
the right-hand side of \eqref{jmv:eq:JM-values}. Since \(J\subset M\),
the same certificate proves its containment in \(\mathcal V_F(M)\).
The upper bounds already established finish both equalities.
\end{proof}

For additional exact checks on the unchanged representatives, the
three individual values are
\[
 F(3v)=14-\frac{63}{8}\sqrt2-\sqrt3,\qquad
 F(p)=-2,\qquad F(q)=\frac{64}{27}.
\]
They agree with the corresponding rows of the complete calculation;
they are not substituted for the mixed-term certificate.

\subsection{Exact inclusions, quotient maps, and the factor of the Albert determinant}

\begin{theorem}[The exact four-to-three chain of value groups]
\label{jmv:thm:chain}
There is a strict chain of additive subgroups of the original real
coefficient field
\[
 \mathcal V_F(N)\ \subsetneq\ \mathcal V_F(J)=\mathcal V_F(M)
                       \ \subsetneq\ \mathcal V_F(\Lambda).
\]
The inclusion matrices in the ordered bases already displayed are
\[
 \mathcal V_F(N)\longrightarrow\mathcal V_F(J)=\mathcal V_F(M):
       (a,b,c)\longmapsto(a,4b,c),
\]
\[
 \mathcal V_F(J)=\mathcal V_F(M)\longrightarrow\mathcal V_F(\Lambda):
       (a,b,c)\longmapsto(a,b,3c).
\]
Their exact quotient maps are
\begin{align*}
 \frac a{27}+\frac{b\sqrt2}{216}+\frac{c\sqrt3}{18}
       &\longmapsto b\pmod4,\\
 \frac a{27}+\frac{b\sqrt2}{216}+\frac{c\sqrt3}{54}
       &\longmapsto c\pmod3.
\end{align*}
In particular
\begin{align*}
 \mathcal V_F(J)/\mathcal V_F(N)&\cong\mathbb Z/4,&
 \mathcal V_F(M)/\mathcal V_F(N)&\cong\mathbb Z/4,\\
 \mathcal V_F(\Lambda)/\mathcal V_F(J)&\cong\mathbb Z/3,&
 \mathcal V_F(\Lambda)/\mathcal V_F(M)&\cong\mathbb Z/3,\\
 \mathcal V_F(M)/\mathcal V_F(J)&=0.&&
\end{align*}
\end{theorem}
\begin{proof}
The first basis comparison uses
\(\sqrt2/54=4(\sqrt2/216)\), with the other basis entries equal.
The second uses \(\sqrt3/18=3(\sqrt3/54)\), with its other
entries equal. Uniqueness of rational radical coefficients proves
that the displayed kernels are exactly the corresponding subgroups.
Surjectivity is witnessed by \(\sqrt2/216\) and \(\sqrt3/54\),
respectively, both actual elements of the proved additive groups.
Their orders modulo the smaller groups are exactly four and three:
the coefficient divisibility conditions are necessary and sufficient.
This also proves strictness. The last quotient is zero by the equality
of the two proved groups. In particular the composite inclusion is
\((a,b,c)\mapsto(a,4b,3c)\), recovering the full quotient
\(\mathbb Z/4\oplus\mathbb Z/3\), with explicit map
\[
 \frac a{27}+\frac{b\sqrt2}{216}+\frac{c\sqrt3}{54}
       \longmapsto (b\bmod4,c\bmod3).
\]
Its identification with \(\mathbb Z/12\) is
\((\bar b,\bar c)\mapsto9b+4c\pmod{12}\), whose inverse is
reduction modulo four and modulo three. Every assertion follows
from the proved cubic coefficient groups, not from an index of
the original lattices.
\end{proof}

The factor two in the original Albert determinant is unchanged:
\[
 \det Q_{\mathbb O}(\lambda;\mathcal LX)
       =\lambda^3-\lambda(X,X)+2F(X).
\]
Thus, if one instead asks for the additive group generated by the
actual cubic term \(2F\), it is exactly \(2\mathcal V_F(L)\).
For \(J,M\) this is
\[
 \frac2{27}\mathbb Z\ \oplus\frac{\sqrt2}{108}\mathbb Z
                           \ \oplus\frac{\sqrt3}{9}\mathbb Z.
\]
Multiplication by two is an injective additive map of the real
coefficient field and restricts to the exact isomorphism
$\mathcal V_F(L)\to\mathcal V_{2F}(L)=2\mathcal V_F(L)$,
$x\mapsto2x$, with inverse $y\mapsto y/2$.
It intertwines the same two inclusion matrices and
the same quotients; no factor has been absorbed into \(F\).

\subsection{Reproduction and scope of the exact certificate}

The accompanying evidence folder is
\begin{flushleft}\small
\texttt{supporting\_materials/workbench/evidence/}\\
\texttt{period\_lattice\_continuation\_20260906/jm\_cubic/}
\end{flushleft}
It contains both full source files \texttt{cubic\_engine.py} and
\texttt{check\_jm\_values.py}. Running the latter with Python's
standard library performs the following finite proof certificate:
all \(20399\) polarization entries for \(N,\Lambda,J,M\), all
displayed denominator and determinant assertions, the integer
Bezout and adjugate axis reconstructions, and the exact inclusion
coefficient comparisons. The output \texttt{EXACT\_RESULTS.json}
retains each selected column, determinant, Bezout coefficient,
and every expanded integer linear combination reconstructing a
value-group axis.

There is additionally an independent coefficient-identity check:
the original 24 simple roots form a real basis of the original
24-dimensional space. The direct Cayley--Dickson implementation
and the rational polynomial implementation are compared on the
24 basis vectors, their 552 ordered-by-index plus/minus pairs,
and their 2024 three-element sums, a total of 2600 exact inputs.
The finite polarization proof shows these evaluations determine
every coefficient of a homogeneous cubic in those 24 variables.
Consequently this is a full polynomial-identity certificate on
the real ambient space, not sampling of lattice points. The
field implementation retains all four rational coefficients in
\(\mathbb Q(\sqrt2,\sqrt3)\), and verifies that the \(\sqrt6\)
coefficient is identically zero. A further 65 direct checks cover
the displayed new generating lists and \(u,p,q\).

This computation proves the stated additive value groups and
their morphisms. It asserts neither equality of individual-value
sets nor that the cubic group reconstructs a lattice or classifies
all cubic-preserving lattice isometries.

% Complete proof module: Wilson construction and retained-coordinate comparison.
\section{A cyclic Wilson Leech lattice in the original triality space}
\label{wcl:sec:construction}

\subsection{The two retained metrics and the octonion convention}

Let $V=\mathbb O_1\oplus\mathbb O_2\oplus\mathbb O_3$ be the original
ordered octonion space. Retain its metric and cubic
\[
\begin{aligned}
 h_0(x,x)&=N(x_1)+N(x_2)+N(x_3),\\
 \tau(x_1,x_2,x_3)&=\operatorname{Re}((x_1x_2)x_3),\qquad
 d(Q_{\mathbb O}(0;x))=2\tau(x_1,x_2,x_3),
\end{aligned}
\]
where
\[
 Q_{\mathbb O}(\lambda;x)=
 \begin{pmatrix}\lambda&x_3&\bar x_2\\
 \bar x_3&\lambda&x_1\\x_2&\bar x_1&\lambda\end{pmatrix}.
\]
Wilson's norm on the same underlying ordered space is
\[
 q_{\rm W}(x)=\tfrac12 h_0(x,x).
\]
The map
\begin{equation}
 \iota:(V,q_{\rm W})\longrightarrow(V,h_0),\qquad
 \iota(x_1,x_2,x_3)=(x_1/\sqrt2,x_2/\sqrt2,x_3/\sqrt2)
 \label{wcl:eq:metric-isometry}
\end{equation}
is an explicit isometry. Both metrics remain specified. In particular,
for the original $\lambda$, which is not rescaled,
\begin{equation}
 d(Q_{\mathbb O}(\lambda;\iota x))
 =\lambda^3-\lambda q_{\rm W}(x)+\frac1{\sqrt2}\tau(x_1,x_2,x_3).
 \label{wcl:eq:cubic-scaling}
\end{equation}
Indeed the quadratic coordinates scale by $1/2$, while the trilinear
term scales by $(1/\sqrt2)^3=1/(2\sqrt2)$ and its original coefficient
is $2$. Thus the off-diagonal cubic after $\iota$ is
$1/(2\sqrt2)$ times its value before $\iota$.

For an explicit realization inside the same $\mathbb O$, choose a
Cayley--Dickson basis
$(1,\mathsf i,\mathsf j,\mathsf k,\ell,\mathsf i\ell,
\mathsf j\ell,\mathsf k\ell)$, with
\[
 (a+b\ell)(c+d\ell)=(ac-\bar d b)+(da+b\bar c)\ell,\qquad
 \mathsf i\mathsf j=\mathsf k .
\]
Wilson's ordered orthonormal basis is realized by
\begin{equation}
 (i_\infty,i_0,i_1,i_2,i_3,i_4,i_5,i_6)
 =(1,\mathsf i,\mathsf j,\ell,\mathsf k,
                  \mathsf j\ell,-\mathsf k\ell,\mathsf i\ell).
 \label{wcl:eq:actual-basis-map}
\end{equation}
The seven oriented Fano triples are $(t,t+1,t+3)$ modulo $7$:
$i_ti_{t+1}=i_{t+3}$, with cyclic products within each triple and
opposite signs on reversal. To verify \eqref{wcl:eq:actual-basis-map},
its seven corresponding products, in order $t=0,\ldots,6$, are
\[
 \mathsf i\mathsf j=\mathsf k,\quad
 \mathsf j\ell=\mathsf j\ell,\quad
 \ell\mathsf k=-\mathsf k\ell,\quad
 \mathsf k(\mathsf j\ell)=\mathsf i\ell,\quad
 (\mathsf j\ell)(-\mathsf k\ell)=\mathsf i,\quad
 (-\mathsf k\ell)(\mathsf i\ell)=\mathsf j,\quad
 (\mathsf i\ell)\mathsf i=\ell.
\]
Alternativity supplies cyclic products on each quaternionic line.
The seven lines cover all pairs of different imaginary basis units;
squares are $-1$ and $1$ is the unit. This proves the algebra
isomorphism on the full multiplication table, hence by bilinearity on
all octonions. It also preserves conjugation and the norm. No
conjugation of a coordinate of $Q_{\mathbb O}$ is performed.

\subsection{The explicit eight-dimensional lattice and two sublattices}

In the displayed Wilson basis write $e_0=i_\infty,e_{t+1}=i_t$ and put
\[
\begin{aligned}
 D_8&=\{(z_0,\ldots,z_7)\in\mathbb Z^8:\textstyle\sum z_i\equiv0\pmod2\},\\
 s&=\tfrac12(-1+i_0+i_1+i_2+i_3+i_4+i_5+i_6),\qquad
 L_{\rm W}=D_8+\mathbb Zs.
\end{aligned}
\]
This is Wilson's lattice $L$, renamed $L_{\rm W}$ to distinguish it
from the project's original rank-24 Niemeier lattice. Since $2s\in
D_8$ and $s\notin D_8$, it is exactly $D_8\cup(s+D_8)$.

\begin{lemma}\label{wcl:lem:e8}
The lattice $L_{\rm W}$ is even and unimodular, has minimum norm $2$,
and has the explicit basis
\[
 b_1=s,\quad b_j=e_{j-1}-e_j\ (2\le j\le7),\quad b_8=e_6+e_7.
\]
Its Gram matrix for the unchanged octonion norm is
\begin{equation}
 G_8=
 \begin{pmatrix}
 2&0&0&0&0&0&0&1\\
 0&2&-1&0&0&0&0&0\\
 0&-1&2&-1&0&0&0&0\\
 0&0&-1&2&-1&0&0&0\\
 0&0&0&-1&2&-1&0&0\\
 0&0&0&0&-1&2&-1&-1\\
 0&0&0&0&0&-1&2&0\\
 1&0&0&0&0&-1&0&2
 \end{pmatrix},\qquad \det G_8=1 .
 \label{wcl:eq:e8-gram}
\end{equation}
Its norm-two vectors are precisely the $112$ vectors $\pm e_i\pm e_j$
with $i<j$ and the $128$ vectors
$\frac12(\pm e_0\pm\cdots\pm e_7)$ with an odd number of minus signs.
\end{lemma}

\begin{proof}
For $z\in D_8$, $N(z)\equiv\sum z_i=0\pmod2$ and
$h(s,z)=\frac12\sum z_i-z_0\in\mathbb Z$.
As $N(s)=2$, every norm in $D_8+\mathbb Zs$ is even.
The covolume of $D_8$ is $2$, and its displayed extension has index
$2$, so $L_{\rm W}$ has covolume $1$. The integral bilinear form follows
by polarizing the even norm, and covolume one proves self-duality.
The columns $b_i$ all belong to $L_{\rm W}$. Their determinant has
absolute value one: expansion along the first coordinate gives
$1/2$ times the determinant $2$ of the indicated $D_7$ basis.
They therefore form a basis. Their scalar products give the displayed
matrix. Positivity and evenness give minimum at least $2$, and
$N(b_i)=2$ attains it. Integer vectors of norm two have exactly two
nonzero entries $\pm1$. Half-integer vectors have norm at least
$8/4=2$, with equality exactly when all eight entries are $\pm1/2$.
Their membership in $s+D_8$ is exactly the stated odd-minus parity.
\end{proof}

Define, with the bars in the positions printed in Wilson,
\[
 \mathsf M=L_{\rm W}\bar s,\qquad \mathsf N=L_{\rm W}s .
\]
All products are right products. Let $B$ be the basis matrix with
columns $b_i$. Right multiplication by $s$ in this basis is the
following integer matrix:
\begin{equation}
 S=B^{-1}R_s B=
 \begin{pmatrix}
 3&0&0&0&0&0&0&2\\
 -2&-1&0&1&-1&1&0&-2\\
 -4&-1&-1&1&0&0&1&-3\\
 -6&-1&-1&0&0&1&0&-4\\
 -8&0&-1&0&-1&1&1&-5\\
 -10&-1&0&0&-1&0&1&-5\\
 -5&0&-1&1&-1&0&0&-2\\
 -7&0&0&0&0&0&0&-4
 \end{pmatrix}.
 \label{wcl:eq:S}
\end{equation}
This is a direct multiplication-table calculation using
\eqref{wcl:eq:actual-basis-map}; it can be checked by multiplying each
column $b_j$ by the displayed $s$. Set $T=-I_8-S$. Since
$s+\bar s=-1$, $s\bar s=2$, and alternativity allows repeated $s$
and $\bar s= -1-s$, one has
\begin{equation}
 T=B^{-1}R_{\bar s}B,\qquad S+T=-I_8,\qquad ST=TS=2I_8,
 \qquad S^{\mathsf T}G_8S=T^{\mathsf T}G_8T=2G_8 .
 \label{wcl:eq:ST-identities}
\end{equation}

\begin{lemma}\label{wcl:lem:complementary}
Both $\mathsf M,\mathsf N$ are sublattices of $L_{\rm W}$ of index
$16$ and minimum norm $4$, and
\[
 2L_{\rm W}\subset\mathsf M,\mathsf N\subset L_{\rm W},\qquad
 \mathsf M+\mathsf N=L_{\rm W},\qquad
 \mathsf M\cap\mathsf N=2L_{\rm W}.
\]
Their images $\overline{\mathsf M},\overline{\mathsf N}$ in
$\mathcal E=L_{\rm W}/2L_{\rm W}$ are complementary
four-dimensional $\mathbb F_2$ subspaces.
\end{lemma}

\begin{proof}
The integer matrices $S,T$ give the inclusions in $L_{\rm W}$;
$ST=2I$ gives the inclusions of $2L_{\rm W}$.
Norm composition shows that each right multiplication scales all
squared lengths by $N(s)=N(\bar s)=2$; thus its determinant has
absolute value $(\sqrt2)^8=16$, proving the indices and minimum norms.
The identity $S+T=-I$ gives $\mathsf M+\mathsf N=L_{\rm W}$.
Modulo $2L_{\rm W}$, each image has size $2^8/16=2^4$.
Their sum is the eight-dimensional space $\mathcal E$, so their
intersection is zero. This proves the asserted intersection upstairs.
\end{proof}

\subsection{The binary glue code and its exact generators}

Wilson's original lattice is exactly
\begin{equation}
\begin{split}
 \Lambda_{\rm W}=\{(x,y,z)\in L_{\rm W}^3:\ &
 x+y,x+z,y+z\in\mathsf M,\quad x+y+z\in\mathsf N\}.
\end{split}
 \label{wcl:eq:Wilson-definition}
\end{equation}
It carries Wilson's norm $q_{\rm W}$, as printed. Define the lattice
inside the original triality metric by
\[
 \Lambda_{\rm cyc}=\iota(\Lambda_{\rm W})\subset(V,h_0).
\]
The simultaneous use of $\mathsf M=L_{\rm W}\bar s$ for pair sums
and $\mathsf N=L_{\rm W}s$ for the total sum is essential.

Write any class in $\mathcal E$ uniquely as $m+n$, with
$m\in\overline{\mathsf M}$ and $n\in\overline{\mathsf N}$.
Then the image of \eqref{wcl:eq:Wilson-definition} in $\mathcal E^3$
is the explicit binary code
\begin{equation}
 \mathcal C=
 \{(m_1+n,m_2+n,m_1+m_2+n):
       m_1,m_2\in\overline{\mathsf M},\
       n\in\overline{\mathsf N}\}.
 \label{wcl:eq:code}
\end{equation}
Indeed the pair conditions say that the three $\mathsf N$ components
are equal; the total-sum condition says that the sum of the three
$\mathsf M$ components is zero. This proves both directions of the
code description. Its dimension is $4+4+4=12$.

For a concrete lift, the first four columns of $S$ and of $T$ are
bases of their images modulo $2$. The determinant of the eight-column
matrix $(T_1,T_2,T_3,T_4,S_1,S_2,S_3,S_4)$ is $-7$, so these columns
are independent modulo $2$; each of the two images already has
dimension four. Put $m_j=b_j\bar s$, $n_j=b_js$ for $1\le j\le4$.
A complete set of $36$ integer generators of $\Lambda_{\rm W}$ is
\begin{equation}
\begin{gathered}
 (2b_i,0,0),\ (0,2b_i,0),\ (0,0,2b_i)\quad(1\le i\le8),\\
 (m_j,m_j,0),\ (m_j,0,m_j),\ (n_j,n_j,n_j)\quad(1\le j\le4).
\end{gathered}
 \label{wcl:eq:generators}
\end{equation}
These vectors are all in \eqref{wcl:eq:Wilson-definition}.
Their reductions span precisely \eqref{wcl:eq:code}; the first
$24$ generate its full kernel $(2L_{\rm W})^3$. This proves that
they generate the entire lattice, not only a sublattice.

\begin{theorem}\label{wcl:thm:Leech}
The explicitly generated lattice $\Lambda_{\rm W}$ with $q_{\rm W}$,
and therefore $\Lambda_{\rm cyc}$ with the unaltered $h_0$, is an
even unimodular rank-24 lattice of minimum norm $4$. In particular
it is rootless. This is the actual Wilson Leech construction, through
the explicit isometry \eqref{wcl:eq:metric-isometry}.
\end{theorem}

\begin{proof}
The code proves
$[L_{\rm W}^3:\Lambda_{\rm W}]=2^{24-12}=2^{12}$.
Since $L_{\rm W}$ has covolume one for its original octonion metric,
$\Lambda_{\rm W}$ has covolume $2^{12}$ for $h_0$.
Changing to the separately recorded metric $q_{\rm W}$ multiplies
volume in dimension $24$ by $(1/\sqrt2)^{24}=2^{-12}$,
so its covolume is exactly one.

For every triple of vectors in a real inner product space,
\[
 N(x)+N(y)+N(z)
 =N(x+y)+N(x+z)+N(y+z)-N(x+y+z).
\]
For a vector of $\Lambda_{\rm W}$ the four terms on the right
are norms in $\mathsf M$ or $\mathsf N$, so each is divisible by
$4$: these are twice norms in the even lattice $L_{\rm W}$.
It follows that $q_{\rm W}(x,y,z)$ is even. Polarization of these
even norms gives an integral bilinear form. An integral lattice of
covolume one equals its dual, proving unimodularity.

If a nonzero vector had $q_{\rm W}=2$, its three coordinate norms
would sum to $4$. Each nonzero coordinate in $L_{\rm W}$ has norm
at least $2$, so one coordinate is zero. The other two then lie in
$\mathsf M$, each of whose nonzero vectors has norm at least $4$.
At most one can be nonzero. That coordinate also lies in $\mathsf N$
by the total-sum condition, so it lies in
$\mathsf M\cap\mathsf N=2L_{\rm W}$ and has norm at least $8$,
a contradiction. Hence all nonzero norms are at least $4$.
The vectors $(2b_i,0,0)$ have $q_{\rm W}=4$, so the minimum is
exactly $4$. The isometry $\iota$ transfers every statement to the
original metric.
\end{proof}

\subsection{Cyclic symmetry, the fixed lattice, and exact three-primary glue}

Retain the canonical cycle $c(x,y,z)=(z,x,y)$. Denote the opposite
cycle used in Wilson coordinates by $c_{\rm W}(x,y,z)=(y,z,x)$.
Direct substitution gives
\[
 c_{\rm W}=c^{-1}=c^2,\qquad c_{\rm W}^2=c,\qquad
 \langle c_{\rm W}\rangle=\langle c\rangle.
\]
Both actions are retained with these directions; no action is renamed
by silently reversing it.
The defining three coordinate memberships, three pair conditions,
and total-sum condition are all unchanged by $c_{\rm W}$; hence both cycles preserve
$\Lambda_{\rm W}$ and $\Lambda_{\rm cyc}$. They preserve their metrics
and the original cubic:
\[
 \operatorname{Re}((xy)z)=\operatorname{Re}((yz)x).
\]
This follows from the octonion adjoint identities
$h(ab,d)=h(b,\bar a d)=h(a,d\bar b)$, which imply real associativity,
and $\operatorname{Re}(ab)=\operatorname{Re}(ba)$.
In the Albert algebra $c_{\rm W}$ is realized by conjugation with the
\emph{real} permutation matrix for $(1\,3\,2)$.
It sends the coordinates $(x_1,x_2,x_3)$ to $(x_2,x_3,x_1)$ with no
conjugations, and permutes the ordered diagonal frame accordingly.
The inverse real permutation matrix, for $(1\,2\,3)$, realizes the
canonical $c$ and sends $(x_1,x_2,x_3)$ to $(x_3,x_1,x_2)$.
Thus the fixed lattices of the two cycles are identical. For the
retained residue isometry $\mathcal Lr=c\mathcal L$, the opposite
action satisfies $\mathcal Lr^{-1}=c_{\rm W}\mathcal L$.

\begin{theorem}\label{wcl:thm:fixed-coinvariant}
Before applying $\iota$, the cyclic fixed lattice and its orthogonal
complement in $\Lambda_{\rm W}$ are exactly
\[
 F=\{(n,n,n):n\in\mathsf N\},\qquad
 C=\{(m_1,m_2,-m_1-m_2):m_1,m_2\in\mathsf M\}.
\]
Their ranks are $8$ and $16$. For the retained Wilson metric, and
therefore after $\iota$ for the original triality metric, their Gram
matrices in the displayed bases are
\begin{equation}
 G_F=3G_8,\qquad
 G_C=\begin{pmatrix}2G_8&G_8\\G_8&2G_8\end{pmatrix}.
 \label{wcl:eq:fixed-coinvariant-grams}
\end{equation}
The fixed-lattice basis is $(b_js,b_js,b_js)$, $1\le j\le8$;
the coinvariant bases are $(b_j\bar s,0,-b_j\bar s)$ and
$(0,b_j\bar s,-b_j\bar s)$. In particular the fixed lattice has
minimum norm $6$ and determinant $3^8$.
The coinvariant lattice has determinant $3^8$, and the map which
negates its second $L_{\rm W}$ parameter identifies its displayed
Gram matrix with the usual $A_2\otimes E_8$ Gram matrix.
\end{theorem}

\begin{proof}
A fixed triple has the form $(x,x,x)$. Its pair condition holds for
every $x\in L_{\rm W}$ because $2L_{\rm W}\subset\mathsf M$.
Its total condition is $3x\in\mathsf N$. As $L_{\rm W}/\mathsf N$
has exponent two, $3x\in\mathsf N$ is equivalent to $x\in\mathsf N$.
This proves the first description. Orthogonality to the whole
diagonal real eight-space says $x+y+z=0$. The pair conditions are
then exactly $x,y,z\in\mathsf M$, and the total condition is automatic.
This proves the description of $C$.

Right multiplication by either $s$ or $\bar s$ multiplies the
octonion metric by $2$. Thus
\[
 q_{\rm W}(as,as,as)=3N(a),
\]
and
\[
 q_{\rm W}(a\bar s,b\bar s,-(a+b)\bar s)
 =N(a)+N(b)+N(a+b)=2N(a)+2N(b)+2h(a,b).
\]
Polarization gives both full Gram matrices. Their determinants are
\[
 \det G_F=3^8\det G_8=3^8,\qquad
 \det G_C=\det\begin{pmatrix}2&1\\1&2\end{pmatrix}^{8}
 (\det G_8)^2=3^8.
\]
Minimum norm $6$ for $F$ follows from the exact
factor $3$ and the minimum $2$ in $L_{\rm W}$.
The sign map $b\mapsto-b$ changes both off-diagonal blocks from
$G_8$ to $-G_8$, proving the stated exact comparison to
$A_2\otimes E_8$.
\end{proof}

There is a complete integral comparison, including the glue:
\begin{equation}
 0\longrightarrow F\oplus C\longrightarrow\Lambda_{\rm W}
 \xrightarrow{\ \Sigma\ }\mathsf N/3\mathsf N
 \longrightarrow0,\qquad
 \Sigma(x,y,z)=x+y+z\pmod{3\mathsf N}.
 \label{wcl:eq:three-primary-exact}
\end{equation}
To prove surjectivity, for every $a\in L_{\rm W}$ the explicit vector
\[
 v(a)=(as,a,-a)
\]
lies in $\Lambda_{\rm W}$ and has sum $as$. Its pair sums are
$a(s+1)=-a\bar s\in\mathsf M$,
$a(s-1)=-a\bar s-2a\in\mathsf M$, and $0$;
its total lies in $\mathsf N$. If a lattice vector has total $3n$,
subtracting $(n,n,n)\in F$ leaves a vector of $C$. This proves the
kernel assertion. Thus
\[
 \Lambda_{\rm W}/(F\oplus C)\cong\mathsf N/3\mathsf N
 \cong L_{\rm W}/3L_{\rm W}\cong(\mathbb Z/3)^8,\qquad
 [\Lambda_{\rm W}:F\oplus C]=3^8.
\]
All maps commute with $c_{\rm W}$ and with its inverse $c$; both
cycles act trivially on the quotient.

The discriminant anti-isometry is also explicit. Use the quadratic
discriminant convention $q_D(u+D)=h_D(u,u)\pmod{2\mathbb Z}$
for an even lattice $D$ with its retained metric. The orthogonal
projections of the lift $v(a)$ are
\begin{align}
 f(a)&=\tfrac13(as,as,as),\nonumber\\
 \gamma(a)&=\left(\tfrac23as,\ a-\tfrac13as,\ -a-\tfrac13as\right).
 \label{wcl:eq:glue-projections}
\end{align}
They belong to $F^*$ and $C^*$, respectively: their pairings with
$F$ or $C$ are the pairings of the integral vector $v(a)$ with the
same lattice, because the other projection is orthogonal.
The classes $f(a)+F$ identify $L_{\rm W}/3L_{\rm W}$ with $F^*/F$.
If $\gamma(a)\in C$, then $f(a)=v(a)-\gamma(a)\in
\Lambda_{\rm W}\cap F_{\mathbb R}=F$, so $a\in3L_{\rm W}$.
Since $C^*/C$ has order $\det G_C=3^8$, the map
\[
 f(a)+F\longmapsto\gamma(a)+C
\]
is an isomorphism of the two discriminant groups. Its quadratic
sign is determined without an unspecified convention:
\[
 q_F(f(a))=\tfrac13N(a),\qquad
 q_C(\gamma(a))=q_{\rm W}(v(a))-\tfrac13N(a)
              =2N(a)-\tfrac13N(a).
\]
As $N(a)$ is even, their sum $2N(a)$ is zero modulo $2\mathbb Z$.
This is the required anti-isometry. Equations
\eqref{wcl:eq:three-primary-exact} and
\eqref{wcl:eq:glue-projections} specify the overlattice itself,
including its actual coset representatives.

\subsection{Finite lattice symmetries and exact cubic tests}

Independent sign changes of the three sectors preserve
$\Lambda_{\rm W}$. Changing the sign of $x$ changes a pair sum by
$-2x\in2L_{\rm W}\subset\mathsf M$ and changes the total by
$-2x\in2L_{\rm W}\subset\mathsf N$. The same argument applies to
each sector. All six sector permutations preserve the defining
conditions too. Under signs $(\epsilon_1,\epsilon_2,\epsilon_3)$
the cubic $\tau$ is multiplied by $\epsilon_1\epsilon_2\epsilon_3$;
hence exactly the even sign changes in this subgroup preserve it.
Cyclic permutations preserve $\tau$ as already proved.

The odd permutation $(x,y,z)\mapsto(y,x,z)$ need not preserve it,
even on the lattice itself. In Wilson's basis put
\[
 a=2(1+i_0),\qquad b=2(1+i_1),\qquad d=2(i_2+i_3).
\]
Every coordinate belongs to $2L_{\rm W}$, so $(a,b,d)\in
\Lambda_{\rm W}$. The real part of
$((1+i_0)(1+i_1))(i_2+i_3)$ is $-1$, because
$i_0i_1=i_3$ and the $i_3^2$ term is the only real contribution.
Reversing the first two factors changes this contribution to $+1$.
Consequently
\begin{equation}
 \tau(a,b,d)=-8,\qquad \tau(b,a,d)=8.
 \label{wcl:eq:transposition-cubic}
\end{equation}
After the specified scalar isometry the original determinant cubic
has values $-4\sqrt2$ and $4\sqrt2$.

The simultaneous Fano-index operations
\[
 A:i_t\longmapsto i_{t+1},\qquad
 B_0:i_t\longmapsto i_{2t},\qquad i_\infty\longmapsto i_\infty
\]
are octonion algebra automorphisms, as the defining oriented Fano
products show. They are orthogonal permutations preserving $D_8$ and
the odd-minus half coset, and they fix the displayed $s$ and $\bar s$.
They consequently preserve $\mathsf M,\mathsf N,\Lambda_{\rm W}$.
Applied simultaneously to all three sectors, they preserve $\tau$.
They satisfy $A^7=B_0^3=1$ and $B_0AB_0^{-1}=A^2$; their permutations
are the distinct maps $t\mapsto 2^j t+k$ with $j=0,1,2$ and
$k\in\mathbb F_7$. Thus they form an explicitly realized
$C_7\rtimes C_3$ of order $21$.

The even sign changes and the cycle $c_{\rm W}$ form
$(\mathbb Z/2)^2\rtimes C_3\cong A_4$. They commute with the
simultaneous Fano operations, and their intersection is trivial:
a simultaneous Fano operation fixes the real unit in every sector
and does not permute sectors. Hence an explicit finite group
\[
 A_4\times(C_7\rtimes C_3)
\]
of order $252$ preserves both the literal lattice and the original
cubic. This is a proved subgroup; no assertion that it is the full
cubic stabilizer is needed.

Wilson also gives the lattice symmetries
\[
 r_t(x,y,z)=(x,yi_t,zi_t)\qquad(t=0,\ldots,6).
\]
For completeness their lattice action can be checked solely from
the explicit finite matrices above. Let
$U_t=B^{-1}R_{i_t}B$. The signed permutation multiplication table
and \eqref{wcl:eq:S} give
\begin{equation}
 U_t\in GL_8(\mathbb Z),\qquad
 T^{-1}(U_t-I_8)\in M_8(\mathbb Z),\qquad
 \tfrac12(U_t-I_8)T\in M_8(\mathbb Z)
 \quad(0\le t\le6).
 \label{wcl:eq:unit-matrix-certificates}
\end{equation}
These are finite exact matrix assertions, not assumptions about
octonionic matrix associativity. The full multiplication rule and
the exact rational algorithm checking every entry for all seven
values are included in the accompanying \texttt{check\_wilson.py}.
Concretely, start with the $64$ products
$1i_t=i_t1=i_t$, $i_t^2=-1$, and the seven oriented Fano triples,
form the eight columns $R_{i_t}(e_j)$, conjugate by the explicitly
displayed $B$, and substitute the integer $T=-I_8-S$.
For each of the three matrices in \eqref{wcl:eq:unit-matrix-certificates}
the reduced denominators of every entry are $1$, for each
$t=0,1,2,3,4,5,6$. Also $\det U_t=1$.
This is a reproducible $7\cdot3\cdot64$-entry certificate.
The first condition is also immediate from the signed permutation:
right multiplication by $i_t$ flips four coordinate signs, preserving
the odd-minus parity of $L_{\rm W}$.

The second matrix assertion says
$L_{\rm W}(i_t-1)\subset\mathsf M$, and the third says
$\mathsf M(i_t-1)\subset2L_{\rm W}$. Thus $x+yi_t$ and $x+zi_t$
differ from the original pair sums by elements of $\mathsf M$.
The pair $(y+z)i_t$ remains in $\mathsf M$, and the transformed total
differs from $x+y+z$ by
$(y+z)(i_t-1)\in2L_{\rm W}\subset\mathsf N$.
All defining conditions are preserved. The inverse follows from
$r_t^2(x,y,z)=(x,-y,-z)$ and the proved sign symmetries.
Norm composition proves that $r_t$ is an isometry.

Nevertheless none of these seven $r_t$ preserves the original cubic.
For $a_t=2(1+i_t)$, the triple $(a_t,a_t,a_t)$ lies in
$(2L_{\rm W})^3\subset\Lambda_{\rm W}$. All products below occur in
the associative real subalgebra generated by $i_t$, and give
\[
 \tau(a_t,a_t,a_t)=8\operatorname{Re}((1+i_t)^3)=-16,
\]
whereas $(1+i_t)i_t=-1+i_t$ and
\[
 \tau(a_t,a_ti_t,a_ti_t)
 =8\operatorname{Re}((1+i_t)(-1+i_t)^2)=16.
\]
Therefore the original determinant cubic on the isometric lattice
$\Lambda_{\rm cyc}$ changes from $-8\sqrt2$ to $8\sqrt2$.
This supplies actual lattice vectors as counterexamples, not only
ambient real vectors outside the lattice.

\subsection{Exact comparison with coordinate lattices and the marked lane}

The construction has the explicit integral sandwich
\begin{equation}
 (\sqrt2L_{\rm W})^3
 \ \subset\ \Lambda_{\rm cyc}\
 \subset\ (L_{\rm W}/\sqrt2)^3,
 \qquad
 [\Lambda_{\rm cyc}:(\sqrt2L_{\rm W})^3]
 =[(L_{\rm W}/\sqrt2)^3:\Lambda_{\rm cyc}]=2^{12}.
 \label{wcl:eq:coordinate-sandwich}
\end{equation}
This is \eqref{wcl:eq:code} transported by the exact isometry
\eqref{wcl:eq:metric-isometry}. The inclusions are literal inclusions
in the fixed three octonion sectors, and each commutes with both
$c_{\rm W}$ and the canonical $c=c_{\rm W}^{-1}$.
Together with the explicit bases, the code, and the three-primary
glue \eqref{wcl:eq:glue-projections}, it determines actual
commensurability maps and cosets rather than only an abstract
isometry class.

Under \eqref{wcl:eq:actual-basis-map}, the single negative column
takes the Wilson odd-minus half roots to even-minus half roots in
the current ordered Cayley--Dickson basis. The element $s$ remains
exactly
\[
 s=\tfrac12(-1+\mathsf i+\mathsf j+\mathsf k+\ell
                     +\mathsf i\ell+\mathsf j\ell-\mathsf k\ell),
\]
with coordinate vector $(-1,1,1,1,1,1,1,-1)/2$.
The signed permutation has determinant $+1$ and all $64$ basis
products agree, as independently checked in the accompanying code.
Thus the lattice, its cyclic action, and the triality cubic live
in one specified multiplication convention.

\subsection{Literal intersections with both retained lattice embeddings}
\label{wcl:sec:literal-intersections}

Let $U=\operatorname{span}_{\mathbb R}(1,\mathsf i,\mathsf j,\mathsf k)$
be the first four coordinates in the \emph{current} ordered
Cayley--Dickson basis, and let
\[
 D_4=\{a=(a_0,a_1,a_2,a_3)\in\mathbb Z^4:
                          a_0+a_1+a_2+a_3\in2\mathbb Z\}\subset U.
\]
This quaternionic $D_4$ denotes the displayed four-coordinate lattice;
the original Niemeier $D$ summand below remains in its original place.

Let $J_{\rm W}$ denote the signed permutation matrix in
\eqref{wcl:eq:actual-basis-map}. If $U_0:\mathbb R^4\to\mathbb R^8$
is the coordinate inclusion in the current basis, then
\begin{equation}
 S^{-1}B^{-1}J_{\rm W}^{-1}U_0=
 \begin{pmatrix}
 1/2&-1/2&-1/2&-1/2\\
 -1/2&0&0&1/2\\
 -1&1/2&0&1/2\\
 -3/2&1&1/2&1/2\\
 -2&1&1&1/2\\
 -5/2&3/2&1&1\\
 -5/4&3/4&3/4&3/4\\
 -7/4&3/4&3/4&3/4
 \end{pmatrix}.
 \label{wcl:eq:quaternion-membership-matrix}
\end{equation}
All matrices on the left have already been explicitly specified.
Thus the equation is a direct rational multiplication certificate.

\begin{lemma}\label{wcl:lem:quaternion-intersection}
In this fixed coordinate frame,
\[
 \mathsf N\cap U=2D_4 .
\]
\end{lemma}
\begin{proof}
Since $\mathsf N\subset L_{\rm W}$, an element of $\mathsf N\cap U$
has integral first four coordinates: the half-integer coset of the
current $E_8$ description has no zero coordinate among its last four.
Write those integral coordinates as $k_0,k_1,k_2,k_3$. Membership in
$\mathsf N$ is exactly integrality of the matrix
\eqref{wcl:eq:quaternion-membership-matrix} applied to $k$.
Subtracting its last two rows gives $k_0/2$, so $k_0$ is even.
The second row then forces $k_3$ even, the third forces $k_1$ even,
and the fourth forces $k_2$ even. Write $k=2a$.
The first six rows are now integral. The seventh and eighth rows
are integral precisely when
$-5a_0+3a_1+3a_2+3a_3$ is even, equivalently
$a_0+a_1+a_2+a_3$ is even. Conversely these conditions make all eight
rows integral. This proves both inclusions with no unstated lattice
identification.
\end{proof}

We now retain the original lattices and their actual map $\mathcal L$.
The original coordinates are $X=(x_{i,j};d_k)$,
$1\le i\le4$, $0\le j\le5$, $1\le k\le4$, with
$\sum_jx_{i,j}=0$. The root lattice is
$R=A_5^{\perp4}\perp D$, with its unchanged Euclidean metric.
Its index-$72$ overlattice $N$ is the inverse image of the original
glue
\[
 (k_1,k_2,k_3,k_4;y)
   =\bigl(3u+2(a+b,a+2b,a,b);\,(u_2+u_4,u_3+u_4)\bigr),
\]
where $u\in\mathbb F_2^4$ has even weight and $a,b\in\mathbb F_3$;
the first four entries are modulo $6$. Here $k_i=-6x_{i,0}\pmod6$
and $y$ is the original $D^\vee/D$ marking, so $d=0$ has $y=00$.
Retain the original neighbour
\[
 \rho_A=(5,3,1,-1,-3,-5)/2,\quad
 \rho_D=(3,2,1,0),\quad \rho=\rho_A^{\oplus4}\perp\rho_D,
\]
\[
 K=\{X\in N:(\rho,X)\in6\mathbb Z\},\qquad
 \alpha=(e_0-e_1,0,0,0;0),\qquad
 \Lambda=K+\mathbb Z(\rho/6-\alpha).
\]
Thus $\rho\in N$, $\rho^2=84$, $(\rho,\alpha)=1$.
Every vector of $N$ or $\Lambda$ has rational original coordinates.

For exact use of the retained isometry, set
\[
 \sigma_i=\sum_{j=0}^5(-1)^jx_{i,j},\qquad
 t_{i,p,j}=x_{i,p+2j}-\frac13\sum_{h=0}^2x_{i,p+2h}.
\]
Write $u_0=1,u_1=\mathsf i,u_2=\mathsf j,u_3=\mathsf k$ and retain
the remaining current Cayley--Dickson basis vectors $u_4,\ldots,u_7$.
The already specified map has the exact coordinate form
\[
 (\mathcal LX)_{j+1,a}=
 \begin{cases}
 t_{i,p,j}+\sqrt2\,\sigma_{a+1}/6,&0\le a\le3,\\
 t_{i,p,j}+\sqrt3\,d_{a-3}/3,&4\le a\le7,
 \end{cases}
 \qquad a=2(i-1)+p.
\]
The placement of $\sigma_{a+1}$ is retained, including its difference
from the index $i$ of the residual. No coordinate is reassigned.

\begin{theorem}\label{wcl:thm:literal-intersection}
The literal intersections in the original metric space are exactly
\begin{equation}
 \Lambda_{\rm cyc}\cap\mathcal L(N)
 =\Lambda_{\rm cyc}\cap\mathcal L(\Lambda)
 =\mathcal I
 :=\{\sqrt2(a,a,a):a\in D_4\subset U\}.
 \label{wcl:eq:literal-common-lattice}
\end{equation}
In original Niemeier coordinates the inverse image of this common
lattice is
\[
 \mathcal L^{-1}(\mathcal I)
 =\{x_{i,j}=(-1)^j a_{i-1},\ d=0:\ a\in D_4\}\subset K.
\]
It has rank four, minimum norm $12$, and Gram matrix
\begin{equation}
 G_{\mathcal I}=6
 \begin{pmatrix}
 2&-1&0&0\\
 -1&2&-1&-1\\
 0&-1&2&0\\
 0&-1&0&2
 \end{pmatrix},
 \qquad\det G_{\mathcal I}=4\cdot6^4=5184.
 \label{wcl:eq:common-gram}
\end{equation}
An explicit basis is $\sqrt2(\delta_j,\delta_j,\delta_j)$,
where
$\delta_1=(1,-1,0,0)$, $\delta_2=(0,1,-1,0)$,
$\delta_3=(0,0,1,-1)$, $\delta_4=(0,0,1,1)$.
\end{theorem}

\begin{proof}
Each coordinate of $\Lambda_{\rm cyc}$ belongs to
$\mathbb Q/\sqrt2=\sqrt2\mathbb Q$ in the current frame.
For a rational original vector $X$, the first four coordinate
formulas above therefore force $t_{i,p,j}=0$ at the corresponding
positions, by linear independence of $1,\sqrt2$ over $\mathbb Q$.
The last four force both $t_{i,p,j}=0$ and $d_{a-3}=0$, by linear
independence of $1,\sqrt2,\sqrt3$. For the latter independence, a
relation with nonzero $\sqrt3$ coefficient would place $\sqrt3$ in
$\mathbb Q(\sqrt2)$; squaring $A+B\sqrt2$ forces $AB=0$ and then
would make $3$ or $3/2$ a rational square, a contradiction.
All residuals therefore vanish and $d=0$. The zero fibre sums give
\[
 x_{i,j}=(-1)^j\sigma_i/6.
\]

First let $X\in N$. Its $A_5^\vee$ integral-difference conditions
are exactly $k_i:=\sigma_i/3\in\mathbb Z$, yielding
$x_{i,j}=(-1)^jk_i/2$. Their discriminant labels are
$-3k_i=3k_i\pmod6$. Their three-primary parts are zero, so the
injectivity of $(a,b)\mapsto(a+b,a+2b,a,b)$ forces $a=b=0$.
Since $d=0$, the binary condition is
$u_2+u_4=u_3+u_4=0$ and $\sum u_i=0$. These equations give
$u_1=u_2=u_3=u_4$. Thus all four $k_i$ have the same parity.
Conversely these conditions satisfy every original glue requirement.
The image is the diagonal triple
\[
 \mathcal LX=(k/\sqrt2,k/\sqrt2,k/\sqrt2),\qquad k\in U.
\]
By the fixed-lattice description of Wilson's construction it
belongs to $\Lambda_{\rm cyc}$ exactly when $k\in\mathsf N$.
The previous lemma makes this equivalent to $k=2a$, $a\in D_4$,
which proves the first intersection equality.

For the second equality, a vector of $\Lambda$ can be written
$X=n+m\rho/6$ with $n\in N$, $m\in\mathbb Z$.
The already proved necessary condition $d_X=0$ gives
$d_n=-m(3,2,1,0)/6\in D^\vee$.
All coordinates of a vector of $D^\vee$ are integral or all
half-integral; the fourth coordinate is zero, so all are integral.
Its third coordinate forces $6\mid m$, and hence $X\in N$.
In fact $\Lambda\cap N=K$: in
$k+m(\rho/6-\alpha)\in N$ the pairing with the original root
$\alpha$ forces $6\mid m$, and
$(\rho,\rho-6\alpha)=84-6=78\in6\mathbb Z$.
Thus an intersection vector belongs to $K$.
For the diagonal-fibre vectors just found,
\[
 (\rho,X)=\frac32\sum_{i=1}^4 k_i.
\]
Consequently the condition $X\in K$ is $\sum k_i\equiv0\pmod4$.
Every $k=2a$ with $a\in D_4$ satisfies it, proving equality of
the two intersections and the displayed inverse image.

Finally
$h_0(\sqrt2(a,a,a),\sqrt2(b,b,b))=6h(a,b)$.
The four $\delta_j$ form a basis of $D_4$, since their determinant
has absolute value two and $D_4$ has index two in $\mathbb Z^4$.
Their scalar products give \eqref{wcl:eq:common-gram}.
The least nonzero norm in $D_4$ is two: a norm-one integer vector
has odd coordinate sum, while $(1,-1,0,0)$ has norm two.
This proves the exact minimum and determinant.
\end{proof}

The two slice indices are also literal. Define the common
four-dimensional real subspace
$\mathcal U=\{(u,u,u):u\in U\}$. The proof gives
\[
 \mathcal L(N)\cap\mathcal U
 =\{(k,k,k)/\sqrt2:k\in P\},\quad
 P=2\mathbb Z^4+\mathbb Z(1,1,1,1),
\]
\[
 \mathcal L(\Lambda)\cap\mathcal U
 =\{(k,k,k)/\sqrt2:k\in P_0\},\quad
 P_0=\{k\in P:\textstyle\sum k_i\equiv0\pmod4\}
     =2D_4+\mathbb Z(1,1,1,1).
\]
The first proof does not need Wilson membership to obtain these
slice equations. The lattice $P$ has basis
$2e_0,2e_1,2e_2,(1,1,1,1)$ and covolume $8$.
The lattice $2D_4$ has covolume $2^4\cdot2=32$, so
$[P:2D_4]=4$. The class of $(1,1,1,1)$ has exact order two
modulo $2D_4$, yielding $[P_0:2D_4]=2$. Therefore
\[
 [\mathcal L(N)\cap\mathcal U:\mathcal I]=4,\qquad
 [\mathcal L(\Lambda)\cap\mathcal U:\mathcal I]=2.
\]
The rank-four intersection, strictly below rank $24$, also proves
that these full lattices are not commensurable \emph{in the retained
embedding}: a common finite-index subgroup would have full rank.
Their exact existing relation is the common lattice with the
explicit embeddings, inverse map, metrics, and indices just proved.

\subsection{The cubic value group on the literal common lattice}
\label{wcl:sec:common-cubic}

Retain $F(y)=\tau(y_1,y_2,y_3)$, so the Albert off-diagonal
determinant is $2F$. For $y=\sqrt2(a,a,a)\in\mathcal I$ write
$a=a_0+\mathbf a$ in the same quaternionic $U$. Since
$\mathbf a^2=-N(\mathbf a)$ and $a_0$ is real, direct multiplication
inside this associative subalgebra gives
\[
 \operatorname{Re}(a^3)
 =a_0^3-3a_0N(\mathbf a)=4a_0^3-3a_0N(a).
\]
Every scalar and factor is therefore
\begin{equation}
 F(\sqrt2(a,a,a))
 =2\sqrt2\bigl(4a_0^3-3a_0N(a)\bigr).
 \label{wcl:eq:common-cubic}
\end{equation}
For $a\in D_4$, $N(a)$ is even, because
$N(a)\equiv\sum a_i=0\pmod2$. The integer in parentheses is
therefore even, so every value belongs to $4\sqrt2\mathbb Z$.
For the actual lattice vector $a=(1,1,0,0)$ it equals
$4-3\cdot2=-2$, giving $F=-4\sqrt2$.
Consequently the \emph{additive subgroup generated by all values}
is exactly
\[
 \bigl\langle F(\mathcal I)\bigr\rangle_{\mathbb Z}
       =4\sqrt2\mathbb Z,\qquad
 \bigl\langle d(Q_{\mathbb O}(0;\mathcal I))\bigr\rangle_{\mathbb Z}
       =8\sqrt2\mathbb Z.
\]
This assertion concerns the generated additive groups; it does not
claim that every element of either group is an individually attained
cubic value.
For an explicit distinction, put
$P(a)=a_0(a_0^2-3(a_1^2+a_2^2+a_3^2))$.
If $3\mid P(a)$, reduction modulo three forces $3\mid a_0$,
and substitution $a_0=3b$ gives
$P(a)=9b(3b^2-a_1^2-a_2^2-a_3^2)$.
Thus $3\mid P(a)$ implies $9\mid P(a)$.
The value $P(a)=6$ is impossible, so $F=12\sqrt2$ is not
attained, although it belongs to the proved additive value group.


\subsection{Primary source and exact reproducibility data}

Wilson's octonionic construction is the source of the pair and total
conditions \cite[Sections~2--4, pp.~2--6]{esn:Wilson2009}. His page~2
uses the odd-minus lattice $L$ and $R=\bar L$; page~3 gives pair sums
in $L\bar s$, the total in $Ls$, and the metric equal to one half of
the sum of octonion norms. The displayed basis map and
\eqref{wcl:eq:metric-isometry} retain these conjugations and the exact
scale in the present coordinates. The distinct coordinate change on
Wilson's page~6 is not substituted for this isometry.

The complete integer-arithmetic program and its output are supplied in the
following directory (the two displayed lines form one path):
\begin{quote}\small
\path{supporting_materials/workbench/evidence/}\\
\path{lattice_foundation_20260906/wilson/}
\end{quote}
The files are \path{check_wilson.py} and \path{EXACT_CHECKS.json}.
They record the generating-lattice index $4096$, Gram determinant $1$,
even diagonal, integral cyclic action, the fixed and coinvariant Gram
matrices, and the index $6561$ of their sum. The same program checks
all $64$ multiplication-table products for
\eqref{wcl:eq:actual-basis-map} and the stated exact cubic tests.
The proofs above establish rootlessness and the global descriptions;
these conclusions are not inferred from a finite search.


% Complete TeX conversion of THREE_LINE_IMAGE_PROOF.md, Sections 1--5.
% Retains the original F=tau from es_niemeier_wilson_cyclic.tex.
\section{The exact cubic image of the original common lattice}
\label{d4i:sec:complete-image}

Retain the coordinates $a_0,a_1,a_2,a_3$ and the original lattice
\[
 D_4=\{a\in\mathbb Z^4:a_0+a_1+a_2+a_3\in2\mathbb Z\}.
 \label{d4i:eq:original-domain}\tag{D4I1}
\]
The polynomial already used in \cref{wcl:sec:common-cubic} is
\[
 P(a)=a_0\bigl(a_0^2-3(a_1^2+a_2^2+a_3^2)\bigr).
 \label{d4i:eq:original-polynomial}\tag{D4I2}
\]
Here and throughout this section $F$ is the original triality cubic
$F=\tau$, with its ordered octonionic arguments as in
\cref{wcl:sec:construction,wcl:sec:common-cubic}. In particular, on the
literal common lattice of \cref{wcl:thm:literal-intersection},
\[
 \mathcal I=\{\sqrt2(a,a,a):a\in D_4\},\qquad
 F(\sqrt2(a,a,a))=2\sqrt2\,P(a).
 \label{d4i:eq:retained-cubic}\tag{D4I3}
\]
No coordinate permutation, scaling of $P$, or alteration of the lattice
is made. The witnesses below are actual points of this original lattice.
The image proof is elementary and requires no external arithmetic
existence theorem.

\subsection{Necessary restrictions in the original coordinates}
\label{d4i:sec:necessary}

For every integer $b$, one has $b^2\equiv b\pmod2$. Consequently,
for $a\in D_4$,
\[
 a_0^2+a_1^2+a_2^2+a_3^2
 \equiv a_0+a_1+a_2+a_3\equiv0\pmod2.
\]
Writing $s=a_1^2+a_2^2+a_3^2$, the factor $a_0^2-3s$ is congruent
to $a_0^2+s\equiv0\pmod2$. Thus $P(a)$ is even.

Modulo three the original polynomial satisfies $P(a)\equiv a_0^3$.
If $3\mid P(a)$, the field $\mathbb Z/3\mathbb Z$ has no nonzero
element with zero cube, so $3\mid a_0$. Substitution of $a_0=3b$
into the original formula gives
\[
 P(a)=3b(9b^2-3s)=9b(3b^2-s).
 \label{d4i:eq:three-divisibility}\tag{D4I4}
\]
Therefore
\[
 P(a)\in2\mathbb Z,\qquad
                 3\mid P(a)\ \Longrightarrow\ 9\mid P(a).
 \label{d4i:eq:necessary-restrictions}\tag{D4I5}
\]

\subsection{Three explicit affine lines and their exact values}
\label{d4i:sec:three-lines}

Define the common direction
\[
 v=(-3,1,1,1).
 \label{d4i:eq:common-direction}\tag{D4I6}
\]
Its coordinate sum is zero, and direct substitution gives
\[
                   P(v)=-3\bigl(9-3(1+1+1)\bigr)=0.
\]
For every $r\in\mathbb Z$ and $\varepsilon\in\{-1,1\}$, put
\[
 L_\varepsilon(r)=(-3r-\varepsilon,r,r,r+\varepsilon)
                =rv+(-\varepsilon,0,0,\varepsilon).
 \label{d4i:eq:two-lines}\tag{D4I7}
\]
Its coordinate sum is
$-3r-\varepsilon+r+r+r+\varepsilon=0$, so it belongs to $D_4$.
The original last-three-coordinate square sum and first-coordinate
square are respectively
\[
 s=r^2+r^2+(r+\varepsilon)^2=3r^2+2r\varepsilon+1,
 \qquad (-3r-\varepsilon)^2=9r^2+6r\varepsilon+1.
\]
Thus the factor in the original cubic is exactly
\[
 a_0^2-3s=(9r^2+6r\varepsilon+1)
                 -(9r^2+6r\varepsilon+3)=-2,
\]
and therefore
\[
 P(L_\varepsilon(r))=(-3r-\varepsilon)(-2)=6r+2\varepsilon.
 \label{d4i:eq:two-line-values}\tag{D4I8}
\]

For every $h\in\mathbb Z$, define the third line
\[
 L_0(h)=(-3h,h+1,h-1,h)=hv+(0,1,-1,0).
 \label{d4i:eq:third-line}\tag{D4I9}
\]
Its coordinate sum is zero, so it also belongs to $D_4$. Here
\[
 s=(h+1)^2+(h-1)^2+h^2=3h^2+2,
\]
and consequently
\[
 P(L_0(h))=(-3h)\bigl(9h^2-3(3h^2+2)\bigr)
                         =(-3h)(-6)=18h.
 \label{d4i:eq:third-line-values}\tag{D4I10}
\]
These identities hold for all signed integer parameters. In particular,
$L_0(0)=(0,1,-1,0)$ has cubic value zero, and the zero vector also
has cubic value zero.

\subsection{The global image and a proved reconstruction map}
\label{d4i:sec:global-image}

Set
\[
 S=\{n\in\mathbb Z:n\in2\mathbb Z\text{ and }
                                  (3\nmid n\text{ or }9\mid n)\}.
 \label{d4i:eq:image-domain}\tag{D4I11}
\]

\begin{theorem}[Complete cubic image and explicit section]
\label{d4i:thm:complete-image}
The original polynomial has the exact image
\[
 P(D_4)=18\mathbb Z\ \cup\ (6\mathbb Z+2)\ \cup\ (6\mathbb Z-2)=S.
 \label{d4i:eq:complete-image}\tag{D4I12}
\]
Equivalently,
\[
 P(D_4)=\{n\in\mathbb Z:
                 n\bmod18\in\{0,2,4,8,10,14,16\}\}.
 \label{d4i:eq:residue-image}\tag{D4I13}
\]
The following map is an actual section $\sigma:S\to D_4$:
\[
 \sigma(n)=
 \begin{cases}
 (-3h,h+1,h-1,h),&n=18h,\\
 (-3r-\varepsilon,r,r,r+\varepsilon),
        &n=2(3r+\varepsilon),\quad\varepsilon\in\{-1,1\}.
 \end{cases}
 \label{d4i:eq:global-section}\tag{D4I14}
\]
Its displayed parameters are uniquely determined in their respective
cases, and $P\circ\sigma=\operatorname{id}_S$.
\end{theorem}

\begin{proof}
Equation~\eqref{d4i:eq:necessary-restrictions} proves
$P(D_4)\subseteq S$. For the reverse inclusion let $n\in S$.
If $3\mid n$, then $9\mid n$, and evenness gives $n=18h$ for
a unique $h\in\mathbb Z$. Equation~\eqref{d4i:eq:third-line-values}
gives $P(L_0(h))=n$.

If $3\nmid n$, write $n=2q$. Then $3\nmid q$, so there is a
unique $\varepsilon\in\{-1,1\}$ with
$q\equiv\varepsilon\pmod3$. The number
$r=(q-\varepsilon)/3$ is consequently a unique integer.
Equation~\eqref{d4i:eq:two-line-values} gives
\[
                    P(L_\varepsilon(r))=6r+2\varepsilon=2q=n.
\]
The two displayed cases are disjoint and cover $S$, and both
witnesses lie in $D_4$ by \cref{d4i:sec:three-lines}. These
calculations prove the asserted section identity and both inclusions
in \eqref{d4i:eq:complete-image}. The first arithmetic progression
contributes residue zero modulo eighteen; the second contributes
$2,8,14$; and the third contributes $16,4,10$. This proves
\eqref{d4i:eq:residue-image}.
\end{proof}

The classification includes every positive, zero, and negative value.
The full excluded set is the union of the odd integers with
$18\mathbb Z+6$ and $18\mathbb Z+12$.
Let
\[
 A=\{a\in\mathbb Z^4:a_0+a_1+a_2+a_3=0\}
 \label{d4i:eq:sum-zero-slice}\tag{D4I15}
\]
be the literal sum-zero sublattice in the original coordinates. The
three witnesses lie in $A\subseteq D_4$, so the two inclusions
$S\subseteq P(A)\subseteq P(D_4)=S$ prove
\[
                              P(A)=P(D_4)=S.
 \label{d4i:eq:slice-image}\tag{D4I16}
\]
The image of the original cubic $F=\tau$ on the common lattice is
therefore exactly
\[
\begin{aligned}
 F(\mathcal I)&=2\sqrt2\,S\\
 &=4\sqrt2\bigl((9\mathbb Z)\cup(3\mathbb Z+1)
                                     \cup(3\mathbb Z-1)\bigr).
\end{aligned}
 \label{d4i:eq:original-cubic-image}\tag{D4I17}
\]
Because the original off-diagonal Albert determinant is $2F$, its
corresponding image is $4\sqrt2\,S$. The additive group generated
by that determinant image is $8\sqrt2\mathbb Z$, while the
additive group generated by $F(\mathcal I)$ is $4\sqrt2\mathbb Z$.
Indeed every $P$ value is even, and the first two lines at parameter
zero attain $P=2$ and $P=-2$. Thus the generated additive group of
the $P$ values is exactly $2\mathbb Z$; applying the two original
scalar factors proves both group assertions without identifying an
image set with its generated group.

\subsection{Exact local images with explicit local sections}
\label{d4i:sec:local-images}

For a prime $p$, let
\[
 D_{4,p}=\{a\in\mathbb Z_p^4:
                            a_0+a_1+a_2+a_3\in2\mathbb Z_p\}.
 \label{d4i:eq:local-domain}\tag{D4I18}
\]
This is precisely $D_4\otimes_{\mathbb Z}\mathbb Z_p$ in the original
coordinates. To verify the equality directly, use the integer basis
\[
\begin{gathered}
 b_0=(1,-1,0,0),\quad b_1=(0,1,-1,0),\\
 b_2=(0,0,1,-1),\quad b_3=(0,0,1,1).
\end{gathered}
 \label{d4i:eq:integral-basis}\tag{D4I19}
\]
For any $a$ in the displayed definition of $D_{4,p}$, the unique
coefficients of these basis vectors are
\[
\begin{aligned}
 c_0&=a_0,&c_1&=a_0+a_1,\\
 c_3&=(a_0+a_1+a_2+a_3)/2,&c_2&=c_3-a_3.
\end{aligned}
 \label{d4i:eq:local-basis-inverse}\tag{D4I20}
\]
All belong to $\mathbb Z_p$, and substitution gives the original
four coordinates. Conversely a $\mathbb Z_p$-linear combination
of the four vectors has coordinate sum $2c_3$. These two direct
calculations prove the tensor-product identification, including
at $p=2$. Taking integral coefficients gives the same basis
description of the original $D_4$.

The three line identities of \cref{d4i:sec:three-lines} hold over
$\mathbb Z_p$, since they are integer polynomial identities. Each
line stays in the sum-zero submodule of $D_{4,p}$.

\begin{theorem}[All local images and the exact local-to-global statement]
\label{d4i:thm:local-images}
For the unchanged polynomial $P$,
\[
\begin{aligned}
 P(D_{4,2})&=2\mathbb Z_2,\\
 P(D_{4,3})&=\mathbb Z_3^{\times}\cup9\mathbb Z_3,\\
 P(D_{4,p})&=\mathbb Z_p\qquad(p\ge5).
\end{aligned}
 \label{d4i:eq:local-images}\tag{D4I21}
\]
An integer $n$ belongs to every local image if and only if
$n\in S$, and hence if and only if $n\in P(D_4)$. The global
section is exactly \eqref{d4i:eq:global-section}; explicit local
sections are given in the proof.
\end{theorem}

\begin{proof}
For $p=2$, the parity calculation of \cref{d4i:sec:necessary}, reduced
modulo two, gives $P(D_{4,2})\subseteq2\mathbb Z_2$. Conversely,
given an arbitrary $n=2q\in2\mathbb Z_2$, put
\[
 r=(q-1)/3\in\mathbb Z_2.
\]
Division by three is valid since three is a unit in $\mathbb Z_2$.
Then $L_1(r)\in D_{4,2}$ and
\[
                         P(L_1(r))=6r+2=2q=n.
\]
Thus $n\mapsto L_1((n/2-1)/3)$ is an explicit section on
$2\mathbb Z_2$, and proves the first equality.

For $p=3$, reduction of the original polynomial modulo three and
the substitution $a_0=3b$ in
\eqref{d4i:eq:three-divisibility} prove that any value divisible by
three belongs to $9\mathbb Z_3$. Every unit
$n\in\mathbb Z_3^{\times}$ is congruent modulo three to exactly
one of $2\varepsilon$, with $\varepsilon\in\{-1,1\}$. Set
\[
 r=(n-2\varepsilon)/6\in\mathbb Z_3.
\]
This quotient is integral because the numerator is divisible by
three and two is a unit. The line $L_\varepsilon(r)$ then has
$P$ value $6r+2\varepsilon=n$.
Every $n\in9\mathbb Z_3$ can be written uniquely as $18h$ with
$h=n/18\in\mathbb Z_3$, again because two is a unit.
The third line has $P(L_0(h))=18h=n$. These two disjoint cases
give the explicit local section
\[
 n\longmapsto
 \begin{cases}
 L_\varepsilon((n-2\varepsilon)/6),
        &n\in\mathbb Z_3^{\times},\quad n\equiv2\varepsilon\pmod3,\\
 L_0(n/18),&n\in9\mathbb Z_3.
 \end{cases}
 \label{d4i:eq:three-adic-section}\tag{D4I22}
\]
They prove the second equality in \eqref{d4i:eq:local-images}.

For every prime $p\ge5$, six is a unit. Given any
$n\in\mathbb Z_p$, choose $r=(n-2)/6\in\mathbb Z_p$.
The point $L_1(r)$ belongs to $D_{4,p}$ and has
$P(L_1(r))=6r+2=n$. Thus $n\mapsto L_1((n-2)/6)$ is a
section on all of $\mathbb Z_p$, proving the third equality.

For integers $n$, the first local condition is precisely evenness;
the second is precisely $3\nmid n$ or $9\mid n$; and the other
primes impose no condition. Membership in every local image is
therefore exactly $n\in S$. The explicit integral section in
\cref{d4i:thm:complete-image} proves that those conditions suffice
for global integral representability. This proves the local-to-global
assertion without approximation or inference from a finite modular
calculation.
\end{proof}

\subsection{The exact factorization and quotient metric on the sum-zero slice}
\label{d4i:sec:factorization}

Define the integral map into the original sublattice $A$ of
\eqref{d4i:eq:sum-zero-slice} by
\[
 \Phi:\mathbb Z^3\longrightarrow A,\qquad
 \Phi(w,u,v)=(-3w-u-v,w+u,w+v,w).
 \label{d4i:eq:slice-map}\tag{D4I23}
\]
Its coordinate sum is zero. For $a\in A$, its proposed inverse is
\[
                 \Phi^{-1}(a)=(a_3,a_1-a_3,a_2-a_3).
 \label{d4i:eq:slice-inverse}\tag{D4I24}
\]
Substitution returns the original second, third, and fourth
coordinates immediately; the first becomes $-a_1-a_2-a_3=a_0$
by the defining equation of $A$. Substituting $\Phi$ into the
inverse returns $(w,u,v)$. These two equalities prove that $\Phi$
is an exact integral-module isomorphism.

For $a=(a_0,y,z,w)\in A$, put $T=y+z+w=-a_0$. Direct expansion
of the original polynomial gives
\[
\begin{aligned}
 P(a)&=-T\bigl(T^2-3(y^2+z^2+w^2)\bigr)\\
     &=2T(y^2+z^2+w^2-yz-yw-zw)\\
     &=T\bigl((y-z)^2+(y-w)^2+(z-w)^2\bigr).
\end{aligned}
 \label{d4i:eq:slice-expansion}\tag{D4I25}
\]
Under $\Phi$, one has $T=3w+u+v$, and the three differences are
$u-v,u,v$. Consequently
\[
                 P(\Phi(w,u,v))=2(3w+u+v)(u^2-uv+v^2).
 \label{d4i:eq:factorization}\tag{D4I26}
\]
This identity explains the three original lines without changing
their coordinates. On $L_\varepsilon(r)$, the quotient differences
are $(u,v)=(-\varepsilon,-\varepsilon)$, so
$u^2-uv+v^2=1$; on $L_0(h)$, they are $(u,v)=(1,-1)$, so
$u^2-uv+v^2=3$. The original first coordinates remain
$-3r-\varepsilon$ and $-3h$, respectively.

The exact quotient map is
\[
 \chi:A\longrightarrow\mathbb Z^2,\qquad
                  \chi(a)=(a_1-a_3,a_2-a_3).
 \label{d4i:eq:quotient-map}\tag{D4I27}
\]
It is onto because $\chi\Phi(0,u,v)=(u,v)$. An element lies in
its kernel exactly when its last three coordinates all equal a
common integer $w$; its first coordinate is then $-3w$.
Thus $\ker\chi=\mathbb Z(-3,1,1,1)$. The section
$(u,v)\mapsto\Phi(0,u,v)$ proves the split exact sequence
\[
 0\longrightarrow\mathbb Z(-3,1,1,1)\longrightarrow A
       \xrightarrow{\ \chi\ }\mathbb Z^2\longrightarrow0.
 \label{d4i:eq:split-exact-sequence}\tag{D4I28}
\]

For $\omega=(-1+\sqrt{-3})/2$, one has
$\omega+\overline\omega=-1$ and $\omega\overline\omega=1$.
Direct multiplication therefore proves
\[
                 (u+v\omega)(u+v\overline\omega)=u^2-uv+v^2.
 \label{d4i:eq:eisenstein-norm}\tag{D4I29}
\]
The quadratic factor in the original cubic is exactly this
Eisenstein norm, with the preceding coefficient two unchanged.

The original Euclidean metric on these same coordinates is,
by expansion,
\[
 \|\Phi(w,u,v)\|^2
            =12w^2+8w(u+v)+2(u^2+uv+v^2).
 \label{d4i:eq:full-slice-metric}\tag{D4I30}
\]
Write $d=(-3,1,1,1)$ for the direction when used in the metric
calculation. Its squared norm is twelve, and the exact pairing is
\[
          \langle\Phi(w,u,v),d\rangle=12w+4(u+v).
 \label{d4i:eq:direction-pairing}\tag{D4I31}
\]
Hence orthogonal projection away from the real line $\mathbb R d$
is given in the original coordinates by
\[
\begin{aligned}
 \pi\Phi(w,u,v)
  &=\Phi(w,u,v)-\left(w+\frac{u+v}{3}\right)d\\
  &=\left(0,\frac{2u-v}{3},\frac{2v-u}{3},-\frac{u+v}{3}\right).
\end{aligned}
 \label{d4i:eq:orthogonal-projection}\tag{D4I32}
\]
This depends only on the quotient coordinate $\chi(a)=(u,v)$.
Its squared norm is
\[
\begin{aligned}
 \|\pi\Phi(w,u,v)\|^2
  &=\frac{(2u-v)^2+(2v-u)^2+(u+v)^2}{9}\\
  &=\frac23(u^2-uv+v^2).
\end{aligned}
 \label{d4i:eq:projected-norm}\tag{D4I33}
\]
The quotient metric on $A/\mathbb Z d$, realized by this exact
projection, thus has Gram matrix
\[
              \frac13\begin{pmatrix}2&-1\\-1&2\end{pmatrix}
 \label{d4i:eq:quotient-gram}\tag{D4I34}
\]
in the displayed quotient coordinates $(u,v)$.
The retained common-lattice map $a\mapsto\sqrt2(a,a,a)$ multiplies
every squared norm by six, as proved in
\cref{wcl:thm:literal-intersection}. Consequently the corresponding
quotient Gram matrix on the image of $A$ in $\mathcal I$ is
\[
                         \begin{pmatrix}4&-2\\-2&4\end{pmatrix}.
 \label{d4i:eq:common-lattice-quotient-gram}\tag{D4I35}
\]
This is the exact quotient-metric calculation on the indicated
rank-three slice. It neither identifies that slice with the full
$D_4$ nor changes the original common-lattice metric.

% Integration-ready proof module; no canonical file has been edited.
% Requires amsmath, amssymb, amsthm, and hyperref in a standalone host.
% All coordinates and quadratic-form conventions are retained explicitly.

\section{The marked Niemeier lattice, its twelve theta indices, and the umbral shadow}
\label{nus:sec:shadow}

The original objects are
\[
 A=\{(x_0,\ldots,x_5)\in\mathbb Z^6:\sum_jx_j=0\},\qquad
 D=\{d\in\mathbb Z^4:\sum_id_i\equiv0\pmod2\},\qquad
 R=A^{\perp4}\perp D.
\]
The metrics are the displayed Euclidean metrics.  The four copies of $A$
retain their original fibre order $f_i289^j$ for $0\le j<6$, with fibre
labels $(00,10,01,11)$.  Their discriminant generator is
$\lambda=(5,-1,-1,-1,-1,-1)/6$.  The discriminant of $D$ retains
$10\mapsto v=(1,0,0,0)$, $01\mapsto s=(1,1,1,1)/2$, and
$11\mapsto c=(-1,1,1,1)/2$.  Consequently the original norm-valued
quadratic refinement, with codomain $\mathbb Q/2\mathbb Z$, is
\[
 q_R(k_1,k_2,k_3,k_4;y)
 =\frac56\sum_i k_i^2+\mathbf1_{y\ne0}\pmod{2\mathbb Z}.
\]
The original glue and Niemeier lattice are
\[
 \begin{split}
 E&=\{u\in\mathbb F_2^4:\sum_i u_i=0\},\qquad
 \tau(u)=(u_2+u_4,u_3+u_4),\\
 j(a,b)&=(a+b,a+2b,a,b),\qquad C_3=j(\mathbb F_3^2),\\
 H&=\{(3u+2j(a,b);\tau(u)):u\in E,\ (a,b)\in\mathbb F_3^2\},\\
 N&=\{x\in R^\vee:x+R\in H\}.
 \end{split}
\]
Thus $|H|=8\cdot9=72$ and $N/R=H$, with the root system
$X=A_5^4D_4$ proved in the retained lattice calculation.  In particular,
no norm-two vectors are introduced by the nonzero glue classes: their
original minima are four or six.  We use this proved literal root system,
rather than a rank-only identification, throughout the construction below.

\subsection{The primary prescription and its complete coefficient matrix}

Write $e(t)=\exp(2\pi\mathrm i t)$ and $q=e(\tau)$ for
$\tau\in\mathbb H$.  Cheng--Duncan--Harvey, arXiv:1307.5793v2,
Section~3.2 \cite[Sections~3.2--3.3]{esn:CDH2014}, source labels
\texttt{thetexp}, \texttt{transf\_theta},
and \texttt{transf\_theta\_2}, use
\[
 \theta_{m,r}(\tau,z)
 =\sum_{k\equiv r\ (2m)}q^{k^2/(4m)}e(kz),
 \qquad r\in\mathbb Z/(2m).
\]
Their Section~3.3, source label \texttt{def:OmegaMatrices}, defines
\[
 \Omega_m(n)_{r,s}
 =\begin{cases}
 1,&r+s\equiv0\pmod{2n},\quad r-s\equiv0\pmod{2m/n},\\
 0,&\text{otherwise},
 \end{cases}
 \qquad n\mid m.
\]
Their ADE table, source label \texttt{ADE1}, assigns
\[
 \Omega^{A_{m-1}}=\Omega_m(1),\qquad
 \Omega^{D_{m/2+1}}=\Omega_m(1)+\Omega_m(m/2),
\]
and extends by addition over components with the same Coxeter number.
These are the primary authors' prescribed maps.  The following calculation
specializes and proves their values in the present marked coordinates; it
does not assert a new ADE classification.

The six-cycle $e_j\mapsto e_{j+1\bmod6}$ on an $A$ fibre has
characteristic polynomial $(x^6-1)/(x-1)$ on its augmentation space.
Its order is six.  In $D$, take the original simple roots
$\alpha_1=e_1-e_2$, $\alpha_2=e_2-e_3$, $\alpha_3=e_3-e_4$,
$\alpha_4=e_3+e_4$.  Their reflection product
$s_{\alpha_1}s_{\alpha_2}s_{\alpha_3}s_{\alpha_4}$ sends
\[
 e_1\mapsto e_2,\quad e_2\mapsto e_3,\quad
 e_3\mapsto-e_1,\quad e_4\mapsto-e_4.
\]
It has order six and characteristic polynomial
\[
 (x^3+1)(x+1)
 =(x^2-x+1)(x+1)^2
 =\frac{(x^2-1)(x^6-1)}{(x-1)(x^3-1)}.
\]
Thus the original five components have Coxeter number $h=m=6$, and the
primary theta prescription has precisely $2h=12$ indices and denominator
$4h=24$ in its exponents.  The root rank remains $4\cdot5+4=24$;
the glue index remains $72$.

Put $J=\Omega_6(3)$.  Its two congruences are
\[
 r+s\equiv0\pmod6,\qquad r-s\equiv0\pmod4.
\]
The simultaneous solutions modulo twelve consist of the single value
$s\equiv5r\pmod{12}$: this value satisfies both congruences, and the
difference of any two solutions is divisible by both six and four, hence
by twelve.  Therefore
\[
 \boxed{\Omega^X=5I_{12}+J,\qquad J_{r,s}=\mathbf1_{s\equiv5r\ (12)}.}
\]
In the full original index order $0,1,\ldots,11$ this is
\begingroup
\csname c@MaxMatrixCols\endcsname=12\relax
\[
\Omega^X=\begin{pmatrix}
6&0&0&0&0&0&0&0&0&0&0&0\\
0&5&0&0&0&1&0&0&0&0&0&0\\
0&0&5&0&0&0&0&0&0&0&1&0\\
0&0&0&6&0&0&0&0&0&0&0&0\\
0&0&0&0&5&0&0&0&1&0&0&0\\
0&1&0&0&0&5&0&0&0&0&0&0\\
0&0&0&0&0&0&6&0&0&0&0&0\\
0&0&0&0&0&0&0&5&0&0&0&1\\
0&0&0&0&1&0&0&0&5&0&0&0\\
0&0&0&0&0&0&0&0&0&6&0&0\\
0&0&1&0&0&0&0&0&0&0&5&0\\
0&0&0&0&0&0&0&1&0&0&0&5
\end{pmatrix}.
\]
\endgroup
The unary series and the actual shadow are
\[
 S_{6,r}(\tau)=\frac1{2\pi\mathrm i}
  \left.\frac{\partial\theta_{6,r}}{\partial z}\right|_{z=0}
 =\sum_{k\equiv r\ (12)}kq^{k^2/24},
 \qquad S^X=\Omega^X S_6.
\]
All series and their derivatives converge normally on compact subsets of
their stated domains, by Gaussian decay.  Substitution $k\mapsto-k$
proves $S_{6,-r}=-S_{6,r}$, so $S_{6,0}=S_{6,6}=0$.
Let $\mathcal A:\mathbb C^5\to\mathbb C^{12}$ insert zero in positions
zero and six and put $(\mathcal Au)_{12-r}=-u_r$ for $1\le r\le5$.
This retains, and explicitly reconstructs, all twelve components.  Then
\[
 \Omega^X\mathcal A=\mathcal A M^X,\qquad
 \boxed{M^X=\begin{pmatrix}
 5&0&0&0&1\\0&4&0&0&0\\0&0&6&0&0\\
 0&0&0&4&0\\1&0&0&0&5
 \end{pmatrix}.}
\]
Indeed its $(r,s)$ entry is
$\Omega^X_{r,s}-\Omega^X_{r,12-s}$; direct substitution gives the matrix.
In particular the fourth A-type multiplicity and the full D-type
contribution have both been retained.  The diagonal entries
$(5,4,6,4,5)$ are exactly the Coxeter exponent multiplicities:
four copies of $(1,2,3,4,5)$ and one copy of $(1,3,3,5)$.

\subsection{Quadratic refinements and the finite theta transformations}

Let $L=\mathbb Z\ell$ with $(\ell,\ell)=12$.  Its dual is
$L^\vee=\mathbb Z\ell/12$.  Hence
\[
 A_L=L^\vee/L=\mathbb Z/12,\qquad
 q_L(r)=r^2/12\pmod{2\mathbb Z},\qquad
 Q_L(r)=q_L(r)/2=r^2/24\pmod{\mathbb Z},
\]
and the associated bilinear form is $B_L(r,s)=rs/12\pmod{\mathbb Z}$.
The symbols $q_L$ and $Q_L$ deliberately have different codomains and
retain the factor two.

Define the unitary Fourier matrices and reflection by
\[
 (F_\pm)_{r,s}=12^{-1/2}e(\pm rs/12),\qquad
 R_{r,s}=\mathbf1_{s\equiv-r\ (12)}.
\]
Finite geometric-series summation gives
\[
 F_\pm^2=R,\qquad F_-=RF_+=F_+R,\qquad F_+F_-=I.
\]
The substitution $k\mapsto k+12$ shows that $e(k^2/24)$ depends only on
$r=k\bmod12$.  Thus, with
$D_{r,s}=e(r^2/24)\delta_{r,s}$,
\[
 \theta_6(\tau+1,z)=D\theta_6(\tau,z).
\]
Poisson summation of the Gaussian gives
\begin{align*}
 \theta_6(-1/\tau,z/\tau)
 &=\sqrt{-\mathrm i\tau}\,e(6z^2/\tau)F_-\theta_6(\tau,z),\\
 \theta_6(-1/\tau,-z/\tau)
 &=\sqrt{-\mathrm i\tau}\,e(6z^2/\tau)F_+\theta_6(\tau,z).
\end{align*}
Here $\sqrt{-\mathrm i\tau}$ is the holomorphic square root positive
when $\tau$ is positive imaginary.  For completeness, apply the Fourier
transform convention $\widehat f(u)=\int_{\mathbb R}f(x)e(-ux)\,dx$
to $f(x)=\exp(\pi\mathrm i\tau x^2/12+2\pi\mathrm i zx)$.
Completing the square evaluates its transform; the real positive Gaussian
integral is obtained by squaring the integral and using polar coordinates,
and holomorphic continuation gives the expression for $\Im\tau>0$.
Periodizing $f(r+12x)$, its $n$-th Fourier coefficient is
$12^{-1}e(nr/12)\widehat f(n/12)$.  Summing these coefficients at zero
and grouping $n$ by its residue modulo twelve gives the first displayed
transformation.  Absolute Gaussian convergence justifies the integration,
grouping, and Fourier evaluation.  Replacing $z$ by $-z$ and using
$\theta_{6,-r}(\tau,z)=\theta_{6,r}(\tau,-z)$ gives the second.
The second, with its plus-sign Fourier matrix, is precisely CDH's
convention; the first uses the usual generator
$S=\left(\begin{smallmatrix}0&-1\\1&0\end{smallmatrix}\right)$
with elliptic argument $z/\tau$.

The matrices carrying the finite multiplier for this first convention are
\[
 U(T)=D,\qquad U(S)=e(-1/8)F_-.
\]
The square-root weight factor has not been included in $U$.
Differentiation of the theta transformation at $z=0$ gives
\[
 S_6(-1/\tau)=\tau\sqrt{-\mathrm i\tau}F_-S_6(\tau),\qquad
 S_6(\tau+1)=D S_6(\tau).
\]
Equivalently, CDH's reflected convention gives
$S_6(-1/\tau)=-\tau\sqrt{-\mathrm i\tau}F_+S_6(\tau)$.
The equality follows from $RS_6=-S_6$.  Thus the unary series has weight
$3/2$ with the same finite multiplier; theta has weight $1/2$.

The permutation $J$ preserves $Q_L$ because
$((5r)^2-r^2)/24=r^2\in\mathbb Z$.  Since $5^{-1}=5$ modulo twelve,
\[
 (JF_\pm)_{r,s}=12^{-1/2}e(\pm5rs/12)
 =(F_\pm J)_{r,s}.
\]
Therefore $\Omega^X$ commutes with both finite multiplier matrices.
The shadow transformation follows by applying this proved intertwiner.
Its component series vanish at infinity since every nonzero summand has
strictly positive exponent; the modular transformation shows the same
cuspidal behavior at every cusp of the full modular group.

There is also an exact comparison with the original A5 discriminant,
without replacing its form.  Reduction
\[
 p:\mathbb Z/12\longrightarrow\mathbb Z/6,\qquad p(r)=r\bmod6
\]
has kernel $\{0,6\}$ and satisfies
\[
 q_A(p(r))=\frac56r^2=10q_L(r)\pmod{2\mathbb Z}.
\]
Thus it is a surjective quadratic similitude with precisely multiplier
ten, and its bilinear forms satisfy the same multiplier relation.
It is not asserted to be an isometry.  Applying $p$ to four coordinates
and the identity to the retained $D$ discriminant gives
\[
 \widetilde D=(\mathbb Z/12)^4\oplus\mathbb F_2^2
 \longrightarrow D_R,
 \qquad (r;y)\longmapsto(p(r_1),\ldots,p(r_4);y).
\]
The pullback quadratic refinement is exactly
$10\sum_iq_L(r_i)+\mathbf1_{y\ne0}$, its radical is
$\{0,6\}^4\oplus0$, and its quotient by that radical is the original
$(D_R,q_R)$.  To prove the radical claim, pairing with each coordinate
generator forces $5r_i/6\in\mathbb Z$, hence $6\mid r_i$; the
nondegenerate $D$ form forces $y=0$.  The vectors so obtained have
quadratic value $30$ in each nonzero coordinate, hence zero modulo two.
The inverse image of the original $H$ has order $16\cdot72=1152$ and
is isotropic for this pullback form.  This gives the exact relation of the
theta indices to all original glue coordinates, with the form and
multiplier specified on both sides.

\subsection{The full Fourier determinant and the retained order-twelve character}
\label{nus:subsec:determinant}

Let $\chi:\mathrm{SL}_2(\mathbb Z)\to\mathbb Z/12$ be the retained
marking
\[
 \chi(T)=-1,\qquad \chi(S)=3.
\]
These values respect the presentation
$S^4=I$, $(ST)^3=S^2$: in the abelian target the required relations
are $4\cdot3=0$ and $3(3-1)=2\cdot3$ modulo twelve.
They therefore define a character, and its value at $T$ makes it onto.
We now calculate the determinant of the entire twelve-dimensional
finite theta representation, rather than that of a reduced component
list.

Put $\zeta=e(1/12)$.  The Vandermonde determinant gives
\[
 \det F_+=12^{-6}\prod_{0\le j<k\le11}(\zeta^k-\zeta^j).
\]
For every such pair,
\[
 \zeta^k-\zeta^j
 =2\mathrm i\sin\!\left(\frac{\pi(k-j)}{12}\right)
       \exp\!\left(\frac{\pi\mathrm i(j+k)}{12}\right),
\]
whose sine is strictly positive.  There are $\binom{12}{2}=66$ pairs,
and each index occurs in eleven of them, so
\[
 \sum_{j<k}(j+k)=11\sum_{j=0}^{11}j=11\cdot66=726.
\]
The phase of the determinant is consequently
\[
 \mathrm i^{66}\exp(726\pi\mathrm i/12)=(-1)\mathrm i=-\mathrm i.
\]
Its remaining factor is positive.  Unitarity, proved by the finite
geometric sum above, makes its absolute value one.  It follows exactly
that
\[
 \boxed{\det F_+=-\mathrm i,\qquad\det F_-=\mathrm i.}
\]
In particular retaining CDH's reflection changes the Fourier determinant
by $\det R=(-1)^5=-1$; it cannot be dropped.

The sum of the squares of all twelve indices is
\[
 \sum_{r=0}^{11}r^2=\frac{11\cdot12\cdot23}{6}=506.
\]
Consequently
\[
 \boxed{\det U(T)=e(506/24)=e(1/12),\qquad
 \det U(S)=e(-12/8)\det F_-=-\mathrm i=e(-3/12).}
\]
For an explicit check that the finite matrices have the relevant
metaplectic relations, let
$G=\sum_{r=0}^{11}e(r^2/24)$.  Listing the squares modulo twenty-four
gives
\[
 G=1+4e(1/24)+2e(1/6)+2e(3/8)+2e(2/3)+e(1/2).
\]
The first and last terms cancel, as do the terms with arguments $1/6$
and $2/3$.  The elementary sine and cosine values at $\pi/12$ and
$3\pi/4$ give
\[
 4e(1/24)+2e(3/8)=\sqrt6(1+\mathrm i)
 =\sqrt{12}\,e(1/8).
\]
Completing the square in the finite sum proves
\[
 (F_-D)^2=e(1/8)D^{-1}F_-,\qquad
 (F_-D)^3=e(1/8)R.
\]
For the first equality, its $(r,s)$ entry before evaluation is
$12^{-1}e(s^2/24)\sum_te(t^2/24-t(r+s)/12)$; shifting $t$ by
$r+s$ produces exactly the stated Gauss sum and matrix entry.
The second follows by multiplication and $F_-^2=R$.
It follows that
\[
 (U(S)U(T))^3=U(S)^2=-\mathrm iR,\qquad U(S)^4=-I_{12}.
\]
Taking determinants kills this metaplectic central sign because the
dimension twelve is even.  The determinant therefore descends to an
ordinary character of $\mathrm{SL}_2(\mathbb Z)$.  Its two computed
generator values, together with generation by $S,T$, prove
\[
 \boxed{\det U(\gamma)=e(-\chi(\gamma)/12).}
\]
The dual determinant is its inverse and is exactly
$e(\chi(\gamma)/12)$, with the retained signs $T\mapsto e(-1/12)$,
$S\mapsto\mathrm i$.

This is a statement about the determinant multiplier, with its domain
and codomain specified.  It is not a statement that a vector of weight
$1/2$ has scalar weight one after taking a determinant.  The separate
Dedekind-eta identities
\[
 \eta(\tau+1)=e(1/24)\eta(\tau),\qquad
 \eta(-1/\tau)=\sqrt{-\mathrm i\tau}\,\eta(\tau)
\]
show that $\eta^2$ has weight one and multiplier $\det U$, whereas
$\eta^{-2}$ has weight minus one and multiplier $(\det U)^{-1}$.
Indeed $\eta^2(S\tau)=(-\mathrm i)\tau\eta^2(\tau)$ and
$\eta^{-2}(S\tau)=\mathrm i\tau^{-1}\eta^{-2}(\tau)$.
The integer-weight factors $(c\tau+d)^{\pm1}$ obey the cocycle identity
by matrix multiplication, so the two generator formulas prove the
statements for every $\gamma$.  At $-I$ the character and odd integer
weight each contribute a minus sign, whose product is one.

\subsection{The residue Fourier map and what the Weyl quotient retains}

The original marking is
$\kappa(289^j)=j\pmod6$, with the full ordered residue list
\[
 (289^j)_{j=0}^5=(1,289,361,169,121,529)\pmod{840}.
\]
On the six-dimensional complex coordinate carrier define
\[
 E_2:\mathbb C[S_{840}]\longrightarrow\mathbb C[\mathbb Z/12],
 \qquad E_2(e_{289^j})=e_{2j},
\]
and on a separately indexed $\mathbb C^6$ define
\[
 B(e_k)=\frac{e_k+e_{k+6}}{\sqrt2},\qquad0\le k<6.
\]
Both are injective and preserve the indicated coordinate Hermitian
products.  Retain the original negative-sign normalized Fourier matrix
$F_6^-=(6^{-1/2}e(-kj/6))_{k,j}$.  Entry-by-entry calculation gives
\[
 \boxed{F_-E_2=B F_6^-.}
\]
Indeed the $(r,j)$ entry of the left side is
$12^{-1/2}e(-rj/6)$, and the right side has the same value, using
$r\bmod6$.  Thus the original unnormalized sample transform is carried
into this formula by its exact factor $\sqrt6$.
Translation $j\mapsto j+1$ intertwines under $E_2$ with translation
$r\mapsto r+2$ modulo twelve.  This gives a literal Fourier and
translation map, not an identification of different quadratic forms.
In particular $Q_L(6)=3/2=1/2\pmod{\mathbb Z}$, so the quotient
$\mathbb Z/12\to\mathbb Z/6$ does not carry $Q_L$ itself to a
quadratic form on the quotient.  Its exact multiplier-ten comparison
with the original A5 form was proved above.

For the original lattice action, $\rho(289)$ is the same six-cycle in
each of the four $A$ coordinate sets and is the identity on $D$.
Each coordinate transposition of $A$ is reflection in the corresponding
root $e_i-e_j$, so the full residue action is in
$W(R)=W(A)^4\times W(D)$.  It acts trivially on every A discriminant
because
$Q_\sigma\lambda-\lambda=e_{\sigma(0)}-e_0\in A$, and it fixes
the D discriminant.  Consequently
\[
 S_{840}\xrightarrow{\rho}O(N)\longrightarrow O(N)/W(R)
\]
is the trivial homomorphism.  This proves exactly which residue-group
action is killed, without discarding its already constructed Fourier map.

The original marked quotient is
\[
 G=O(N)/W(R)\cong\mathrm{GL}_2(\mathbb F_3).
\]
Here its actual coordinate map uses
$\ell_1=(1,1)$, $\ell_2=(1,2)$, $\ell_3=(1,0)$,
$\ell_4=(0,1)$.  For $A\in\mathrm{GL}_2(\mathbb F_3)$ determine
the unique signs and permutation by
\[
 \ell_iA=\varepsilon_i\ell_{p(i)},\qquad
 (P_pu)_i=u_{p(i)},\qquad B_p\tau(u)=\tau(P_pu).
\]
The four row lines exhaust the projective line, so the signs and
permutation exist uniquely.  The binary equation defines $B_p$ uniquely
since $\tau$ is onto.  Then $(k;y)\mapsto(\varepsilon_i k_{p(i)};B_py)$
preserves precisely the original two glue codes, and its ternary
restriction is $M_gj=jA$.  Conversely a marked glue stabilizer gives this
same $A$.  Multiplication of these identities proves the group
isomorphism, with the original row/column convention retained.

To compare its chosen signed lifts with CDH's simple-root permutation
section, let $Q_{\rm rev}e_j=e_{5-j}$ in an $A$ fibre.  When
$\varepsilon_i=-1$, replace the representative $-I$ by
$-Q_{\rm rev}$; their quotient is $Q_{\rm rev}\in W(A)$ and
$-Q_{\rm rev}(e_j-e_{j+1})=e_{4-j}-e_{5-j}$.
Thus the resulting map permutes the original simple roots by diagram
reversal.  When $\varepsilon_i=1$ no change is needed.
For D, put $\Delta=(\alpha_1\ \alpha_2\ \alpha_3\ \alpha_4)$.
Associate the original classes $v,s,c$ with outer nodes $1,4,3$,
respectively; their fundamental weights are
$e_1$, $(1,1,1,1)/2$, and $(1,1,1,-1)/2$, the last of which differs
from the retained $c$ by $(1,0,0,-1)\in D$.
For $B_p$, let $P_{B_p}$ permute these outer nodes exactly as $B_p$
permutes $v,s,c$ and fix node two.  Then
\[
 d_{B_p}=\Delta P_{B_p}\Delta^{-1}
\]
is the exact original-coordinate diagram representative.  Its
orthogonality follows because permuting the outer nodes preserves the
original Gram matrix.  It preserves $D$, and its action on fundamental
weights proves its discriminant action is exactly $B_p$.
Its quotient with the originally chosen D lift has trivial discriminant
action, hence belongs to the proved Weyl kernel $W(D)$.
This specifies the actual change of section on every summand.

Define the following four class functions using those retained maps:
\[
 a_g=\#\{i:p(i)=i\},\quad
 b_g=\sum_{p(i)=i}\varepsilon_i,\quad
 s_g=\operatorname{sgn}(B_p|_{\{v,s,c\}}),\quad
 c_g=\#\operatorname{Fix}(B_p|_{\{v,s,c\}})-1.
\]
The first two are the traces on the four-component permutation and
signed permutation representations.  The last is the trace on the
two-dimensional sum-zero subspace of the three outer-node permutation
representation, and $s_g$ is its sign character.  Thus these are precisely
CDH's $(\bar\chi^{X_A},\chi^{X_A},\chi^{X_D},\check\chi^{X_D})$;
the remaining $\bar\chi^{X_D}$ is one because the single D component
is fixed.  The comparison of sections above proves this identification
for the actual marked quotient.

CDH Section~5.1 \cite[Section~5.1]{esn:CDH2014} defines the full
twisted coefficient matrix by
\[
 (\Omega_g^{X_A})_{r,s}
 =\bigl(b_g\mathbf1_{r\ \text{even}}+
          a_g\mathbf1_{r\ \text{odd}}\bigr)\delta_{r,s},
\]
and, for its exceptional $m=6$ D4 case, by
\[
 (\Omega_g^{X_D})_{r,s}=
 \begin{cases}
 c_g\delta_{r,s},&r\equiv0,3\pmod6,\\
 \delta_{r,s}+s_gJ_{r,s},&r\equiv1,5\pmod6,\\
 s_g\delta_{r,s}+s_gJ_{r,s},&r\equiv2,4\pmod6.
 \end{cases}
\]
All deltas in this formula have indices modulo twelve.  The exact maps
from the marked glue quotient to the prescribed twisted shadows are
\[
 g\longmapsto\Omega_g^X=\Omega_g^{X_A}+\Omega_g^{X_D},\qquad
 S_g^X=\Omega_g^XS_6.
\]
They are class-function constructions from traces of representations;
the assignment $g\mapsto\Omega_g^X$ is not asserted to be a group
representation.  On the odd five-component presentation the explicit
matrix is
\[
 \boxed{M_g^X=\begin{pmatrix}
 a_g+1&0&0&0&s_g\\
 0&b_g&0&0&0\\
 0&0&a_g+c_g&0&0\\
 0&0&0&b_g&0\\
 s_g&0&0&0&a_g+1
 \end{pmatrix}.}
\]
This follows by the same signed component reconstruction $\mathcal A$:
in rows two and four the two D terms cancel, in row three the contribution
is $c_g$, and rows one and five pair with one another.
For reference, all possible character tuples are
\[
\begin{array}{c|rrrr|r}
 [g]&a_g&b_g&s_g&c_g&\text{number of elements}\\\hline
 1A&4&4&1&2&1\\
 2A&4&-4&1&2&1\\
 2B&2&0&-1&0&12\\
 4A&0&0&1&2&6\\
 3A&1&1&1&-1&8\\
 6A&1&-1&1&-1&8\\
 8A,8B&0&0&-1&0&12
\end{array}
\]
with the CDH class names of their table
\texttt{tab:chars:eul:6}.  The formulas apply to each actual $A$, so
their validity does not depend on a choice of conjugacy-class names.
The finite checker obtains the table from all $48$ invertible matrices
over $\mathbb F_3$ and the original row forms.

The distinguished residue action has quotient $g=1$ and therefore
produces the untwisted shadow, while a product $\rho(r)L_g$ produces
the same $S_g^X$ for every $r$.  The four-component signed action and
the D4 triality quotient survive and enter the displayed matrix.  The
central element $-I\in\mathrm{GL}_2(\mathbb F_3)$, for example,
changes $b_g$ from four to minus four while fixing $a_g,s_g,c_g$;
it changes the even-index odd-space entries with exactly that sign.

The already marked composite
\[
 \mathrm{SL}_2(\mathbb Z)\xrightarrow{\chi}\mathbb Z/12
 \longrightarrow\mathbb Z/6\xrightarrow{j\mapsto289^j}S_{840}
\]
is onto, has kernel $\chi^{-1}(\{0,6\})$, and sends $T$ to
$289^5=529$.  Following it by the Weyl quotient kills it, by the
proved residue calculation.  Following $\chi$ into the dual determinant
multiplier retains its complete order twelve, by the determinant theorem.
These are explicit different outgoing morphisms from the same marking,
including their kernels and signs.

\subsection{The shadow-to-Fricke map with the original scalar seventy-two}

The Coxeter calculation above gives the original formal frame exponent
\[
 \pi_X=-5[1]+[2]-[3]+5[6].
\]
Let $\mathcal E(\sum a_d[d])=\prod_d\eta(d\tau)^{a_d}$.
It is a group homomorphism from addition of formal exponents to
multiplication of nonzero holomorphic functions, by the exponent laws.
The retained scalar function is
\[
 \boxed{T_6(\tau)=\mathcal E(-\pi_X)
 =\frac{\eta(\tau)^5\eta(3\tau)}
        {\eta(2\tau)\eta(6\tau)^5}.}
\]
The minus sign on $\pi_X$ is essential: $\mathcal E(\pi_X)=T_6^{-1}$.

Here is a direct coefficient map from the actual shadow matrix to this
eta function.  Its diagonal difference vector is
$\alpha=(5,4,6,4,5)$, and set $\alpha_0=0$ modulo six.
Define the Lambert-series map at this marked matrix by
\[
 \mathscr L(\Omega^X)(\tau)
 =1+\sum_{n\ge1}\alpha_{n\bmod6}
                    \frac{nq^n}{1-q^n}.
\]
Convergence is absolute and locally uniform since $|q|<1$ and
$\alpha$ is bounded.  The coefficient of $q^N$ is exactly
$\sum_{d\mid N}d\alpha_{d\bmod6}$; thus this specifies every
coefficient, not only its initial terms.
For all positive integers $n$ the retained multiplicities satisfy
\[
 \alpha_{n\bmod6}
 =5-\mathbf1_{2\mid n}+\mathbf1_{3\mid n}-5\mathbf1_{6\mid n}.
\]
Checking the six residues proves this identity.  The product definition
$\eta(\tau)=q^{1/24}\prod_{n\ge1}(1-q^n)$ and termwise logarithmic
differentiation then give
\[
 \boxed{f^X:=\mathscr L(\Omega^X)
 =-q\frac{d}{dq}\log T_6
 =5\lambda_6+\lambda_2-\lambda_3,}
\]
where the original functions on the right are
\[
 \lambda_d=q\frac{d}{dq}\log\frac{\eta(d\tau)}{\eta(\tau)}
 =\frac{d-1}{24}+\sum_{n\ge1}\sigma_1(n)(q^n-dq^{dn}).
\]
The constant is
$[5(6-1)+(2-1)-(3-1)]/24=1$, and the other coefficients agree by
the six-residue identity.  This proves the shadow-to-eta calculation
directly for every coefficient.

The precise primary connecting map is CDH
Section~4.2 \cite[Section~4.2]{esn:CDH2014}:
\[
 \sigma^X(\tau,z)=\sum_{r\ (12)}\overline{S_r^X(\tau)}
                                \theta_{6,r}(\tau,z),\qquad
 \mathscr S_{1,1}(\sigma^X)=f^X.
\]
Its source labels are \texttt{eqn:sigmaX}, \texttt{def:weight2},
\texttt{weight2:Aseries}, and \texttt{weight2\_EZ}.
The displayed Lambert map is its explicit specialization at this
root-system shadow.  It verifies the resulting function independently
of the general Shimura-lift theorem.  In particular the shadow and the
retained Fricke function are connected by this complete coefficient map.

The original modular subgroup here is $\Gamma_0(6)$, and
$[\mathrm{SL}_2(\mathbb Z):\Gamma_0(6)]=12$.  To compute the index,
reduction modulo six gives
$\mathrm{SL}_2(\mathbb Z/6)\cong
\mathrm{SL}_2(\mathbb F_2)\times\mathrm{SL}_2(\mathbb F_3)$ of
order $6\cdot24=144$.  Reduction is onto: elementary upper and lower
unipotent matrices generate over each field, and CRT choices of their
entries realize each factor while acting trivially on the other.
The upper triangular subgroup modulo six has $\varphi(6)\cdot6=12$
elements (choose its first unit diagonal entry and its upper-right entry;
the other diagonal entry is its inverse).  Its inverse image is
$\Gamma_0(6)$, proving index $144/12=12$.
This modular-group index and the twelve theta indices have their
specified constructions; neither replaces the lattice index seventy-two.

Under the exact Fricke transformation $W_6\tau=-1/(6\tau)$,
the eta inversion formula sends each $\eta(dW_6\tau)$ to
$\sqrt{-\mathrm i(6/d)\tau}\,\eta((6/d)\tau)$.
The retained exponent vector at $d=1,2,3,6$ is $(5,-1,1,-5)$;
its reversed vector is its negative and its exponent sum is zero.
The $\tau$ factors therefore cancel.  Since each $6/d$ is positive,
the consistent square-root branch gives the positive scalar
\[
 \prod_{d\mid6}(6/d)^{a_d/2}
 =6^{5/2}3^{-1/2}2^{1/2}=72.
\]
Consequently
\[
 \boxed{T_6(W_6\tau)=72/T_6(\tau).}
\]
For the weight-two slash action of the determinant-six representative,
\[
 (f|_2W_6)(\tau)=(6\tau^2)^{-1}f(-1/(6\tau)).
\]
Differentiating the last identity and using
$d(W_6\tau)/d\tau=(6\tau^2)^{-1}$ gives
\[
 \boxed{f^X|_2W_6=-f^X.}
\]
On the Jacobi side, the primary Eichler--Zagier operator
$\mathcal W_6(6)$ acts on theta by $\Omega_6(6)=R$.
This is verified by its defining congruences, $r+s=0\pmod{12}$,
$r-s=0\pmod2$.  On the odd shadow and on the corresponding
skew-Jacobi form its eigenvalue is minus one.  The specialized coefficient
map just proved carries that negative eigenline to the negative Fricke
eigenline of $f^X$.  This proves an exact equivariant linear comparison
on these lines; it does not evaluate an $\mathrm{SL}_2(\mathbb Z)$
character on the determinant-six matrix.

Finally, the original sheet coordinates remain
\[
 J_6=T_6+5+72/T_6,\qquad Y_6=T_6-72/T_6.
\]
Direct substitution and differentiation prove
\[
 J_6(W_6\tau)=J_6(\tau),\quad
 Y_6(W_6\tau)=-Y_6(\tau),\quad
 Y_6^2=(J_6-5)^2-288,\quad
 dJ_6=Y_6\,dT_6/T_6=-2\pi\mathrm iY_6 f^X\,d\tau.
\]
Thus the original scalar $72$, the additive shift $5$, the branch
constant $288$, and the differential factor $-2\pi\mathrm i$ all
survive the complete shadow-to-sheet calculation.

\subsection{Established scope and reproducibility}

The constructed maps are the ADE coefficient specialization, the exact
finite-theta intertwiner, the multiplier determinant and its dual, the
quadratic similitude to the original discriminants, the marked glue-group
trace prescription, the original six-to-twelve Fourier map, and the
coefficient-level shadow-to-Fricke map.  Their kernels, weights, signs,
and constants were specified and proved above.  They establish concrete
relations even where a direct isometry or an identical group action is
absent.  They do not construct a new graded moonshine module or a new
moonshine theorem.  The name ``umbral shadow'' here refers to precisely
the CDH prescription, with its source provenance retained.

The standard-library program \texttt{check\_shadow.py} uses exact
arithmetic in $\mathbb Q[z]/(z^8-z^4+1)$ with $z=e(1/24)$.
It verifies both full and odd matrices, the Fourier determinant by the
Vandermonde product, the Gauss sum, the character signs, the Fourier
embedding, all $48$ original-code character records, the first $240$
Lambert coefficients, and the squared positive Fricke scalar $72^2$.
The proofs above establish the unbounded identities independently of
those finite checks.


% END ES-NIEMEIER HOST-ONLY INPUT.

\begingroup
\sloppy
\raggedright
\begin{thebibliography}{99}

\bibitem{Milnor1965}
J.~W. Milnor, \emph{Topology from the Differentiable Viewpoint},
University Press of Virginia, 1965, \S7, Theorem B and Lemma 4,
pp.~42--50.
\url{https://www.maths.ed.ac.uk/~aar/papers/milnortop.pdf}.



\bibitem{AlpogeS6_2026}
L.~Alp\"oge, with the disclosed Fable/Claude contribution,
\emph{The \((3,4,\infty)\) modular family of \(2\)-tori, completed at its
three special points, is a complex structure on \(S^6\)}, together with
the citable canonical audit and source bundle, Zenodo record 22238779
(2026), \href{https://doi.org/10.5281/zenodo.22238779}
{doi:10.5281/zenodo.22238779}.  Local audit locators are below
\texttt{supporting\_materials/workbench/research/}; the one-copy quotient
and framing discussion are in
\texttt{cusp\_sphere\_route\_strengthening.tex} and
\texttt{cusp\_s12\_carrier\_homotopy.tex}.

\bibitem{ConradSL2Z}
K.~Conrad, \emph{SL\(_2(\Z)\)}, author-hosted expository notes,
Theorem~1.1, pp.~1--2; Theorem~C.1, pp.~18--19,
\href{https://kconrad.math.uconn.edu/blurbs/grouptheory/SL(2,Z).pdf}
{author's PDF}, accessed 5 September 2026.

\bibitem{BorelSerre1953}
A.~Borel and J.-P.~Serre,
\emph{Groupes de Lie et puissances r\'eduites de Steenrod},
Amer. J. Math. \textbf{75} (1953), no.~3, 409--448;
Sections~12 and~15,
\href{https://doi.org/10.2307/2372495}{doi:10.2307/2372495}.

\bibitem{HatcherSpectralSequences}
A.~Hatcher,
\emph{Spectral Sequences in Algebraic Topology}, Chapter~5 draft,
Section~5.1; Theorem~5.41(b), p.~573, and the suspension calculation,
pp.~577--578,
\href{https://pi.math.cornell.edu/~hatcher/AT/ATch5.pdf}
{author-hosted chapter PDF}.

\bibitem{Hatcher2002}
A.~Hatcher,
\emph{Algebraic Topology},
Cambridge University Press, Cambridge, 2002; the opening
low-dimensional homotopy-group table, p.~339,
\href{https://pi.math.cornell.edu/~hatcher/AT/ATpage.html}
{author's electronic edition}.

\bibitem{Husemoller1994}
D.~Husemoller,
\emph{Fibre Bundles}, third edition,
Graduate Texts in Mathematics, vol.~20, Springer, New York, 1994;
Chapter~17, Section~7, Theorem~7.2, pp.~274--275,
\href{https://doi.org/10.1007/978-1-4757-2261-1}
{doi:10.1007/978-1-4757-2261-1}.

\bibitem{IsaksenWangXu2023}
D.~C.~Isaksen, G.~Wang, and Z.~Xu,
\emph{Stable homotopy groups of spheres: from dimension 0 to 90},
Publ. Math. IH\'ES \textbf{137} (2023), 107--243; Table~1, p.~113,
\href{https://doi.org/10.1007/s10240-023-00139-1}
{doi:10.1007/s10240-023-00139-1}.

\bibitem{KomoriMatsumotoTsumura2023}
Y.~Komori, K.~Matsumoto, and H.~Tsumura,
\emph{The Theory of Zeta-Functions of Root Systems},
Springer Monographs in Mathematics, Springer, Singapore, 2023;
Section~10.1, pp.~213 and~216,
\href{https://doi.org/10.1007/978-981-99-0910-0}
{doi:10.1007/978-981-99-0910-0}.

\bibitem{MareWillems2013}
A.-L.~Mare and M.~Willems,
\emph{Topology of the octonionic flag manifold},
M\"unster J. Math. \textbf{6} (2013), 483--523,
\href{https://www.uni-muenster.de/FB10/mjm/vol_6/mjm_vol_6_14.pdf}
{published article};
\href{https://arxiv.org/abs/0809.4318}{arXiv:0809.4318}.

\bibitem{Toda1962}
H.~Toda,
\emph{Composition Methods in Homotopy Groups of Spheres},
Annals of Mathematics Studies \textbf{49},
Princeton University Press, Princeton, NJ, 1962;
Proposition~3.1,
\href{https://doi.org/10.1515/9781400882625}
{doi:10.1515/9781400882625}.

\bibitem{retbib:Canonical}
L.~Alp\"oge, with the disclosed Fable/Claude contribution,
\emph{The \((3,4,\infty)\) modular family of \(2\)-tori, completed at
its three special points, is a complex structure on \(S^6\)},
canonical audit and reconstruction source bundle, Zenodo record 22238779
(2026), original manuscript \S3.7 and \S\S4.1--4.5.
\url{https://doi.org/10.5281/zenodo.22238779}.

\bibitem{retbib:ConeRoutes}
\emph{Real \(S^{12}\) and \(S^{24}\) Routes from Alp\"oge's
Candidate \((3,4,\infty)\) Two-Torus Fibration},
companion research draft in the preceding canonical source bundle;
native quotient, localized cusp map, and framed-cusp calculation.
Current proof-bearing edition:
\url{https://www.overleaf.com/read/rtmyqxyrzprn}.

\bibitem{retbib:Milnor1965}
J.~W.~Milnor, \emph{Topology from the Differentiable Viewpoint},
University Press of Virginia, 1965, \S7, Theorem B and Lemma 4.
\url{https://www.maths.ed.ac.uk/~aar/papers/milnortop.pdf}.

\bibitem{retbib:Pontryagin1950}
L.~S.~Pontryagin,
\emph{Homotopy classification of mappings of an \((n+2)\)-dimensional
sphere on an \(n\)-dimensional one},
Dokl. Akad. Nauk SSSR \textbf{70} (1950), 957--959.
\url{https://www.sas.rochester.edu/mth/sites/doug-ravenel/otherpapers/pont3.pdf}.

\bibitem{retbib:Pontryagin1955}
L.~S.~Pontryagin,
\emph{Smooth manifolds and their applications in homotopy theory},
Trudy Mat. Inst. Steklov \textbf{45} (1955);
English translation, Amer. Math. Soc. Transl. (2) \textbf{11} (1959).
\url{https://www.sas.rochester.edu/mth/sites/doug-ravenel/otherpapers/pont4.pdf}.

\bibitem{retbib:Johnson}
D.~Johnson, \emph{Spin structures and quadratic forms on surfaces},
J. London Math. Soc. (2) \textbf{22} (1980), 365--373.

\bibitem{retbib:DebrayGunningham}
A.~Debray and S.~Gunningham,
\emph{The Arf--Brown TQFT of pin-minus surfaces},
arXiv:1803.11183, especially \S3.1.1 and Proposition 5.1.
\url{https://adebray.github.io/papers/arf_brown.pdf}.

\bibitem{retbib:Hatcher}
A.~Hatcher, \emph{Algebraic Topology}, Cambridge University Press,
2002; group table p.~339, Corollary 4.24, p.~360, and pp.~384--385.
\url{https://pi.math.cornell.edu/~hatcher/AT/AT.pdf}.

\bibitem{hbridgebib:HBRetained}
\emph{The retained-angle cusp quotient has nonzero top-cell class},
complete derivative-framing theorem and its based
Pontryagin--Thom proof, 5 September 2026;
companion to the canonical cusp/sphere-route draft.
Current cumulative source collection:
\url{https://www.overleaf.com/read/rtmyqxyrzprn}.

\bibitem{hbridgebib:HBHigher}
\emph{The higher torus wrap: obstruction, exact local lift, and the
Jordan--triality rung},
canonical higher-rung research draft, section on peripheral
modulus and the mod--24 boundary, 5 September 2026.
Current cumulative source collection:
\url{https://www.overleaf.com/read/rtmyqxyrzprn}.

\bibitem{hbridgebib:HBHatcher}
A.~Hatcher, \emph{Algebraic Topology}, Cambridge University Press,
2002, Corollary~4.24, p.~360, stable homotopy groups and the
integral relation \(\eta^3=12\nu\), pp.~384--385,
and Proposition~4.56 with proof on p.~388.
\url{https://pi.math.cornell.edu/~hatcher/AT/AT.pdf}.

% Host-only helper: the workbench supplies it in its central bibliography.
\bibitem{Yokota2009}
I.~Yokota, \emph{Exceptional Lie groups}, 2009;
Sections~1.10, 1.14, 1.16 and Theorem~2.7.1,
\url{https://arxiv.org/abs/0902.0431}.

% Bibliography entries for es_niemeier_shared.tex.
% Input once inside the host's existing thebibliography environment.
% This file deliberately does not open or close that environment.

\bibitem{esn:CDH2014}
M.~C.~N.~Cheng, J.~F.~R.~Duncan, and J.~A.~Harvey,
\emph{Umbral Moonshine and Niemeier Lattices},
Research in the Mathematical Sciences \textbf{1} (2014), Article~3;
Section~2.4, definition of the Weyl quotient and Table~3,
\url{https://arxiv.org/abs/1307.5793v2}.

\bibitem{esn:MareWillems2013}
A.-L.~Mare and M.~Willems,
\emph{Topology of the octonionic flag manifold},
M\"unster Journal of Mathematics \textbf{6} (2013), 483--523;
Definition~2.2, Propositions~2.4 and~2.7, pp.~491--492,
and equations (5.4)--(5.6), pp.~510--511,
\url{https://www.uni-muenster.de/FB10/mjm/vol_6/mjm_vol_6_14.pdf};
\url{https://arxiv.org/abs/0809.4318}.

\bibitem{esn:Residual2026}
The Clankers,
\emph{Residual divisor shells, support-homogeneous signs, and the
\(A_5^4D_4\) Niemeier refinement of the modulo-\(840\)
Erd\H{o}s--Straus datum},
April 2026, maintained-source snapshot of 5 August 2026;
theorems \texttt{thm:reduced-shell}, \texttt{thm:isotropic},
\texttt{thm:root-preservation}, and \texttt{thm:two-channel}.
The exact source title, statement excerpts, snapshot digest, and
counting correction are included in the accompanying
\emph{ES--Niemeier source provenance} \cite{esn:Provenance}.

\bibitem{esn:ESMain2026}
\emph{Exact support--glue--Fricke transport},
theorem \nolinkurl{thm:exact-support-glue-Fricke-transport}
in the maintained Erd\H{o}s--Straus Foundation source,
inspected 5 September 2026.
The exact earlier-glue generators and transporter, with their original
equation labels, are reproduced in the accompanying
\emph{ES--Niemeier source provenance} \cite{esn:Provenance}.

\bibitem{esn:Provenance}
\emph{ES--Niemeier source provenance},
companion document to \emph{The modulo-840 lattice, triality, and exact
shell maps}, 5 September 2026;
\nolinkurl{ES_NIEMEIER_PROVENANCE.md}, included with this module.
It supplies exact passages and labels from the maintained residual
note, the support--glue theorem, and the canonical
\emph{Higher torus wrapping and the octonionic 24-dimensional continuation}.


\bibitem{esn:ConwaySloane1982}
J.~H.~Conway and N.~J.~A.~Sloane,
\emph{Twenty-three constructions for the Leech lattice},
Proc. R. Soc. Lond. A \textbf{381} (1982), 275--283, pp.~277--280,
\url{https://doi.org/10.1098/rspa.1982.0071}.

\bibitem{esn:Borcherds1985}
R.~E.~Borcherds, \emph{The Leech lattice},
Proc. R. Soc. Lond. A \textbf{398} (1985), 365--376,
Sections~6--7,
\url{https://math.berkeley.edu/~reb/papers/leech/leech.pdf}.

\bibitem{esn:Wilson2009}
R.~A.~Wilson, \emph{Octonions and the Leech lattice},
Journal of Algebra \textbf{322} (2009), 2186--2190,
\url{https://doi.org/10.1016/j.jalgebra.2009.03.021}.
Author version submitted 18 December 2008, revised 20 March 2009,
Sections~2--4, preprint pp.~2--6,
\url{https://webspace.maths.qmul.ac.uk/r.a.wilson/pubs_files/octoLeech1rev.pdf}.

\end{thebibliography}
\endgroup

\end{document}
